\documentclass[11pt,oneside]{report}
\usepackage[margin=1.1in]{geometry}
\usepackage{amsmath,amssymb,amsthm,mathtools}
\usepackage{booktabs,threeparttable,graphicx,setspace,enumitem,longtable,array}
\usepackage{xcolor}
\usepackage{etoolbox}
\usepackage{placeins}
\AtBeginEnvironment{tabular}{\footnotesize}
\usepackage[authoryear,round]{natbib}
\usepackage[colorlinks=true,linkcolor=blue!50!black,citecolor=blue!50!black,urlcolor=blue!50!black]{hyperref}
\onehalfspacing
\emergencystretch=6em
\Urlmuskip=0mu plus 2mu\relax
\setcounter{tocdepth}{1}

\newtheorem{theorem}{Theorem}[section]
\newtheorem{proposition}[theorem]{Proposition}
\newtheorem{lemma}[theorem]{Lemma}
\newtheorem{corollary}[theorem]{Corollary}
\newtheorem{conjecture}[theorem]{Conjecture}
\theoremstyle{definition}
\newtheorem{definition}[theorem]{Definition}
\newtheorem{assumption}[theorem]{Assumption}
\theoremstyle{remark}
\newtheorem{remark}[theorem]{Remark}
\newtheorem{example}[theorem]{Example}
\newcommand{\E}{\mathbb{E}}
\newcommand{\R}{\mathbb{R}}
\newcommand{\KL}{D_{\mathrm{KL}}}
\newcommand{\lab}[1]{\textsc{[#1]}}

% status tags used in the new chapters
\newcommand{\tagbox}[2]{\fcolorbox{#1}{#1!7}{\textsf{\scriptsize #2}}}
\newcommand{\Proved}{\tagbox{green!45!black}{PROVED IN THIS THESIS}}
\newcommand{\ProvedIn}[1]{\tagbox{green!45!black}{PROVED IN #1}}
\newcommand{\Target}{\tagbox{blue!70!black}{TARGET: STATEMENT FIXED, PROOF TO BE WRITTEN}}
\newcommand{\Conj}{\tagbox{orange!85!black}{CONJECTURE}}
\newcommand{\Pending}{\tagbox{red!70!black}{PENDING DATA}}
\newcommand{\pendbox}[2]{\par\medskip\noindent\fcolorbox{red!70!black}{red!4}{\parbox{0.96\linewidth}{\small\textbf{\color{red!70!black}Pending --- #1.} #2}}\par\medskip}

% each chapter's references appear as a section at the end of that chapter
\renewcommand{\bibsection}{\section*{References}}

% numbers generated from committed result files (used in Chapters 4 and 5)
% generated by export_results_extra.py from the committed result files; do not edit
\newcommand{\EoneSaturated}{28}
\newcommand{\EoneSmallIid}{16}
\newcommand{\EonePholmOne}{0.98}
\newcommand{\DonePholmOne}{0.031}
\newcommand{\DonePholmTwo}{0.008}
\newcommand{\GoneTheta}{0}
\newcommand{\GoneThetaTLRzero}{383}
\newcommand{\GoneThetaTLRone}{461}
\newcommand{\GoneStressGuard}{$+$0.044}
\newcommand{\GoneStressGuardCI}{[$-$0.14,0.23]}
\newcommand{\GoneStressP}{0.71}
\newcommand{\GoneStressDates}{46}
\newcommand{\GoneAllGuard}{$+$0.42}
\newcommand{\GoneAllGuardCI}{[0.35,0.48]}
\newcommand{\GoneAllHours}{16{,}560}
\newcommand{\GoneAllRatio}{1.52}
\newcommand{\VonePairs}{20}
\newcommand{\VoneIncluded}{3}
\newcommand{\VoneCWp}{0.36}
\newcommand{\VoneDRtwoCI}{[$-$0.29,0.19]}
\newcommand{\VoneTestHours}{2{,}319}
\newcommand{\PilotSymbols}{30}
\newcommand{\PilotMarginX}{29}
\newcommand{\PilotMarginXsafe}{29}
\newcommand{\PilotMarginRet}{11}
\newcommand{\PilotNoFb}{27}
\newcommand{\PilotSame}{28}
\newcommand{\PilotFbRatio}{0.24}
\newcommand{\PilotFbNeg}{22}
\newcommand{\PilotIdentified}{0}
\newcommand{\PilotSaturated}{30}
\newcommand{\PilotCurveTen}{16}
\newcommand{\PilotCurveTwenty}{22}
\newcommand{\PilotCurveForty}{28}
\newcommand{\PilotIdMax}{2}
\newcommand{\PilotDrift}{13}
\newcommand{\PilotWithinSlope}{$-$0.34}
\newcommand{\PilotWithinIn}{24}
\newcommand{\PilotAcrossIn}{18}

% generated by b2_llm.py; do not edit
\newcommand{\BtwoControl}{15}\newcommand{\BtwoRamp}{30}\newcommand{\BtwoRampCrisis}{30}
\newcommand{\BtwoFA}{0}
\newcommand{\BtwoEarly}{30}\newcommand{\BtwoLead}{3.5}
\newcommand{\BtwoLater}{27}\newcommand{\BtwoEarlier}{1}\newcommand{\BtwoTies}{2}\newcommand{\BtwoDelay}{3.0}
\newcommand{\BtwoFixedOnly}{15}\newcommand{\BtwoLLMOnly}{0}\newcommand{\BtwoCrisesFast}{15}\newcommand{\BtwoCrisesLLMV}{0}
\newcommand{\BtwoBelow}{1}\newcommand{\BtwoAbove}{29}
\newcommand{\BtwoDiscFA}{14}\newcommand{\BtwoCalls}{252{,}798}\newcommand{\BtwoParseFail}{0}
\newcommand{\BtwoPone}{passed}\newcommand{\BtwoPtwo}{failed}\newcommand{\BtwoPthree}{passed}\newcommand{\BtwoPfour}{passed}

% generated by b2_extra_numbers.py from evidence/B2_results.json; do not edit
\newcommand{\BtwoLeadFast}{5.0}\newcommand{\BtwoLeadSlow}{2.0}\newcommand{\BtwoTwinLead}{3.0}
\newcommand{\BtwoDiscTtzFA}{15}\newcommand{\BtwoDiscLostFA}{4}
\newcommand{\BtwoSwitchMin}{24}\newcommand{\BtwoSwitchMax}{24}\newcommand{\BtwoFinalMin}{-4}\newcommand{\BtwoFinalMax}{-3}
\newcommand{\BtwoLossCut}{0.40}\newcommand{\BtwoChangeRate}{1.8}
\newcommand{\BtwoShare}{0.66}\newcommand{\BtwoStatic}{0.60}\newcommand{\BtwoTwinsAbove}{28}\newcommand{\BtwoTwinsBelow}{2}
\newcommand{\BtwoMaxDelay}{15}\newcommand{\BtwoHours}{3.6}

\begin{document}
\pagenumbering{roman}
\begin{titlepage}
\centering
\vspace*{1.2in}
{\LARGE\bfseries From Human Markets to Machine Markets\par}
\vspace{0.4cm}
{\Large Certify, Don't Identify: Safety and Identification\\ in Markets with Adaptive Trading Agents\par}
\vspace{1.2cm}
{\large Pratik Niroula\par}
\vspace{0.3cm}
{\normalsize Department of Computer Information Science, Minnesota State University, Mankato\par}
\vspace{1.2cm}
{\normalsize Thesis working draft, revised 29 October 2026\par}
\vfill
\noindent\fcolorbox{black!60}{black!3}{\parbox{0.92\textwidth}{\small
\textbf{What this document is.} The intended final form of the thesis. Papers~1 and~2 are working drafts and appear as Chapters~\ref{ch:p1} and~\ref{ch:p2} (the second condensed); the theory of Chapter~\ref{ch:p3}, the AI-agents chapter, the registered tests, the introduction and the conclusion are new. Every result in the new chapters carries a tag: \Proved{} (with its proof), \ProvedIn{CH.~2/3} (proved in a paper chapter), \Target{} (the statement the thesis will prove, with a proof plan), \Conj{}, or \Pending{} (an empirical test specified but not yet run). No number in the new chapters is invented; Appendix~\ref{app:status} lists the status of every main result.}}
\end{titlepage}

\chapter*{Abstract}
\addcontentsline{toc}{chapter}{Abstract}
Electronic markets are feedback systems: adaptive trading agents, now increasingly AI agents, respond to prices that their own trades move. Supervisors respond by asking to see inside the agents. This thesis asks whether that knowledge is needed for safety and answers that, in the closed-loop markets studied here, it is not: certify, don't identify. Prices reveal the agents only through a merged measure of their memories and weights. Separating two agents whose memory rates differ by $\delta$ costs trading capital of order $\delta^{-4}$, separating $K$ agents costs experiments of order $(\Delta\eta)^{-(4K-2)}$, and with general memory accuracy improves only as $1/\log N$. Under the positive-cohort fixed-inventory bounded-shock contract, worst-case loss depends on the summed price response, which is certified with $\sigma^2\epsilon^{-2}\log(1/\epsilon)$ experiments whatever their number and spacing, from passive data at a cost of order one over the stability margin, and not at all without a bound on memory; a pre-trade guard built on the certificate is the largest safe admission rule. A general theory explains the pattern: every question is a functional of the merged measure, and its exact cost is the value of a polynomial approximation problem; the safety question sits on the boundary of the parametric regime and pays one logarithm. Applied to AI agents, the model shows that adoption lowers the stability margin when the common strategy has stronger positive feedback than the strategies it replaces, and it gives an instrument that needs from deployers only a memory bound and a sensitivity bound: certify the margin daily from public data and trigger on its trend. In a registered simulation the trend warned about two weeks before instability in 96\% of markets with no false alarm, while a threshold computed from a model of the agents erred by nearly a factor of two; with language-model desks the trend again warned before every crisis with no false alarm, but only three to four days ahead. In data, prices leave the agents' memories unrestricted while disclosed trader-class positions restrict them; with enough days public data decide the safety question in new markets whose agents they do not identify; and in thin books the safety statistic forecasts the size of adverse excursions beyond a full volatility baseline. The evidence is narrower than the theory: decided margins rest mainly on price reversion in markets far from the boundary, the margin does not forecast adverse moves, and a guard estimated from public data does no better than a volatility limit. Of the 23 primary hypotheses of the thesis's registered tests, 11 pass; all failures are reported.

\chapter*{Statement of co-authorship}
\addcontentsline{toc}{chapter}{Statement of co-authorship}
Chapter~\ref{ch:p1} reproduces the paper \emph{Auditable Memory and Financial Identification Limits in Electronic Limit Order Books}, joint with Prashant Niroula. Chapter~\ref{ch:p2} condenses the paper \emph{From Human Markets to Machine Markets: Transient Amplification, an Exact Depth Condition, and a Maximal Pre-Trade Guard}, joint with Suboh Alkhushayni, Yuhao Chen, Mansi Bhavsar and Prashant Niroula. \textit{[To be completed with the co-authors: the candidate's contribution to each chapter, and permission to include the papers.]} Chapters~\ref{ch:p3}, \ref{ch:ai} and~\ref{ch:tests}, the introduction and the conclusion are the candidate's own work.

\chapter*{Use of AI tools}
\addcontentsline{toc}{chapter}{Use of AI tools}
The candidate used Claude (Anthropic), an AI assistant, to write and debug analysis and simulation code, to check derivations, numerical results and Lean formalizations, and to draft and edit text. Every number in the thesis is produced by code that was executed and archived, and an automated check compares the printed numbers with the files that produced them. The candidate is responsible for reviewing the content. The 29 October development revision additionally used OpenAI Codex to propose source corrections, analysis code and an experiment notebook; these changes require author and independent expert review.

\tableofcontents
\clearpage
\pagenumbering{arabic}

\chapter{Introduction}\label{ch:intro}

\section{The question}

Many of the systems on which modern economies depend are feedback systems populated by adaptive agents. In an electronic market, trading algorithms watch prices, trade on what they see, and move the prices they watch. In a power grid, inverters and flexible loads respond to frequency and in doing so change it. In a stream of traffic, each adaptive cruise controller responds to the car in front and becomes the input to the car behind. In each case the operator who is responsible for safety faces agents whose rules are private, heterogeneous and changing.

A natural response is to ask to see inside. Securities regulators require firms to flag algorithmic orders and identify the algorithm behind each decision; the US Consolidated Audit Trail records who stands behind every order; grid operators ask aggregators for device-level models. Behind these requirements sits an implicit claim: that safety cannot be assured without knowing who the agents are.

This thesis tests that claim. Its central question is:
\begin{quote}
\emph{When can a feedback system with unknown adaptive agents be certified safe without identifying the agents, and what does each kind of knowledge cost?}
\end{quote}

\section{The answer in one paragraph}

For the closed-loop systems studied here, the information that certifies safety is generically far smaller and far cheaper than the information that identifies the agents. Prices reveal the agents only through a merged measure of their memory rates and weights (Chapter~\ref{ch:p1}). Recovering the agents from that measure is a moment problem whose cost explodes as agents become similar: the trading capital needed to separate two memory rates $\delta$ apart grows as $\delta^{-4}$, the number of experiments needed to separate $K$ agents as $(\Delta\eta)^{-(4K-2)}$, and with general memory no method does better than accuracy $1/\log N$. Safety needs much less. The worst-case transient loss of a position depends on the agents only through one number, the summed price response $S$, and $S$ is certified at the parametric rate whatever the number and spacing of the agents, provided their memory is bounded (Chapter~\ref{ch:p2}). The framework of Chapter~\ref{ch:framework} explains why: every question an operator can ask is a functional of the merged measure, and its cost is decided by how that functional behaves as memory becomes long. Smooth questions such as ``is the system safe?'' are cheap; questions that need the fine structure of the measure, such as ``who is trading?'', are expensive; and questions that are singular at infinite memory cannot be answered from finite data at all.

\section{Contributions}

\paragraph{Identification limits (Chapter~\ref{ch:p1}).} In a linear cohort model whose reduced form arises from Kalman filtering and regularized online learning, prices identify only the merged rate--weight measure. Separating a known pair costs executed quantity of order $\delta^{-2}$ and capital of order $\delta^{-4}$, with explicit constants and a finite-horizon certificate; $K$ unknown cohorts cost $(\Delta\eta)^{-(4K-2)}$ experiments, matched by Prony's method; general completely monotone memory gives the logarithmic rate; passive data cost the same polynomial order. Class-tagged disclosure removes the barrier and is worth a factor $(\Delta\eta)^{-(4K-4)}$ in samples. Trader-class positions from Binance and the CFTC pin down memories that prices leave unrestricted.

\paragraph{Safety without identification (Chapter~\ref{ch:p2}).} In a limit order book with adaptive cohorts, cumulative transient risk has the closed form $S=\tilde\kappa/(m-q\tilde\kappa)$, a book overshoots if and only if depletion outweighs impact-weighted reversion, and a pre-trade guard on the summed impulse response is the largest safe admission rule. $S$ is certified with $O(\sigma^2\epsilon^{-2}\log(1/\epsilon))$ experiments; a finite window certifies it only under a memory bound; the certified guard is minimax-optimal in lost capacity up to a logarithmic factor; disclosure has a computable value for safe capacity; and reversion is under-supplied in equilibrium. In a pre-registered empirical program the statistic forecasts adverse excursions in thin books, while mechanism tests fail.

\paragraph{A framework that contains both (Chapter~\ref{ch:framework}).} Questions are functionals of the merged measure and observation contracts are noise models on its moments. For linear questions the thesis proves an upper bound on the cost of certification in terms of two quantities of the question, its tail at long memory and its coefficient energy (Theorem~\ref{thm:dich}(i), \Proved{}), and states the matching lower bound as its main open technical target. The rates of Chapters~\ref{ch:p1} and~\ref{ch:p2} appear as cases of one picture.

\paragraph{Beyond linear and stationary agents (Chapter~\ref{ch:beyondlinear}).} A comparison theorem shows that the summed-gain guard stays safe for agents with sector-bounded nonlinear rules (Theorem~\ref{thm:sector}, \Proved{}); a second shows that when agents keep learning, a certificate has a lifetime equal to its safety margin divided by the speed of learning (Theorem~\ref{thm:lifetime}, \Proved{}). Whether the guard survives agents who adapt to it is stated as a conjecture.

\paragraph{Beyond markets (Chapter~\ref{ch:beyondmarkets}).} The framework maps to string stability of vehicle platoons with unknown commercial cruise controllers and to frequency safety in grids with unknown device populations (\Target{}, \Pending{}).

\paragraph{Evidence program (Chapter~\ref{ch:evidence}).} Two tests are specified that would test the thesis directly: a ground-truth test that the certificate computed from prices agrees with the one computed from disclosed classes while the classes themselves are not recoverable from prices, and a prospective test on liquidation data recorded after registration (\Pending{}).

\section{Status of the results}

\begin{table}[h]\centering\small
\begin{threeparttable}
\caption{Main results of the thesis and their status}\label{tab:intro_status}
\begin{tabular}{p{0.52\textwidth}p{0.2\textwidth}p{0.18\textwidth}}
\toprule
Result & Where & Status\\
\midrule
Prices identify only the merged measure & Thm~\ref{p1:thm:merged} & proved\\
Capital of order $\delta^{-4}$; $(\Delta\eta)^{-(4K-2)}$ experiments; $1/\log N$ & Ch.~\ref{ch:p1}, \S\ref{p1:sec:resources}--\ref{p1:sec:mixtures} & proved\\
Safety depends only on $S$; guard is maximal & Lemma~\ref{p2:thm:guard} & proved\\
$S$ certified at the parametric rate; impossible without a memory bound & Thms~\ref{p2:thm:window}, \ref{p2:thm:certcost} & proved\\
Upper bound on certification cost for any linear question & Thm~\ref{thm:dich}(i) & proved here\\
Matching lower bound (the dichotomy) & Thm~\ref{thm:dich}(ii) & target\\
Guard safe for sector-bounded nonlinear agents & Thm~\ref{thm:sector} & proved here\\
Certificate lifetime under learning & Thm~\ref{thm:lifetime} & proved here\\
Guard under strategic adaptation & Conj.~\ref{conj:strategic} & conjecture\\
Platoons and grids & Ch.~\ref{ch:beyondmarkets} & target, pending data\\
Ground-truth and prospective tests & Ch.~\ref{ch:evidence} & pending data\\
\bottomrule
\end{tabular}
\end{threeparttable}
\end{table}

\section{Relation to existing work}

The thesis connects four literatures that rarely meet. \emph{System identification} \citep{t:ljung1999} asks how well a system's parameters can be recovered from input--output data; Chapter~\ref{ch:p1} shows that in closed-loop markets the relevant parameters are zeros of the transfer function and are recoverable only at polynomial or logarithmic rates. \emph{Mixture and moment problems} in statistics give the closest mathematical relatives: the optimal rate for estimating a $k$-atom mixing distribution is $n^{-1/(4k-2)}$ \citep{t:heinrich2018,t:wu2020}, super-resolution has the same Vandermonde geometry \citep{t:candes2014,t:moitra2015}, and the estimation of smooth \emph{functionals} of a distribution can be far easier than estimating the distribution itself \citep{t:lepski1999,t:cai2011}. The dichotomy of Chapter~\ref{ch:framework} is the counterpart of that last observation for safety questions in feedback systems. \emph{Robust and positive control} \citep{t:megretski1997,t:farina2000,t:rantzer2015} supplies the comparison arguments of Chapter~\ref{ch:beyondlinear} and the $\ell_1$ gain that bounds worst-case loss. \emph{Partial identification} in econometrics \citep{t:manski2003} provides the stance: when the structure is not identified, ask which decisions are nonetheless determined. In finance, the closest work is on transient impact and resilience, liquidity spirals and automated-market crashes, reviewed in Chapters~\ref{ch:p1} and~\ref{ch:p2} \citep{t:brunnermeier2009,t:kirilenko2017}.

\section{Reading guide}

Chapter~\ref{ch:framework} can be read first or last: it states the common framework, and each of its results points to the chapter where the special case is proved. Chapters~\ref{ch:p1} and~\ref{ch:p2} are self-contained papers and keep their own notation, appendices and references; Appendix~\ref{app:status} maps the notation across chapters. Chapters~\ref{ch:beyondlinear}--\ref{ch:evidence} build on Chapter~\ref{ch:p2}.

\chapter{Auditable memory and identification limits}\label{ch:p1}

This chapter, the first paper of the thesis, reproduces the paper \emph{Auditable Memory and Financial Identification Limits in Electronic Limit Order Books}, joint with Prashant Niroula. It establishes the identification half of the thesis: what prices reveal about the agents of a closed-loop market, and what it costs to learn more. Its notation is its own; Appendix~\ref{app:status} maps it to the notation of Chapter~\ref{ch:p3}.

\begin{quote}\small\textbf{Summary.} Markets are increasingly driven by algorithms that learn from recent prices and, by trading, move the prices they learn from. Supervisors, exchanges and trading desks often want to know which group of traders is behind a price move and how fast each group reacts. We ask how much of this can be learned from public prices and order flow, and what it costs to learn more. In a model where several groups of traders follow prices with exponentially fading memory, prices reveal only one combined summary of the groups: groups with the same memory look like a single group, however much data is collected. Groups with different memories can be told apart, but the cost rises steeply as their memories get closer. For $K$ unknown groups whose memories are $\Delta\eta$ apart, the number of experiments needed grows as $(\Delta\eta)^{-(4K-2)}$, with matching lower and upper bounds on suitably rich classes: for two groups, halving the spacing multiplies the data needed by 64. Watching without trading needs at least as many observations, and if memory is not exponential, accuracy improves only with the logarithm of the data. To tell apart a known pair of groups, an observer who trades against them needs capital that rises sixteen-fold each time their difference is halved. Records that label positions or trades by trader class remove this barrier: each group becomes an ordinary estimation problem whose difficulty does not depend on how similar the groups are. In a simple decision model, such records are worth buying for groups that are neither too similar nor too different, and only when their fixed cost is shared across enough markets. In a hypothetical benchmark, re-learning two unknown groups each week by test trades would cost an outside participant at least \$15 million to \$356 million a year per market; this is a shadow price of the information, not a social cost. In tests specified before the data were opened, CFTC trader-class reports pin down the persistence of the faster class's positions in all eight futures markets studied, where weekly prices do not, and Binance's class data do so in 29 of 29 new markets. We do not show that these persistences are learning memories; the Binance persistence is close to a random walk once a within-day bias is removed, and class data did not improve return forecasts.\end{quote}

\paragraph{Reproducibility.}Every number in the paper is produced by a script or notebook cell that writes a result file from public data or from a closed form; Appendix~\ref{p1:app:repro} maps each table, figure and number to its source. The specifications of the October 2026 follow-up tests (Section~\ref{p1:sec:learners} and Section~\ref{p1:sec:disc_data}) were committed to the repository before their test data were downloaded; because each specification reached the public repository in the same push as its results, that order is documented by the commit history rather than by an independent timestamp. ``Registered'' in this paper means specified and committed to the public repository before the test data were opened; it is not registration with an independent registry. The proofs have been checked by the authors and against numerical computations; they have not yet been independently reviewed or formally verified, and the numerical checks reproduce results with the same code unless stated otherwise.

\section{Introduction}

Markets are increasingly run by algorithms. Trend-following funds, market-making programs, execution algorithms and, more recently, AI trading agents all read recent prices, decide what to trade, and by trading move the prices they read. When something goes wrong, a sharp sell-off, a sudden loss of liquidity, a crowded trade unwinding, the first questions supervisors and risk managers ask are who was trading and how. Which group of traders drove the move? How quickly do they react, and how long do they remember what prices did? Could the same group do it again?

This chapter asks how much of that can be learned from what everyone can see, prices and anonymous order flow, and, when public data are not enough, what it costs to find out in other ways.

The question matters now. Trading is passing to learning algorithms and AI agents \citep{p1:colliard2022,p1:dou2025}, and regulators have begun to require firms to say which algorithm sent each order. Under MiFID II, EU investment firms must flag algorithmic orders and identify the algorithm behind each investment and execution decision \citep{p1:esma2016}. The US Consolidated Audit Trail (CAT) records who stands behind every order in US stocks and listed options, at an operating budget of about \$147 million for 2026 \citep{p1:sec2012,p1:catbudget2026}. Such systems serve many purposes, among them enforcement, rebuilding market events after the fact, and surveillance. One question they raise, and the one we study, is how much of the information about who drives price moves could instead be learned from public data, and at what cost.

\paragraph{The short answer.} Prices tell an observer less than one might hope, and sharpening what they do tell is expensive.
\begin{itemize}[itemsep=1pt]
\item Prices reveal one combined summary of all the trader groups, not the groups themselves. Two groups that react to prices with the same memory look exactly like one larger group, however much price data is collected (Theorem~\ref{p1:thm:merged}).
\item Groups with different memories can be told apart in principle, but the data needed grow very fast as their memories get closer. For two unknown groups, halving the difference between their memories multiplies the data needed by 64; for three groups, by 1{,}024 (Theorems~\ref{p1:thm:kcohort} and~\ref{p1:prop:esprit}).
\item Records that label each position or trade by trader class remove this problem. With labels, each group becomes an ordinary estimation problem whose difficulty does not depend on how similar the groups are (Proposition~\ref{p1:prop:tagged}).
\end{itemize}
We keep three statements apart throughout: what cannot be learned at all, what can be learned but only at high cost, and where estimates become unstable.

\paragraph{Why prices hide the groups.} The reason is the feedback loop. Algorithms trade on recent prices, their orders move prices, and those prices feed back into the algorithms. We model each group as tracking the price with a fading memory: recent prices count most, older prices count less, and one number, the group's memory rate, sets how fast the past fades. A desk behaves exactly this way if it uses a Kalman filter or a regularized online learning rule (Lemma~\ref{p1:lem:micro}). Because every group reaches the price only through the sum of all their trades, the price records the total and loses the parts. Think of several people pushing the same door: from the door's movement one can learn the total push and how it builds and fades, but two people who push in the same rhythm look like one stronger person.

\paragraph{Three ways to learn more, and what each costs.} An observer that wants to know more, whether a supervisor, an exchange's surveillance team or a large trading desk, has three options (Section~\ref{p1:sec:observer}).
\begin{enumerate}[label=(\alph*),itemsep=1pt]
\item \emph{Wait and watch prices.} This costs time. A passive observer needs at least as many observations, in order of magnitude, as an active one needs experiments (Theorem~\ref{p1:prop:passive}).
\item \emph{Trade against the groups} and watch how the price responds, much as one taps a wall to find a hidden beam. Every test trade moves the price against the trader, so learning this way costs money. To tell one group from two groups whose memory rates differ by $2\delta$, the quantity traded grows like $1/\delta^{2}$ and the capital at risk like $1/\delta^{4}$: halve the difference and the capital needed rises sixteen-fold (Theorems~\ref{p1:thm:quantity} and~\ref{p1:thm:cap}). If memory does not fade exponentially, accuracy improves only with the logarithm of the data, so halving the error squares the data needed (Theorem~\ref{p1:thm:general}).
\item \emph{Acquire labelled records}, such as audit-trail data or class-level position reports. This costs whatever the reporting system costs, but it removes the explosion in data: labels save a factor of order $(\Delta\eta)^{-(4K-2)}$ at fixed accuracy, where $K$ is the number of groups and $\Delta\eta$ the spacing of their memory rates (Theorem~\ref{p1:thm:discvalue}).
\end{enumerate}
Proposition~\ref{p1:prop:choice} says when each option is the cheapest.

\paragraph{How large are the costs?} To give a sense of scale we price the trading option in a hypothetical benchmark. The model parameters are illustrative; only the size of the measurement noise is estimated, from Bitcoin data. An outside participant that re-learns two unknown algorithms each week with test trades would lose at least \$15 million to \$356 million a year per market; with three unknown algorithms it would lose more than \$40 million a year per market even when their memories are far apart (Section~\ref{p1:sec:disc_dollars}). These are lower bounds for one stated way of learning, not estimates of what any regulator spends: regulators do not trade to learn. They are best read as the price of the information that labelled records deliver without any trading.

\paragraph{Why this matters for finance.} The results point in a constructive direction for several groups of readers.
\begin{itemize}[itemsep=1pt]
\item \emph{Supervisors and policymakers} get a way to decide what data to require. If the goal is only to watch overall price risk, public prices are enough, because the combined price response is learned at the ordinary rate (Remark~\ref{p1:rem:riskneutral}). Labelled records earn their cost when supervision must be targeted, for example when one algorithm must be paused while the others keep supplying liquidity, and when the reporting system covers enough markets to share its fixed cost (Proposition~\ref{p1:prop:choice}). This is an economies-of-scale argument for market-wide regimes such as CAT, and a reason to keep a stable account identifier even when personal data are removed (Section~\ref{p1:sec:disclosure}).
\item \emph{Exchanges and trading venues}, which already see account-level flow, learn which questions about their participants public prices cannot answer, and why publishing class-level positioning, as some crypto venues already do, adds information that prices do not contain.
\item \emph{Risk managers and trading desks} get a realistic limit on what can be inferred about other traders or about crowding from market data, and a warning that common feedback measures computed from public data mostly reflect filtering and timing rather than true feedback (Section~\ref{p1:sec:empirics}).
\item \emph{Researchers} who estimate models with several trader types get one reason why the parameters that separate the types are often imprecise, and a number that says how precise they can be before any estimation (Section~\ref{p1:sec:lit}).
\end{itemize}

\paragraph{What the data show.} Two real reporting systems let us check the main prediction, that class records pin down what prices leave open. The U.S.\ Commodity Futures Trading Commission requires large traders in futures markets to report their positions and publishes them weekly by trader class. In a test whose specification we committed before downloading any data, in eight futures markets (CME Bitcoin and Ether, the E-mini S\&P 500, 10-year Treasury notes, the euro, the yen, crude oil and natural gas), the class reports pin down how persistent the faster class's positions are to a range 9 to 71 times narrower than weekly prices allow, while prices alone leave it almost anywhere in the admissible range (Section~\ref{p1:sec:cot}). Binance's five-minute trader-class ratios show the same in all 30 perpetual futures markets we studied and, in a second test on markets never used before, in 29 of 29 (Section~\ref{p1:sec:disc_data}). We also report what we did not expect: leveraged funds sell rather than buy after price rises in three financial futures markets; the Binance persistence, once a within-day bias is removed, is close to a random walk; and class data did not improve return forecasts. The data therefore show what prices cannot pin down. They do not show that this persistence is the learning memory of the model.

\paragraph{Two further results.} Markets with the same eigenvalues, the same one-step volatility and the same long-run price variance can still have very different worst-case short-run losses, so watching eigenvalues or volatility cannot replace knowing who trades (Theorem~\ref{p1:thm:blind}). And standard feedback measures computed from public data mostly reflect filtering and clock alignment: two fee changes that moved trading volume by factors of 2.6 and 4.9 left measured feedback unchanged (Section~\ref{p1:sec:empirics}).

\paragraph{Contribution.} A large literature asks how transparency affects liquidity and trading costs \citep{p1:pagano1996,p1:boehmer2005,p1:bessembinder2006}. We ask a different question: what disclosure lets an observer learn about the traders behind price moves, and what that information would cost to obtain otherwise. We answer in data, in trading capital and in dollars, with matching upper and lower bounds where we can prove them. The mathematical tools, moment matching, two-point testing and Prony-type estimation, are known from mixture estimation and signal processing; what is new is their closed-loop, passive and financial form (Section~\ref{p1:sec:lit}).

\paragraph{How to read the paper.} Every main result is stated formally and then explained in a short paragraph headed ``In words'', which says what it means for a trader or a supervisor; proofs are in the appendix. Section~\ref{p1:sec:worked} works through two desks that use Kalman filters, in numbers, and Table~\ref{p1:tab:map} maps each question to its answer. Readers who want the finance first can read this introduction, Section~\ref{p1:sec:observer}, the worked example and Section~\ref{p1:sec:disclosure}, and take the statistical results as given. Table~\ref{p1:tab:glossary} explains the main symbols.

\begin{table}[!htbp]\centering\small
\begin{threeparttable}
\caption{The main symbols in plain words}\label{p1:tab:glossary}
\begin{tabular}{>{\raggedright\arraybackslash}p{0.17\textwidth}>{\raggedright\arraybackslash}p{0.75\textwidth}}
\toprule
Symbol & Meaning\\
\midrule
$K$ & Number of trader groups (``cohorts'')\\
$\eta_i$ & Group $i$'s learning rate: the weight it puts on the newest price. Large $\eta_i$ means a short memory and fast reactions; small $\eta_i$ means a long memory\\
$r_i=1-\eta_i$ & Group $i$'s memory rate: the share of last period's view that survives\\
$q_i$ & Group $i$'s trading intensity: how strongly it trades on its view. Positive $q_i$ means it buys when the price is above value, as trend followers do\\
$\Delta\eta$, $2\delta$ & How far apart the groups' learning rates are ($\Delta\eta$ for unknown groups, $2\delta$ for a known pair)\\
$a$, $\kappa$ & How persistent pricing errors are, and how far one unit of net order flow moves the price\\
$M$ (merged measure) & The combined summary of all groups that prices reveal: where the memories are and how much trading sits at each\\
$N$ & Number of experiments (test trades from a fresh start) or observations\\
Par units & The model's own unit of quantity and money, before conversion to USDT (Table~\ref{p1:tab:capital})\\
\bottomrule
\end{tabular}
\end{threeparttable}
\end{table}

\paragraph{What we do not claim.} The identification results hold for the stated linear model and for the stated observation contracts, that is, for what the observer is assumed to see. They are not general limits on surveillance with participant-level records, which can be far more informative (Propositions~\ref{p1:prop:disclosure} and~\ref{p1:prop:tagged}). The matching upper bound holds on a class with separated rescaled rates and weights bounded away from zero. Capital figures depend on the observation-noise calibration, which we estimate rather than assume. The empirical results are reduced-form diagnostics; they do not identify structural learning parameters. Our pre-registered real-data test of the identification boundary failed its criterion on one of its three parts: on 42 days of data the merged response was precisely estimated in only 5 of 30 markets (Section~\ref{p1:sec:disc_data}). Binance's class ratios are positions, not tagged order flow, so they are the closest public analogue to the disclosure of Proposition~\ref{p1:prop:tagged}, not an instance of it. We hold the algorithms' memories and intensities fixed across disclosure regimes; disclosure that reveals positions may lead informed desks to trade less or to split orders across accounts, which we do not model.

\subsection{Related literature}\label{p1:sec:lit}

\paragraph{Algorithms, AI and supervision.} Recent work in finance asks how learning algorithms and AI agents behave in markets: reinforcement-learning market makers and informed traders can sustain wide spreads or collusive outcomes \citep{p1:colliard2022,p1:dou2025}, prices formed against Q-learners have distinctive dynamics \citep{p1:banerjee2026}, and algorithmic markets can be manipulated \citep{p1:tremacoldi2025}. These papers model or simulate the algorithms. We take the supervisor's side and ask what can be learned about them from market data, and what disclosure is worth. A related literature measures the cost of regulation and of regulatory reporting \citep{p1:trebbi2024}; our dollar comparison with the Consolidated Audit Trail is a statement about the benefit side: what the audit trail buys that prices cannot.

\paragraph{Price formation and impact.} Our price equation is a discrete linear impact model in the tradition of \citet{p1:kyle1985}, \citet{p1:glosten1985} and \citet{p1:hasbrouck1991}, in which order flow moves prices and prices revert toward value. Transient impact and the cost of trading follow \citet{p1:almgren2001}, \citet{p1:obizhaeva2013}, \citet{p1:gatheral2010} and \citet{p1:bouchaud2004}; \citet{p1:gatheral2010} shows that absence of dynamic arbitrage requires a positive-definite decay kernel, \citet{p1:alfonsi2012} that excluding transaction-triggered manipulation requires convex, nonincreasing decay, and \citet{p1:alfonsi2016} extend positive definiteness to several assets. Positive definiteness is the property Lemma~\ref{p1:lem:toeplitz} delivers in our cohort economy. \citet{p1:bouchaud2018} document order splitting and long memory in order flow, which is the main confounder in Section~\ref{p1:sec:empirics}. Relative to this literature we treat the impact kernel not as a primitive but as the closed-loop output of adaptive traders, and ask what it reveals about them.

\paragraph{Algorithmic trading and feedback.} Algorithms that react to recent prices can amplify or dampen moves. \citet{p1:brunnermeier2009} and \citet{p1:shleifer1992} study liquidity spirals and fire sales, and \citet{p1:kirilenko2017} document the behaviour of automated intermediaries in the 2010 Flash Crash. \citet{p1:sornette2023} show that non-normal interactions can generate transient bubbles. Our contribution is on the measurement side: we show what a closed-loop observer can and cannot learn about the feedback, and that standard reduced-form indices are dominated by filtering and timing.

\paragraph{Learning algorithms and what can be inferred about them.} Empirical game-theoretic analysis treats the choice among simulated trading strategies as a game \citep{p1:wellman2006}. Empirically, algorithmic trading has been found to improve liquidity and price discovery on average \citep{p1:hendershott2011,p1:brogaard2014}. These studies simulate the algorithms or observe them through regulatory data; we ask what an outside observer can infer about them from market data alone, and answer in samples and in capital. Our stance is that of partial identification \citep{p1:manski2003}: when the structure is not identified, characterize what the data do determine, here the merged measure and, through it, the intensity-weighted distribution of the desks' information quality (Proposition~\ref{p1:prop:primitives}).
\paragraph{Control without identification.} We do not claim that acting on a system without identifying it is new. \citet{p1:moallemi2008} study strategic execution against an observer who learns from prices, which bears directly on our assumption that cohorts do not react to being observed, and \citet{p1:parkvanroy2015} study execution that learns price impact while trading, with regret guarantees; our object is different, the latent trader cohorts behind the impact rather than the aggregate impact coefficients. Active identification of linear systems has matching bounds and optimal experiment design \citep{p1:wagenmaker2020}, and task-directed exploration asks what must be learned for a decision \citep{p1:wagenmaker2021}; our information matrix differs because the unknowns are zeros of a closed-loop response and the input is paid for in trading losses. and \citet{p1:vanwaarde2020} characterize when control is possible without unique system identification. Our results concern a specific market observation contract and its costs in executed quantity, capital and money.
\paragraph{Closest work.} The closest literature is super-resolution and moment-based mixture estimation, which gives rates for recovering clustered point sources and mixing measures from noisy moments \citep{p1:moitra2015,p1:heinrich2018,p1:batenkov2021,p1:liu2023,p1:wu2020,p1:kong2017}. Our $K$-cohort rate is, in that language, the square of the amplitude condition number $(\Delta\eta)^{-(2K-1)}$ for $K$ clustered nodes \citep{p1:heinrich2018,p1:wu2020}, and we do not claim the exponent as new. What changes when the mixture sits inside a closed-loop financial market is fourfold (Table~\ref{p1:tab:novelty}): first, in mixture estimation and open-loop system identification the unknown rates are locations or poles of the observed law, whereas in a closed-loop market the cohorts feed back into the price, making the memory rates transmission \emph{zeros} of the observable transfer function $H(z)$ that must first be extracted from the loop (Theorem~\ref{p1:thm:merged}); with $a$ and $\kappa$ known this extraction is a smooth invertible map, so it decides which method is valid (pole fitting on prices is not) rather than the rate; second, passive price observations face at least the same $(\Delta\eta)^{-(4K-2)}$ sample-complexity barrier as active probing, uniformly in the spacing, because $H(z)$ is minimum phase and the Whittle likelihood inherits the moment gap (Theorems~\ref{p1:thm:minphase} and~\ref{p1:prop:passive}); third, in a market statistical sample size translates into executed quantity $\Theta(\delta^{-2})$ and round-trip trading loss or margin capital $\Theta(\delta^{-4})$ via the self-impact Toeplitz quadratic form $\tfrac12 v^\top T_h v$ (Theorems~\ref{p1:thm:quantity}--\ref{p1:thm:rh}); and fourth, class-tagged disclosure breaks the feedback loop cohort by cohort and cuts the sample requirement by $(\Delta\eta)^{-(4K-4)}$ (Proposition~\ref{p1:prop:tagged} and Theorem~\ref{p1:thm:discvalue}). Beyond the $K$-cohort case, we prove matching lower and upper bounds of exact logarithmic order $\Theta(1/\log N)$ for general completely monotone memory (Theorem~\ref{p1:thm:general}) and price the resulting disclosure savings against audit-trail budgets. Several ingredients are known in other settings, and we claim only their closed-loop, passive and financial forms: the $\delta^{-4}$ law for telling one component from two closely spaced ones is classical in resolution analysis \citep{p1:shahram2005}; logarithmic rates are known for mixtures of exponentials \citep{p1:rebafka2015} and for general mixing distributions, by the same moment-matching and polynomial-approximation technique \citep{p1:wu2020}; and labelled samples are known to be far more informative than unlabelled ones in mixture problems \citep{p1:castelli1996}.

\paragraph{Estimating heterogeneous-agent models from prices.} A large literature models markets as populations of trader types that switch between forecasting rules \citep{p1:brock1998} and estimates such models from price data. Estimates of the parameters that separate the types are often imprecise: \citet{p1:amilon2008} reports that the estimated model reproduces properties of returns but that its fit is generally poor, \citet{p1:franke2012} compare switching mechanisms by the method of simulated moments, and \citet{p1:kukacka2017} find the switching coefficient statistically insignificant by simulated maximum likelihood while the belief parameters are significant. These papers take the identification of the types from prices as given. Our results suggest one reason the type-separating parameters are hard to pin down: around a steady state, prices determine only the merged measure, and separating types whose memories differ by $\Delta\eta$ takes a number of samples of order $(\Delta\eta)^{-(4K-2)}$ for $K$ types (Theorems~\ref{p1:thm:merged} and~\ref{p1:thm:kcohort}). The models these papers estimate are nonlinear, so this is a statement about their linearization, not a claim that their estimates are wrong; it does give a quantity, the separation of the types' memories, that bounds the precision one can expect before estimating.

\paragraph{Closed-loop system identification.} Engineering has long studied identification under feedback \citep{p1:ljung1999,p1:forssell1999}. There the object is the plant, the difficulty is that feedback correlates the input with the noise, and the remedy is external excitation that makes the experiment informative. Our problem is the reverse in two ways. The object is the feedback law (the traders), not the plant (the price), and the memory rates enter as zeros of the observed transfer function. And external excitation does not remove the difficulty: with a probe the merged measure is identified exactly (Theorem~\ref{p1:thm:merged}), so the barrier is not qualitative identifiability but the conditioning of the map from memories to responses, which we measure in samples, quantity and capital. In finance, \citet{p1:leach1993} show that a market maker can gain by setting prices to learn about demand, so learning through trading has an equilibrium counterpart; our probing observer is a reduced-form version of that idea without the equilibrium. In econometric terms, as $\Delta\eta\to0$ the cohort parameters pass from strong to weak identification \citep{p1:andrews2012}, and our rates state that passage in samples.

\paragraph{Inferring trader types from market data.} Market microstructure has a tradition of inferring hidden trader populations from public data. The probability of informed trading \citep{p1:easley1996} fits a mixture model to daily counts of buys and sells to estimate the share of informed orders, and what it captures has been debated \citep{p1:duarte2009}. With regulatory account data, \citet{p1:kirilenko2017} classify traders by their holding and trading patterns. Our results speak to the case where types differ in how long they remember prices: from prices alone the separation is costly, and account or class records of the kind \citet{p1:kirilenko2017} use remove most of the cost (Proposition~\ref{p1:prop:tagged}).

\begin{table}[t]\centering\small
\begin{threeparttable}
\caption{Comparison with finite-mixture estimation and open-loop identification}\label{p1:tab:novelty}
\setlength{\tabcolsep}{4pt}
\begin{tabular}{*{4}{>{\raggedright\arraybackslash}p{0.225\textwidth}}}
\toprule
Feature & Finite-mixture estimation \citep{p1:heinrich2018,p1:wu2020} & Open-loop exponential fitting \citep{p1:roy1989,p1:batenkov2021} & Closed-loop market identification (this chapter)\\
\midrule
Where rates live & Component locations / moments of density & Poles of impulse response & Transmission zeros of closed-loop $H(z)$ (Thm~\ref{p1:thm:merged})\\
Active vs.\ passive data & i.i.d.\ draws only & Designed input vs.\ output spectrum & At least the same $(\Delta\eta)^{-(4K-2)}$ order, uniformly, via the minimum-phase Whittle gap (Thm~\ref{p1:prop:passive})\\
Economic unit of cost & Sample size $N$ & Signal-to-noise ratio & Quantity $\Theta(\delta^{-2})$, trading loss and capital $\Theta(\delta^{-4})$ via $\tfrac12v^\top T_hv$ (Thms~\ref{p1:thm:quantity}--\ref{p1:thm:rh})\\
Effect of class tagging & Observed component labels ($\Theta(1)$) & Separate sensor per mode & Breaks closed loop cohort-by-cohort; saves $\Theta((\Delta\eta)^{-(4K-2)})$ at fixed accuracy, $\Theta((\Delta\eta)^{-(4K-4)})$ for rates (Thm~\ref{p1:thm:discvalue})\\
\bottomrule
\end{tabular}
\end{threeparttable}
\end{table}

\begin{table}[!htbp]\centering\footnotesize
\begin{threeparttable}
\caption{Closest results and how this chapter differs}\label{p1:tab:closest}
\setlength{\tabcolsep}{3pt}
\begin{tabular}{*{5}{>{\raggedright\arraybackslash}p{0.172\textwidth}}}
\toprule
Existing result & Observations & Target & Conclusion & Difference from this chapter\\
\midrule
Mixture estimation \citep{p1:heinrich2018,p1:wu2020} & i.i.d.\ draws from a $K$-component mixture & Mixing measure & Rates that slow sharply as components merge; $n^{-1/(4K-2)}$ for Gaussian location mixtures & Same exponent; here the measure is observed through a feedback loop, as zeros of $H$, actively or passively, and the cost is also priced in quantity and capital\\
Super-resolution and Prony \citep{p1:moitra2015,p1:batenkov2021,p1:katz2023} & Noisy Fourier or exponential samples & Point-source locations & Stability governed by Vandermonde conditioning and node separation; Prony attains the min-max bound for clustered exponents & Real nodes extracted from a closed loop; our matching upper bound (Theorem~\ref{p1:prop:esprit}) adapts this known result\\
Closed-loop system identification \citep{p1:ljung1999,p1:forssell1999} & Input--output data under feedback & The plant & Consistency requires informative (externally excited) experiments & Target is the feedback law (the traders); excitation gives exact identification of $M$, and the barrier is conditioning, measured in samples and capital\\
Heterogeneous-agent estimation \citep{p1:brock1998,p1:franke2012,p1:kukacka2017} & Prices and returns & Type and switching parameters & Type-separating parameters often imprecise & Bounds that precision by the types' memory separation (in the linearization)\\
Probability of informed trading \citep{p1:easley1996} & Daily buy and sell counts & Share of informed orders & Mixture estimate of trader population & Types that differ in memory; price-only versus class-tagged data\\
Data informativity \citep{p1:vanwaarde2020} & Finite data, often noiseless & Whether a property is determined & Qualitative conditions & Quantitative minimax costs under noise\\
\bottomrule
\end{tabular}
\end{threeparttable}
\end{table}

\paragraph{Statistical identification.} The identification results use two-point testing arguments \citep{p1:lecam1986,p1:tsybakov2009}. The moment-matching construction of Section~\ref{p1:sec:mixtures} is the dynamical analogue of the slow rates in finite-mixture estimation \citep{p1:manole2022} and of super-resolution of near-colliding sources \citep{p1:liu2023,p1:batenkov2021}; the general-memory result combines best-approximation duality \citep{p1:bernstein1914,p1:rivlin1981} with the moment comparisons used by \citet{p1:kong2017} and \citet{p1:wu2020}, and its logarithmic rate parallels supersmooth deconvolution \citep{p1:fan1991}; recovery of exponentials connects to ESPRIT and matrix-pencil methods \citep{p1:roy1989,p1:hua1990}. Minimum-phase spectral factorization and the Whittle likelihood are classical \citep{p1:whittle1953}. What is new here is the mapping from these statistical rates into executed quantity and money through a financially specified account, the passive equivalence, and the value of class records; that the relevant rates are zeros, not poles, of the observable transfer function matters for which estimator is valid, not for the rate.

\paragraph{Non-normal dynamics and supervision.} Transient growth of stable non-normal systems is classical \citep{p1:trefethen2005,p1:kreiss1962}. We use it to show that supervision through eigenvalues and second moments of prices is blind to transient risk (Theorem~\ref{p1:thm:blind}), which motivates the identification-free control of Chapter~\ref{ch:p2}. On the regulatory side, pre-trade risk controls for broker-dealers with market access are governed by SEC Rule 15c3-5 \citep{p1:sec2010}, and participant-level audit trails are the natural disclosure regime for Proposition~\ref{p1:prop:tagged}. A large literature studies how pre- and post-trade transparency affects liquidity and trading costs \citep{p1:pagano1996,p1:boehmer2005,p1:bessembinder2006}; we ask a different question, what disclosure lets a supervisor learn about the algorithms that trade, and give the answer in samples. Crypto venues, whose market structure is studied by \citet{p1:makarov2020}, are a natural laboratory for this question because they publish class-level positioning every five minutes, while the CFTC publishes it weekly for regulated futures.

Partial fractions, Gaussian testing, convolution inequalities and positive-system gains are established tools. Our contribution is to connect them into a quantitative chain from memory resolution, through probing resources, to money, with matching orders, explicit constants and a supervisory interpretation, and then to measure the relevant quantities in data.

{\footnotesize\begin{longtable}{p{0.34\textwidth}p{0.36\textwidth}p{0.18\textwidth}}
\caption{Map of results}\label{p1:tab:map}\\
\toprule
Question & Answer & Where\\
\midrule
What do prices identify? & The merged rate--weight measure only & Thm~\ref{p1:thm:merged}\\
Quantity to separate a known pair & $\Theta(\delta^{-2})$, exact error formula & Thm~\ref{p1:thm:quantity}\\
Capital to separate a known pair & $\Theta(\delta^{-4})$, counterparty or own funds & Thm~\ref{p1:thm:cap}, Lem.~\ref{p1:thm:capital}\\
Constant at the calibration & $[61.41,64.48]$ par units; $\approx10^5$ USDT & Thms~\ref{p1:thm:spectral}--\ref{p1:thm:rh}, Table~\ref{p1:tab:capital}\\
Does adaptivity help? & At most by a factor $\approx0.49$ in expected loss & Thm~\ref{p1:thm:adaptive}\\
Unknown mixtures & $(\Delta\eta)^{-(4K-2)}$, tight for every $K$ & Thms~\ref{p1:thm:kcohort}, \ref{p1:prop:esprit}\\
General memory & Accuracy $\asymp1/\log N$, lower and upper bounds & Lem.~\ref{p1:lem:altpair}, Thm~\ref{p1:thm:general}\\
Passive data & Identifies $M$; at least the same polynomial cost for every $K$ & Thms~\ref{p1:thm:minphase}, \ref{p1:prop:passive}, Prop.~\ref{p1:prop:passive0}\\
Nonlinear and adaptive rules? & Settled online learners obey $\delta^{-4}$; fast adaptation or thresholds break it & Table~\ref{p1:tab:learners}\\
What removes the barrier? & Tagged flow, audited equity & Props~\ref{p1:prop:disclosure}, \ref{p1:prop:tagged}\\
What is tagging worth? & Factor $(\Delta\eta)^{-(4K-2)}$ (fixed accuracy) or $(\Delta\eta)^{-(4K-4)}$ (rates); $\approx333\times$ for BTCUSDT (illustrative) & Thm~\ref{p1:thm:discvalue}\\
Disclosure in dollars? & One market's annual probing loss exceeds the 2026 CAT budget once $K\ge3$ at $\Delta\eta\le0.10$, $K=4$ at $\Delta\eta=0.20$, or for a known pair at $\Delta\eta=0.04$ (both noise calibrations) & Table~\ref{p1:tab:cat_dollars}\\
Do class data restrict class persistence? & Yes: 30 of 30, and 29 of 29 new markets (registered); whether it is memory is open & Table~\ref{p1:tab:disc}\\
Can eigenvalues certify safety? & No & Thm~\ref{p1:thm:blind}\\
What do feedback indices measure? & Mostly filtering and timing; small real feedback & Section~\ref{p1:sec:empirics}\\
Does feedback track activity? & No: fees move volume 2.6--4.9$\times$, not $\hat Q$ & Section~\ref{p1:sec:empirics}\\
\bottomrule
\end{longtable}}


\subsection{Outline}
Section~\ref{p1:sec:model} sets out the model. Section~\ref{p1:sec:ident} gives the identification boundary. Section~\ref{p1:sec:resources} links memory resolution to quantity and money for a known pair. Section~\ref{p1:sec:account} computes account-level capital. Section~\ref{p1:sec:mixtures} treats unknown mixtures and passive data. Section~\ref{p1:sec:learners} checks the results against stylized nonlinear and adaptive trading rules. Section~\ref{p1:sec:disclosure} studies disclosure and compares its cost in dollars with the U.S. Consolidated Audit Trail. Section~\ref{p1:sec:blind} gives the spectral blind spot. Section~\ref{p1:sec:empirics} reports the empirical evidence. Section~\ref{p1:sec:conclusion} concludes. Proofs are in Appendix~\ref{p1:app:proofs}; Appendix~\ref{p1:app:ia} collects additional results, and Appendix~\ref{p1:app:repro} maps every number to its source.

\section{Model and microfoundations}\label{p1:sec:model}

\paragraph{The model in words.} The market has one price and $K$ groups of traders. Each period three things happen. First, the price moves: it drifts back toward its fundamental value, and it is pushed by the groups' trades, by any test trade and by news. Second, each group updates its view by mixing the newest price into its running average, with weight $\eta_i$ on the newest price and $1-\eta_i$ on its old view. Third, each group trades in proportion to its view, with intensity $q_i$. A group with a large $\eta_i$ has a short memory and reacts quickly, like a short-horizon trend follower; a group with a small $\eta_i$ has a long memory and reacts slowly. The equations below write this down.

Time is discrete. Let $p_t$ be the price, $p^*$ the fundamental value and $x_t=p_t-p^*$ the pricing error. Competitive dealers absorb net order flow with price impact $\kappa>0$, and arbitrage pulls the error toward zero with persistence $a\in[0,1)$:
\begin{equation}\label{p1:eq:price}
x_t = a x_{t-1} + \kappa \sum_{i=1}^K q_i y_{i,t-1} + \kappa v_t + w_t,
\end{equation}
where $v_t$ is an executed probe quantity and $w_t$ a price innovation. Cohort $i$ holds a state $y_{i,t}$ with exponential memory and trades $q_i y_{i,t}$:
\begin{equation}\label{p1:eq:cohort}
y_{i,t} = (1-\eta_i)\,y_{i,t-1} + \eta_i x_t,\qquad \eta_i\in(0,1).
\end{equation}
The ordering is \emph{price first}: the price responds to the probe, then cohorts update. Write $r_i=1-\eta_i$ for the memory rate and $w_i=q_i\eta_i$ for the merged weight. The feedback kernel and its transform are
\begin{equation}\label{p1:eq:kernel}
k_m=\sum_i w_i r_i^{\,m},\qquad G(z)=\sum_i \frac{w_i z}{z-r_i},
\end{equation}
and the transfer function from probes to prices, $H(z)=\sum_{j\ge0}h_jz^{-j}$, is
\begin{equation}\label{p1:eq:H}
H(z)=\frac{\kappa z}{z-a-\kappa G(z)}=\frac{\kappa zD(z)}{(z-a)D(z)-\kappa N(z)},\qquad G=\frac ND,\quad D(z)=\prod_{i}(z-r_i),
\end{equation}
with impulse response $h_0=H(\infty)=\kappa,h_1,h_2,\ldots$; in particular $H(1)=\kappa/(1-S)=\sum_jh_j$. Stacking $z_t=(x_t,y_{1,t},\ldots,y_{K,t})^\top$ gives $z_t=Jz_{t-1}+bv_t$, where, because cohorts update on the same period's price, the probe enters every state on impact through
\begin{equation}\label{p1:eq:b}
b=\kappa\,(1,\eta_1,\ldots,\eta_K)^\top ,
\end{equation}
and
\begin{equation}\label{p1:eq:J}
J=\begin{pmatrix} a & \kappa q_1 & \cdots & \kappa q_K\\ \eta_1 a & 1-\eta_1+\eta_1\kappa q_1 & \cdots & \eta_1\kappa q_K\\ \vdots & & \ddots & \vdots\\ \eta_K a & \eta_K\kappa q_1 & \cdots & 1-\eta_K+\eta_K\kappa q_K\end{pmatrix},\qquad S\equiv a+\kappa\sum_i q_i .
\end{equation}
In the \emph{positive economy} ($q_i\ge0$), $J$ is entrywise nonnegative, its price row sums to $S$, and the economy is Schur stable iff $S<1$.

\begin{lemma}[Filtering and learning microfoundations]\label{p1:lem:micro}
\begin{enumerate}[label=(\roman*)]
\item \emph{Kalman memory.} A desk that tracks a latent random-walk trend $v_{t+1}=v_t+\zeta_{t+1}$, $\zeta\sim N(0,\sigma_\zeta^2)$, from signals $r_t=v_t+\nu_t$, $\nu\sim N(0,\sigma_\nu^2)$, using the steady-state Kalman filter, updates as $\hat v_{t+1}=(1-\eta)\hat v_t+\eta r_t$ with
\begin{equation}\label{p1:eq:kgain}
\eta=\frac{-\sigma_\zeta^2+\sqrt{\sigma_\zeta^4+4\sigma_\zeta^2\sigma_\nu^2}}{2\sigma_\nu^2}\in(0,1).
\end{equation}
\item \emph{Gradient memory.} Online gradient descent with $L_2$ regularization $\gamma\in(0,1)$, step size $\alpha>0$ and loss $\tfrac12H(x_t-\theta_t)^2$ gives $\theta_{t+1}=(1-\eta)\theta_t+\eta q x_t$ with $\eta=\gamma+\alpha H\in(0,1)$ and $q=\alpha H/\eta$. A tabular value update with learning rate $\eta$ at a fixed state has the same form; general reinforcement-learning policies need not.
\end{enumerate}
\end{lemma}

\begin{remark}[Why exponential memory]
Lemma~\ref{p1:lem:micro} shows that two different design philosophies, Bayesian filtering and regularized online learning, produce the same observable reduced-form belief state $y_{i,t}$. When desk $i$ chooses order flow $u_{i,t}$ subject to $I_{i,t}=(1-\delta_{\rm inv})I_{i,t-1}+u_{i,t}$ and an inventory cap $|I_{i,t}|\le\bar I_i$ to maximize expected continuation return $\beta_i y_{i,t}u_{i,t}$ net of quadratic execution cost $\tfrac12\lambda_{{\rm tc},i}u_{i,t}^2$ and quadratic inventory holding cost $\tfrac12\gamma_{{\rm inv},i}I_{i,t}^2$, its interior first-order condition gives $u_{i,t}^*=q_iy_{i,t}-\psi_{{\rm inv},i}I_{i,t-1}$ with market-response weight $q_i=\beta_i/(\lambda_{{\rm tc},i}+\gamma_{{\rm inv},i})$ and inventory mean-reversion $\psi_{{\rm inv},i}=\gamma_{{\rm inv},i}(1-\delta_{\rm inv})/(\lambda_{{\rm tc},i}+\gamma_{{\rm inv},i})$, which shrinks toward zero as trading costs or inventory penalties tighten and destabilizes the market when the closed-loop spectral radius reaches one (\texttt{microfoundation\_and\_baselines.py}). This filtering-and-control recursion specifies an individual agent's decision rule, not a complete economic equilibrium (dealer--trader equilibrium is analysed in Chapter~\ref{ch:p2}). Two further simplifications matter for reading the results: the observer is taken to know the fundamental value $p^*$, so the pricing error $x_t$ is observed, and the dealer sector is a fixed linear impact $\kappa$ rather than an equilibrium quote. The paragraph before Section~\ref{p1:sec:cot} states what the data tests lose when neither holds. The memory rate encodes the desk's signal-to-noise ratio or regularization, which is the private information a supervisor would like to recover; Section~\ref{p1:sec:ident} shows why prices reveal it only in merged form.
\end{remark}

\begin{lemma}[Positive price-response geometry]\label{p1:lem:cm}
In the positive, stable economy the impulse response is a nonnegative mixture of decaying exponentials,
\[
h_j=\sum_\ell \omega_\ell\lambda_\ell^j,\qquad \omega_\ell\ge0,\quad \sum_\ell\omega_\ell=\kappa,\quad 0\le\lambda_\ell<1 .
\]
It is completely monotone and nonincreasing; its peak is the direct impact $\kappa$, even though $J$ is non-normal. In price/memory coordinates the price-only gain is $\sup_t\|e_x^\top J^t\|_2=1$.
\end{lemma}

Numerically, $(-1)^n\Delta^nh_j\ge0$ for $n\le4$ held in all 892 stable random positive economies tested. Lemma~\ref{p1:lem:cm} matters twice: it makes probing costs positive (Section~\ref{p1:sec:resources}), and it shows that in this economy transient amplification of the price response cannot arise from cohort feedback alone. Chapter~\ref{ch:p2} locates amplification in order-book depth instead.

\subsection{Discussion of the assumptions}\label{p1:sec:assumptions}

\paragraph{Linearity.} The model is linear in prices, flows and memory states. Linearity is what makes the merged-measure equivalence exact and the testing problems Gaussian. Nonlinear strategies (for example threshold rules or position limits) can break the equivalence class, but they do so in regions of the state space that a small probe does not visit; for small probes the linearization around the steady state is the relevant object.

\paragraph{Price-first timing.} We order the price update before the cohort update. Under the alternative lagged timing the transition matrix becomes $\begin{psmallmatrix}a&1\\ \eta Q&1-\eta\end{psmallmatrix}$ and its eigenvalues differ, so stability statements must be computed under the timing actually used. The identification results depend on the transfer function and survive either convention.

\paragraph{Exponential memory.} Lemma~\ref{p1:lem:micro} justifies exponential memory for filtering and regularized learning. Neural and reinforcement-learning policies are not exponential filters. Rules whose orders are smooth functions of an exponential-memory state are exponential filters to first order, so the results hold for small probes; Section~\ref{p1:sec:learners} checks this in stylized simulations and finds where it fails, and how closely production algorithms follow such rules remains an empirical question we do not settle. Mixtures of exponentials can approximate any completely monotone kernel, so a power-law memory corresponds to a continuum of cohorts, for which Theorem~\ref{p1:thm:general} shows that accuracy improves only logarithmically in the data.

\paragraph{Algorithms that react to being probed.} The probing results assume that the cohorts keep their memories while they are probed. An algorithm that detects a probe could pause, shorten its memory or trade against the probe. Four observations limit what this does to our conclusions. First, a reaction that all cohorts share does not separate them, and an algorithm that hides by mimicking another rate makes the hypotheses closer, so neither lowers the costs we derive; a reaction that differs across cohorts is itself a signature, but one the observer must model, which is a new identification problem rather than a free one. Second, the passive route needs no probe: Theorem~\ref{p1:prop:passive} shows that natural price variation carries information about the same separation at a rate of the same polynomial order, so an observer who fears detection pays in time, not in rate. Third, detection is itself a testing problem against the same noise: the generalized-eigenvector probe of Section~\ref{p1:sec:account} trades at most 1.25 par units in any period, and a probe sized within normal flow variation is hard to tell from noise, so the threat bites mainly for large single pulses, which our ledger shows are not optimal anyway. Fourth, what we cannot rule out is an algorithm designed to defeat surveillance, which turns identification into a game between observer and observed; our results are the benchmark against which such a game would be measured, and we leave it to future work.

\paragraph{Known primitives.} The resource theorems take $\kappa$, $a$ and the noise law as known. Unknown primitives can only make separation harder, so the lower bounds remain valid. The upper bounds as stated assume the primitives known; estimating them adds a cost we do not analyse, which we expect to be of lower order because they enter the response at first order.

\paragraph{Observation contracts.} Each result names what the observer sees. A price-only observer, an observer with executed quantities, one with noisy statements, one with audited equity and one with tagged flow face different experiments. The paper's central message is that these differences, not the observer's statistical skill, determine what can be learned.
\subsection{The observer, its options and its costs}\label{p1:sec:observer}
The results are stated from the side of one observer outside the trader groups: a supervisor, an exchange's surveillance unit or a large desk. Its \emph{target} is attribution, recovering the merged weight and memory rate of each group to a stated accuracy, or deciding between two candidate populations at a stated error. Before it observes, it knows the price persistence $a$, the impact $\kappa$ and the law of the observation noise, and, for the $K$-group results, the number of groups; it does not know their rates or weights. These are the conditions under which the lower bounds are statements about the data rather than about the observer's ignorance; unknown primitives can only enlarge the set of candidate populations.

The observer has three options, and each result prices one of them.
\begin{enumerate}[label=(\alph*),itemsep=1pt]
\item \emph{Wait.} Record stationary prices driven by natural innovations, without trading. The cost is time: the number of observations needed (Theorems~\ref{p1:thm:minphase} and~\ref{p1:prop:passive}).
\item \emph{Probe.} Execute a sequence $v_t$ of trades against the groups and record the price response. Probes execute at the average price of a linear book after the uncontrolled state update; the prober bears the execution loss $\tfrac12v^\top T_hv$ plus fees (Lemma~\ref{p1:lem:toeplitz}) and must fund margin throughout. The counted resource is the number of independent experiments, each starting from the reset state, or the executed quantity and capital of a single experiment (Theorems~\ref{p1:thm:quantity}--\ref{p1:thm:spectral}).
\item \emph{Acquire disclosure.} Receive each class's flow or position, observed with noise, at an operating cost per supervisory cycle and a compliance cost per market (Proposition~\ref{p1:prop:tagged}, Theorem~\ref{p1:thm:discvalue}, Proposition~\ref{p1:prop:choice}).
\end{enumerate}
Within each period the order is: the probe (if any) executes and the price updates, the groups update their states on the new price, and the observer records the price with noise. The groups do not change their rules in response to being observed; Section~\ref{p1:sec:assumptions} discusses what happens if they do. Probing experiments start from the reset state. A probe sent into a running market adds the free response of the current state, which by linearity superposes with the forced response; if the state is unknown it acts as a nuisance term, which can only make separation harder, so the lower bounds still hold, while the upper bounds as stated need the reset state or a known one. Regulators do not, in practice, trade to learn; the probing option describes an outside participant, or prices the information a supervisor would otherwise have to buy by trading.

Three money quantities appear and are kept apart: the \emph{capital required}, the equity an account must hold so that its margin constraint is never breached while it probes (Theorem~\ref{p1:thm:cap} and Lemma~\ref{p1:thm:capital}); the \emph{round-trip trading loss}, the execution loss of a probe that ends flat, which is what the account actually spends (Theorems~\ref{p1:thm:spectral}, \ref{p1:thm:adaptive} and~\ref{p1:thm:rh}); and the \emph{annualised loss} of renewing that probe weekly, used for the comparison with audit-trail budgets (Section~\ref{p1:sec:disc_dollars}). The lower bounds hold for the stated probe classes, deterministic probes and, in Theorem~\ref{p1:thm:adaptive}, non-anticipating sequential probes; they are not lower bounds on information obtained through other channels, which is the point of Section~\ref{p1:sec:disclosure}. Two rates recur and measure different things: $\delta^{-4}$ is the capital (and the average-power horizon) for a \emph{known} symmetric pair whose rates differ by $2\delta$, while $(\Delta\eta)^{-(4K-2)}$, which is $(\Delta\eta)^{-6}$ for $K=2$, is the number of experiments for $K$ \emph{unknown} groups with spacing $\Delta\eta$.

\section{The identification boundary}\label{p1:sec:ident}

A central question in algorithmic market surveillance is whether an observer of executions and prices can identify the individual strategies active in the market.

\begin{theorem}[Merged rate--weight equivalence]\label{p1:thm:merged}
Let $\kappa>0$ and $a\in[0,1)$ be known. Use zero initial states and a common probe contract, and observe the modeled price state $x_t$.
\begin{enumerate}[label=(\roman*)]
\item \emph{Equivalence class.} Two economies $\mathcal E^A=\{(q_i^A,\eta_i^A)\}$ and $\mathcal E^B=\{(q_j^B,\eta_j^B)\}$ produce identical price responses for every probe sequence if and only if their merged measures coincide:
\begin{equation}\label{p1:eq:merged}
\sum_{i:\eta_i^A=\eta} q_i^A\eta \;=\; \sum_{j:\eta_j^B=\eta} q_j^B\eta\qquad\forall\eta\in(0,1].
\end{equation}
Write $a'=a+\kappa M(\{0\})$, so that cohorts with $\eta=1$ (rate $0$) act as part of the anchor. For $a'\neq0$ the minimal realization has dimension $1+|\mathrm{supp}(M)\cap(0,1)|$, counting distinct positive rates with nonzero merged weight after signed cancellation; for $a'=0$ the price carries no state of its own and the dimension is $|\mathrm{supp}(M)\cap(0,1)|$. A cohort with $\eta=1$ merges into the anchor $a$ without adding a pole.
\item \emph{Separation complexity under average power.} Consider testing a single cohort with rate $\bar\eta$ against two cohorts with rates $\bar\eta\pm\delta$ and the same total weight, on a compact family that is uniformly Schur stable with rates bounded away from $0$ and $1$. Observe $\tilde x_t=x_t+\nu_t$ with $\nu_t\sim N(0,\sigma_\nu^2)$ and use deterministic probes with $\frac1T\sum_{t<T}v_t^2\le P$. Any test with worst-case error at most $\alpha\in(0,1/4)$ needs
\begin{equation}\label{p1:eq:Tpower}
T\;\ge\;\frac{4\sigma_\nu^2(1-2\alpha)^2}{C_{\mathcal F}^2P\,\delta^4}=\Omega(\delta^{-4}),
\end{equation}
where $C_{\mathcal F}$ depends only on the family.
\end{enumerate}
\end{theorem}

\paragraph{In words.} If two trading desks update on prices with the same memory, no amount of price data can tell them apart or say how much each trades: prices reveal only how much trading happens at each memory speed. Desks with different memories can be told apart in principle, and the rest of the paper prices how hard that is.


The memory rates are poles of $G$ and \emph{zeros} of $H$, not closed-loop poles. This has a practical consequence we verify in Section~\ref{p1:sec:mixtures}: pole-fitting methods applied directly to the price impulse response recover closed-loop eigenvalues ($0.4256$ and $0.8224$ for $K=1$, $\eta=0.3$, $a=0.5$, $\kappa=0.8$, $q=0.2$) and miss the memory rate $0.70$.

\begin{remark}[Scope]\label{p1:rem:scope}
The result concerns the stated model and observation contract. Cohorts sharing an exact rate enter prices through their merged weight; this does not determine identities or wealth allocations. Nearby rates can be distinguishable with enough information. The $\delta^{-4}$ statement in (ii) is a lower bound on a declared family, with family-dependent constants. Individual algorithms may choose their rates strategically, but any unilateral change maps into the same observer-facing merged measure, which reinforces rather than circumvents the boundary.
\end{remark}

\section{Probing resources, memory resolution and money}\label{p1:sec:resources}

Observation length alone cannot buy identification when the lifetime quantity of controlled executions is fixed. This section connects the merged-memory experiment to a financially specified input/output contract. Throughout, $\Phi$ is the standard normal distribution function.

\subsection{Quantity}

\begin{theorem}[Sharp probing-resource orders for a known symmetric pair]\label{p1:thm:quantity}
Fix $0\le a<1$, $\kappa,q>0$, $S=a+\kappa q<1$ and $\eta\in(0,1)$. Compare a single cohort of weight $q$ and rate $\eta$ with two cohorts of weight $q/2$ and rates $\eta\pm\delta$, both in $[\epsilon,1-\epsilon]$. All parameters are known. The price receives executed quantity through $\kappa v_t$, starts from zero, has no process noise, and is observed with independent $N(0,\sigma^2)$ noise. Probes are deterministic and common to both alternatives. Let $d_T=(h^\delta_j-h^0_j)_{j<T}$ and $K_*=2\kappa^2q/[\epsilon^3(1-S)^2]$.
\begin{enumerate}[label=(\roman*)]
\item For every probe with $\sum_t|v_t|\le V$,
\begin{equation}\label{p1:eq:klV}
\KL(P_0^v\|P_\delta^v)\le\frac{K_*^2V^2\delta^4}{2\sigma^2},\qquad e_{\rm worst}\ge\frac12\Bigl(1-\frac{K_*V\delta^2}{2\sigma}\Bigr)_+ .
\end{equation}
\item The optimal worst-case error over these probes and all classifiers is exactly
\begin{equation}\label{p1:eq:exacterr}
e^*_T(V)=\Phi\!\left(-\frac{V\|d_T\|_2}{2\sigma}\right),
\end{equation}
attained by a pulse at the first input and the likelihood-ratio test. Waiting adds information but $\|d_T\|_2\uparrow\|d_\infty\|_2<\infty$, so it cannot make the error arbitrarily small at fixed $V$. Error at most $\alpha<1/2$ requires $V\ge2\sigma(1-2\alpha)/(K_*\delta^2)$ independently of $T$.
\item The first two taps coincide and $d_2=-\kappa^2q\delta^2$. A single pulse of size
\begin{equation}\label{p1:eq:pulse}
A=\frac{2\sigma\Phi^{-1}(1-\alpha)}{\kappa^2q\delta^2}
\end{equation}
and the third observation achieve error exactly $\alpha$. The necessary and sufficient quantity order is $\delta^{-2}$.
\item Under average power $\sum_{t<T}v_t^2\le PT$, the necessary horizon order is $\delta^{-4}$, attained by the pulse whenever $T\ge\max\{3,\lceil A^2/P\rceil\}$.
\item For a fixed monetary price exposure $\lambda\neq0$, additive reserve $L$ and cost envelope $c\sum|v_t|$, the sharp robust quantity allowance under budget $B>L$ is $V_B=(B-L)/(|\lambda|\kappa+c)$.
\end{enumerate}
\end{theorem}

\paragraph{In words.} To tell one desk from two desks whose memories differ slightly, the observer must trade against them, and the quantity it must trade grows like $1/\delta^2$: halving the difference quadruples the trading needed. Waiting longer after a fixed trade does not help, because the trade's echo in the price dies out.


The identity in (iii) is verified to machine precision for $\delta\in\{0.02,0.05,0.10\}$: the difference is $-1.2\times10^{-4}$, $-7.5\times10^{-4}$ and $-3.0\times10^{-3}$ at $\kappa=1$, $q=0.3$, and $-3.2\times10^{-4}$ at the calibration $\kappa=0.8$, $q=0.2$, $\delta=0.05$. The theorem does not solve estimation of an unknown mixture (Section~\ref{p1:sec:mixtures}).

\begin{figure}[t]\centering
\includegraphics[width=\textwidth]{p1figs/fig_impulse.pdf}
\caption{Left: impulse responses of the single cohort and the split pair ($\delta=0.05$) are nearly indistinguishable. Right: the difference scaled by $\delta^2$ collapses onto one curve, with first nonzero value $-\kappa^2q$ at lag 2 (dotted). Calibration $a=0.5$, $\kappa=0.8$, $q=0.2$, $\eta=0.4$.}\label{p1:fig:impulse}
\end{figure}

\paragraph{A Monte Carlo check of the exact error.} Formula \eqref{p1:eq:exacterr} is exact for the likelihood-ratio test. We simulated 200,000 noisy paths per hypothesis at the calibration ($\sigma=0.01$, $\delta=0.05$) for several pulse sizes and horizons (Table~\ref{p1:tab:mcerr}); simulated errors match the formula to Monte Carlo precision under both hypotheses.

\begin{table}[t]\centering
\begin{threeparttable}
\caption{Exact known-pair error: formula versus simulation}\label{p1:tab:mcerr}
\begin{tabular}{rrccc}
\toprule
Pulse $V$ & Horizon $T$ & $\Phi(-V\|d_T\|_2/2\sigma)$ & Simulated error, $H_0$ & Simulated error, $H_\delta$\\
\midrule
20 & 3 & 0.3745 & 0.3734 & 0.3762\\
40 & 3 & 0.2611 & 0.2599 & 0.2603\\
40 & 10 & 0.0029 & 0.0030 & 0.0027\\
80 & 20 & $1.8\times10^{-10}$ & 0 & 0\\
\bottomrule
\end{tabular}
\begin{tablenotes}\footnotesize\item 200,000 paths per hypothesis; $a=0.5$, $\kappa=0.8$, $q=0.2$, $\eta=0.4$. \end{tablenotes}
\end{threeparttable}
\end{table}

Figure~\ref{p1:fig:impulse} shows why the problem is hard. The two economies have the same direct impact and the same first-lag response; they first differ at lag two, by an amount proportional to $\delta^2$, and the whole difference sequence scales with $\delta^2$. A test therefore needs the signal-to-noise ratio $V\|d\|/\sigma$ to grow like $\delta^{-2}$, which is the quantity order, and because the adverse mark of a probe grows with the square of its size, capital grows like $\delta^{-4}$.

\subsection{Execution cost}

A probe changes inventory as well as price. We therefore specify cash and execution timing. Price has units of currency per asset, and $\kappa$ currency per squared quantity.

\begin{lemma}[Fair execution]\label{p1:lem:toeplitz}
For a quote schedule refreshed after the uncontrolled state update, execution at the average price of a linear book costs, relative to the constant fundamental mark,
\begin{equation}\label{p1:eq:cost}
C(v)=\tfrac12v^\top T_hv+c\|v\|_1,\qquad (T_h)_{ij}=h_{|i-j|}.
\end{equation}
In the positive stable economy $T_h$ is positive definite, so no nonzero round trip produces a deterministic gain.
\end{lemma}

We verify the ledger identity \eqref{p1:eq:cost} against an explicit cash account to $5.7\times10^{-13}$ on 1,000 random probes, and the smallest eigenvalue of $T_h$ over 1,000 random stable positive economies is $0.141$. A direct-only proxy $\sum\tfrac12\kappa v_t^2$ omits cross-time impact: for $v=(+1,-1)$ it gives $0.80$ against a true Toeplitz loss of $h_0-h_1=0.3488$.

\subsection{Capital supplied by a counterparty}

\begin{theorem}[Sharp robust quantity cap from cash and margin]\label{p1:thm:cap}
Use the positive, stable, reset economy with constant fundamental price $P_*>0$ and no process noise. A designated dealer fills only nonnegative probes $v_t\ge0$ at the average price $P^0_{t+1}+\kappa v_t/2$ of a linear schedule, holds inventory $I_{t+1}=I_t-v_t$ and cash $K_{t+1}=K_t+v_t(P^0_{t+1}+\kappa v_t/2)-cv_t$, starts with equity $B>0$, and must satisfy $\mu|I_t|P_t\le E_t$ at every time. For the robust contract admitting every probe with $\sum_tv_t\le V$, the largest allowance guaranteeing margin is
\begin{equation}\label{p1:eq:VB}
V_B=\frac{2B}{\beta+\sqrt{\beta^2+4(\mu+\tfrac12)\kappa B}},\qquad\beta=c+\mu P_*,
\end{equation}
equivalently $(\mu+\tfrac12)\kappa V^2+(c+\mu P_*)V\le B$. The cap holds across all positive cohort representations with $S<1$, without resolving individual rates. An observer that receives only noisy prices and executed quantities has best error $\Phi[-V_B\|d_T\|_2/(2\sigma)]$ within this allowance; fixed error $\alpha<1/2$ requires capital of order at least $\delta^{-4}$, and a single positive pulse attains that order.
\end{theorem}

\paragraph{In words.} A probe moves the price against the account that makes it, so the margin it needs grows with the square of the probe. Because the probe grows like $1/\delta^2$, the capital grows like $1/\delta^4$: halving the difference multiplies the capital needed by 16.


Even with zero initial inventory, the probe creates a quadratic adverse mark that a fixed initial price exposure would miss. Figure~\ref{p1:fig:knownpair} illustrates the experiment.

\begin{figure}[t]\centering
\includegraphics[width=\textwidth]{p1figs/fig_knownpair.pdf}
\caption{Known-pair experiment funded by the counterparty cash/margin contract. (a) Exact Gaussian error at fixed capital 100 and 1,000, separation $\delta=0.05$. (b) The sharp robust quantity allowance \eqref{p1:eq:VB}. (c) Capital for error $0.1$ with 200 observations. Parameters: $a=0.5$, $\kappa=0.8$, $q=0.2$, $\eta=0.4$, observation s.d.\ $0.003$, reference price 100, margin $0.25$, cost $0.02$. No process noise.}\label{p1:fig:knownpair}
\end{figure}

\begin{proposition}[Exact financial disclosure changes the experiment]\label{p1:prop:disclosure}
In this noiseless-state account, exact disclosure of the third equity mark after a single positive pulse $0<A\le V_B$ distinguishes the known pair with zero error for every $\delta>0$.
\end{proposition}

\paragraph{In words.} If the observer can see an account's exact equity, a single test trade settles the question, however close the two candidate markets are. The hard part of the problem is therefore not the market itself but what the observer is allowed to see.

Exact marked equity is extra information, not a free observation assumed by the price-only theorem. Real account marks have their own precision: with reports rounded to $0.01$, a guaranteed change requires a larger trade (Section~\ref{p1:sec:disclosure}).

\subsection{Capital paid by the probe trader}

\begin{lemma}[Time-domain Toeplitz coercivity: capital order for funded probes]\label{p1:thm:capital}
Use the same economy and execution law. The statistical observer receives only noisy prices and executed quantities. The probe trader starts with cash $B>0$ and zero inventory; a finite signed sequence $v$ changes inventory by $+v_t$ and cash by $K_{t+1}=K_t-v_t(P^0_{t+1}+\kappa v_t/2)-c|v_t|$, after which trading stops and the trader must remain solvent as the price returns to $P_*$. Put
\[
r=1-\epsilon(1-S)<1,\qquad m=\kappa\frac{1-r}{1+r}>0,\qquad K_*=\frac{2\kappa^2q}{\epsilon^3(1-S)^2}.
\]
Every solvent strategy satisfies $\KL(P_0^v\|P_\delta^v)\le K_*^2B\delta^4/(\sigma^2m)$, so error at most $\alpha<1/2$ requires
\begin{equation}\label{p1:eq:Bmin}
B\;\ge\;\frac{2\sigma^2m(1-2\alpha)^2}{K_*^2\delta^4},
\end{equation}
without a quantity cap and at any horizon. Conversely, for any margin fraction $0<\mu\le1$, the pulse $A$ in \eqref{p1:eq:pulse} achieves error $\alpha$ and remains margin-solvent at every later mark with capital $B_A=\tfrac{\kappa}{2}A^2+(c+\mu P_*)A$. Necessary and sufficient capital orders are $\delta^{-4}$.
\end{lemma}

Lemma~\ref{p1:thm:capital} is an order statement: an elementary time-domain coercivity argument for Toeplitz forms that establishes the $\delta^{-4}$ rate before the exact frequency-domain optimization of Section~\ref{p1:sec:account}. Its explicit lower constant is correspondingly loose: at the calibrated parameters it evaluates to $4.9\times10^{-4}$ par units at $\delta=0.05$ with $\epsilon=0.35$, five orders of magnitude below any design we can execute, and because $m/K_*^2\propto\epsilon^7(1-S)^5$ it collapses near the boundary of the family. We therefore state it as a lemma; Section~\ref{p1:sec:account} computes the constant that matters, within 5\% of a feasible design at every separation.

\subsection{A worked example: two Kalman desks}\label{p1:sec:worked}

To make the orders concrete, consider two desks that track the same trend with Kalman filters (Lemma~\ref{p1:lem:micro}(i)) and differ only in their signal-to-noise ratios: $\sigma_\zeta^2/\sigma_\nu^2=0.09$ and $0.16$. Their gains are $\eta_1=0.2584$ and $\eta_2=0.3279$, so $\bar\eta=0.2931$ and $\delta=0.0348$. With equal weights, the economy is observationally the split pair of Theorem~\ref{p1:thm:quantity} around $\bar\eta$, and the response difference first appears at lag two, $d_2=-\kappa^2q\delta^2=-1.55\times10^{-4}$ (verified).

A supervisor who wants to tell this market from a single desk with gain $\bar\eta$ at 5\% error using a single pulse and the third observation needs $A=2\sigma\Phi^{-1}(0.95)/(\kappa^2q\delta^2)=212$ units of executed quantity, and, under a 0.5\% margin and 4 basis points of fees, single-pulse capital of $18{,}166$ par units. The pulse is very inefficient: by the $\delta^{-4}$ scaling of Theorem~\ref{p1:thm:spectral}, evaluating Theorem~\ref{p1:thm:spectral} at this pair ($\bar\eta=0.2931$, $\delta=0.0348$) gives about 77 par units for the optimal flat probe (rescaling the $\bar\eta=0.4$ constant by $\delta^{-4}$ alone would give 262, because the constant also depends on $\bar\eta$). Converting with the measured Roll noise (factor $(0.40/0.01)^2$), the optimal probe costs roughly $1.2\times10^5$ USDT and the naive pulse about $3\times10^7$ USDT. Two desks whose signal-to-noise ratios differ by almost a factor of two are therefore expensive, but not impossible, to tell apart by probing; three or more desks are much harder (Theorem~\ref{p1:thm:kcohort}).

\subsection{What prices hide, in economic terms}\label{p1:sec:primitives}

The worked example states the barrier in signal-to-noise ratios. That is the general picture. Write $s=\sigma_\zeta^2/\sigma_\nu^2$ for a desk's signal-to-noise ratio, so that by \eqref{p1:eq:kgain} its memory rate is $1-\eta(s)$ with
\[
\eta(s)=\tfrac12\bigl(\sqrt{s^2+4s}-s\bigr),\qquad \eta'(s)=\tfrac12\Bigl(\frac{s+2}{\sqrt{s^2+4s}}-1\Bigr)>0 ,
\]
and $\eta$ is increasing and concave. With no inventory penalty, the first-order condition of the Remark on exponential memory gives the trading intensity $q_i=\beta_i/\lambda_{{\rm tc},i}$: what the desk expects to earn per unit of belief, over what it pays to trade.

\begin{proposition}[Identification in economic primitives]\label{p1:prop:primitives}
Consider $K$ Kalman desks (Lemma~\ref{p1:lem:micro}(i)) with signal-to-noise ratios $s_i\in[\underline s,\bar s]\subset(0,\infty)$ and intensities $q_i=\beta_i/\lambda_{{\rm tc},i}$, and write $c_-=\eta'(\bar s)$ and $c_+=\eta'(\underline s)$.
\begin{enumerate}[label=(\roman*)]
\item Prices and aggregate order flow depend on the desks only through the measure
\[
\sum_iq_i\,\eta(s_i)\,\delta_{1-\eta(s_i)}.
\]
Desks are identified through their signal-to-noise ratios and intensities only; expected returns and trading costs enter only through their ratio, and two populations with the same intensity-weighted distribution of signal-to-noise ratios are observationally equivalent.
\item Fix a centre $s_c$ and let the ratios lie in $s_c+\Delta s\,[-1,1]\subset[\underline s,\bar s]$. Recovering the intensity-weighted distribution of signal-to-noise ratios to Wasserstein accuracy of order $\Delta s$, or its weights $q_i\eta(s_i)$, from $N$ noisy reset experiments with worst-case error probability below $1/4$ requires $N\ge c\,\sigma^2(\Delta s)^{-(4K-2)}$; on separated configurations, Prony's method of Theorem~\ref{p1:prop:esprit} composed with $\eta^{-1}$ achieves this order. The constants are those of Theorems~\ref{p1:thm:kcohort} and~\ref{p1:prop:esprit} multiplied by powers of $c_-$ and $c_+$.
\end{enumerate}
\end{proposition}

\paragraph{In words.} What prices reveal about Kalman-filter desks is the mix of how good their information is (their signal-to-noise ratios), weighted by how hard they trade. Expected returns and trading costs matter only through their ratio. Pinning down that mix carries the same steep cost as separating memory rates.

\begin{proof}
(i) By Lemma~\ref{p1:lem:micro}(i) desk $i$ is a cohort with memory rate $1-\eta(s_i)$ and weight $q_i$, so Theorem~\ref{p1:thm:merged} applies, and $\beta_i,\lambda_{{\rm tc},i}$ enter the economy only through $q_i$. (ii) On $[\underline s,\bar s]$, $\eta$ is a $C^1$ bijection onto its image with $c_-\le\eta'\le c_+$. Lower bound: take the two economies of Theorem~\ref{p1:thm:kcohort} with centre $r_c=1-\eta(s_c)$ and rate spacing $\Delta\eta=c_-\Delta s$. By the mean value theorem their gains are images of ratios in $s_c+\Delta s\,[-1,1]$, so they are two populations of Kalman desks in the class, and the bound $N\ge c\,\sigma^2(\Delta\eta)^{-(4K-2)}=c\,c_-^{-(4K-2)}\sigma^2(\Delta s)^{-(4K-2)}$ applies to recovering their weights. Because $\eta$ is $c_+$-Lipschitz, the Wasserstein distance between the two distributions of ratios is at least $1/c_+$ times that between the two distributions of rates, which Theorem~\ref{p1:thm:kcohort} bounds below by a constant times $\Delta\eta=c_-\Delta s$. Upper bound: ratios in $s_c+\Delta s\,[-1,1]$ with spacing at least $\gamma\Delta s$ have gains in $\eta(s_c)+c_+\Delta s\,[-1,1]$ with spacing at least $c_-\gamma\Delta s$, a class of the form required by Theorem~\ref{p1:prop:esprit} with rescaled separation $\gamma c_-/c_+$; Prony's method recovers the rates and weights at that theorem's rate, and $\eta^{-1}$, which is $1/c_-$-Lipschitz, maps rate errors to ratio errors.
\end{proof}

The private information a supervisor would most like to know, how good each desk's information is, is exactly what prices hide. In the worked example the two desks' ratios differ by a factor of almost two, and separating them already costs hundreds of thousands of dollars of probing capital. The proposition takes intensities to vary independently of information quality, because they also reflect costs and capital. The next proposition asks what happens when the economics ties them together.

\begin{proposition}[Intensity tied to information quality]\label{p1:prop:tied}
Let each desk's intensity be a known function of its signal-to-noise ratio, $q=\varphi(s)$ with $\varphi$ positive and $C^2$ on $[\underline s,\bar s]$ (for a mean--variance desk, $\varphi(s)$ increases with $s$), and let $n_i\ge0$ denote the mass of desks with ratio $s_i$ (a divisible population, not an integer desk count). Write $s(\eta)$ for the inverse gain map and $w(\eta)=\varphi(s(\eta))\,\eta$ for the weight per desk as a function of its gain.
\begin{enumerate}[label=(\roman*)]
\item \emph{Unknown numbers of desks.} The merged measure is $\sum_in_i\,w(\eta(s_i))\,\delta_{1-\eta(s_i)}$. Every positive measure with finitely many atoms in $[1-\eta(\bar s),1-\eta(\underline s)]$ is of this form, with $n_i=m_i/w(\eta_i)$ for an atom of mass $m_i$ at $1-\eta_i$. So when the numbers of desks of each type are not known, tying intensity to information quality removes no economy, and every lower bound of Proposition~\ref{p1:prop:primitives}(ii) holds unchanged.
\item \emph{Known numbers of desks.} If the $n_i$ are also known, the set of economies is smaller, but the single-versus-pair problem keeps its order: one type at gain $\bar\eta$ with $2n$ desks and two types at $\bar\eta\pm\delta$ with $n$ desks each have taps that differ by $n\,\delta^2\,\partial_\eta^2[w(\eta)(1-\eta)^m]_{\bar\eta}+o(\delta^2)$, so separating them needs $N\ge c\,\sigma^2\delta^{-4}$ reset experiments, the order the symmetric pair has when intensities vary freely.
\item \emph{Self-confirming configurations.} Suppose the ratio each desk perceives is the one implied by the closed-loop dynamics it trades in, through a map $T$ from configurations $(s_1,\dots,s_K)$ to perceived ratios that is continuous on a box $\prod_i[\underline s,\bar s]$ and maps it into itself. Then a configuration with $T(s)=s$ exists, and parts (i)--(ii) and Proposition~\ref{p1:prop:primitives} apply at it.
\end{enumerate}
\end{proposition}

\begin{proof}
(i) Desk type $i$ is a cohort with gain $\eta_i=\eta(s_i)$, memory rate $1-\eta_i$ and total weight $n_i\varphi(s_i)$, so it contributes $n_i\varphi(s_i)\eta_i=n_iw(\eta_i)$ to the atom at $1-\eta_i$. Since $w>0$, any atom mass $m_i\ge0$ is attained with $n_i=m_i/w(\eta_i)$. (ii) Write $g_m(\eta)=w(\eta)(1-\eta)^m$; the tap difference is $n[g_m(\bar\eta+\delta)+g_m(\bar\eta-\delta)-2g_m(\bar\eta)]=n\delta^2g_m''(\bar\eta)+o(\delta^2)$, with $\sum_m\sup|g_m''|^2<\infty$ on gains bounded away from zero, so the Euclidean distance between the two tap vectors is $O(n\delta^2)$ and the Kullback--Leibler divergence over $N$ experiments is $O(Nn^2\delta^4/\sigma^2)$; Le Cam's two-point argument gives the bound. (iii) is Brouwer's fixed-point theorem.
\end{proof}

So the barrier does not depend on intensities varying freely: it survives any known link between intensity and information quality as long as the supervisor does not also know how many desks of each kind there are, and even then the simplest separation problem keeps its $\delta^{-4}$ order. Whether the full $(\Delta\eta)^{-(4K-2)}$ exponent survives with known numbers of desks is open. Part~(iii) shows only that self-confirming configurations exist under continuity; it does not derive the map $T$ from a model of the desks' learning.

\subsection{Natural excitation and the cost of ambiguity}

Natural price innovations can also separate the pair with zero controlled quantity: an exact finite-sample test on independent three-step blocks needs order $\delta^{-4}$ blocks (Internet Appendix, Proposition~\ref{p1:prop:passive0}), the same order as the passive Whittle bound of Section~\ref{p1:sec:mixtures}. Non-identification also has a price on the quote side: when the relevant set of economies includes costlier ones, a robust dealer widens its quote, which Chapter~\ref{ch:p2} analyses; cohort non-identification alone does not establish that the worst relevant cost exceeds the informed one.
\section{Account-level capital}\label{p1:sec:account}

Lemma~\ref{p1:thm:capital} gives the order but not the constant. This section computes what a trader actually needs under a calibrated linear-impact environment, then converts it to Bitcoin dollars.

\paragraph{Numerical values.} Decimal optima and brackets below are numerical evaluations of the stated bounds.

\subsection{Observation design and the tri-region partition}

An observer distinguishes two candidate mechanisms with rates $\theta_0$ and $\theta_1=\theta_0+\delta$ over $N$ periods, submitting probes $V=(V_1,\ldots,V_N)^\top$ that generate inventory $I_t=\sum_{s\le t}V_s$. Under contract $I_0$ (active probing, Gaussian observation noise) $p_t=(h_i*V)_t+\epsilon_t$, $\epsilon_t\sim N(0,\sigma^2)$, so
\begin{equation}\label{p1:eq:klconv}
\KL(P_0^{\otimes N}\|P_1^{\otimes N})=\frac1{2\sigma^2}\sum_{t=1}^N\bigl((h_1*V)_t-(h_0*V)_t\bigr)^2 .
\end{equation}
The observer minimizes capital $C(N,V)=\sum_t(\lambda V_t^2+\phi|V_t|)+\max_t\mu p_0|I_t|$ subject to (i) separation $\KL\ge\tfrac12(z_\alpha+z_\beta)^2$, (ii) caps $|V_t|\le V_{\rm eff}$ and $|I_t|\le I_{\max}$, and (iii) a solvency budget $\max_t{\rm Loss}_t(V)\le B$. Here $V_{\rm eff}=\min(V_{\max},D_{\max}/\kappa)$ combines resting depth and a price-displacement limit.

\begin{proposition}[Tri-region partition and finite-session resolution floor]\label{p1:prop:tri}
Under observation horizon $T$ per pulse, $(\delta,B)$ space splits into three regions.
\begin{enumerate}[label=(\roman*)]
\item \emph{Feasible single pulse.} For $\delta\ge\delta_{\min}(T)$ and $B\ge C^*(\delta)$ one pulse within capacity suffices. For a simple rate shift ($\eta_1=\eta_0+\delta$) the response difference is $\Theta(\delta)$ and $C^*=\Theta(\delta^{-2})$; for symmetric splitting the first-order variation cancels, the difference is $\Theta(\delta^2)$ and $C^*=\Theta(\delta^{-4})$.
\item \emph{Sequential bounded probing.} For $\delta<\delta_{\min}(T)$ the required quantity exceeds $V_{\rm eff}$. Capped pulses each supply divergence $V_{\rm eff}^2\|d_T\|^2/(2\sigma^2)$, so
\begin{equation}\label{p1:eq:npulses}
N_{\rm pulses}(T,\delta)=\Bigl\lceil\frac{4\sigma^2z_\alpha^2}{V_{\rm eff}^2\|d_T\|_2^2}\Bigr\rceil .
\end{equation}
\item \emph{Infeasible with isolated pulses.} If the budget cannot finance sequential costs, peak inventory breaches $I_{\max}$, or a session deadline binds, the finite-session floor $\delta_{\min}(T_{\rm session})$ is the root of $V_{\rm eff}\|d_{T_{\rm cycle}}(\delta)\|_2=2\sigma z_\alpha/\sqrt{\lfloor T_{\rm session}/T_{\rm cycle}\rfloor}$, with small-$\delta$ approximation $\delta_{\min}=\Theta(T_{\rm session}^{-1/4})$.
\end{enumerate}
\end{proposition}

\paragraph{In words.} A single test trade works only when the groups are far enough apart and the account has enough capital. Closer groups need many capped test trades, and a longer observation window after each trade makes every test trade more informative.

At the calibration $\kappa=0.8$, $q=0.2$, $\sigma=0.01$, $V_{\max}=50$, $D_{\max}=5$ ($V_{\rm eff}=6.25$), the single-tap floor is $\delta_{\min}(3)=\sqrt{2\sigma z_\alpha/(V_{\rm eff}\kappa^2q)}=0.2028$; extended observation lowers it to $0.1376$ ($T=4$), $0.0973$ ($T=10$), $0.0914$ ($T=20$) and $0.0905$ ($T=100$). At $\delta=0.05$ the simplified one-leg benchmark (cost $\tfrac12\kappa V_{\rm eff}^2+\phi P_*V_{\rm eff}=15.875$ per pulse plus initial margin $3.125$) needs 271 pulses and $4{,}305.25$ par units with $T=3$, 58 pulses and $923.88$ with $T=4$, but only 12 pulses and $193.63$ with $T=20$ and 11 pulses and $177.75$ with $T=100$; optimizing the pulse size (6.06 rather than 6.25) lowers the $T=20$ figure to $182.03$. Most of the saving comes from observing the transient response rather than a single tap (Figure~\ref{p1:fig:phase}).

\begin{figure}[t]\centering
\includegraphics[width=0.62\textwidth]{p1figs/fig_phase.pdf}
\caption{Capital against separation: isolated-pulse benchmarks with three and twenty observations per pulse, and the spectral lower bound of Theorem~\ref{p1:thm:spectral}. The dotted line is the single-tap resolution floor $\delta_{\min}(3)=0.2028$. Recomputed in this revision.}\label{p1:fig:phase}
\end{figure} These are costs within an isolated-pulse design class, not minimum capital; the next subsection replaces the formula by a ledger.

\subsection{Account-level capital and the gap to the lower bound}

We run each candidate design through an explicit synthetic ledger at reference price $P_*=100$: fills at the average linear-book price, taker fees of 4 basis points, marks at the prevailing price, and 0.5\% maintenance margin. Initial capital is $C_0=\max_t({\rm margin}_t-{\rm equity}_t)$, maximized over both hypotheses, and the test uses the full price path including unwind legs.

\begin{table}[t]\centering\small
\begin{threeparttable}
\caption{Account-level initial capital at $\delta=0.05$ (par units)}\label{p1:tab:ledger}
\setlength{\tabcolsep}{4pt}
\begin{tabular}{lrrrrrr}
\toprule
Design & $C_0$ & Error & End inv. & Peak margin & Max $|v_t|$ & Trades\\
\midrule
12 buys, $T=20$ apart, never unwound & 229.13 & 0.0578 & 75.00 & 39.38 & 6.25 & 12\\
6 buy/sell round trips & 189.45 & 0.0308 & 0.00 & 3.28 & 6.25 & 12\\
5 buy/sell round trips & 157.89 & 0.0442 & 0.00 & 3.28 & 6.25 & 10\\
12 buys, each unwound after 20 steps & 378.82 & 0.0040 & 0.00 & 3.28 & 6.25 & 24\\
Generalized-eigenvector probe ($L=120$) & 65.06 & 0.0500 & 0.00 & 2.74 & 1.24 & 120\\
Generalized-eigenvector probe ($L=160$) & 64.86 & 0.0500 & 0.00 & 2.17 & 1.09 & 160\\
Fee-refined probe ($L=240$) & 64.48 & 0.0500 & 0.00 & --- & 1.30 & 240\\
\bottomrule
\end{tabular}
\begin{tablenotes}\footnotesize\item The pulse designs place a pulse every 40 steps and close it 20 steps later; the fee-refined probe is re-optimized by sequential quadratic programming over $L=240$. The numerically evaluated lower bound of Theorem~\ref{p1:thm:spectral} is 61.41; the explicit constant of Lemma~\ref{p1:thm:capital} is $4.9\times10^{-4}$.\end{tablenotes}
\end{threeparttable}
\end{table}

Table~\ref{p1:tab:ledger} gives three conclusions. First, twelve buys that are never unwound fail the 5\% target once the test uses the joint price path (divergence $4.95<5.41$) and end with 75 units of inventory. Second, if every unwind leg is itself informative, five or six round trips meet the target with zero terminal inventory. Third, pulses are not optimal: a generalized-eigenvector probe, which maximizes divergence per unit of fee-free round-trip cost, needs about 65 par units with maximum trade 1.25 and maximum displacement 2.23, well inside limits.

For flat fee-paying policies $v\in\R^L$ with $\sum_tv_t=0$, the capital problem is
\begin{equation}\label{p1:eq:capprob}
\min_v\ \max_{i\in\{0,1\}}L_i(v)+\sum_t\phi P_*|v_t|+\max_t\mu P_*|I_t|\quad\text{s.t.}\quad\|(H_1-H_0)v\|_2\ge2\sigma z_{1-\alpha},
\end{equation}
where $L_i(v)=\tfrac12v^\top T_{h_i}v$ is the exact Toeplitz round-trip loss. Zero-padding embeds horizon $L$ into $L+1$ at identical cost and information, so the minimum is non-increasing in $L$.

\begin{theorem}[Spectral dissipation lower bound]\label{p1:thm:spectral}
Let $\hat H_i(z)=\sum_{j\ge0}h_{i,j}z^{-j}$, $\Delta\hat H=\hat H_1-\hat H_0$, and define the dissipation kernel
\begin{equation}\label{p1:eq:gdiss}
g_i(\omega)=\mathrm{Re}\,\hat H_i(e^{i\omega})-\kappa/2>0\qquad\forall\omega\in[-\pi,\pi].
\end{equation}
\begin{enumerate}[label=(\roman*)]
\item \emph{Toeplitz energy identity.} For any flat round trip, $L_i(v)=\frac1{2\pi}\int_{-\pi}^{\pi}g_i(\omega)|V(e^{i\omega})|^2d\omega$.
\item \emph{Lower bound.} $\|(H_1-H_0)v\|_2^2\le\sup_\omega\frac{|\Delta\hat H(e^{i\omega})|^2}{g_i(\omega)}L_i(v)$, so every deterministic round trip with error at most $\alpha$ has
\begin{equation}\label{p1:eq:Cdet}
L_i(v)\;\ge\;4\sigma^2z_{1-\alpha}^2\inf_\omega\frac{g_i(\omega)}{|\Delta\hat H(e^{i\omega})|^2}\;\equiv\;C^*_{\rm det,quad}(\delta).
\end{equation}
\item \emph{Bracket.} At $\delta=0.05$, $C^*_{\rm det,quad}=61.14$ at $\omega^*=0.2526$; adding the unavoidable fee $\phi P_*\cdot2\sigma z/\sup_\omega|\Delta\hat H|$ gives the numerical lower estimate $C_{\rm det}\ge61.41$. A feasible 240-step sequence attains 64.48, so the fee-paying optimum lies in $[61.41,64.48]$.
\item \emph{Process noise.} Let the price also receive the innovations $w_t\sim N(0,\sigma_w^2)$ of \eqref{p1:eq:price}, independent of observation noise $\nu_t\sim N(0,\sigma_\nu^2)$, so that from the reset state the observer sees $\tilde x_t=(h_i*v)_t+\kappa^{-1}(h_i*w)_t+\nu_t$, $t<T$. For the known pair of Theorem~\ref{p1:thm:quantity}, put $\sigma_{\rm eff}^2=\sigma_\nu^2+\sigma_w^2/4$, $\bar M=1/(1-S)$ and $C_w=2\sigma_w^4\bar M^2K_*^2(\sigma_\nu^2+\sigma_w^2\bar M^2)/(\kappa^2\sigma_\nu^2\sigma_{\rm eff}^4)$. Every deterministic round trip whose test has errors $\alpha$ and $\beta$ under the two hypotheses satisfies
\begin{equation}\label{p1:eq:Cnoise}
L_i(v)\;\ge\;2\sigma_{\rm eff}^2\bigl[{\rm kl}(\alpha,1-\beta)-C_wT\delta^4\bigr]_+\inf_\omega\frac{g_i(\omega)}{|\Delta\hat H(e^{i\omega})|^2},
\end{equation}
with ${\rm kl}$ as in Theorem~\ref{p1:thm:adaptive}. In-loop noise therefore leaves the $\delta^{-4}$ capital order intact for every horizon $T\le{\rm kl}(\alpha,1-\beta)/(2C_w\delta^4)$; at horizons of order $\delta^{-4}$ and longer, the noise itself separates the pair without a probe (Theorem~\ref{p1:prop:passive}).
\end{enumerate}
\end{theorem}

\paragraph{In words.} This turns the order into a dollar floor. No round trip that tells the two candidate markets apart at 5\% error can lose less than this amount, and a concrete probe comes within 5\% of it.


\begin{figure}[t]\centering
\includegraphics[width=\textwidth]{p1figs/fig_dissipation.pdf}
\caption{Left: the dissipation kernel $g_0(\omega)$ is strictly positive. Right: the capital ratio in \eqref{p1:eq:Cdet}; dots mark its minimum, near $\omega^*\approx0.25$ for every separation. Recomputed in this revision.}\label{p1:fig:dissip}
\end{figure}

Positivity of $g_i$ follows from Lemma~\ref{p1:lem:cm}: $g_i(\omega)=\tfrac12\sum_\ell\omega_\ell\frac{1-\lambda_\ell^2}{1-2\lambda_\ell\cos\omega+\lambda_\ell^2}>0$. We compute $C^*_{\rm det,quad}$ across separations (Table~\ref{p1:tab:capbounds}): $61.14$ at $\delta=0.05$, with $\delta^4C^*_{\rm det,quad}$ between $3.67\times10^{-4}$ and $3.84\times10^{-4}$, confirming the $\delta^{-4}$ law. The fee term, the feasible upper bounds and the adaptive bound of Theorem~\ref{p1:thm:adaptive} are reported in the same table.

The infimum over $\omega$ is evaluated on a fine frequency grid; a local refinement agrees to four digits ($61.1377$ at $\omega^*=0.2526$), and the reported constant is the larger of the two hypotheses' bounds. ``Certified'' here means a numerical evaluation of the stated bound, not an interval-arithmetic enclosure.

\begin{theorem}[Adaptive and sequential probing]\label{p1:thm:adaptive}
Let $(v_t)_{t\le\tau}$ be any non-anticipating policy adapted to the price filtration, with an almost surely finite stopping time and $\sum_{t\le\tau}v_t=0$, followed by a terminal test $\psi_\tau\in\{0,1\}$ with $P_0(\psi_\tau=1)\le\alpha$ and $P_1(\psi_\tau=0)\le\beta$. Then
\begin{equation}\label{p1:eq:adaptive}
\E_0[L_0(v)]\;\ge\;2\sigma^2\,{\rm kl}(\alpha,1-\beta)\inf_\omega\frac{g_0(\omega)}{|\Delta\hat H(e^{i\omega})|^2},
\end{equation}
where ${\rm kl}(\alpha,1-\beta)=\alpha\log\tfrac{\alpha}{1-\beta}+(1-\alpha)\log\tfrac{1-\alpha}{\beta}$.
Hence adaptation can lower expected fee-free loss by at most the factor ${\rm kl}(\alpha,1-\beta)/(2z_{1-\alpha}^2)\approx0.4897$ at $\alpha=\beta=0.05$, and cannot alter the $\delta^{-4}$ order.
\end{theorem}

\paragraph{In words.} Letting the observer adapt its test trades as data arrive, and stop as soon as it is confident, helps only a little: at the calibration it can cut the expected loss by at most about half. The steep cost is a property of the market, not of a clumsy probing plan.

At $\delta=0.05$ the adaptive bound is $29.94$, or $30.22$ with the minimal fee. Pairing it with the feasible deterministic 64.48 gives a numerical ratio of at least $0.464$. In 20{,}000 simulated sequential probability-ratio tests, which split the optimal probe into ten copies of equal information, mean loss was $39.5$ under $H_0$, between the two bounds, but losses vary with the stopping time: the 95th and 99th percentiles of required account capital were $91.7$ and $128.4$, and funding a trader with only the mean loss gives a 37.5\% insolvency rate before the test ends. Expected-loss savings are not solvency savings.

\begin{table}[t]\centering\small
\begin{threeparttable}
\caption{Capital lower bounds and feasible upper bounds across separations (par units)}\label{p1:tab:capbounds}
\begin{tabular}{lcccccc}
\toprule
$\delta$ & $\omega^*$ & $C^*_{\rm det,quad}$ & $\delta^4C^*_{\rm det,quad}$ & Certified $C_{\rm det}$ & Feasible UB & Adaptive LB\\
\midrule
0.02 & 0.2546 & 2397.85 & $3.84\times10^{-4}$ & 2399.58 & 2428.25 & 1176.04\\
0.03 & 0.2541 & 473.19 & $3.83\times10^{-4}$ & 473.97 & 484.50 & 232.51\\
0.05 & 0.2526 & 61.14 & $3.82\times10^{-4}$ & 61.41 & 64.48 & 30.22\\
0.08 & 0.2488 & 9.26 & $3.79\times10^{-4}$ & 9.37 & 10.38 & 4.64\\
0.10 & 0.2453 & 3.77 & $3.77\times10^{-4}$ & 3.83 & 4.44 & 1.91\\
0.15 & 0.2331 & 0.73 & $3.67\times10^{-4}$ & 0.76 & 1.01 & 0.38\\
\bottomrule
\end{tabular}
\begin{tablenotes}\footnotesize\item Calibration $a=0.5$, $\kappa=0.8$, $q=0.2$, $\eta=0.4$, $\sigma=0.01$, $\alpha=0.05$. From \texttt{recompute\_theory.py} and \texttt{recompute\_p1\_capital.py}; the numerical-bound column adds the fee lower bound to the fee-free infimum; the feasible column is the best flat probe found over $L\le240$; the adaptive column adds the fee lower bound to the adaptive bound.\end{tablenotes}
\end{threeparttable}
\end{table}

\begin{theorem}[Finite-horizon Rayleigh--H\"older duality]\label{p1:thm:rh}
For a horizon of $L\ge4$ trading periods, let $E\in\R^{L\times(L-1)}$ be the first-difference basis of flat probes ($\sum_tv_t=0$), $Q_i=\tfrac12T_{h_i}$ (so that $L_i(v)=v^\top Q_iv$), $A_{\rm diff}=(H_1-H_0)^\top(H_1-H_0)$, and let $G_i=E^\top Q_iE\succ0$ and $M=E^\top A_{\rm diff}E\succeq0$ be the reduced execution-loss and information matrices, with $G_i$-orthonormal generalized eigenvectors $e_j$ and eigenvalues $0\le\lambda_1\le\cdots\le\lambda_{L-3}<\lambda_{L-2}\le\lambda_{L-1}$; put $c_{\rm req}=4\sigma^2z_{1-\alpha}^2$ and $c_{\rm fee}=\phi P_*$. For $\theta\in[0,\pi)$ let $d(\theta)=\cos\theta\,e_{L-2}+\sin\theta\,e_{L-1}$, $\lambda(\theta)=\lambda_{L-2}\cos^2\theta+\lambda_{L-1}\sin^2\theta$, $s_\theta={\rm sign}(Ed(\theta))$, and
\begin{equation}\label{p1:eq:rhdefs}
\begin{gathered}
F_{\rm dom}(\theta)=c_{\rm fee}\sqrt{\frac{c_{\rm req}}{\lambda(\theta)}}\,\|Ed(\theta)\|_1,\qquad W_\perp(\theta)=c_{\rm fee}^2\sum_{j\le L-3}\frac{(s_\theta^\top Ee_j)^2}{1-\lambda_j/\lambda(\theta)},\\
\gamma(\theta)=\frac{\lambda(\theta)\,\lambda_{L-3}}{(\lambda(\theta)-\lambda_{L-3})\,c_{\rm req}} .
\end{gathered}
\end{equation}
Then every flat probe meeting the separation constraint of \eqref{p1:eq:capprob} satisfies
\begin{equation}\label{p1:eq:rh}
L_i(v)+c_{\rm fee}\|v\|_1\;\ge\;\min_{\theta}\Bigl\{\frac{c_{\rm req}}{\lambda(\theta)}+\Psi(\theta)\Bigr\}\;\equiv\;C_{\rm RH}(L)\;\ge\;\frac{c_{\rm req}}{\lambda_{L-1}},
\end{equation}
\[
\Psi(\theta)=\min_{x\ge0}\Bigl\{x^2+\Bigl[F_{\rm dom}(\theta)\bigl(1-\gamma(\theta)x^2\bigr)_+-\sqrt{W_\perp(\theta)}\,x\Bigr]_+\Bigr\}\;\ge0 ,
\]
a piecewise-quadratic minimization in closed form; when $\gamma F_{\rm dom}<1$ and the minimizer has $\gamma x^2\le1$, $\Psi\ge\max\bigl(0,F_{\rm dom}-W_\perp/(4(1-\gamma F_{\rm dom}))\bigr)$, which is what the bound uses.
\end{theorem}

\paragraph{In words.} This is the exact cheapest probe over a finite number of trading periods, fees included. It confirms that the dollar floor of Theorem~\ref{p1:thm:spectral} is not an artefact of assuming an infinitely long probe.

The bound is computable for any horizon. Table~\ref{p1:tab:rh} reports its quadratic part $c_{\rm req}/\lambda_{L-1}$, which decays toward the infinite-horizon value 61.14 at rate $O(L^{-1})$, the full bound $C_{\rm RH}(L)$, and the best feasible fee-paying probe we find at each horizon. At $L=80$ the full bound is $64.72$ against a best probe of $64.76$ on the same loss-plus-fee objective (Table~\ref{p1:tab:rh} reports $65.02$ for ledger capital, which also includes margin), so the dual is essentially tight at short horizons; at longer horizons the fee term of the dual weakens and the bracket widens to about 5\%. The fee of unpadded full-support eigenvectors grows as $\Theta(\sqrt L)$, which is why zero-padded or $L_1$-localized probes are needed at long horizons.

\begin{table}[t]\centering\small
\begin{threeparttable}
\caption{Finite-horizon capital at $\delta=0.05$ (par units)}\label{p1:tab:rh}
\setlength{\tabcolsep}{4pt}
\begin{tabular}{lcccc}
\toprule
Horizon $L$ & Quadratic part $c_{\rm req}/\lambda_{L-1}$ & Full bound $C_{\rm RH}(L)$ & Best feasible probe & Gap to $C_{\rm det}=61.41$\\
\midrule
80 & 62.55 & 64.72 & 65.02 & 5.9\%\\
120 & 61.84 & 63.39 & 64.59 & 5.2\%\\
160 & 61.54 & 62.46 & 64.50 & 5.0\%\\
240 & 61.32 & 61.76 & 64.48 & 5.0\%\\
320 & 61.25 & --- & $\le64.48$ & $\le5.0\%$\\
480 & 61.19 & --- & $\le64.48$ & $\le5.0\%$\\
\bottomrule
\end{tabular}
\begin{tablenotes}\footnotesize\item Quadratic bound $c_{\rm req}/\lambda_{L-1}$ from the generalized eigenproblem of Theorem~\ref{p1:thm:rh} on flat probes (equivalently $c_{\rm req}/(2\lambda'_L)$ with $\lambda'$ computed against $T_{h_0}$ rather than $\tfrac12T_{h_0}$); it converges to the infinite-horizon value 61.14. The full bound minimizes \eqref{p1:eq:rh} over 3{,}600 values of $\theta$ (\texttt{hardening/rh\_bound.py}); at $L=80$ and $120$ a direct numerical minimization of loss plus fee under the separation constraint gives $64.76$ and $64.43$, above the bound. Feasible probes: top generalized eigenvector refined by sequential quadratic programming on the full ledger; at $L=320$ and $480$ the $L=240$ probe padded with zeros is feasible at the same cost, so the best feasible value can only be lower. \end{tablenotes}
\end{threeparttable}
\end{table}

\subsection{Robustness: heteroskedastic impact and probing protocols}

\begin{corollary}[Invariance under impact-induced heteroskedasticity]\label{p1:cor:garch}
Suppose probing inflates innovation variance as $\sigma^2_{i,t}(v)=\sigma^2(1+\xi[(h_i*v)_t]^2)$ and slippage scales as $\kappa_{\rm eff}=\kappa(1+\xi_0|v_t|)$. For capacity-bounded probes, $\KL(P_0^v\|P_1^v)\le\frac{1+O(\xi^2V_{\rm eff}^2)}{2\sigma^2}\|(H_1-H_0)v\|_2^2=O(\|v\|_2^2\delta^4)$, and convex impact only raises dissipation. Required capital remains $\Theta(\delta^{-4})$.
\end{corollary}

Four probing protocols (unconstrained single pulse, repeated capped pulses, a sequential test with dynamic throttling, and a predetermined generalized-eigenvector trajectory) all show capital slopes near $-4$ in $\log\delta$: $-3.90$, $-3.84$, $-3.47$ and $-3.81$ over $\delta\in[0.03,0.12]$, where the sequential slope is flattened by the 1.7 pulses needed at $\delta=0.12$ (Appendix Table~\ref{p1:tab:protocols}). At equal observation horizon, sequential stopping saves 14--49\% of capital relative to repeated capped pulses, but most of the apparent saving from the single-tap baseline comes from observing the transient response, not from adaptivity.

\subsection{From par units to Bitcoin dollars}\label{p1:sec:dollars}

Capital scales as $\sigma^2\delta^{-4}$, so the dollar magnitude hinges on the observation noise. Treating the 4 USDT ten-second price volatility as observation noise would give 9.8 million USDT at $\delta=0.05$, but that conflates price volatility with measurement error. We now estimate observation noise from Binance BTCUSDT one-second closes on January 10--20, 2023, with the Roll covariance estimator $\hat\sigma=\sqrt{-{\rm cov}(\Delta p_t,\Delta p_{t-1})}$ \citep{p1:roll1984}. Daily estimates range from $0.09$ to $0.77$ USDT, median $0.40$ USDT.

\begin{table}[t]\centering
\begin{threeparttable}
\caption{Capital lower bound under alternative observation-noise calibrations (USDT)}\label{p1:tab:capital}
\begin{tabular}{lrrr}
\toprule
Noise calibration & $\delta=0.02$ & $\delta=0.05$ & $\delta=0.10$\\
\midrule
Roll estimator, median (0.40 USDT) & 3{,}874{,}696 & 99{,}193 & 6{,}200\\
One-second price-change s.d., median (1.98 USDT) & 93{,}688{,}369 & 2{,}398{,}422 & 149{,}902\\
Earlier assumption (4 USDT, upper scenario) & 383{,}812{,}500 & 9{,}825{,}600 & 614{,}100\\
\bottomrule
\end{tabular}
\begin{tablenotes}\footnotesize
\item $C(\sigma,\delta)=61.41\times(\sigma/0.01)^2\times(0.05/\delta)^4$, with $\sigma$ the unrounded median. Noise estimated on 11 days of checksum-verified Binance BTCUSDT one-second klines. \end{tablenotes}
\end{threeparttable}
\end{table}

\begin{table}[t]\centering\small
\begin{threeparttable}
\caption{Daily observation-noise estimates, Binance BTCUSDT one-second closes}\label{p1:tab:roll}
\begin{tabular}{lccc}
\toprule
Date & Roll $\hat\sigma$ (USDT) & One-second change s.d.\ (USDT) & Mean price (USDT)\\
\midrule
2023-01-10 & 0.152 & 0.679 & 17{,}288\\
2023-01-11 & 0.092 & 0.891 & 17{,}456\\
2023-01-12 & 0.232 & 2.292 & 18{,}343\\
2023-01-13 & 0.348 & 1.762 & 19{,}075\\
2023-01-14 & 0.773 & 3.511 & 20{,}818\\
2023-01-15 & 0.482 & 1.790 & 20{,}778\\
2023-01-16 & 0.629 & 2.357 & 21{,}036\\
2023-01-17 & 0.457 & 1.976 & 21{,}194\\
2023-01-18 & 0.695 & 2.541 & 21{,}122\\
2023-01-19 & 0.402 & 1.316 & 20{,}850\\
2023-01-20 & 0.316 & 2.038 & 21{,}306\\
\midrule
Median & 0.402 & 1.976 & \\
\bottomrule
\end{tabular}
\begin{tablenotes}\footnotesize\item The Roll estimator measures the negative first-order autocovariance of price changes, which reflects bid-ask bounce and recording noise rather than price volatility.\end{tablenotes}
\end{threeparttable}
\end{table}

\begin{figure}[t]\centering
\includegraphics[width=0.62\textwidth]{p1figs/fig_capital.pdf}
\caption{Capital lower bound against memory separation for three observation-noise calibrations.}\label{p1:fig:capital}
\end{figure}

At realistic noise, separating rates that differ by $0.05$ costs about $9.9\times10^4$ USDT, which is not prohibitive for a well-capitalized desk (Table~\ref{p1:tab:capital}). The barrier becomes large only for finer separation ($3.9$ million USDT at $\delta=0.02$) and, through Theorem~\ref{p1:thm:kcohort}, when several cohorts are unknown. Two factors we do not model would raise these figures: walking several levels of a real book costs more than the linear top-of-book impact assumed here, and impact-induced heteroskedasticity raises the noise a probe must overcome. The capital figures are therefore lower bounds for the stated model.
\section{Unknown mixtures and passive data}\label{p1:sec:mixtures}

Theorem~\ref{p1:thm:quantity} concerns a \emph{known} pair. A supervisor faces an unknown number of cohorts with unknown rates and weights. This section asks how fast the cost of recovering them grows.

\subsection{Why naive pole fitting fails}

Because the memory rates are zeros of $H$ (Theorem~\ref{p1:thm:merged}), extracting poles directly from the price impulse response by matrix-pencil or Prony methods targets closed-loop eigenvalues. A four-step inversion repairs this: (1) decouple $g_t=h_{t+1}-ah_t$; (2) form a Hankel matrix of $g$ and compute a minimal realization by Kung's eigensystem realization; (3) compute its transmission zeros, which are the roots of the denominator of $G$; (4) project reconstructed flow onto the recovered memory filters to obtain the weights (Appendix~\ref{p1:app:era}). Without noise this recovers $K=0$, $K=1$ and well-separated two-cohort systems exactly (Appendix~\ref{p1:app:era}); with noise, per-cohort weights become unreliable once rates are close. The reason is information-theoretic, as the rest of this section shows.

\subsection{Local information: the Cram\'er--Rao benchmark}

\begin{proposition}[Cram\'er--Rao benchmark]\label{p1:prop:crb}
Let $\theta$ collect the weights $w_1,\ldots,w_K$ and rates $\eta_1,\ldots,\eta_K$. Suppose the impulse response $h_t(\theta)$, $t\le T$, is observed with i.i.d.\ Gaussian error from $N$ reset probes of size $V_{\rm probe}$ against innovation noise $\sigma_p$, so the error has standard deviation $\sigma_h=\sigma_p/(V_{\rm probe}\sqrt N)$. With Jacobian $J_{tj}=\partial h_t/\partial\theta_j$ and $F_{\rm unit}=J^\top J$, any regular unbiased estimator has ${\rm Var}(\hat w_k)\ge\sigma_h^2[F_{\rm unit}^{-1}]_{kk}$, and a 95\% interval of half-width ${\rm tol}$ requires
\begin{equation}\label{p1:eq:Ncrb}
N^*\ge\Bigl(\frac{1.96}{\rm tol}\Bigr)^2\Bigl(\frac{\sigma_p}{V_{\rm probe}}\Bigr)^2[F_{\rm unit}^{-1}]_{kk}.
\end{equation}
For $K=2$ with spacing $\Delta\eta$, $[F_{\rm unit}^{-1}]_{kk}=\Theta((\Delta\eta)^{-6})$, so $N^*=\Theta((\Delta\eta)^{-6})$.
\end{proposition}

\paragraph{In words.} This gives the number of test trades any well-behaved estimator needs to pin down one group's trading weight to a stated precision. Each extra group with a nearby memory raises the number sharply, as Table~\ref{p1:tab:crb} shows.

\begin{table}[t]\centering\small
\begin{threeparttable}
\caption{Cram\'er--Rao sample sizes for multi-cohort weight recovery ($T=120$ taps)}\label{p1:tab:crb}
\begin{tabular}{lcrccc}
\toprule
Scenario & $K$ & $N^*$ probes & $\sigma_{\rm crit}$ & ML success ($f=1$) & RMSE/CRB\\
\midrule
$\eta=0.30$ & 1 & 0.1 & $6.05\times10^{-3}$ & 94\% & 1.02\\
$\Delta\eta=0.30$ & 2 & 47.8 & $2.31\times10^{-4}$ & 94\% & 1.01\\
$\Delta\eta=0.20$ & 2 & 597 & $6.55\times10^{-5}$ & 96\% & 0.98\\
$\Delta\eta=0.10$ ($\delta=0.05$ pair) & 2 & 40{,}374 & $7.96\times10^{-6}$ & 93\% & 1.03\\
$\Delta\eta=0.05$ & 2 & 2{,}620{,}320 & $9.88\times10^{-7}$ & 95\% & 1.03\\
spacing 0.10 & 3 & 17{,}349{,}791 & $3.84\times10^{-7}$ & 98\% & 1.22\\
spacing 0.20 & 3 & 27{,}038 & $9.73\times10^{-6}$ & 98\% & 0.88\\
\bottomrule
\end{tabular}
\begin{tablenotes}\footnotesize\item Tolerance $\pm0.01$ at 95\%. $\sigma_{\rm crit}$ is the single-probe noise at which one probe suffices; operational single-probe noise is $1.6\times10^{-3}$. ML columns: least squares over all weights and rates from six generic starts (never the true values), 200 replications per row at $\sigma_{\rm crit}$; success means every weight within $\pm0.01$. Fitted log--log slope of $N^*$ in $\Delta\eta$: $-6.04$. \end{tablenotes}
\end{threeparttable}
\end{table}

At the critical noise level maximum likelihood attains the bound: its RMSE for the hardest weight is between 0.88 and 1.22 times the Cram\'er--Rao value and every weight is within tolerance in 93\% to 98\% of replications (Table~\ref{p1:tab:crb}), so the critical noise levels, not estimator failure, are what make close cohorts expensive. Under ordinary surveillance budgets ($N\le50$) unbiased estimators cannot resolve close weights.

\subsection{A global two-point bound for any estimator}

Proposition~\ref{p1:prop:crb} applies to regular unbiased estimators. Could biased, ridge, shrinkage or Bayesian estimators escape the $(\Delta\eta)^{-6}$ scaling? No.

\begin{theorem}[Global two-point bound and matching upper bound, $K=2$]\label{p1:thm:twocohort}
In this theorem $w_k$ denotes cohort $k$'s trading intensity $q_k$, not the merged weight $q_k\eta_k$ of Section~\ref{p1:sec:model}, and $\Delta\eta$ is the actual spacing of the two rates. Fix total intensity $q=0.20$, mean rate $\bar\eta=0.40$, target separation $\tau=0.01$, and $\Delta\eta\in(0,8/11)$. Let $p_\mp=\tfrac12\mp\tau/q$ and
\begin{equation}\label{p1:eq:pair}
\theta_0=(qp_-,qp_+,\bar\eta-p_+\Delta\eta,\bar\eta+p_-\Delta\eta),\qquad\theta_1=(qp_+,qp_-,\bar\eta-p_-\Delta\eta,\bar\eta+p_+\Delta\eta).
\end{equation}
\begin{enumerate}[label=(\roman*)]
\item Label cohorts in increasing order of rate, so that weight vectors are compared rate by rate (as unordered sets the two weight vectors coincide). Both economies have spacing $\Delta\eta$, weight separation $\|w(\theta_1)-w(\theta_0)\|_\infty=2\tau$ in this labelling, and equal rate moments $M_m=\sum_kw_k\eta_k^m$ for $m\in\{0,1,2\}$, while
\begin{equation}\label{p1:eq:dM3}
M_3(\theta_1)-M_3(\theta_0)=4\tau p_-p_+(\Delta\eta)^3 .
\end{equation}
Hence $\|h(\theta_1)-h(\theta_0)\|_2=\Theta(\tau(\Delta\eta)^3)$.
\item For any measurable estimator (including biased, ridge, James--Stein or MAP) that reports weights in the same rate order, the supremum being over the class of two-cohort economies with total weight $q$, mean rate $\bar\eta$ and spacing $\Delta\eta$,
\begin{equation}\label{p1:eq:lecam2}
\sup_\theta P_\theta(\|\hat w-w(\theta)\|_\infty\ge\tau)\;\ge\;1-\Phi\!\Bigl(\frac{\sqrt N\|h(\theta_1)-h(\theta_0)\|_2}{2\sigma_h}\Bigr),
\end{equation}
so failure probability at most $0.05$ requires $N\ge4\sigma_h^2z_{0.95}^2/\|h(\theta_1)-h(\theta_0)\|_2^2=\Omega((\Delta\eta)^{-6})$. The same bound holds without any labelling convention for the permutation-invariant loss
\[
d_\infty(\hat\theta,\theta)=\min_{\pi\in S_2}\max_{k}\bigl(|\hat w_{\pi(k)}-w_k|+|\hat\eta_{\pi(k)}-\eta_k|\bigr)
\]
in place of $\|\hat w-w(\theta)\|_\infty$, whenever $\Delta\eta\ge2\tau q/(q+2\tau)$ ($\approx0.018$ at the stated values; for smaller spacings the unordered weights coincide and the relevant permutation-invariant statement is the $W_1$ bound of Theorem~\ref{p1:thm:kcohort}, at accuracy proportional to $\Delta\eta$), because then $d_\infty(\theta_0,\theta_1)=\min\{2\tau(1+\Delta\eta/q),\,(1+2\tau/q)\Delta\eta\}\ge2\tau$.
\item Conversely, the matching upper bound $N=O((\Delta\eta)^{-6})$ on rich separated two-cohort classes is Theorem~\ref{p1:prop:esprit} with $K=2$. On the separated class with weights and rates bounded away from the boundary, Toeplitz deconvolution followed by a Prony--Hankel root step has Hankel determinant $g_1g_3-g_2^2=w_1w_2\eta_1\eta_2r_1r_2(\Delta\eta)^2>0$ for the deconvolved kernel $g_m=\sum_iw_i\eta_ir_i^m$, and perturbation analysis gives the same order.
\end{enumerate}
\end{theorem}

\paragraph{In words.} For two unknown groups, no method can tell the two candidate markets apart reliably with fewer than a number of experiments of order $(\Delta\eta)^{-6}$, and a practical method (Prony's method, Theorem~\ref{p1:prop:esprit}) achieves this order. Halving the spacing multiplies the data needed by 64.

We recomputed the two-point quantities exactly (Table~\ref{p1:tab:lecam}): the moment gap equals $4\tau p_-p_+(\Delta\eta)^3$ to machine precision, and the Le Cam requirement at $\Delta\eta=0.10$ is $5{,}927.6$ resets. The worst-case failure rates of four re-implemented estimators all lie above the lower bound (Table~\ref{p1:tab:lecam}).

\begin{table}[t]\centering\small
\begin{threeparttable}
\caption{Two-point sample complexity for two-cohort weight recovery}\label{p1:tab:lecam}
\begin{tabular}{lcccc}
\toprule
$\Delta\eta$ & $\Delta M_3$ & $\|\Delta h\|_2$ & $N^*_{\rm LeCam}$ (95\%) & $N^*_{\rm CRB}$\\
\midrule
0.30 & $2.67\times10^{-4}$ & $1.98\times10^{-3}$ & 7.1 & 47.8\\
0.20 & $7.92\times10^{-5}$ & $5.60\times10^{-4}$ & 88.2 & 597.0\\
0.14 & $2.72\times10^{-5}$ & $1.89\times10^{-4}$ & 775.3 & 5{,}267.9\\
0.10 & $9.90\times10^{-6}$ & $6.84\times10^{-5}$ & 5{,}927.6 & 40{,}374.2\\
0.08 & $5.07\times10^{-6}$ & $3.49\times10^{-5}$ & 22{,}740.5 & 155{,}048.9\\
0.06 & $2.14\times10^{-6}$ & $1.47\times10^{-5}$ & 128{,}333.0 & 875{,}732.3\\
0.05 & $1.24\times10^{-6}$ & $8.50\times10^{-6}$ & 383{,}859.1 & 2{,}620{,}321.2\\
0.04 & $6.34\times10^{-7}$ & $4.35\times10^{-6}$ & 1{,}466{,}361.5 & 10{,}012{,}662.4\\
\midrule
\multicolumn{5}{p{0.85\textwidth}}{\emph{Worst-case failure at $\Delta\eta=0.10$, $N=50$, 1,000 trials per hypothesis}: bound 44.1\%; ML 94.5\%, ridge 47.5\%, positive-part shrinkage 100\%, heavy prior 73.9\%.}\\
\bottomrule
\end{tabular}
\begin{tablenotes}\footnotesize\item $\tau=0.01$, $\sigma_h=1.6\times10^{-3}$, $a=0.5$, $\kappa=0.8$, 121 taps. Recomputed columns from \texttt{recompute\_theory.py}. Estimators re-implemented in \texttt{recompute\_p1\_estimators.py}: ML is least squares from four starts; ridge and heavy prior are posterior modes with prior standard deviations $(0.05,0.10)$ and $(0.005,0.02)$ for (weights, rates) around the symmetric centre of the two hypotheses; positive-part shrinkage pulls the ML weights toward their mean. Failure means some weight is off by at least $\tau$.\end{tablenotes}
\end{threeparttable}
\end{table}

\begin{remark}[The hard directions carry no risk]\label{p1:rem:riskneutral}
The two economies of Theorem~\ref{p1:thm:twocohort} have the same total weight $\sum_iq_i=q$, so their summed price response $\sum_th_t=\kappa/(1-a-\kappa\sum_iq_i)$ is identical, $2.353$ at the calibration of Table~\ref{p1:tab:lecam}: the directions in which prices are least informative about the cohorts are directions in which cumulative price risk does not move. The summed response is a linear functional of the response and is estimated at the parametric rate whatever the spacing of the memories; Chapter~\ref{ch:p2} turns this into a certification theorem, a minimax-optimal pre-trade guard and a decision value of disclosure (Chapter~\ref{ch:p2}). Identification and safety ask different questions of the same data.
\end{remark}

\subsection{Many unknown cohorts}

\begin{theorem}[$K$-cohort two-point lower bound]\label{p1:thm:kcohort}
Fix $K\ge2$, a center $r_c\in(0,1)$, total merged weight $W>0$ and $0<\Delta\eta<\min(r_c,1-r_c)$, so that all rates lie in $(0,1)$. There exist two $K$-cohort economies with rates in $r_c+\Delta\eta[-1,1]$, positive weights, and merged measures whose moments of order $0,\ldots,2K-2$ coincide while the moment of order $2K-1$ differs by a constant times $(\Delta\eta)^{2K-1}$. Their impulse responses differ by $\|h^A-h^B\|_2\asymp(\Delta\eta)^{2K-1}$, so recovering the weights from $N$ independent noisy reset experiments with worst-case error probability below $1/4$ requires
\[
N\;\ge\;c\,\sigma^2(\Delta\eta)^{-(4K-2)} .
\]
The same requirement holds for the permutation-invariant loss $W_1(\hat M,M)$ between estimated and true merged measures at accuracy $c_KW\Delta\eta$, where $2c_K>0$ is the Wasserstein distance between the two normalized quadrature rules, so no relabelling of atoms escapes it.
\end{theorem}

The construction uses the $K$-point Gauss and Gauss--Radau rules of the same weight function: both are positive, have $K$ atoms and integrate polynomials of degree $2K-2$ exactly. For $K=2$ it reproduces the order of Theorem~\ref{p1:thm:twocohort}. Table~\ref{p1:tab:radau} lists the two rules: their low moments agree to machine precision and their moment of order $2K-1$ differs by $0.222$, $0.0533$ and $0.0131$ for $K=2,3,4$. The shrinking gap explains why constants deteriorate with $K$ on top of the exponent.

\paragraph{What counts as an experiment.} In Theorem~\ref{p1:thm:kcohort} an experiment is a unit impulse from the reset state whose response is observed with noise of variance $\sigma^2$ per tap. The order in $\Delta\eta$ does not depend on that choice. For any probe $v$ the mean difference between the two economies is at most $\sup_\omega|\Delta\hat H(e^{i\omega})|\,\|v\|_2$, and $\sup_\omega|\Delta\hat H|=O((\Delta\eta)^{2K-1})$ because $\Delta\hat H=z^{-1}\hat H_A\hat H_B\,\Delta G$ with both responses bounded on the circle. Experiments with probe energy $E$ each therefore need $NE\gtrsim\sigma^2(\Delta\eta)^{-(4K-2)}$: a larger probe buys the same information with fewer experiments, and Section~\ref{p1:sec:account} prices the capital that a larger probe requires. The bound does not cover information obtained without trading against the groups, such as class records (Section~\ref{p1:sec:disclosure}).

\begin{table}[t]\centering\small
\begin{threeparttable}
\caption{Gauss and Gauss--Radau rules on $[-1,1]$ used in Theorem~\ref{p1:thm:kcohort}}\label{p1:tab:radau}
\begin{tabular}{cp{0.42\textwidth}cc}
\toprule
$K$ & Gauss nodes / Gauss--Radau nodes & Max gap, orders $<2K-1$ & Gap at order $2K-1$\\
\midrule
2 & $\pm0.5774$ \;/\; $-1,\ 0.3333$ & $6\times10^{-17}$ & 0.2222\\
3 & $0,\ \pm0.7746$ \;/\; $-1,\ -0.2899,\ 0.6899$ & $2\times10^{-16}$ & 0.0533\\
4 & $\pm0.3400,\ \pm0.8611$ \;/\; $-1,\ -0.5753,\ 0.1811,\ 0.8228$ & $4\times10^{-16}$ & 0.0131\\
\bottomrule
\end{tabular}
\begin{tablenotes}\footnotesize\item All Gauss--Radau weights are positive. \end{tablenotes}
\end{threeparttable}
\end{table} Table~\ref{p1:tab:kcohort} and Figure~\ref{p1:fig:kcohort} confirm the order numerically.

\begin{table}[t]\centering
\begin{threeparttable}
\caption{Numerical check of Theorem~\ref{p1:thm:kcohort}}\label{p1:tab:kcohort}
\begin{tabular}{lccc}
\toprule
 & $K=2$ & $K=3$ & $K=4$\\
\midrule
Theoretical exponent of $N^*$ & $-6$ & $-10$ & $-14$\\
Fitted log--log slope (small $\Delta\eta$) & $-5.96$ & $-9.96$ & $-14.01$\\
\bottomrule
\end{tabular}
\begin{tablenotes}\footnotesize
\item Le Cam bound $N^*\propto\sigma^2/\|h^A-h^B\|_2^2$ with $\sigma=1$, $a=0.5$, $\kappa=1$, total merged weight $0.10$, rates centered at $0.70$, responses over 400 lags. \end{tablenotes}
\end{threeparttable}
\end{table}

\begin{figure}[t]\centering
\includegraphics[width=0.62\textwidth]{p1figs/fig_kcohort.pdf}
\caption{Two-point lower bound on the number of reset experiments needed to separate $K$-cohort economies, as a function of the rate spacing.}\label{p1:fig:kcohort}
\end{figure}

\begin{theorem}[Matching upper bound for every $K$]\label{p1:prop:esprit}
Fix $K\ge2$, a compact set $\mathcal S\subset\{s\in[-1,1]^K:s_1<\cdots<s_K\}$ of ordered rescaled rates with $\min_{i\ne j}|s_i-s_j|\ge\gamma>0$, and weights in $[\omega_{\min},\omega_{\max}]$ with $\omega_{\min}>0$. Consider the class of $K$-cohort economies with rates $r_c+\Delta\eta\,s$, $s\in\mathcal S$, and $N$ independent reset experiments whose responses carry Gaussian noise with variance $\sigma^2$. There are constants $c_0,C>0$ depending only on $(K,\mathcal S,\gamma,\omega_{\min},\omega_{\max},a,\kappa,r_c)$ such that, whenever $\epsilon_N\equiv\sigma N^{-1/2}(\Delta\eta)^{-(2K-1)}\le c_0/\sqrt{\log(2K/\beta)}$, Prony's method applied to the first $2K$ estimated kernel values, with its output atoms listed in increasing order of rate, recovers the weights with
\[
\max_i|\hat w_i-w_i|\;\le\;C\sqrt{\log(2K/\beta)}\;\epsilon_N
\]
with probability at least $1-\beta$. Hence $N=C'\sigma^2(\Delta\eta)^{-(4K-2)}$ experiments suffice for any fixed accuracy and confidence, On any such class that contains the Gauss and Gauss--Radau pair of Theorem~\ref{p1:thm:kcohort} (a rich class), together with that theorem the minimax sample requirement is $N^*\asymp\sigma^2(\Delta\eta)^{-(4K-2)}$. On smaller classes the upper bound need not be tight: if the rates are known, for example, the weights need only order $\sigma^2(\Delta\eta)^{-2}$ experiments.
\end{theorem}

\begin{proof}
The kernel values $k_m=\sum_iw_ir_i^m$ are the power moments of the merged measure. The closed-loop recursion $h_t=ah_{t-1}+\kappa\sum_mk_mh_{t-1-m}$ (plus $\kappa$ at $t=0$) can be solved for $k_0,\ldots,k_{2K-1}$ from $h_0,\ldots,h_{2K}$ with $h_0=\kappa\ne0$; the map is smooth with Lipschitz constant independent of $\Delta\eta$, so the estimated kernel values have error $O_P(\sigma N^{-1/2})$. Centring at $r_c$ is a fixed triangular linear map, $\mu_m=\sum_{j\le m}\binom mj(-r_c)^{m-j}k_j$, and $\mu_m=(\Delta\eta)^m t_m$ with $t_m=\sum_iw_is_i^m$. The map $F:(s,\omega)\mapsto(t_0,\ldots,t_{2K-1})$ does not depend on $\Delta\eta$; its Jacobian is a confluent Vandermonde matrix with determinant $\prod_i\omega_i\prod_{i<j}(s_i-s_j)^4$, bounded away from zero on the compact class. By the quantitative inverse function theorem there are $\rho_0>0$ and $L<\infty$, uniform over the class, such that $F^{-1}$ exists and is $L$-Lipschitz on the ball of radius $\rho_0$ around every point of $F(\text{class})$. An error $e$ in $\mu_m$ is an error $e(\Delta\eta)^{-m}$ in $t_m$, largest at $m=2K-1$. The observed response errors are Gaussian; after nonlinear deconvolution the kernel errors need not be exactly Gaussian. On the compact family, a local Lipschitz bound for that smooth map and Gaussian concentration for the response coordinates give, for sufficiently large $N$, $\|\hat t-t\|_\infty\le C_1\sqrt{\log(2K/\beta)}\,\epsilon_N$ with probability at least $1-\beta$. When this is below $\rho_0$, $\hat t$ lies in the ball where $F^{-1}$ is $L$-Lipschitz. The moment map is invariant to relabelling the atoms, so it is injective only on ordered nodes; on the class of ordered, $\gamma$-separated nodes it is injective, and the inverse function theorem applies there. Prony's method computes the algebraic inverse of the moment map up to the order of the atoms, so, listing its output in increasing order of rate (which the error bound below keeps unambiguous once $C\sqrt{\log(2K/\beta)}\,\epsilon_N<\gamma/2$), its output is $F^{-1}(\hat t)$ and the stated bound follows with $C=LC_1$ and $c_0=\min(\rho_0,\gamma/(2L))/C_1$, which also keeps $C\sqrt{\log(2K/\beta)}\,\epsilon_N<\gamma/2$. No separate root-perturbation bound is needed: on the ball where $F^{-1}$ is $L$-Lipschitz, Prony's roots and weights are the coordinates of $F^{-1}(\hat t)$, so the conditioning of the root step is part of $L$ (for an explicit step-by-step analysis of Prony's method with clustered exponents see \citealp{p1:katz2023}). The matching statement requires the class to contain the Gauss and Gauss--Radau pair of Theorem~\ref{p1:thm:kcohort} (ordered nodes with separation at least the smaller of their two minimum gaps, and weights within $[\omega_{\min},\omega_{\max}]$ after scaling to total weight $W$); for any such class the orders match, with constants that depend on the class.
\end{proof}

\paragraph{In words.} With $K$ unknown groups whose memories lie within $\Delta\eta$ of each other, the number of experiments needed grows like $(\Delta\eta)^{-(4K-2)}$. For two groups, halving the spread multiplies the data needed by 64; for three groups, by 1{,}024. Theorem~\ref{p1:thm:kcohort} shows that no method does better, and this theorem shows that Prony's method achieves it.


Table~\ref{p1:tab:kupper} confirms the order numerically. For numerical stability the estimator works in centred and rescaled coordinates (the estimator is the same; only its floating-point evaluation changes). The argument parallels the stability analysis of super-resolution for clustered complex nodes \citep{p1:batenkov2021}, and \citet{p1:katz2023} show that Prony's method attains the min-max bound for clustered exponents; here the nodes are real and are first extracted from the closed loop, and it reproduces the ESPRIT slope of $-6.15$ reported for $K=2$ in the earlier version.

\begin{table}[t]\centering\small
\begin{threeparttable}
\caption{Sample requirement of Prony's method against the lower-bound order}\label{p1:tab:kupper}
\begin{tabular}{lcccc}
\toprule
 & $K=2$ & $K=3$ & $K=4$ & $K=5$\\
\midrule
Lower-bound order $-(4K-2)$ & $-6$ & $-10$ & $-14$ & $-18$\\
Fitted slope of $N^*$ in $\Delta\eta$ (Prony) & $-5.94$ & $-10.00$ & $-13.93$ & $-18.34$\\
\bottomrule
\end{tabular}
\begin{tablenotes}\footnotesize\item Rates $0.7+\Delta\eta\,s$ with $s$ the $K$-point Gauss nodes, total weight 0.1, $\Delta\eta\in[0.03,0.2]$; $N^*$ is the number of experiments for a weight RMSE of 10\% of the smallest weight, from 400 simulated estimates per point. \end{tablenotes}
\end{threeparttable}
\end{table}

\subsection{Near-colliding rates: a uniform rate for the merged measure}\label{p1:sec:uniform}

Theorem~\ref{p1:prop:esprit} needs the rates to stay $\gamma$-separated and every weight to stay away from zero. Neither restriction can be dropped for weight recovery: two atoms of weights $w$ and $W-w$ at rates $r$ and $r+\Delta$ and two of weights $W-w$ and $w$ at the same rates have merged measures that converge to the same measure as $\Delta\to0$, so no estimator recovers the individual weights uniformly over near-colliding economies. What can be recovered uniformly is the merged measure itself, in the permutation-invariant Wasserstein distance.

\begin{theorem}[Uniform rate without separation]\label{p1:thm:uniformw1}
Fix $K\ge1$, a memory bound $\bar r<1$ and $0<m_{\min}\le B$. Let $\mathcal M_K$ be the merged measures with at most $K$ atoms in $[0,\bar r]$ and total mass $k_0\in[m_{\min},B]$, with no separation condition and no lower bound on individual weights, and write $\bar M=M/k_0$ for the normalized memory distribution. From $N$ reset experiments as in Theorem~\ref{p1:prop:esprit}, let $\hat M$ be any element of $\mathcal M_K$ whose first $2K$ kernel values are closest in the sup norm to the estimated ones. There is $C$ depending only on $(K,\bar r,a,\kappa,B)$ such that, with probability at least $1-\beta$,
\[
W_1(\hat{\bar M},\bar M)\;\le\;C\Bigl(\frac{\sigma\sqrt{\log(4K/\beta)}}{m_{\min}\sqrt N}\Bigr)^{1/(2K-1)}\qquad\text{uniformly over }\mathcal M_K .
\]
Hence $N=O(\sigma^2m_{\min}^{-2}\epsilon^{-(4K-2)})$ experiments give accuracy $\epsilon$, and with the Wasserstein statement of Theorem~\ref{p1:thm:kcohort} the minimax rate for the merged memory distribution is $N^{-1/(4K-2)}$, uniformly over near-colliding economies.
\end{theorem}

\paragraph{In words.} Even with no lower limit on how close the groups' memories are or how small a group is, the combined memory distribution can still be estimated, and no faster rate is possible. What is lost when groups nearly coincide is the ability to tell them apart, not the ability to describe the combined picture.

\begin{proof}
As in the proof of Theorem~\ref{p1:prop:esprit}, the kernel values $k_0,\ldots,k_{2K-1}$ are a Lipschitz function of the first $2K+1$ responses on the bounded set of responses the class produces, so a union bound gives $\max_{t<2K}|\bar k_t-k_t|\le e=C_1\sigma\sqrt{\log(4K/\beta)/N}$ with probability at least $1-\beta$ (the Lipschitz constant is for a bounded neighbourhood of the class; when the noise is too large for that, the right-hand side of the theorem exceeds $\bar r\ge W_1$ once $C$ is enlarged, so the bound holds trivially). The class is compact and the moment map continuous, so a closest element $\hat M$ exists, and $\max_t|k_t(\hat M)-k_t(M)|\le2e$ by the triangle inequality. Since all atoms lie in $[0,1]$, $|k_t|\le k_0$, and the normalized moments $m_t=k_t/k_0$ satisfy $|m_t(\hat M)-m_t(M)|\le2e/k_0+(|\hat k_t|/\hat k_0)(2e/k_0)\le4e/m_{\min}$. The affine map $r\mapsto2r/\bar r-1$ carries both normalized measures to distributions on $[-1,1]$ with at most $K$ atoms, multiplies $W_1$ by $2/\bar r$ and the differences of moments of order below $2K$ by at most $(1+2/\bar r)^{2K-1}$. The moment comparison inequality for $K$-atomic distributions on $[-1,1]$ \citep{p1:wu2020}, $W_1\le c_K\delta^{1/(2K-1)}$ when the first $2K-1$ moments differ by at most $\delta$, then gives the bound with $C=c_K(\bar r/2)(1+2/\bar r)(4C_1)^{1/(2K-1)}$.
\end{proof}

The theorem separates two questions that Theorem~\ref{p1:prop:esprit} answers together. Which memory distribution do the participants have? That has the same polynomial cost whether or not their rates collide. Which participant has which weight? That becomes arbitrarily expensive as rates approach each other (Theorem~\ref{p1:thm:kcohort}) and impossible when they coincide (Theorem~\ref{p1:thm:merged}). A supervisor who needs the first can have it at the rate of Theorem~\ref{p1:thm:kcohort}; one who needs the second has to be told the number of cohorts and a separation, or has to obtain tagged data (Section~\ref{p1:sec:disclosure}).

\subsection{General memory: the minimax rate is logarithmic}\label{p1:sec:general}

Exponential memory is a modelling choice. Suppose instead that the feedback kernel is completely monotone, $k_m=\int_{[0,\rho]}r^m\,d\mu(r)$ with $\rho<1$, where $\mu$ is any nonnegative measure of total mass $W$: a finite or infinite mixture of exponential memories with merged weight $W$. Write $\mathcal M(\rho,W)$ for this class, assume the closed loop is uniformly stable, $s\equiv a+\kappa W/(1-\rho)<1$, and measure accuracy by
\[
d(\mu,\mu')=\sup\Bigl\{\textstyle\int f\,d(\mu-\mu'):\ f\ \text{1-Lipschitz on }[0,\rho],\ f(0)=0\Bigr\},
\]
which equals the Wasserstein-1 distance between measures of equal mass and also compares estimates whose mass is not exactly $W$. Let $E_L(|x|)$ denote the error of the best uniform approximation of $|x|$ on $[-1,1]$ by polynomials of degree $L$, and let $c_*=\inf_{L\ge1}L\,E_L(|x|)$. \citet{p1:bernstein1914} showed that $L\,E_L(|x|)$ converges to a positive constant ($0.2802$; \citealp{p1:varga1985}), so $c_*>0$; numerically $c_*=1/4$, attained at $L=2$ (Table~\ref{p1:tab:general}).

\begin{lemma}[Moment-matched cohort pairs]\label{p1:lem:altpair}
For every $n\ge2$ there are two probability measures on $[0,\rho]$, each with exactly $n$ atoms, whose moments of orders $0,\ldots,2n-2$ coincide and whose Wasserstein-1 distance is at least $\rho\,E_{2n-2}(|x|)\ge c_*\rho/(2n-2)$.
\end{lemma}

The atoms are the points where the best approximation of $|r-\rho/2|$ by polynomials of degree $2n-2$ attains its maximal error with alternating signs, and the weights are those of the divided difference on these points, which annihilates every polynomial of degree $2n-2$ (Appendix~\ref{p1:app:proofs}). Each measure is an $n$-atom merged measure in the class; one atom may sit at rate $0$ (a cohort with $\eta=1$), which the class allows.

\begin{theorem}[Minimax rate for general memory]\label{p1:thm:general}
Let $N$ independent reset experiments observe the impulse response with Gaussian noise of variance $\sigma^2$ per tap, and let $B(\rho)=(1+\sqrt2)(1+2/\rho)$ and $C_h=\kappa^2/(1-s)^2$.
\begin{enumerate}[label=(\roman*)]
\item \emph{Lower bound.} Let $\Lambda_N=\log(N/\sigma^2)+\log\bigl(2C_h^2W^2/((1-\rho)\rho)^2\bigr)$. Whenever $\Lambda_N>4\log(1/\rho)$, every estimator $\hat\mu$ satisfies
\[
\sup_{\mu\in\mathcal M(\rho,W)}P_\mu\Bigl(d(\hat\mu,\mu)\ge\frac{c_*\,\rho W\log(1/\rho)}{\Lambda_N}\Bigr)\;\ge\;\frac1{4e},
\]
and the bound already holds over $n$-cohort economies, for a finite $n$ that grows with $N$.
\item \emph{Upper bound.} With $a$ and $\kappa$ known, invert the averaged response into kernel values $\hat k_0,\ldots,\hat k_L$, with $L+1=\lfloor\log(N/\sigma^2)/(4\log B(\rho))\rfloor$, and let $\hat\mu$ be any measure on $[0,\rho]$ minimizing $\max_{j\le L}|\int r^jd\hat\mu-\hat k_j|$ (a linear program). For every $\beta\in(0,1)$ there is $N_0$ depending only on $(\rho,W,a,\kappa,\beta)$ such that for $N/\sigma^2\ge N_0$,
\[
\sup_{\mu\in\mathcal M(\rho,W)}P_\mu\Bigl(d(\hat\mu,\mu)>\frac{3\pi\rho W\log B(\rho)}{\log(N/\sigma^2)}\Bigr)\;\le\;\beta .
\]
\end{enumerate}
Hence the minimax accuracy is of exact order $1/\log(N/\sigma^2)$, and reaching accuracy $w$ requires and suffices $N/\sigma^2=\exp(\Theta(1/w))$.
\end{theorem}

\begin{proof}[Proof sketch]
(i) Put the measures of Lemma~\ref{p1:lem:altpair}, scaled to mass $W$, into the model. Their difference annihilates polynomials of degree $2n-2$, so the generating function of the kernel difference is $\int(1-rz)^{-1}d(\mu_A-\mu_B)=z^{2n-1}\int r^{2n-1}(1-rz)^{-1}d(\mu_A-\mu_B)$; on $|z|=1$ it is at most $2W\rho^{2n-1}/(1-\rho)$, and Parseval bounds $\|k^A-k^B\|_2$ by the same quantity. The response difference is $H_AH_B$ times the shifted kernel difference, so $\|h^A-h^B\|_2\le C_h\|k^A-k^B\|_2$. Choosing the smallest $n$ for which the Kullback--Leibler divergence of the $N$-sample experiments is at most one, the Bretagnolle--Huber form of Le Cam's lemma gives the bound with half the Wasserstein separation of the lemma. (ii) For a 1-Lipschitz $f$ with $f(0)=0$, Jackson's theorem gives a polynomial $q$ of degree $L$ with $\sup|f-q|\le\pi\rho/(4(L+1))$, and expanding $q$ in Chebyshev polynomials of $2r/\rho-1$ bounds the sum of its monomial coefficients by $\rho B(\rho)^{L+1}$. Hence $d(\hat\mu,\mu)\le\pi\rho(W+\delta)/(2(L+1))+2\delta\rho B(\rho)^{L+1}$, where $\delta$ is the largest kernel error. The kernel is recovered from the response by a series reciprocal, $K(z)=(1-az-\kappa/H(z))/(\kappa z)$, whose error is at most $2((1+s)/\kappa)^2(L+2)$ times the largest response error, so $\delta=O_P(\sigma N^{-1/2}L\sqrt{\log L})$. The choice of $L$ makes $\delta B(\rho)^{L+1}=O_P((\sigma^2/N)^{1/4}L\sqrt{\log L})$, negligible against the first term. Full proofs are in Appendix~\ref{p1:app:proofs}.
\end{proof}

\paragraph{In words.} If trader memory is a general mixture of exponentials rather than a few clean ones, accuracy improves only with the logarithm of the data: to halve the error, the data must be squared. The few-groups model of the previous results is therefore the optimistic case.


Theorem~\ref{p1:thm:general} replaces a polynomial barrier with a logarithmic one. With exponential memory and $K$ known cohorts, halving the error multiplies the required number of experiments by a constant; with general memory, halving the error squares it. The lower bound holds already for finite mixtures, so knowing that the memory is a finite mixture of exponentials does not help unless the number of cohorts is also known (Theorem~\ref{p1:thm:uniformw1}). The constants of the two bounds differ, by a factor that grows as $\rho\to1$, so the dependence on how long memory lasts is not pinned down. The source of the gap is visible in the proof: as $\rho\to1$ the moment-matched pairs separate their kernels more slowly ($\rho^{2n-1}$ decays slowly), which weakens the lower bound through $\log(1/\rho)$, while the upper bound pays $\log B(\rho)$, which stays bounded away from zero, for the conditioning of polynomial approximation on $[0,\rho]$. Logarithmic rates of this kind are familiar from deconvolution with supersmooth errors \citep{p1:fan1991} and from moment methods for mixing distributions \citep{p1:kong2017,p1:wu2020}; here they arise from the closed loop of a market.

Table~\ref{p1:tab:general} checks the construction and the estimator. The pairs of Lemma~\ref{p1:lem:altpair} have exactly $n$ atoms, agree in all moments below order $2n-1$ to $10^{-61}$, and remain at Wasserstein distance about $0.4/n$, while the distance between their kernels shrinks by a factor of about $0.2$ per atom at $\rho=0.9$. The linear-programming estimator applied to a continuous memory distribution has error times $\log(N/\sigma^2)$ between $0.18$ and $0.28$ for $N$ from $10^4$ to $10^{24}$: a $10^{20}$-fold increase in data halves the error. The exponential-memory model of the previous subsections, where the cost is polynomial in $1/\Delta\eta$, is therefore the optimistic case: relaxing the parametric form makes identification from prices much harder, not easier.

\begin{table}[t]\centering\small
\begin{threeparttable}
\caption{General memory: moment-matched cohort pairs and the moment estimator}\label{p1:tab:general}
\setlength{\tabcolsep}{3.5pt}
\begin{tabular}{lcccccc}
\toprule
\multicolumn{7}{l}{\emph{A. Best approximation of $|x|$}}\\
$L$ & 2 & 4 & 6 & 10 & 20 & 30\\
$L\,E_L(|x|)$ & 0.2500 & 0.2705 & 0.2756 & 0.2785 & 0.2797 & 0.2800\\
\midrule
\multicolumn{7}{l}{\emph{B. Pairs of Lemma~\ref{p1:lem:altpair}}}\\
Atoms per measure $n$ & 2 & 4 & 6 & 8 & 10 & 12\\
$W_1(\mu_A,\mu_B)$, $\rho=0.9$ & 0.2250 & 0.1037 & 0.0669 & 0.0494 & 0.0391 & 0.0324\\
$n\,W_1$, $\rho=0.9$ & 0.450 & 0.415 & 0.402 & 0.395 & 0.391 & 0.388\\
$\|k^A-k^B\|_2$, $\rho=0.9$ & $4.41\times10^{-2}$ & $1.03\times10^{-3}$ & $4.08\times10^{-5}$ & $1.96\times10^{-6}$ & $1.05\times10^{-7}$ & $5.95\times10^{-9}$\\
$\|k^A-k^B\|_2$, $\rho=0.99$ & $7.41\times10^{-2}$ & $5.84\times10^{-3}$ & $1.07\times10^{-3}$ & $2.76\times10^{-4}$ & $8.33\times10^{-5}$ & $2.76\times10^{-5}$\\
\midrule
\multicolumn{7}{l}{\emph{C. Linear-programming estimator, continuous memory distribution, $\rho=0.9$}}\\
$\log_{10}N$ & 4 & 8 & 12 & 16 & 20 & 24\\
Moments used ($L+1$) & 2 & 3 & 5 & 6 & 7 & 8\\
Error $d(\hat\mu,\mu)$ & 0.00878 & 0.00759 & 0.00432 & 0.00445 & 0.00382 & 0.00399\\
Error $\times\log(N/\sigma^2)$ & 0.202 & 0.245 & 0.179 & 0.226 & 0.229 & 0.276\\
\bottomrule
\end{tabular}
\begin{tablenotes}\footnotesize\item A: Remez algorithm; the limit is Bernstein's constant $0.2802$. B: moments of orders $0,\ldots,2n-2$ agree to $10^{-61}$ (60-digit arithmetic); every $W_1$ exceeds $\rho E_{2n-2}(|x|)$. C: true merged measure Beta(2,3) on $[0,0.9]$ with mass $0.1$; kernel values observed with noise $\sigma N^{-1/2}$, $\sigma=10^{-3}$; $L$ from the rule in Theorem~\ref{p1:thm:general}(ii); measure fitted on a 401-point grid; 20 replications. Kind: recomputed (\texttt{general\_minimax.py}).\end{tablenotes}
\end{threeparttable}
\end{table}

\subsection{Passive data}

\begin{theorem}[Passive identification: minimum phase]\label{p1:thm:minphase}
In a stable economy with $0<a<1$, distinct $0<r_i<1$ and nonzero merged weights, the zeros of $H$ are the memory rates $r_i$ together with a structural zero at $z=0$, and $H$ is minimum phase. Degenerate anchor and zero-rate cases require the cancellations described in the merged-equivalence proof. With a single price innovation, the price spectrum identifies $H$ up to scale, hence the memory rates; with $\kappa$ and $a$ known, as in Theorem~\ref{p1:thm:merged}, it identifies the merged measure.
\end{theorem}

The result holds for signed weights (mean-reverting cohorts): in 1,733 stable random economies with weights of both signs, the nonstructural zeros equalled the memory rates and lay inside the unit disk. With a second independent shock in the cohorts' memories, the Jacobian of the price spectrum with respect to six model parameters has full rank, with singular values from $34.51$ down to $0.0698$. Passive noise helps identification rather than hindering it.

\begin{theorem}[Passive cost]\label{p1:prop:passive}
Suppose the price is observed passively as the stationary output $x_t=H_\theta(L)\xi_t$ of i.i.d.\ Gaussian innovations $\xi_t\sim N(0,\sigma^2)$, $\sigma^2>0$ known, without observation noise, for $\theta$ in a compact, uniformly Schur-stable family with rates in $[\eta_{\min},\eta_{\max}]\subset(0,1)$, and let $D(\theta_1\|\theta_0)$ be the Kullback--Leibler divergence rate per observation.
\begin{enumerate}[label=(\roman*)]
\item For a single cohort at $\bar\eta$ against the symmetric pair at $\bar\eta\pm\delta$ with the same total weight, $D=\Theta(\delta^4)$.
\item For the moment-matched two-cohort economies of Theorem~\ref{p1:thm:twocohort}, $D=\Theta(\tau^2(\Delta\eta)^6)$.
\item For every $K\ge2$, for the $K$-cohort Gauss and Gauss--Radau economies of Theorem~\ref{p1:thm:kcohort}, $D=\Theta((\Delta\eta)^{4K-2})$.
\end{enumerate}
Hence any test based on $T$ passive observations with worst-case error at most $\alpha<1/4$ needs $T=\Omega(\delta^{-4})$, $T=\Omega(\tau^{-2}(\Delta\eta)^{-6})$ and $T=\Omega((\Delta\eta)^{-(4K-2)})$ respectively, uniformly as the separation shrinks: passive data cost at least the same polynomial order as active probing for every number of cohorts. For each fixed pair, at fixed type-I error the likelihood-ratio test's type-II error decays at exponential rate $D$ (Stein's lemma), and both errors decay at the Chernoff-information rate, which is of the same order in the separation; so the orders are the right ones; we do not claim that this attainment is uniform as the separation shrinks.
\end{theorem}

\paragraph{In words.} Watching the market without trading does not escape the barrier. A passive observer needs at least the same polynomial order of observations as a prober needs experiments; it pays in time rather than in trading losses.


The proof (Appendix~\ref{p1:app:proofs}) combines the Whittle form of the divergence rate with the resolvent identity $H_1-H_0=z^{-1}H_1H_0(G_1-G_0)$: the two spectra differ by the leading unmatched moment, of order $\delta^2$ for the pair and $\tau(\Delta\eta)^3$ for the moment-matched economies, and the divergence rate is the square of that difference. Numerically, the fitted slopes of the required sample length are $-4.21$ and $-5.93$, against the theoretical $-4$ and $-6$; for part (iii) the fitted slopes of the Whittle divergence rate in $\Delta\eta$ are $5.91$, $9.98$ and $14.05$ for $K=2,3,4$, against $6$, $10$ and $14$ (\texttt{hardening/p1\_checks.py}).

Passive surveillance is possible in principle but needs at least the same polynomial order of observations as active probing needs experiments. Proposition~\ref{p1:prop:passive0} is the exact finite-sample version for the known pair.

\section{Robustness to nonlinear and adaptive trading rules}\label{p1:sec:learners}

The theorems of Sections~\ref{p1:sec:ident}--\ref{p1:sec:mixtures} use the linear exponential-memory recursion~\eqref{p1:eq:cohort}. Real desks cap positions, trade only beyond thresholds, adapt their gains and learn online. A small probe of size $A$ moves the state by $O(A)$, so any rule whose order is a twice-differentiable function of an exponential-memory state is linear to first order and differs from it by $O(A^2)$; the theorems should then hold for small probes and fail for large ones. We check this, and locate where it fails, in simulations of five stylized rules at $a=0.5$, $\kappa=0.8$, total weight $q=0.2$ and mean rate $\bar\eta=0.4$, over 240 periods after a pulse probe.

The rules are: (1) \emph{momentum with a position limit}, the cohort's order a smooth saturating ($\tanh$) function of its memory state $y_{i,t}$, normalized to unit slope at the operating point; (2) a \emph{threshold rule} that trades only when $|y_{i,t}|>\tau=0.05$, in two versions, an exact dead zone and one averaged over background state variation of standard deviation $0.04$; (3) an \emph{adaptive-gain rule} in the form of recursive least squares with forgetting factor $1-\eta_i$, whose gain falls when the regressor is large; (4) an \emph{online neural readout}, a linear readout of a fixed two-layer $\tanh$ feature map with eight hidden units, trained by regularized stochastic gradient steps of size $\eta_i$; and (5) a \emph{tabular value rule} that updates a 21-bin table at rate $\eta_i$, in two versions, kernel-smoothed across bins and a single-bin lookup. These are stylized: each is built so that its small-signal dynamics are an exponential filter with rate $\eta_i$, and they are not calibrated to any production system.

Two checks follow the theory. \emph{Test~A (merged measure):} one cohort with weight $q$ and position limit $1$ against two cohorts with weights $q/2$ at the same rate and limits $0.55$ and $1.65$, so the two economies have the same merged measure; the raw response gap $\|X^{(1)}(A)-X^{(2)}(A)\|_2$ should be $O(A^2)$ (slope $2$ in $\log A$) and the probe-normalised gap $\|X^{(1)}(A)-X^{(2)}(A)\|_2/A$ should be $O(A)$ (slope $1$ in $\log A$) over $A\in[0.002,0.064]$. \emph{Test~B ($\delta^{-4}$ law):} one cohort at $\bar\eta$ against the symmetric pair at $\bar\eta\pm\delta$, $\delta\in[0.015,0.12]$, at probe size $A=0.005$; the normalized response difference $\|\Delta h\|_2$ should scale as $\delta^2$ and the certified capital of Theorem~\ref{p1:thm:spectral} as $\delta^{-4}$. The saturation threshold $A_{\rm sat}$ is the smallest probe at which $\|\Delta h\|_2$ at $\delta=0.05$ departs by more than 10\% from its small-probe value.

To check that these conclusions do not depend on hand-built rules, we pre-registered a second experiment (\texttt{P1\_\allowbreak trained\_\allowbreak learners.json}, SHA-256 \texttt{3e8bf30cbfb6\ldots}) in which online learners are trained from scratch inside the closed-loop market for $T_{\rm burn}=5{,}000$ periods of flow and noise innovations; their learned parameters are saved and reloaded into the probing simulator without rescaling. Four learning rules are used: (i) recursive least squares with forgetting factor $\lambda=0.995$, learning its return-forecast slope; (ii) a two-layer $\tanh$ network with eight hidden units whose readout is trained by regularized online stochastic gradient descent on trading profit; (iii) a scalar policy-gradient rule trained on trading profit; and (iv) a tabular Q-learner with 21 states and 21 actions, TD(0) Bellman backups ($\gamma=0.90$) and Boltzmann action choice. Each learner acts on an exponential-memory state $y_{i,t}$ with rate $\eta_i$: training sets its response gain and curvature, not its memory rate, so the experiment asks whether learned policies preserve the merged-measure and $\delta^{-4}$ predictions, not whether learning produces exponential memory. The specification fixed four predictions: (a) a single exponential fits each trained learner's small-probe impulse response with $R^2\ge0.90$; (b) two trained populations with the same merged measure differ by at most $1\%$ in relative $L_2$ response at $A=0.005$; (c) the certified capital to separate trained populations whose memory rates differ by $2\delta$ scales with log--log slope in $[-4.15,-3.85]$; and (d) we report where equivalence and the $\delta^{-4}$ law fail when learning continues during the probe. A first implementation of this experiment did not pass the trained parameters to the probing simulator; the results below come from the corrected implementation (\texttt{validate\_\allowbreak trained\_\allowbreak learners\_\allowbreak hardened.py}), which reloads the saved parameters.

\begin{table}[t]\centering\footnotesize
\begin{threeparttable}
\caption{Merged-measure equivalence and the $\delta^{-4}$ law under stylized nonlinear rules and closed-loop trained online learners}\label{p1:tab:learners}
\setlength{\tabcolsep}{1.5pt}
\begin{tabular}{@{}lcccccc@{}}
\toprule
Rule & Test A raw & Test A norm.\ & Test B slope & Test B slope & Capital at & $A_{\rm sat}$ /\\
 & ($O(A^2)$) & ($O(A)$) & ($\|\Delta h\|_2$) & (capital) & $\delta=0.05$ & $R^2_{\rm exp}$\\
\midrule
\multicolumn{7}{@{}l@{}}{\emph{Panel A. Stylized nonlinear and adaptive rules (fixed operating point)}}\\
Linear benchmark & --- & --- & $2.007$ & $-4.001$ & $61.42$ & ---\\
1. Momentum ($\tanh$ limit) & $2.015$ & $1.015$ & $2.007$ & $-4.001$ & $61.47$ & $0.834$\\
2. Threshold rule, averaged & $2.043$ & $1.043$ & $2.007$ & $-4.001$ & $61.42$ & $0.049$\\
3. Adaptive gain (RLS) & $2.064$ & $1.064$ & $2.007$ & $-4.001$ & $61.46$ & $0.751$\\
4. Online neural readout & $2.011$ & $1.011$ & $2.007$ & $-4.001$ & $61.48$ & $0.751$\\
5. Tabular value, smoothed & $2.039$ & $1.039$ & $2.007$ & $-4.001$ & $61.44$ & $1.271$\\
\midrule
\multicolumn{7}{@{}p{0.94\textwidth}@{}}{\emph{Panel B. Learners trained in the closed loop, saved and reloaded (frozen during the probe; shared training stream)}}\\
6. Recursive least squares & $3.000$ & $2.000$ & $2.008$ & $-4.014$ & $54.58$ & $R^2=0.993$\\
7. Two-layer network (SGD) & $1.878$ & $0.878$ & $2.007$ & $-4.012$ & $65.19$ & $R^2=0.994$\\
8. Policy gradient on profit & $3.000$ & $2.000$ & $2.004$ & $-4.006$ & $58.05$ & $R^2=0.995$\\
9. Tabular Q-learner (TD(0)) & $2.108$ & $1.108$ & $1.994$ & $-3.973$ & $186{,}821$ & $R^2=1.000$\\
\midrule
\multicolumn{7}{@{}l@{}}{\emph{Panel C. Where the trained-learner results change (range over the four learners)}}\\
Independent training streams & \multicolumn{2}{c}{rel.\ gap $\le10^{-5}$} & $0.13$ to $0.92$ & $-1.75$ to $-0.38$ & --- & (c) fails\\
Learning during the probe & \multicolumn{2}{c}{rel.\ gap $\le1.1\times10^{-3}$} & $1.99$ to $2.01$ & $-4.02$ to $-3.98$ & --- & (c) holds\\
\midrule
\multicolumn{7}{@{}p{0.82\textwidth}@{}}{Exact dead zone ($\tau=0.05$): $\|\Delta h\|_2=0$ at $A=0.005$; the probe moves the state past $\tau$ only for $A>0.139$.}\\
\multicolumn{7}{@{}p{0.82\textwidth}@{}}{Single-bin value lookup: $\|\Delta h\|_2=0$ at $A=0.005$; the probe leaves the central bin only for $A>0.149$.}\\
\bottomrule
\end{tabular}
\begin{tablenotes}\footnotesize\item Panel A: deterministic simulations around the operating point; Test A reports the raw response-gap slope in $\log A$ ($O(A^2)$) and the probe-normalised gap slope $\|X^{(1)}(A)-X^{(2)}(A)\|_2/A$ ($O(A)$, plotted in Figure~\ref{p1:fig:learners}). Panel B: learners trained for $5{,}000$ periods, parameters saved and reloaded without rescaling and frozen during the probe; the two compared populations are trained on the same innovation stream; relative equivalence gaps at $A=0.005$ are at most $10^{-5}$. A raw Test A slope above $2$ means the gap is of higher order than $A^2$. Panel C: the same learners trained on independent innovation streams (first row) or updated at learning rate $0.35$ during the probe (second row). Capital in par units, certified bound of Theorem~\ref{p1:thm:spectral} ($61.42$ for the linear model); the Q-learner settles on a small response gain ($0.022$), so separating it costs far more capital. Kind: recomputed (Panel A: \texttt{recompute\_p1\_learners.py} $\to$ \texttt{R\_P1\_learners.json}; Panels B--C: \texttt{validate\_trained\_learners\_hardened.py} $\to$ \texttt{R\_P1\_trained\_learners\_hardened.json}).\end{tablenotes}
\end{threeparttable}
\end{table}

\begin{figure}[t]\centering
\includegraphics[width=0.92\textwidth]{p1figs/fig_learners.pdf}
\caption{Left: Test~B, probe-normalised response difference $\|\Delta h\|_2$ against the rate half-spacing $\delta$ at probe size $A=0.005$, five stylized rules and the linear benchmark. Right: Test~A, probe-normalised response gap $\|X^{(1)}(A)-X^{(2)}(A)\|_2/A$ between two economies with the same merged measure against the probe size $A$, with an $O(A)$ reference line (corresponding to an $O(A^2)$ raw gap).}\label{p1:fig:learners}
\end{figure}

Table~\ref{p1:tab:learners} and Figure~\ref{p1:fig:learners} show both where the linear reduction holds and where it breaks down. In Panel~A, for small probes every smooth or dithered rule behaves like the linear model: economies with the same merged measure have raw response gap $O(A^2)$ (slopes $2.011$ to $2.064$) and probe-normalised gap $O(A)$ (slopes $1.011$ to $1.064$), while the response difference and certified capital scale as $\delta^2$ (slope $2.007$) and $\delta^{-4}$ (slope $-4.001$), with capital at $\delta=0.05$ within $0.1\%$ of the linear $61.42$. In Panel~B, predictions (a)--(c) pass for all four trained learners: a single exponential explains $R^2\ge0.993$ of each small-probe impulse response, populations with the same merged measure differ by at most $10^{-5}$ in relative response at $A=0.005$, and certified capital scales with slopes $-4.014$ to $-3.973$, inside the registered band. The level of capital differs across learners because training sets the response gain.

Panel~C shows where this changes. When the two populations are trained on independent innovation streams, prediction (c) fails: the capital slopes are $-1.75$, $-0.97$, $-1.62$ and $-0.38$. Independently trained populations end with learned gains that differ by sampling error of order $T_{\rm burn}^{-1/2}$, so they differ in merged weight as well as in memory rate, and a difference in weight is cheaper to detect than an $O(\delta^2)$ difference in rates. The $\delta^{-4}$ law therefore prices the separation of memory rates between populations whose feedback weights match; it is not a claim about telling apart any two trained populations. When learning continues during the probe at rate $0.35$, the law survives (slopes $-4.02$ to $-3.98$) and the merged-measure gap rises only to $1.1\times10^{-3}$, so the breakdown we anticipated under prediction (d) did not occur in this design.

The linear theory fails in two further places. First, saturation, adaptive gains and the neural readout leave the small-signal regime once probes exceed $A_{\rm sat}\in[0.751,1.271]$ par units, and the dithered threshold rule leaves it already at $A_{\rm sat}=0.049$, the threshold itself. Second, rules with hard discontinuities are different in kind: an exact dead zone ($\tau=0.05$) or a single-bin lookup does not respond at all to probes smaller than $0.139$ and $0.149$ par units ($\|\Delta h\|_2=0$), concealing those cohorts from small probes altogether. Neither smooth nonlinearity nor hard thresholds offers a cheaper route to separating memory rates than the linear benchmark; populations whose learned weights differ are told apart by their weights, not their rates.

\section{What disclosure buys}\label{p1:sec:disclosure}

Because active probing is costly and passive data are slow, the natural policy question is what a supervisor gains from disclosure. We compare six observation contracts at $\delta=0.05$.
\begin{itemize}[itemsep=1pt]
\item $I_0$: public prices only.
\item $I_1$: prices and aggregated executed flow.
\item $I_2$: prices and one independent noisy report of dealer marked equity per reset ($\sigma_E=0.5$).
\item $I_3$: prices and exact audited equity (Proposition~\ref{p1:prop:disclosure}), also with reports rounded to $0.01$.
\item $I_4$: algorithm update timestamps without the update values.
\item $I_5$: a supervisory audit registry with access to internal algorithm states.
\end{itemize}

\begin{table}[t]\centering\small
\begin{threeparttable}
\caption{Disclosure contracts at $\delta=0.05$ (pulse-design capital, par units)}\label{p1:tab:contracts}
\begin{tabular}{llrrr}
\toprule
Contract & Channel & Pulses & Benchmark & Saving vs.\ $T=20$\\
\midrule
$I_0$: prices, $T=3$ & public prices & 271 & 4{,}305.25 & $-2{,}123.5\%$\\
$I_0^*$: prices, $T=20$ & full price path & 12 & 193.62 & 0.0\%\\
$I_1$: prices and signed flow & public flow & 12 & 193.62 & 0.0\%\\
$I_2$: noisy equity ($\sigma_E=0.5$) & statement & 12 & 193.62 & 0.0\%\\
$I_3$: exact equity, 5-unit lot & audited equity & 1 & 12.70 & 93.4\%\\
$I_3$: exact equity, 0.01 rounding & rounded report & 1 & 15.52 & 92.0\%\\
$I_4$: timestamps only & no likelihood & 12 & 193.62 & 0.0\%\\
$I_5$: audit registry & internal states & 0 & 0.00 & 100.0\%\\
\bottomrule
\end{tabular}
\begin{tablenotes}\footnotesize\item Closed-form values. Noisy equity adds information linearly in the number of resets, $\KL=N\Delta_E^2/(2\sigma_E^2)$, and saves only 1.5\% against the $T=3$ baseline. The equity difference is $\Delta E=\kappa^2q\delta^2V^2$; with reports rounded to $0.01$, a guaranteed change needs $V\ge5.59$.\end{tablenotes}
\end{threeparttable}
\end{table}

Table~\ref{p1:tab:contracts} shows that disclosure is not one thing. Timestamps without update values carry no information about $\delta$ in this model. Noisy periodic statements barely help. Exact audited equity collapses the requirement to a single trade, limited only by lot sizes and report precision, and an audit registry removes the trading requirement altogether.

The cleanest form of useful disclosure is flow tagged by participant class.

\begin{proposition}[Tagged flow removes the barrier]\label{p1:prop:tagged}
Let the price be driven by i.i.d.\ innovations of variance $\sigma_\xi^2$ in a positive, stable economy, with stationary autocovariances $\gamma_x(s)=\sum_\ell c_\ell\lambda_\ell^{|s|}$ ($c_\ell\ge0$, $0\le\lambda_\ell<1$ the closed-loop modes). Suppose the observer receives prices and each cohort's flow $\tilde u_{i,t}=u_{i,t}+\varepsilon_{i,t}$, $u_{i,t}=q_iy_{i,t}$, with i.i.d.\ measurement noise of variance $\sigma_f^2$ independent of prices. Then $\tilde u_{i,t}=r_i\tilde u_{i,t-1}+w_ix_t+e_{i,t}$ with $e_{i,t}=\varepsilon_{i,t}-r_i\varepsilon_{i,t-1}$, and the just-identified IV estimator of $\beta_i=(r_i,w_i)$ with instruments $Z_t=(x_{t-1},x_t)^\top$ satisfies $\sqrt T(\hat\beta_i-\beta_i)\to N(0,V_i)$,
\begin{equation}\label{p1:eq:ivvar}
V_i=\Sigma_i^{-1}\Omega_i\Sigma_i^{-\top},\qquad \Sigma_i=\begin{pmatrix}w_iA_0&\gamma_x(1)\\ w_iA_1&\gamma_x(0)\end{pmatrix},\qquad \Omega_i=\sigma_f^2\bigl[(1+r_i^2)\Gamma_0-r_i(\Gamma_1+\Gamma_1^\top)\bigr],
\end{equation}
where $A_m=\sum_{s\ge0}r_i^s\gamma_x(s+m)$, $\Gamma_0=\begin{psmallmatrix}\gamma_x(0)&\gamma_x(1)\\ \gamma_x(1)&\gamma_x(0)\end{psmallmatrix}$ and $\Gamma_1=\E[Z_tZ_{t-1}^\top]$. The Jacobian determinant is
\begin{equation}\label{p1:eq:ivdet}
\det\Sigma_i=w_i\sum_{\ell,m}c_\ell c_m\frac{1-\lambda_\ell\lambda_m}{1-r_i\lambda_m}\;>\;0 .
\end{equation}
$V_i$ depends only on cohort $i$'s own $(w_i,r_i)$, $\sigma_f^2$ and the price autocovariances, which are bounded and nondegenerate uniformly on a uniformly stable family whatever the rate spacing; so the merged weights are estimated at rate $T^{-1/2}$ with a constant that does not depend on $|r_i-r_j|$.
\end{proposition}

\paragraph{In words.} If each group's trading is reported separately, each group becomes an ordinary regression whose precision does not depend on how similar the groups are.


Each tagged flow follows $u_{i,t}=r_iu_{i,t-1}+w_ix_t$, a regression whose information about $w_i$ does not involve the other cohorts' rates; the closed form \eqref{p1:eq:ivdet} matches the simulated Jacobian ($0.1330$ against $0.1322$ in a two-cohort economy with $4\times10^5$ periods; \texttt{hardening/p1\_checks.py}). In simulation ($T=5{,}000$, 60 replications) the root-mean-square error of the recovered weights is flat in $\Delta\eta$ (fitted slope $0.011$, error between $5.8\times10^{-4}$ and $6.6\times10^{-4}$). This is the formal content of a disclosure regime such as a consolidated audit trail: tagging flow by participant class turns a polynomial barrier into an ordinary estimation problem.

\subsection{How much disclosure is worth}\label{p1:sec:discvalue}

Proposition~\ref{p1:prop:tagged} says that tagging removes the barrier; the next result says by how much.

\begin{theorem}[Value of class-tagged disclosure]\label{p1:thm:discvalue}
Take a rich class of Theorem~\ref{p1:prop:esprit}: rates $r_c+\Delta\eta\,s$ with $s$ in a compact set of $\gamma$-separated ordered points that has nonempty interior and contains the Gauss and Gauss--Radau pair of Theorem~\ref{p1:thm:kcohort}, and merged weights in a nondegenerate interval $[\omega_{\min},\omega_{\max}]$, $0<\omega_{\min}<\omega_{\max}$. In each reset experiment the observer applies a unit impulse and records the price response with independent $N(0,\sigma^2)$ noise per tap; under tagging it also records each class's flow response $f_{i,t}=w_i\sum_{s\le t}r_i^{\,t-s}h_s$ with independent $N(0,\sigma_f^2)$ noise per tap. Let $N_p$ and $N_t$ be the numbers of experiments needed, without and with tagging, for a fixed accuracy and confidence. There are constants $0<c\le C<\infty$, depending on $(\sigma,\sigma_f)$ and the class but not on $\Delta\eta$, such that
\begin{enumerate}[label=(\roman*),itemsep=1pt]
\item to recover the weights to fixed accuracy, $c\le N_t\le C$, so $N_p/N_t\asymp(\Delta\eta)^{-(4K-2)}$;
\item to recover the rates to accuracy proportional to $\Delta\eta$, $c(\Delta\eta)^{-2}\le N_t\le C(\Delta\eta)^{-2}$, so $N_p/N_t\asymp(\Delta\eta)^{-(4K-4)}$.
\end{enumerate}
\end{theorem}

\begin{proof}
$N_p\asymp\sigma^2(\Delta\eta)^{-(4K-2)}$ by Theorems~\ref{p1:thm:kcohort} and~\ref{p1:prop:esprit}. \emph{Upper bounds.} Average the $N$ price responses to estimate $h$ with error of order $\sigma N^{-1/2}$ per tap. Given $h$, class $i$'s mean flow is a two-parameter family $F(w_i,r_i;h)$ whose Jacobian columns $\sum_sr_i^{t-s}h_s$ and $w_i\sum_s(t-s)r_i^{t-s-1}h_s$ are linearly independent, with a Gram determinant bounded below uniformly on the class because $w_i\ge\omega_{\min}$, the rates lie in a compact subset of $(0,1)$ and $h_0=\kappa>0$; the bound does not involve the other classes or $\Delta\eta$. Least squares on class $i$'s flows with the plug-in $h$ therefore estimates $(w_i,r_i)$ with error $O_P(N^{-1/2})$, uniformly on the class, which gives the upper bounds in (i) and in (ii) for accuracy $\epsilon\Delta\eta$. \emph{Lower bounds.} At a configuration in the interior of the class, perturb one class's weight by $\epsilon$, or its rate by $\epsilon\Delta\eta$; for small $\epsilon$ both perturbed configurations remain in the class because the weight interval is nondegenerate and the rate set has nonempty interior. The mean differences of both the price and the flow records are of order $\epsilon$, respectively $\epsilon\Delta\eta$, so the Kullback--Leibler divergence of $N$ experiments is at most $CN\epsilon^2$, respectively $CN\epsilon^2(\Delta\eta)^2$, and Le Cam's two-point lemma gives the lower bounds.
\end{proof}

\paragraph{In words.} This is the value of disclosure measured in data. Tagging replaces a requirement that explodes as the groups become similar with one that does not grow at all (for the weights) or grows only like $(\Delta\eta)^{-2}$ (for the rates).


Proposition~\ref{p1:prop:tagged} and Theorem~\ref{p1:thm:discvalue} use different observation arrangements: the proposition concerns passive, stationary data with prices observed without noise, where the price is a valid instrument; the theorem concerns reset experiments with noisy prices. With noisy prices, regressing flow on the observed price is an errors-in-variables problem, which is why the theorem works from the known impulse instead.

With the same noise and lags for both observers, the Cram\'er--Rao requirements give fitted slopes of $-6.08$ for prices and $-4.34$ for the ratio at $K=2$ (theory $-6$ and $-4$). Calibrated to the class memories estimated below, prices alone would need about 333 times as many days as the published class ratios to separate the top-trader and all-account memories in BTCUSDT (rates $0.9889$ and $0.9986$ per five minutes; about 6{,}000 days against 18), and about 19 times as many in ETHUSDT (rates $0.9793$ and $0.9972$). These rates come from day-demeaned regressions, which bias persistence down (Section~\ref{p1:sec:disc_data}), so the factors are illustrative.

\subsection{Disclosure in dollars}\label{p1:sec:disc_dollars}

Theorem~\ref{p1:thm:discvalue} states the value of class-tagged disclosure in samples. Theorems~\ref{p1:thm:spectral} and~\ref{p1:thm:kcohort} let us translate that value into a flow of round-trip trading losses and set it beside the official operating budgets of the U.S.\ Consolidated Audit Trail (CAT). Two adjustments make the comparison exact and like-for-like:
\begin{enumerate}[label=(\roman*),itemsep=1pt]
\item \emph{Exact $K$-cohort moment-gap constant.} Rather than extrapolating from Cram\'er--Rao ratios, we evaluate the spectral round-trip trading loss bound of Theorem~\ref{p1:thm:spectral},
\[
L_{\rm rt}^*(K,\Delta\eta)\;=\;\frac{4\sigma^2 z_{1-\alpha}^2}{\sup_\omega\bigl(|\Delta\hat H_{K,\Delta\eta}(e^{i\omega})|^2/g_0(\omega)\bigr)},
\]
directly for the $K$-point Gauss versus Gauss--Radau economies of Theorem~\ref{p1:thm:kcohort} and Table~\ref{p1:tab:radau}, whose moments below order $2K-1$ coincide and whose $(2K-1)$-th moment gap has the exact constant $\Delta_2=0.222222$, $\Delta_3=0.053333$, and $\Delta_4=0.013061$ at spacing $\Delta\eta$ (with $\Delta\eta=2\delta$ for the symmetric known pair).
\item \emph{Flow against flow.} A single probe loss $L_{\rm rt}^*$ is a one-time certification cost, whereas an audit trail is an ongoing annual expenditure. Because algorithm populations turn over and regulatory surveillance is periodic, we compare annual CAT operating costs against the annualised round-trip trading loss $52\times L_{\rm rt}^*$ under weekly certification renewal (matching the CFTC weekly reporting cycle of Section~\ref{p1:sec:cot}). We benchmark against three primary official figures: the CAT LLC 2025 operating budget of $\$248.8$~million in the November 2024 filing (SEC Release No.\ 34-101901, Dec.\ 12, 2024, p.~4), revised to $\$228.3$~million in May 2025, and the CAT LLC 2026 Financial and Operating Budget of 31 March 2026, with total expenses of $\$147.0$~million, which we use \citep{p1:catbudget2026}; the SEC's 2016 CAT NMS Plan Approval Order estimate of $\$2.4$~billion in initial industry build costs and $\$1.7$~billion per year in ongoing industry costs \citep[Release No.\ 34-79318, Nov.\ 15, 2016, 81 FR 84696 at 84801, p.~363]{p1:sec2016}; and the SEC's Rule 613 proposing and adopting releases estimating $\$4.0$~billion initial and $\$2.1$~billion annual ongoing costs \citep[Release No.\ 34-62174, May 26, 2010, 75 FR 32556 at 32598, 32600, pp.~147, 153; Release No.\ 34-67457, July 18, 2012, 77 FR 45722 at 45790--45795, pp.~220--238]{p1:sec2012}.
\end{enumerate}

\begin{table}[t]\centering\footnotesize
\begin{threeparttable}
\caption{Price-only round-trip probing loss per market and annualised weekly-renewal cost against official CAT budgets}\label{p1:tab:cat_dollars}
\setlength{\tabcolsep}{1.5pt}
\begin{tabular}{@{}>{\scriptsize}l*{6}{>{\scriptsize}r}@{}}
\toprule
 & \multicolumn{2}{c}{\scriptsize Loss / probe (USDT)} & \multicolumn{2}{c}{\scriptsize Annual ($52\times$, USDT)} & \multicolumn{2}{c}{\scriptsize Mkts ($\$147.0$M / $\$1.7$B)}\\
Economy ($\Delta\eta$) & Roll & 1-sec & Roll & 1-sec & Roll & 1-sec\\
\midrule
$K=2$, known ($0.20$) & 6{,}200 & 149{,}901 & 322{,}400 & 7{,}794{,}852 & 456.1 / 5{,}273.0 & 18.86 / 218.09\\
$K=2$, known ($0.10$) & 99{,}192 & 2{,}398{,}422 & 5{,}157{,}984 & 124{,}717{,}944 & 28.5 / 329.6 & 1.18 / 13.63\\
$K=2$, known ($0.04$) & 3{,}874{,}696 & 93{,}688{,}369 & 201{,}484{,}192 & 4{,}871{,}795{,}188 & 0.7 / 8.4 & 0.03 / 0.35\\
\midrule
$K=2$, Radau ($0.20$) & 4{,}292 & 103{,}788 & 223{,}184 & 5{,}396{,}976 & 658.9 / 7{,}617.0 & 27.25 / 314.99\\
$K=2$, Radau ($0.10$) & 283{,}488 & 6{,}854{,}612 & 14{,}741{,}376 & 356{,}439{,}824 & 10.0 / 115.3 & 0.41 / 4.77\\
$K=3$, Radau ($0.20$) & 804{,}403 & 19{,}450{,}107 & 41{,}828{,}956 & 1{,}011{,}405{,}564 & 3.5 / 40.6 & 0.15 / 1.68\\
$K=3$, Radau ($0.10$) & 1{,}137{,}325{,}880 & 27{,}500{,}017{,}407 & 59{,}140{,}945{,}760 & 1{,}430{,}000{,}905{,}164 & 0.0 / 0.0 & 0.00 / 0.00\\
$K=4$, Radau ($0.20$) & 156{,}632{,}086 & 3{,}787{,}291{,}895 & 8{,}144{,}868{,}472 & 196{,}939{,}178{,}540 & 0.0 / 0.2 & 0.00 / 0.01\\
\bottomrule
\end{tabular}
\begin{tablenotes}\footnotesize\item Separation variable $\Delta\eta$ throughout ($\Delta\eta=2\delta$ for the symmetric known pair). Gauss--Radau rows use the exact $K$-cohort quadrature pair of Theorem~\ref{p1:thm:kcohort} and Table~\ref{p1:tab:radau} (moment gaps $\Delta_2=0.222222$, $\Delta_3=0.053333$, $\Delta_4=0.013061$ at $W=0.08$, $r_c=0.60$). Loss per probe is the spectral round-trip trading loss bound of Theorem~\ref{p1:thm:spectral} at the median BTCUSDT Roll noise ($0.40$ USDT) and median one-second price-change standard deviation ($1.98$ USDT); annual loss multiplies by $52$ weekly renewals. The final columns divide the CAT LLC 2026 Financial and Operating Budget of 31 March 2026 (total expenses $\$147.0$M; earlier vintages: $\$248.8$M for 2025 in the November 2024 filing, SEC Release No.\ 34-101901, p.~4, revised to $\$228.3$M in May 2025) and the SEC 2016 annual industry estimate ($\$1.7$B, Release No.\ 34-79318, p.~363) by the annual probing loss per market. Probing cost is a lower bound on the cost of learning from prices by trading, not on passive learning and not a social cost. Kind: recomputed (\texttt{recompute\_p1\_cat\_costs.py} $\to$ \texttt{R\_P1\_cat\_disclosure\_dollars.json}).\end{tablenotes}
\end{threeparttable}
\end{table}

The figures are a shadow cost of learning by trading, the price an outside participant would pay; a supervisor without an audit trail does not probe, and its cost is detection delay, which Proposition~\ref{p1:prop:choice} models. The cost results require a positive economy: with signed weights the execution form $\tfrac12v^\top T_hv$ can be indefinite and a round trip can earn money, so none of the dollar figures applies there. Probing cost is a lower bound on the cost of learning from prices by trading, not on passive learning and not a social cost. It is also a lower bound for deterministic probes only: sequential probing can lower the expected loss by a factor of up to $0.49$ (Theorem~\ref{p1:thm:adaptive}). Only the observation noise is estimated from data; the structural parameters are illustrative, and the constant is sensitive to them. At $\delta=0.05$ the certified constant of Theorem~\ref{p1:thm:spectral} is $358.8$, $179.0$, $61.1$ and $13.9$ par units for $a=0$, $0.25$, $0.5$ and $0.7$, and $3.9$, $19.6$, $61.1$ and $305.2$ for $\bar\eta=0.2$, $0.3$, $0.4$ and $0.6$; at $\Delta\eta=0.20$ instead of $0.10$ the unknown-pair figures fall by a factor of about $64$. The Gauss--Radau pair is used in the orientation that gives the larger, still valid, lower bound. The table's ``known'' rows have rates $\bar\eta\pm\delta$ (spread $2\delta=\Delta\eta$), while the Radau rows place rates at $r_c+\Delta\eta\,s$ with $s\in[-1,1]$, a spread of about $1.33\Delta\eta$ for $K=2$, so the two kinds of row are not at identical spreads. Within that interpretation, Table~\ref{p1:tab:cat_dollars} shows how rapidly the polynomial exponent $(\Delta\eta)^{-(4K-2)}$ and the shrinking Gauss--Radau moment-gap constant $\Delta_K$ turn the comparison in favour of disclosure. For a known pair at coarse separation $\Delta\eta=0.20$ ($\delta=0.10$), weekly price-only renewal costs $\$0.32$~million per market per year at Roll noise and $\$7.79$~million at one-second noise, so matching the $\$147.0$~million 2026 CAT LLC budget requires $456.1$ or $18.86$ markets. At $\Delta\eta=0.10$ ($\delta=0.05$) the unknown $K=2$ Gauss--Radau pair costs $\$14.74$~million (Roll) or $\$356.44$~million (one-second) per market per year, matching the CAT LLC budget in $10.0$ or $0.41$ markets and the SEC's $\$1.7$~billion industry figure in $115.3$ or $4.77$ markets. Once a supervisor must separate $K=3$ unknown cohorts at $\Delta\eta=0.20$, annual probing loss rises to $\$41.83$~million (Roll) and $\$1.01$~billion (one-second) per market ($3.5$ and $0.15$ markets at the CAT LLC budget; $40.6$ and $1.68$ markets at $\$1.7$~billion), and at $K=3, \Delta\eta=0.10$ or $K=4, \Delta\eta=0.20$ a single market's annual probing loss exceeds both the CAT LLC budget and the SEC's total annual industry estimate; with the 2026 budget, a single market also exceeds the CAT budget for a known pair at $\Delta\eta=0.04$ under both noise calibrations. The renewal frequency is an assumption: break-even market counts scale inversely with it, so monthly renewal multiplies every count in the last two columns by $4.3$ and renewal on each of 252 trading days divides them by $4.8$.

\paragraph{When is disclosure worth its cost?} The dollar comparison prices disclosure; it does not say when a supervisor should buy it. Remark~\ref{p1:rem:riskneutral} shows that a supervisor who only needs to bound aggregate price risk does not need disclosure at all: the summed response is identified from public prices at the ordinary rate. Attribution matters when supervision is selective: one algorithm must be stopped or sanctioned while the others keep supplying liquidity, and the law requires attributing the conduct to it at a given error level. A supervisor without an audit trail does not run trading probes; it waits for public data to accumulate, and the cost of waiting is the harm done in the meantime.

Consider $M$ markets, a supervisory cycle of $\bar T$ periods, and in each market one algorithm whose conduct does harm $\ell>0$ per period until it is attributed. With prices alone, attribution at the required error level takes $T_p(\Delta\eta)=c_p(\Delta\eta)^{-(4K-2)}$ periods; with class-tagged records it takes $T_t(\Delta\eta)=c_t(\Delta\eta)^{-2}$ periods, the orders of Theorems~\ref{p1:prop:passive} and~\ref{p1:thm:discvalue} with constants $c_p,c_t>0$. We take the attribution times to be of exactly these orders, which is a modelling assumption: Theorem~\ref{p1:prop:passive} gives only a lower bound for prices, and if price-only attribution is slower, $G$ below is larger and disclosure is chosen at least on the band the proposition describes. $T_t\propto(\Delta\eta)^{-2}$ is the rate target; for weights at fixed accuracy $T_t$ does not grow as $\Delta\eta$ shrinks, and the lower edge of the band disappears. Attribution that would take longer than the cycle does not happen within it. An audit trail costs $F$ per cycle to operate and $c$ per market to comply with, including any information leakage. The supervisor's expected loss is therefore $\mathcal L_{\rm price}=M\ell\min\{\bar T,T_p\}$ without disclosure and $\mathcal L_{\rm disc}=F+Mc+M\ell\min\{\bar T,T_t\}$ with it. Write $\Delta\eta_t=(c_t/\bar T)^{1/2}$ and $\Delta\eta_p=(c_p/\bar T)^{1/(4K-2)}$ for the spacings at which tagged and price-only attribution just fit within the cycle.

\begin{proposition}[When an audit trail is chosen]\label{p1:prop:choice}
Let $K\ge2$ and $\Delta\eta_t<\Delta\eta_p$, and let $G(\Delta\eta)=\min\{\bar T,T_p(\Delta\eta)\}-\min\{\bar T,T_t(\Delta\eta)\}$ be the attribution delay that disclosure saves.
\begin{enumerate}[label=(\roman*),itemsep=1pt]
\item The supervisor strictly prefers disclosure if and only if $F/M+c<\ell\,G(\Delta\eta)$.
\item $G$ vanishes for $\Delta\eta\le\Delta\eta_t$, increases on $(\Delta\eta_t,\Delta\eta_p]$ to its maximum $G^*=\bar T-T_t(\Delta\eta_p)>0$, decreases wherever it is positive beyond $\Delta\eta_p$, and is negative for $\Delta\eta>(c_p/c_t)^{1/(4K-4)}$, where it returns toward zero from below. Hence the spacings at which disclosure is chosen form an interval around $\Delta\eta_p$, possibly empty: disclosure is never chosen for algorithms so similar that even tagged records cannot attribute them within the cycle, nor for algorithms so different that prices attribute them quickly.
\item The interval is non-empty if and only if the audit trail covers more than $M^*=F/(\ell G^*-c)$ markets (with $M^*=\infty$ if $\ell G^*\le c$), and it widens as $M$ grows.
\end{enumerate}
\end{proposition}

Part (i) restates the comparison of losses. Part (ii) says that disclosure is valuable in a middle band, where algorithms are close enough that prices cannot attribute them within the cycle but far enough apart that tagged records can; the polynomial gap between the two delays, $(\Delta\eta)^{-(4K-2)}$ against $(\Delta\eta)^{-2}$, is what makes the band wide. Part (iii) is an economies-of-scale result: an audit trail with a fixed operating cost pays only when it spans enough markets, which is consistent with the design of market-wide regimes such as CAT rather than venue-by-venue reporting. A supervisor that can trade has a second option, active probing, whose cost is the frontier of Table~\ref{p1:tab:cat_dollars}; because that cost grows at the same polynomial order, it shortens price-only attribution only for well-separated algorithms and leaves the comparison for close ones unchanged. The model takes the harm $\ell$ as given and holds the algorithms' behaviour fixed across regimes; Section~\ref{p1:sec:conclusion} returns to both.

\paragraph{What the data can and cannot test.} The theorems concern experiments that no outside researcher can run: reset initial states, known price persistence and impact, and either controlled probes or stationary data from a known innovation law. Public data meet none of these exactly. Order flow is endogenous, the fundamental value is unobserved and replaced by returns, observations overlap in time, parameters drift, and the class reports are positions aggregated over many accounts rather than tagged flow. The tests below therefore check the theory's observable prediction, that class records restrict persistence which prices leave unrestricted, under a stated mapping from class positions to cohort states. They do not estimate the theorems' rates, and a pass does not show that class persistence is learning memory.

\subsection{Mandated disclosure: the CFTC's trader-class reports}\label{p1:sec:cot}

The cleanest test uses disclosure mandated by law. The U.S. Commodity Futures Trading Commission requires large traders in regulated futures markets to report their positions and publishes them weekly by class: leveraged funds, asset managers and dealers in financial futures, and managed money, producer/merchants and swap dealers in commodities. This is class-tagged disclosure mandated by law. We registered a test (SHA-256 \texttt{6f1d6f35ae88\ldots}) before downloading any data, in eight markets that span asset classes: CME Bitcoin and Ether, the E-mini S\&P 500, 10-year Treasury notes, the euro, the yen, WTI crude oil and Henry Hub natural gas, using every weekly report through October 2026 (282 to 1{,}039 weeks per market). For each class, $n_{c,t}=(\text{long}-\text{short})/\text{open interest}$ follows $n_{c,t}=\alpha_c+r_cn_{c,t-1}+w_c\,{\rm ret}_t$, the weekly version of the cohort recursion. The price-only alternative fits a two-exponential kernel to twelve weekly return lags. Intervals come from a 13-week moving-block bootstrap. The registered hypotheses were: (H1, primary) the fast class's memory rate is pinned to an interval narrower than $0.25$ while both price-only intervals are wider, in at least 75\% of markets; (H2) the fast class trend-follows, $w_c>0$, in at least 75\%; (H3) the fast class has the shorter memory in at least half.

\begin{table}[t]\centering\small
\begin{threeparttable}
\caption{Fitted persistence of trader-class positions, CFTC Commitments of Traders, eight regulated futures markets}\label{p1:tab:cot}
\setlength{\tabcolsep}{3pt}\footnotesize
\begin{tabular}{lrccccr}
\toprule
Market & Weeks & Fast $\hat r$ [95\% CI] & Slow $\hat r$ & $\hat r_{\rm fast}-\hat r_{\rm slow}$ [95\% CI] & $\hat w_{\rm fast}$ & Ratio\\
\midrule
Bitcoin (CME) & 434 & 0.951 [0.928, 0.969] & 0.986 & $-0.035$ [$-0.062$, $-0.010$] & $-$ & 23.7\\
Ether (CME) & 282 & 0.890 [0.826, 0.932] & 0.960 & $-0.070$ [$-0.143$, $-0.011$] & $-$ & 9.1\\
E-mini S\&P 500 & 1{,}039 & 0.914 [0.884, 0.937] & 0.954 & $-0.040$ [$-0.071$, $-0.012$] & $-$ & 18.3\\
10-year T-note & 1{,}039 & 0.992 [0.985, 0.998] & 0.993 & $-0.000$ [$-0.006$, $+0.007$] & (0) & 71.4\\
Euro & 1{,}035 & 0.961 [0.947, 0.972] & 0.996 & $-0.034$ [$-0.051$, $-0.022$] & $+$ & 37.3\\
Yen & 1{,}035 & 0.955 [0.939, 0.968] & 0.984 & $-0.029$ [$-0.047$, $-0.013$] & $+$ & 34.0\\
WTI crude oil & 1{,}039 & 0.961 [0.946, 0.974] & 0.999 & $-0.037$ [$-0.053$, $-0.023$] & $+$ & 35.2\\
Natural gas & 1{,}039 & 0.981 [0.969, 0.990] & 0.980 & $+0.001$ [$-0.014$, $+0.017$] & $+$ & 44.0\\
\midrule
\multicolumn{7}{l}{\emph{Pre-registered hypotheses}: H1 (primary) 8 of 8, pass; H2 4 of 8, fail; H3 6 of 8, pass.}\\
\bottomrule
\end{tabular}
\begin{tablenotes}\footnotesize\item Fast class: leveraged funds (financial futures) or managed money (commodities); slow class: asset managers or producer/merchants. $\hat w_{\rm fast}$: sign of the response to returns, $+$ or $-$ when the 95\% interval excludes zero, (0) otherwise. Ratio: mean width of the two price-only rate intervals over the width of the fast-class interval; the price-only intervals span 0.95--0.96 of the admissible range $[0,0.99]$ in every market. Weekly returns from Binance spot (Bitcoin, Ether) and Yahoo Finance (others; the Treasury return is minus the change in the 10-year yield). Specification SHA-256 \texttt{6f1d6f35ae88\ldots}; test opened once. Kind: real data (Part 11).\end{tablenotes}
\end{threeparttable}
\end{table}

The primary hypothesis holds in all eight markets (Table~\ref{p1:tab:cot}). The class reports pin the fitted persistence rate of each fast class's positions to an interval between $0.013$ and $0.106$ wide, while the price-only kernel leaves both rates anywhere in the admissible range; the width ratios run from 9 to 71. The ratio partly reflects how little weekly returns reveal, which is the point of the theory rather than an artefact of it. A fitted persistence of aggregated positions is not a structural learning memory: it also reflects how slowly positions turn over and how accounts enter and leave a class. The classes differ as the cohort model allows. Leveraged funds and managed money have the less persistent positions in six of eight markets, with half-lives of 6 to 18 weeks against months to years for asset managers and producers; in Treasuries and natural gas the two persistences are indistinguishable.

The trend-following hypothesis fails. The fast class buys after rises in the euro, the yen, crude oil and natural gas, but sells after rises in Bitcoin, Ether and the S\&P 500 futures. The first half matches the finding of \citet{p1:moskowitz2012} that speculators' positions tend to go with past returns. We did not predict this. One candidate explanation is that leveraged-fund positions in financial futures are dominated by the short futures leg of cash--futures basis trades, which a single memory per class cannot represent \citep{p1:barth2021}. A post-hoc check is consistent with it but not decisive: weekly position changes of leveraged funds and asset managers are strongly negatively correlated in Bitcoin, Ether, the S\&P 500 and Treasuries ($-0.37$ to $-0.67$) but not in the euro or yen ($-0.04$ and $-0.05$), and holding the asset managers' change fixed removes the negative response in Bitcoin and halves it in Ether and the S\&P 500. Because class positions sum to zero across all reporting categories, some negative correlation is mechanical, so this is a signature, not a test. The failure sharpens rather than weakens the case for class tagging: class labels are not strategies, and a supervisor who wants to know how a class responds to prices needs its positions, which the price process does not reveal. Negative estimated $w_c$ also lies outside the positive economy of Section~\ref{p1:sec:model}, so the dollar bounds of Section~\ref{p1:sec:disc_dollars} do not apply to these classes; the registered restriction test does not use positivity.

Together with Table~\ref{p1:tab:disc}, the evidence now covers two disclosure regimes, one voluntary and one mandated by law, in crypto and in regulated equity-index, rate, currency and commodity futures. In both, class-tagged positions restrict persistence rates that prices leave unrestricted; whether that persistence is learning memory or slow turnover of positions, these data do not say.

\subsection{Voluntary disclosure: Binance's trader-class ratios}\label{p1:sec:disc_data}

Binance publishes voluntarily, for one venue but every five minutes, long/short ratios for its top-trader group, by position and by account, and for all accounts. These are class-level positions, the closest public analogue to the tagged data of Proposition~\ref{p1:prop:tagged}, but not homogeneous cohorts: the top-trader position ratio is built from position sizes, the all-account ratio from counts of long and short accounts, the top-trader group is a subset of all accounts, and its membership can change. We treat a class's reported net position as the observable counterpart of its cohort state: it obeys the same one-lag recursion class by class, so the block-diagonal information structure of Proposition~\ref{p1:prop:tagged} carries over. The mapping is an assumption, because the reports are positions rather than flows and returns stand in for the unobserved deviation $x_t$. For each class $c$ we estimate $n_{c,t}=r_cn_{c,t-1}+w_c\,{\rm ret}_t$, with $n=(R-1)/(R+1)$ from the ratio $R$. The price-only alternative regresses signed five-minute taker flow on twelve lagged five-minute returns and fits two exponentials to the kernel; it estimates the memory rates without class information. Intervals come from a day bootstrap with 200 draws.

\begin{table}[t]\centering\small
\begin{threeparttable}
\caption{Fitted persistence rates with and without class-tagged data}\label{p1:tab:disc}
\setlength{\tabcolsep}{3pt}
\begin{tabular}{lcccp{0.1\textwidth}p{0.08\textwidth}}
\toprule
Symbol & Days & Top position $\hat r$ [95\% CI] & All accounts $\hat r$ [95\% CI] & Width ratio & Classes differ\\
\midrule
BTCUSDT & 48 & 0.9889 [0.9842, 0.9935] & 0.9986 [0.9955, 1.0005] & 121.7 & yes\\
ETHUSDT & 48 & 0.9793 [0.9616, 0.9907] & 0.9972 [0.9934, 0.9997] & 55.6 & yes\\
CELRUSDT & 48 & 0.9884 [0.9830, 0.9922] & 0.9882 [0.9804, 0.9914] & 81.6 & no\\
CVXUSDT & 33 & 0.9822 [0.9732, 0.9908] & 0.9829 [0.9694, 0.9900] & 45.9 & no\\
\midrule
\multicolumn{6}{l}{\emph{Pre-registered test on fresh days, 30 USD-M perpetuals, 2025--2026}}\\
\multicolumn{3}{l}{(i) width ratio $>10$} & \multicolumn{3}{l}{30 of 30 (range 27.6--124.0)}\\
\multicolumn{3}{l}{(ii) price-only rate interval wider than 0.25} & \multicolumn{3}{l}{30 of 30 (range 0.65--0.99)}\\
\multicolumn{3}{l}{(iii) merged response $\sum_mk_m$ identified} & \multicolumn{3}{l}{5 of 30}\\
\multicolumn{3}{l}{Classes differ} & \multicolumn{3}{l}{11 of 30}\\
\bottomrule
\end{tabular}
\begin{tablenotes}\footnotesize\item Top panel: training days 2023--2024; width ratio is the mean price-only rate-interval width over the mean tagged width. The price-only intervals for both rates cover most of $[0,0.99]$. Bottom panel: specification registered before the fresh days were downloaded (SHA-256 \texttt{65acda361861\ldots}); a symbol passes if all three of the first rows hold, and the test required 80\% of symbols to pass. \end{tablenotes}
\end{threeparttable}
\end{table}

Table~\ref{p1:tab:disc} gives three results. First, tagged class data restrict persistence rates that prices and flow leave essentially unrestricted: with the one-hour price-only kernel, the price-only intervals are 28 to 124 times wider in every one of 30 symbols on fresh days, and cover most of the admissible range. Part of this ratio is mechanical, because a series of positions identifies its own persistence directly while the price-only route must fit two exponentials to a weakly identified kernel. Restricting attention to the five symbols in which the price-only fit does identify the merged response (ACH, API3, OGN, ZEN and DASH) does not change the conclusion: the price-only rate intervals are still $0.64$ to $0.92$ wide and the ratios run from $27.6$ to $62.5$ (median $47.5$). We therefore read the ratios as a measure of how little prices reveal about class memory, not as a horse race between two estimators of equal standing. The ratios also overstate the gap for a second reason: the price-only kernel spans one hour (twelve five-minute lags), while the tagged half-lives are several hours, so the price-only fit cannot see memories of that length. An exploratory check on the 30 symbols and 105 days of the follow-up test refits the price-only kernel with 48 and 96 lags (four and eight hours) and compares both routes in one unit, the half-life. With 96 lags the price-only slow half-life (median 5.6 hours) reaches the tagged scale (median 4.2 hours). The tagged intervals remain narrower in all 30 symbols, but by a median factor of 7.5 in log half-life rather than 28 to 124, and the price-only interval still reaches a bound of the fit in 25 of 30. The horizon-matched comparison was then registered as a primary hypothesis of a test on 30 symbols never used in this project, before its data were downloaded, and passed: the tagged interval was narrower in 29 of 29 evaluable symbols (median factor 7.3, smallest 2.6), and the price-only interval reached a bound of the fit in 26. A registered secondary of the same test qualifies what the tagged estimate measures. Day-demeaned regressions bias persistence down; corrected by a split-panel jackknife, the median tagged rate is 1.0008 per five minutes, and the corrected interval lies below one in only 7 of 87 class fits. Within a day the class ratios are close to a random walk, so the half-lives of a few hours reported here are largely produced by removing each day's mean. The evidence therefore supports a narrower claim than identification of class memory: positioning data restrict a persistence parameter that prices cannot bound. Whether that persistence reflects learning memory or slow turnover of positions, which theory's tagged contract (class-tagged flow) would distinguish, is not established by these data. Second, the classes differ where theory says the information should be: in BTCUSDT and ETHUSDT the top-trader group's positioning has a half-life of about three to five hours against about one to two days for all accounts, while in thin symbols the classes are often indistinguishable. Third, the pre-registered test fails its 80\% criterion, on the part of the prediction that concerns what prices \emph{do} reveal: the merged response is significantly different from zero in only 5 of 30 symbols. Theorem~\ref{p1:thm:minphase} says the merged measure is identified; on 42 days of thin-market data it is not estimated precisely. This is a failure of the prediction as registered, and we report it as one. A post hoc reading, not tested here, is that identification is an asymptotic statement and 42 days of passive thin-market data are far from that regime. Two follow-ups were registered before their data were downloaded, and each ran once. On 105 untouched days in the same 30 markets, price-only intervals were again at least ten times wider than class-tagged ones in all 30 symbols (median ratio 74.9 with the one-hour kernel), but the merged response was significantly different from zero in only 9 symbols, and pinning it down was at least ten times cheaper than separating the memory rates in only 10 of 30 (23 were required): that prediction failed. On 58 further days, the advantage of tagging grew with the sample in 22 of 29 symbols, the number required, but the tagged intervals shrank more slowly than the parametric rate (20 of 29 in the predicted range, 22 required), consistent with drift in class persistence over 21 months. Both results are reported in full in the companion thesis.

To test directly whether the memory rates identified by class-tagged positions improve out-of-sample return predictions after large order-flow shocks, we pre-registered a dedicated out-of-sample test (\texttt{P1\_disclosure\_oos.json}, SHA-256 \texttt{f0182ce50978\ldots}) before downloading the 22nd-of-month data for January 2025--September 2026 across 19 thin perpetuals ($6{,}137$ top-decile five-minute flow shocks on 2024 training days; $4{,}579$ shocks on untouched 2025--2026 test days). On the training sample we estimate each symbol's two memory rates $\hat r_1,\hat r_2$ either from the class-tagged positions or from the price-only two-exponential fit, construct the filtered memory states $y_{k,t}=\hat r_k y_{k,t-1}+(1-\hat r_k){\rm ret}_t$, and forecast the 30-minute post-shock cumulative return $R_{t+1\to t+6}$ out of sample relative to an AR(1) return-and-flow baseline. The primary hypothesis fails: on the untouched test shocks, the price-only memory states yield $\Delta R^2=-0.199$~percentage points relative to the baseline, the class-tagged memory states yield $\Delta R^2=-1.288$~points, and the tagged-minus-price-only difference is $\Delta R^2=-1.089$~points (95\% date-block bootstrap CI $[-2.605,+0.376]$, Clark--West $t=0.96$, $p=0.1684$). Likewise, adding the level and one-hour change of the top-trader and all-account positions to the hourly tail-risk controls of Chapter~\ref{ch:p2} leaves out-of-sample $R^2$ unchanged ($-0.19$~points, CI $[-0.99,0.83]$).

After this test failed, we registered a follow-up (\texttt{P1\_disclosure\_state\_oos.json}, SHA-256 \texttt{5136cc6b846d\ldots}) on untouched 2nd-and-9th-of-month data for January 2025--September 2026 across the same 19 thin perpetuals ($12{,}046$ top-decile five-minute flow shocks on 24 training dates in 2024; $9{,}875$ shocks on 42 test dates; test opened once). It asks the narrower question that Theorem~\ref{p1:thm:merged} and Proposition~\ref{p1:prop:tagged} answer directly: can prices and flow track the class states themselves? A two-exponential price-and-flow filter ($\alpha_{\rm fast}=0.50$, $\alpha_{\rm slow}=0.10$ on five-minute returns and signed flow, plus shock return, shock flow and symbol fixed effects) is compared with a filter that adds the current class states and their one-bar changes. Both registered primaries pass by a wide margin (Table~\ref{p1:tab:disc_state_oos}), which confirms that prices cannot reconstruct class position levels; as shown below, this reflects the near-unit-root persistence of positions rather than predictive identification of strategy dynamics: for the 30-minute-ahead divergence between top-trader and all-account positions, out-of-sample $R^2$ is $-4.8\%$ with prices and flow and $46.0\%$ with class data; for the top-trader state it is $4.4\%$ and $54.8\%$.

\begin{table}[!htbp]\centering\footnotesize
\begin{threeparttable}
\caption{Pre-registered out-of-sample cohort-state trajectory test on untouched 2nd-and-9th-of-month days (\texttt{P1\_disclosure\_state\_oos}), 2025--2026}\label{p1:tab:disc_state_oos}
\setlength{\tabcolsep}{1.4pt}
\begin{tabular}{lccccc}
\toprule
30-min post-shock target (19 thin) & $R^2_{\rm pf}$ & $R^2_{\rm tag}$ & $\Delta R^2$ (pts) [95\% CI] & CW $t$ ($p$) & Res.\\
\midrule
Prim.\ 1: Divergence $d_{t+6}$ ($9{,}875$ shocks) & $-4.766$ & $45.962$ & $+50.728$ [$+45.291$, $+55.918$] & $21.27$ ($<10^{-16}$) & pass\\
Prim.\ 2: Top state $n_{{\rm top},t+6}$ ($9{,}875$ shocks) & $4.441$ & $54.845$ & $+50.404$ [$+47.226$, $+53.724$] & $21.85$ ($<10^{-16}$) & pass\\
Sec.: Divergence $d_{t+6}$, all ($210{,}144$ bars) & $-0.586$ & $37.053$ & $+37.639$ [$+36.232$, $+39.107$] & $72.75$ ($<10^{-16}$) & ---\\
Sec.: Top state $n_{{\rm top},t+6}$, all ($210{,}144$ bars) & $0.186$ & $46.779$ & $+46.593$ [$+45.555$, $+47.643$] & $80.44$ ($<10^{-16}$) & ---\\
\bottomrule
\end{tabular}
\begin{tablenotes}\footnotesize\item Out-of-sample $R^2$ in percent. Trained on 24 dates in 2024 ($12{,}046$ shocks, $120{,}384$ bars); evaluated once on 42 untouched test dates (2nd and 9th of each month, Jan 2025--Sep 2026); 2{,}000-draw date-block bootstrap 95\% CIs; Holm $p<10^{-16}$ across both primaries. Specification SHA-256 \texttt{5136cc6b846d\ldots} (\texttt{R\_P1\_disclosure\_state\_oos.json}).\end{tablenotes}
\end{threeparttable}
\end{table}

The gain is almost entirely persistence. In an exploratory decomposition on the same test shocks (\texttt{hardened\_empirical\_audits.py} $\to$ \texttt{R\_hardened\_empirical\_audits.json}), an AR(1) on the current level alone accounts for $100.5\%$ and $99.1\%$ of the improvement, and when the target is the 30-minute \emph{change} in positions the class data add $-0.26$ and $+1.41$ points. Because the class filter observes the series it forecasts, this is a weak test: it confirms that prices and flow cannot reconstruct where the classes are positioned, which Theorem~\ref{p1:thm:merged} implies, but it does not show that class data reveal how their positions will move, and neither registered test finds that class data improve return forecasts.

\section{A spectral blind spot}\label{p1:sec:blind}

Even without identification, could a supervisor monitor risk through eigenvalues or volatility?

\begin{theorem}[Spectral blindness]\label{p1:thm:blind}
For any eigenvalues $0<\lambda_2<\lambda_1<1$, one-step price variance $s_1>0$ and stationary price variance $V$ large enough, there is a family of two-state markets $J_\theta=\begin{psmallmatrix}\lambda_1&\theta\\0&\lambda_2\end{psmallmatrix}$ with shock covariance ${\rm diag}(s_1,s_2(\theta))$ such that all members share the eigenvalues, the one-step price variance and the stationary price variance, while the worst-case transient gain $\sup_k\|e_1^\top J_\theta^k\|$ grows without bound in $\theta$.
\end{theorem}

\paragraph{In words.} Two markets can look identical by the usual risk measures, the same eigenvalues and the same price variance, and yet one of them can produce far larger losses in the short run. Monitoring eigenvalues or volatility alone therefore cannot certify that a market is safe.

At $\lambda=(0.8,0.7)$, $s_1=1$ and $V=5$, the gain is $1.77$, $5.10$, $16.95$, $50.85$ and $169.50$ for $\theta=1,3,10,30,100$, with stationary variance $5.000$ in every case; the cumulative gain rises from $18.3$ to $1{,}667.7$. No monitor that depends only on eigenvalues and second moments of the price can be uniformly safe (Figure~\ref{p1:fig:blind}).

\begin{figure}[t]\centering
\includegraphics[width=0.6\textwidth]{p1figs/fig_blind.pdf}
\caption{Price-row transient gain of the spectral-blindness family $J_\theta$: eigenvalues $(0.8,0.7)$, one-step price variance 1 and stationary price variance 5 are the same for every $\theta$.}\label{p1:fig:blind}
\end{figure}

Two clarifications prevent over-reading. First, the cohort economy of Section~\ref{p1:sec:model} is itself non-normal, yet its price output does not amplify (Lemma~\ref{p1:lem:cm}). Large weighted gains can be manufactured by output weights alone: with $C_w=(1,160)$, a price-and-memory functional peaks at $160.003$ at $t=0$, which is the static weight, not dynamic growth (Appendix~\ref{p1:app:coordw}). Second, in a network of $M$ identical cohort economies coupled by symmetric nonnegative cross-impact, the coupled operator stays symmetric in the coordinates of Lemma~\ref{p1:lem:cm}; with cross-impact $0.05$, the spectral radius rises from $0.822$ ($M=1$) to $0.973$ ($M=10$) while the Kreiss constant stays at one, as it must for a stable symmetric operator (\texttt{recompute\_oct20.py}). A stable normal contraction cannot exhibit Euclidean state-norm transient growth. This does not exclude scalar input-output overshoot relative to initial impact, which also depends on the input and output maps. The positive-cohort no-overshoot result is model-specific; Chapter~\ref{ch:p2} studies a depth-depletion mechanism.
\section{Empirical feedback and its measurement}\label{p1:sec:empirics}

The theory says what prices can reveal. The empirical work has two roles, and neither is a test of the identification bounds themselves, which concern controlled experiments no outside researcher can run. Section~\ref{p1:sec:disc_data} tests the theory's policy prediction, that class-tagged data pin down memories that prices leave unrestricted. This section asks what the reduced-form feedback indices a supervisor would compute from public data actually measure; the full tables are in the published version of the paper.

Empirical work replaces the unobserved fundamental by a trailing exponential moving average. Under a martingale price the filtered state $x^{\rm EMA}$ is then an AR(1) with mechanical persistence $1-\alpha$, equal to $29/31\approx0.9355$ for a 300-second filter on ten-second bars, so measured persistence must be compared with that value, not with zero. With a VAR(1) in $(x^{\rm EMA},f)$, $f$ the signed flow, the loop gain is $\hat Q=\hat J_{12}\hat J_{21}/(1-\hat J_{22})$. Five findings follow, on a Binance spot panel of eight pairs on 96 dates in 2021--2024 (6.6 million ten-second transitions), on 72 days of one-second BTCUSDT data and on the 30 perpetuals of Chapter~\ref{ch:p2}. First, timing dominates: pairing the state with next-bar flow, a look-ahead, gives $\hat Q=0.0081$; aligning flow with the bar in which it is observed lowers it by more than an order of magnitude, to $0.00037$; raw returns give an interval that includes zero; and a placebo random walk reproduces the mechanical persistence ($0.9351$) with $\hat Q$ indistinguishable from zero. Second, feedback that the filter cannot produce is detectable but small: in BTCUSDT $\hat Q=0.0045$ [$0.0034$, $0.0056$] against a placebo of zero, while price persistence ($0.9364$) is almost entirely mechanical. Third, order splitting biases lagged-flow regressions (in simulation OLS never covers the truth, while Cochrane--Orcutt with the true autocorrelation, 2SLS and LIML have correct size), although at ten-second bars the residual serial correlation is negligible and controlling for 30 lags of flow lowers $\hat Q$ further. Fourth, measured feedback rises with depth across the 30 perpetuals (Spearman correlation with log depth $+0.70$). Fifth, when Binance set BTC spot fees to zero and later reinstated them, BTCUSDT volume rose by a factor of $2.6$ and then fell by a factor of $4.9$ relative to seven control pairs, while $\hat Q$ did not move (pooled effect $-0.00016$, interval $[-0.00087,0.00055]$). Measured feedback in crypto markets is real but small, is a feature of deep books, does not track trading activity, and is easily confused with filtering and timing artefacts.

\section{Implications for supervision and market design}\label{p1:sec:policy}

\paragraph{What a supervisor can learn from public data.} Prices and anonymous flow reveal the merged measure $M$ and the closed-loop impulse response, hence the transfer function, and nothing finer. A supervisor who wants to know \emph{which} algorithms drive a feedback loop cannot get there from public data at reasonable cost when their memories are close. This is not a failure of econometric technique: Theorem~\ref{p1:thm:twocohort} holds for every estimator, biased or not.

\paragraph{What probing costs.} At realistic observation noise, an exchange or a large desk could separate two well-separated memory rates for about $10^5$ USDT of capital. Probing is therefore feasible for coarse questions and infeasible for fine ones. Because capital grows as $\delta^{-4}$ and experiments as $(\Delta\eta)^{-(4K-2)}$, halving the separation multiplies the cost by 16 for a known pair and by $2^{4K-2}$ for $K$ unknown cohorts.

\paragraph{What disclosure buys.} Disclosure regimes differ sharply in value (Table~\ref{p1:tab:contracts}). Timestamps without update values are worthless for this purpose. Noisy periodic statements are nearly worthless. Audited equity and, above all, flow tagged by participant class remove the barrier entirely (Proposition~\ref{p1:prop:tagged}). This gives a quantitative argument for audit-trail regimes that record the participant class of each order, and a quantitative argument against relying on periodic statements for the same purpose. Calibrated to Binance's published class ratios, separating the memories of the top-trader group and of all accounts in BTCUSDT would take prices alone about 333 times as long (Theorem~\ref{p1:thm:discvalue}). The class reports that the CFTC already requires have the same property in regulated markets (Table~\ref{p1:tab:cot}), so the argument concerns the granularity and frequency of existing disclosure, not a new kind of reporting.

\paragraph{When audit trails already exist.} The US Consolidated Audit Trail (SEC Rule 613) gives the SEC and the self-regulatory organizations a record of every order in NMS stocks and listed options from origination to execution or cancellation, linked to an anonymized customer identifier and to the broker-dealer that handled it; since 2025 it no longer collects the names, addresses and years of birth of natural persons \citep{p1:sec2012,p1:sec2016,p1:sec2024cat,p1:sec2025}. Our results bear on such regimes in three ways. First, they price what the record is worth for the question of who drives a feedback loop: participant-level records remove a barrier that, for an observer who learns by probing, costs round-trip trading loss of order $\delta^{-4}$, and they save samples by a factor of order $(\Delta\eta)^{-(4K-2)}$ at fixed accuracy or $(\Delta\eta)^{-(4K-4)}$ for rates resolved to a fraction of their spacing (Proposition~\ref{p1:prop:tagged}, Theorem~\ref{p1:thm:discvalue}); set against the CAT LLC 2026 budget ($\$147.0$~million in March 2026; the 2025 budget was $\$248.8$~million in November 2024, revised to $\$228.3$~million in May 2025) and the SEC's 2016 annual industry estimate ($\$1.7$~billion), the annualised weekly-renewal probing loss that disclosure makes unnecessary is small for two well-separated classes and exceeds the entire 2026 CAT budget once $K\ge3$ at $\Delta\eta\le0.10$, $K=4$ at $\Delta\eta=0.20$, or for a known pair at $\Delta\eta=0.04$ (Section~\ref{p1:sec:disc_dollars}, Table~\ref{p1:tab:cat_dollars}). For this purpose a stable account identifier is what matters and names are not needed, which is a quantitative reason to keep the account linkage while removing personal data. Second, an audit trail records accounts, not decision rules. One account can run several algorithms and one algorithm can trade through several accounts, so an observer with account-level records who wants rule-level memories faces our problem within each account's flow, with fewer cohorts but the same exponent; MiFID II's algorithm identifier \citep{p1:esma2016} is the closer analogue of the tagged flow in Proposition~\ref{p1:prop:tagged}. Third, the audit trail does not reach the markets we study. Crypto perpetual futures are outside it, futures regulated by the CFTC have separate reporting, and its data are not available to exchanges, other participants or researchers. For these observers the price-only and coarse-disclosure contracts of this chapter are the relevant ones.

\paragraph{What to monitor instead.} Because transient loss can be arbitrarily large at fixed eigenvalues and volatility (Theorem~\ref{p1:thm:blind}), supervision that wants to control risk rather than identify participants should monitor quantities that bound transient loss directly: the summed absolute impulse response and confidence sets for the transfer function, both of which are identified from public data. These are cheap where identification is expensive (Remark~\ref{p1:rem:riskneutral}); Chapter~\ref{ch:p2} shows that the summed response is certified at the parametric rate from controlled experiments, and only under a declared bound on memory length from a finite window (Chapter~\ref{ch:p2}).

\paragraph{How to measure feedback.} Empirical feedback indices should (i) compare persistence with the mechanical value of the filter, (ii) align flow with the bar in which it is observed, (iii) control for order-splitting persistence or use instruments with a defended exclusion, and (iv) report placebo pipelines. Each of these changes the sign or the order of magnitude of published-style estimates in our data.

\section{Conclusion and an open problem}\label{p1:sec:conclusion}

In a linear market with exponential-memory cohorts and known impact parameters, prices reveal a merged summary of algorithmic feedback. Recovering the algorithms behind it costs executed quantity of order $\delta^{-2}$ and trading capital of order $\delta^{-4}$ for a known pair, with a numerically evaluated constant of about 61 par units at the calibration, and a number of experiments of exact order $(\Delta\eta)^{-(4K-2)}$ for $K$ unknown cohorts, and at least that order of passive observations; in simulations of stylized nonlinear rules and of online learners trained in the closed loop the same costs hold when the competing populations' feedback weights match, populations whose learned weights differ are told apart by those weights instead, and hard thresholds hide cohorts from small probes altogether. Without the exponential form, accuracy improves only as $1/\log N$, with matching lower and upper bounds. Class-tagged disclosure removes the barrier, and its value can be stated in samples and in dollars: a sample factor of order $(\Delta\eta)^{-(4K-4)}$ for resolving rates (and $(\Delta\eta)^{-(4K-2)}$ at fixed accuracy), about 333 for the trader classes Binance already publishes in BTCUSDT, and, for an observer who learns by probing, an annualised weekly-renewal loss that matches the $\$147.0$~million 2026 CAT LLC budget across $10.0$ (Roll) or $0.41$ (one-second) markets for an unknown two-cohort Gauss--Radau pair at $\Delta\eta=0.10$ and exceeds it in a single market once $K\ge3$ at $\Delta\eta=0.10$; the CFTC's class reports restrict persistence rates that prices leave unrestricted in all eight regulated markets we examined, while a pre-registered out-of-sample test finds that class-tagged memory states do not improve 30-minute post-shock return forecasts ($\Delta R^2=-1.09$ points). At realistic noise the capital barrier for a moderate separation is about $10^5$ USDT. Measured feedback in crypto markets is real but small, rises with depth, does not move when fees change trading volume severalfold, and is easily confused with filtering and timing artefacts. Finally, eigenvalues and volatility cannot certify safety, because transient loss can be arbitrarily large at fixed spectrum and variance.

This leaves an open question for supervision: \emph{what can be measured and controlled without identifying the algorithms?} The impulse response and the transfer function, and finite-sample confidence sets for them, are identified from public data even when cohorts are not. The structural matrix $J$ is not: economies with the same merged measure can have different state dimensions, so $J$ is identified only up to the equivalence class of minimal realizations of the same transfer function. A well-posed next question is: for a stated class of stable closed-loop markets, given an estimated transfer function and a finite-sample confidence set, which intervention policy of least cost guarantees a stated bound on future price displacement without recovering individual cohorts? Chapter~\ref{ch:p2} uses exactly this interface to build an identification-free pre-trade control.

\paragraph{Data and code availability.} All data are public: Binance archives (\url{https://data.binance.vision}), verified against published checksums; the CFTC Commitments of Traders reports; and Yahoo Finance prices for the CFTC markets other than Bitcoin and Ether. The pre-registered specification of the disclosure test, its SHA-256 hash and its test-set unlock counter are in the replication package. Every number in this chapter is produced by a notebook cell or script that writes a result file (Appendix~\ref{p1:app:repro}).

\paragraph{Use of AI tools.} The authors used Claude (Anthropic), an AI assistant, to write and debug analysis code, run checks of derivations and numerical results, and draft and edit text. Every number in the paper is produced by code the authors executed and archived in the replication package. The authors reviewed all content and take full responsibility for it.

\paragraph{Disclosures.} The authors declare no relevant financial interests.

\begin{thebibliography}{99}\setlength{\itemsep}{2pt}\small
\bibitem[Alfonsi et~al.(2012)]{p1:alfonsi2012} Alfonsi, A., A. Schied, and A. Slynko (2012). Order book resilience, price manipulation, and the positive portfolio problem. \emph{SIAM Journal on Financial Mathematics} 3. doi:10.1137/110822098.
\bibitem[Alfonsi et~al.(2016)]{p1:alfonsi2016} Alfonsi, A., A. Schied, and F. Kl\"ock (2016). Multivariate transient price impact and matrix-valued positive definite functions. \emph{Mathematics of Operations Research} 41, 914--934.
\bibitem[Almgren and Chriss(2001)]{p1:almgren2001} Almgren, R., and N. Chriss (2001). Optimal execution of portfolio transactions. \emph{Journal of Risk} 3, 5--39.
\bibitem[Amilon(2008)]{p1:amilon2008} Amilon, H. (2008). Estimation of an adaptive stock market model with heterogeneous agents. \emph{Journal of Empirical Finance} 15, 342--362.
\bibitem[Batenkov et~al.(2021)]{p1:batenkov2021} Batenkov, D., G. Goldman, and Y. Yomdin (2021). Super-resolution of near-colliding point sources. \emph{Information and Inference} 10, 515--572.
\bibitem[Bernstein(1914)]{p1:bernstein1914} Bernstein, S. (1914). Sur la meilleure approximation de $|x|$ par des polyn\^omes de degr\'es donn\'es. \emph{Acta Mathematica} 37, 1--57.
\bibitem[Bessembinder et~al.(2006)]{p1:bessembinder2006} Bessembinder, H., W. Maxwell, and K. Venkataraman (2006). Market transparency, liquidity externalities, and institutional trading costs in corporate bonds. \emph{Journal of Financial Economics} 82, 251--288.
\bibitem[Boehmer et~al.(2005)]{p1:boehmer2005} Boehmer, E., G. Saar, and L. Yu (2005). Lifting the veil: An analysis of pre-trade transparency at the NYSE. \emph{Journal of Finance} 60, 783--815.
\bibitem[Bouchaud et~al.(2004)]{p1:bouchaud2004} Bouchaud, J.-P., Y. Gefen, M. Potters, and M. Wyart (2004). Fluctuations and response in financial markets: The subtle nature of ``random'' price changes. \emph{Quantitative Finance} 4, 176--190.
\bibitem[Bouchaud et~al.(2018)]{p1:bouchaud2018} Bouchaud, J.-P., J. Bonart, J. Donier, and M. Gould (2018). \emph{Trades, Quotes and Prices: Financial Markets Under the Microscope}. Cambridge University Press.
\bibitem[Brock and Hommes(1998)]{p1:brock1998} Brock, W.~A., and C.~H. Hommes (1998). Heterogeneous beliefs and routes to chaos in a simple asset pricing model. \emph{Journal of Economic Dynamics and Control} 22, 1235--1274.
\bibitem[Brogaard et~al.(2014)]{p1:brogaard2014} Brogaard, J., T. Hendershott, and R. Riordan (2014). High-frequency trading and price discovery. \emph{Review of Financial Studies} 27, 2267--2306.
\bibitem[Brunnermeier and Pedersen(2009)]{p1:brunnermeier2009} Brunnermeier, M. K., and L. H. Pedersen (2009). Market liquidity and funding liquidity. \emph{Review of Financial Studies} 22, 2201--2238.
\bibitem[Castelli and Cover(1996)]{p1:castelli1996} Castelli, V., and T.~M. Cover (1996). The relative value of labeled and unlabeled samples in pattern recognition with an unknown mixing parameter. \emph{IEEE Transactions on Information Theory} 42, 2102--2117.
\bibitem[Cochrane and Orcutt(1949)]{p1:cochrane1949} Cochrane, D., and G. H. Orcutt (1949). Application of least squares regression to relationships containing auto-correlated error terms. \emph{Journal of the American Statistical Association} 44, 32--61.
\bibitem[Colliard et~al.(2022)]{p1:colliard2022} Colliard, J.-E., T. Foucault, and S. Lovo (2022). Algorithmic pricing and liquidity in securities markets. CEPR Discussion Paper 17606.
\bibitem[Banerjee and Szydlowski(2026)]{p1:banerjee2026} Banerjee, S., and M. Szydlowski (2026). Trading against algorithms: Price dynamics and risk-sharing in a market with Q-learners. Working paper.
\bibitem[Moallemi et~al.(2008)]{p1:moallemi2008} Moallemi, C.~C., B. Park, and B. Van Roy (2008). Strategic execution in the presence of an uninformed arbitrageur. arXiv preprint arXiv:0801.3001.
\bibitem[Tremacoldi-Rossi(2025)]{p1:tremacoldi2025} Tremacoldi-Rossi, P. (2025). Manipulating algorithmic markets. Working paper.
\bibitem[Trebbi et~al.(2024)]{p1:trebbi2024} Trebbi, F., M.~B. Zhang, and M. Simkovic (2024). The cost of regulatory compliance in the United States. Working paper.
\bibitem[van Waarde et~al.(2020)]{p1:vanwaarde2020} van Waarde, H.~J., J. Eising, H.~L. Trentelman, and M.~K. Camlibel (2020). Data informativity: A new perspective on data-driven analysis and control. \emph{IEEE Transactions on Automatic Control} 65, 4753--4768.
\bibitem[Dou et~al.(2025)]{p1:dou2025} Dou, W.~W., I. Goldstein, and Y. Ji (2025). AI-powered trading, algorithmic collusion, and price efficiency. NBER Working Paper 34054.
\bibitem[Duarte and Young(2009)]{p1:duarte2009} Duarte, J., and L. Young (2009). Why is PIN priced? \emph{Journal of Financial Economics} 91, 119--138.
\bibitem[Easley et~al.(1996)]{p1:easley1996} Easley, D., N.~M. Kiefer, M. O'Hara, and J.~B. Paperman (1996). Liquidity, information, and infrequently traded stocks. \emph{Journal of Finance} 51, 1405--1436.
\bibitem[ESMA(2016)]{p1:esma2016} European Securities and Markets Authority (2016). Guidelines on transaction reporting, order record keeping and clock synchronisation under MiFID II. ESMA/2016/1452.
\bibitem[Fan(1991)]{p1:fan1991} Fan, J. (1991). On the optimal rates of convergence for nonparametric deconvolution problems. \emph{Annals of Statistics} 19, 1257--1272.
\bibitem[Leach and Madhavan(1993)]{p1:leach1993} Leach, J.\,C., and A.\,N. Madhavan (1993). Price experimentation and security market structure. \emph{Review of Financial Studies} 6, 375--404.

\bibitem[Andrews and Cheng(2012)]{p1:andrews2012} Andrews, D.\,W.\,K., and X. Cheng (2012). Estimation and inference with weak, semi-strong, and strong identification. \emph{Econometrica} 80, 2153--2211.

\bibitem[Barth and Kahn(2021)]{p1:barth2021} Barth, D., and R.\,J. Kahn (2021). Hedge funds and the Treasury cash-futures disconnect. Office of Financial Research Working Paper 21-01.

\bibitem[Forssell and Ljung(1999)]{p1:forssell1999} Forssell, U., and L. Ljung (1999). Closed-loop identification revisited. \emph{Automatica} 35, 1215--1241.
\bibitem[Franke and Westerhoff(2012)]{p1:franke2012} Franke, R., and F. Westerhoff (2012). Structural stochastic volatility in asset pricing dynamics: Estimation and model contest. \emph{Journal of Economic Dynamics and Control} 36, 1193--1211.
\bibitem[Gatheral(2010)]{p1:gatheral2010} Gatheral, J. (2010). No-dynamic-arbitrage and market impact. \emph{Quantitative Finance} 10, 749--759.
\bibitem[Glosten and Milgrom(1985)]{p1:glosten1985} Glosten, L. R., and P. R. Milgrom (1985). Bid, ask and transaction prices in a specialist market with heterogeneously informed traders. \emph{Journal of Financial Economics} 14, 71--100.
\bibitem[Hasbrouck(1991)]{p1:hasbrouck1991} Hasbrouck, J. (1991). Measuring the information content of stock trades. \emph{Journal of Finance} 46, 179--207.
\bibitem[Heinrich and Kahn(2018)]{p1:heinrich2018} Heinrich, P., and J. Kahn (2018). Strong identifiability and optimal minimax rates for finite mixture estimation. \emph{Annals of Statistics} 46, 2844--2870.
\bibitem[Hendershott et~al.(2011)]{p1:hendershott2011} Hendershott, T., C.~M. Jones, and A.~J. Menkveld (2011). Does algorithmic trading improve liquidity? \emph{Journal of Finance} 66, 1--33.
\bibitem[Hua and Sarkar(1990)]{p1:hua1990} Hua, Y., and T. K. Sarkar (1990). Matrix pencil method for estimating parameters of exponentially damped/undamped sinusoids in noise. \emph{IEEE Transactions on Acoustics, Speech, and Signal Processing} 38, 814--824.
\bibitem[Katz et~al.(2023)]{p1:katz2023} Katz, R., N. Diab, and D. Batenkov (2023). On the accuracy of Prony's method for recovery of exponential sums with closely spaced exponents. arXiv preprint arXiv:2302.05883.
\bibitem[Kirilenko et~al.(2017)]{p1:kirilenko2017} Kirilenko, A., A. S. Kyle, M. Samadi, and T. Tuzun (2017). The Flash Crash: High-frequency trading in an electronic market. \emph{Journal of Finance} 72, 967--998.
\bibitem[Kong and Valiant(2017)]{p1:kong2017} Kong, W., and G. Valiant (2017). Spectrum estimation from samples. \emph{Annals of Statistics} 45, 2218--2247.
\bibitem[Kreiss(1962)]{p1:kreiss1962} Kreiss, H.-O. (1962). \"Uber die Stabilit\"atsdefinition f\"ur Differenzengleichungen die partielle Differentialgleichungen approximieren. \emph{BIT Numerical Mathematics} 2, 153--181.
\bibitem[Kukacka and Barunik(2017)]{p1:kukacka2017} Kukacka, J., and J. Barunik (2017). Estimation of financial agent-based models with simulated maximum likelihood. \emph{Journal of Economic Dynamics and Control} 85, 21--45.
\bibitem[Kyle(1985)]{p1:kyle1985} Kyle, A. S. (1985). Continuous auctions and insider trading. \emph{Econometrica} 53, 1315--1335.
\bibitem[Le Cam(1986)]{p1:lecam1986} Le Cam, L. (1986). \emph{Asymptotic Methods in Statistical Decision Theory}. Springer.
\bibitem[Liu and Ammari(2022)]{p1:liu2023} Liu, P., and H. Ammari (2022). Super-resolution of positive near-colliding point sources. arXiv preprint arXiv:2212.00536.
\bibitem[Ljung(1999)]{p1:ljung1999} Ljung, L. (1999). \emph{System Identification: Theory for the User}, 2nd ed. Prentice Hall, Upper Saddle River, NJ.
\bibitem[Makarov and Schoar(2020)]{p1:makarov2020} Makarov, I., and A. Schoar (2020). Trading and arbitrage in cryptocurrency markets. \emph{Journal of Financial Economics} 135, 293--319.
\bibitem[Manole and Ho(2022)]{p1:manole2022} Manole, T., and N. Ho (2022). Refined convergence rates for maximum likelihood estimation under finite mixture models. \emph{Proceedings of the 39th International Conference on Machine Learning}, PMLR 162.
\bibitem[Manski(2003)]{p1:manski2003} Manski, C.~F. (2003). \emph{Partial Identification of Probability Distributions}. Springer, New York.
\bibitem[Niroula et~al.(2026)]{p1:niroula2026b} Niroula, P., S. Alkhushayni, Y. Chen, M. Bhavsar, and P. Niroula (2026). From human markets to machine markets: Transient amplification, an exact depth condition, and a maximal pre-trade guard. Working paper, Minnesota State University, Mankato.
\bibitem[Moitra(2015)]{p1:moitra2015} Moitra, A. (2015). Super-resolution, extremal functions and the condition number of Vandermonde matrices. In \emph{Proceedings of the 47th Annual ACM Symposium on Theory of Computing (STOC)}, 821--830.
\bibitem[Moskowitz et~al.(2012)]{p1:moskowitz2012} Moskowitz, T. J., Y. H. Ooi, and L. H. Pedersen (2012). Time series momentum. \emph{Journal of Financial Economics} 104, 228--250.
\bibitem[Obizhaeva and Wang(2013)]{p1:obizhaeva2013} Obizhaeva, A. A., and J. Wang (2013). Optimal trading strategy and supply/demand dynamics. \emph{Journal of Financial Markets} 16, 1--32.
\bibitem[Pagano and R\"oell(1996)]{p1:pagano1996} Pagano, M., and A. R\"oell (1996). Transparency and liquidity: A comparison of auction and dealer markets with informed trading. \emph{Journal of Finance} 51, 579--611.
\bibitem[Park and Van Roy(2015)]{p1:parkvanroy2015} Park, B., and B. Van Roy (2015). Adaptive execution: Exploration and learning of price impact. \emph{Operations Research}, doi:10.1287/opre.2015.1415.
\bibitem[Rebafka and Roueff(2015)]{p1:rebafka2015} Rebafka, T., and F. Roueff (2015). Nonparametric estimation of the mixing density using polynomials. arXiv preprint arXiv:1002.4516.
\bibitem[Rivlin(1981)]{p1:rivlin1981} Rivlin, T. J. (1981). \emph{An Introduction to the Approximation of Functions}. Dover.
\bibitem[Roll(1984)]{p1:roll1984} Roll, R. (1984). A simple implicit measure of the effective bid-ask spread in an efficient market. \emph{Journal of Finance} 39, 1127--1139.
\bibitem[Roy and Kailath(1989)]{p1:roy1989} Roy, R., and T. Kailath (1989). ESPRIT---estimation of signal parameters via rotational invariance techniques. \emph{IEEE Transactions on Acoustics, Speech, and Signal Processing} 37, 984--995.
\bibitem[CAT LLC(2026)]{p1:catbudget2026} Consolidated Audit Trail, LLC (2026). Consolidated Audit Trail, LLC 2026 Financial and Operating Budget, 31 March 2026 (total expenses \$147{,}044{,}869). \url{https://www.catnmsplan.com/cat-financial-and-operating-budget}.
\bibitem[SEC(2010)]{p1:sec2010} U.S. Securities and Exchange Commission (2010). Risk management controls for brokers or dealers with market access. Release No.\ 34-63241.
\bibitem[SEC(2012)]{p1:sec2012} U.S. Securities and Exchange Commission (2012). Consolidated Audit Trail; Final Rule (Release No.\ 34-67457, July 18, 2012, 77 FR 45722 at 45790--45795, pp.\ 220--238; proposing release No.\ 34-62174, May 26, 2010, 75 FR 32556 at 32598, 32600, pp.\ 147, 153).
\bibitem[SEC(2016)]{p1:sec2016} U.S. Securities and Exchange Commission (2016). Joint Industry Plan; Order Approving the National Market System Plan Governing the Consolidated Audit Trail. Release No.\ 34-79318, November 15, 2016, 81 FR 84696 at 84801, p.\ 363.
\bibitem[SEC(2024)]{p1:sec2024cat} U.S. Securities and Exchange Commission (2024). Notice of Filing of Proposed Rule Change Relating to the 2025 CAT Budget. Release No.\ 34-101901, December 12, 2024, p.\ 4 (revised May 2025).
\bibitem[SEC(2025)]{p1:sec2025} U.S. Securities and Exchange Commission (2025). Order granting exemptive relief from certain requirements of the National Market System Plan governing the Consolidated Audit Trail. Release No.\ 34-102386, February 10, 2025.
\bibitem[Shahram and Milanfar(2005)]{p1:shahram2005} Shahram, M., and P. Milanfar (2005). On the resolvability of sinusoids with nearby frequencies in the presence of noise. \emph{IEEE Transactions on Signal Processing} 53, 2579--2588.
\bibitem[Shleifer and Vishny(1992)]{p1:shleifer1992} Shleifer, A., and R. W. Vishny (1992). Liquidation values and debt capacity: A market equilibrium approach. \emph{Journal of Finance} 47, 1343--1366.
\bibitem[Sornette et~al.(2023)]{p1:sornette2023} Sornette, D., S. Lera, J. Lin, and K. Wu (2023). Non-normal interactions create socio-economic bubbles. \emph{Communications Physics} 6, 261.
\bibitem[Staiger and Stock(1997)]{p1:staiger1997} Staiger, D., and J. H. Stock (1997). Instrumental variables regression with weak instruments. \emph{Econometrica} 65, 557--586.
\bibitem[Trefethen and Embree(2005)]{p1:trefethen2005} Trefethen, L. N., and M. Embree (2005). \emph{Spectra and Pseudospectra: The Behavior of Nonnormal Matrices and Operators}. Princeton University Press.
\bibitem[Tsybakov(2009)]{p1:tsybakov2009} Tsybakov, A. B. (2009). \emph{Introduction to Nonparametric Estimation}. Springer.
\bibitem[Varga and Carpenter(1985)]{p1:varga1985} Varga, R. S., and A. J. Carpenter (1985). On the Bernstein conjecture in approximation theory. \emph{Constructive Approximation} 1, 333--348.
\bibitem[Wagenmaker and Jamieson(2020)]{p1:wagenmaker2020} Wagenmaker, A., and K. Jamieson (2020). Active learning for identification of linear dynamical systems. In \emph{Proceedings of the 33rd Conference on Learning Theory}, PMLR 125, 3487--3582.
\bibitem[Wagenmaker et~al.(2021)]{p1:wagenmaker2021} Wagenmaker, A.~J., M. Simchowitz, and K. Jamieson (2021). Task-optimal exploration in linear dynamical systems. In \emph{Proceedings of the 38th International Conference on Machine Learning}, PMLR 139, 10641--10652.
\bibitem[Wellman(2006)]{p1:wellman2006} Wellman, M.~P. (2006). Methods for empirical game-theoretic analysis. In \emph{Proceedings of the 21st National Conference on Artificial Intelligence (AAAI-06)}, 1552--1555.
\bibitem[Whittle(1953)]{p1:whittle1953} Whittle, P. (1953). Estimation and information in stationary time series. \emph{Arkiv f\"or Matematik} 2, 423--434.
\bibitem[Wu and Yang(2020)]{p1:wu2020} Wu, Y., and P. Yang (2020). Optimal estimation of Gaussian mixtures via denoised method of moments. \emph{Annals of Statistics} 48, 1981--2007.
\end{thebibliography}

\section{Appendix: Proofs}\label{p1:app:proofs}

\paragraph{Lemma~\ref{p1:lem:micro}.} (i) For $v_{t+1}=v_t+\zeta_{t+1}$ and $r_t=v_t+\nu_t$, the steady-state prior variance $P$ solves $P=P\sigma_\nu^2/(P+\sigma_\nu^2)+\sigma_\zeta^2$, i.e.\ $P^2-\sigma_\zeta^2P-\sigma_\zeta^2\sigma_\nu^2=0$, whose positive root is $P=\bigl(\sigma_\zeta^2+\sqrt{\sigma_\zeta^4+4\sigma_\zeta^2\sigma_\nu^2}\bigr)/2$. The gain is $\eta=P/(P+\sigma_\nu^2)\in(0,1)$, which simplifies to \eqref{p1:eq:kgain}, and the update $\hat v_{t+1}=\hat v_t+\eta(r_t-\hat v_t)$ has the stated form. (ii) $\nabla_\theta\tfrac12H(x_t-\theta)^2=-H(x_t-\theta)$, so $\theta_{t+1}=(1-\gamma)\theta_t+\alpha H(x_t-\theta_t)=(1-\eta)\theta_t+\eta qx_t$ with $\eta=\gamma+\alpha H$ and $q=\alpha H/\eta$.

\paragraph{Lemma~\ref{p1:lem:cm}.} Omit zero-weight cohorts. For $a>0$ put $b_i=\kappa q_i$ and
\[
T={\rm diag}\bigl(\sqrt a,\sqrt{b_1/\eta_1},\ldots\bigr),\qquad u=\bigl(\sqrt a,\sqrt{b_1\eta_1},\ldots\bigr)^\top .
\]
The price-first transition becomes the symmetric matrix $J'=TJT^{-1}={\rm diag}(0,r_1,\ldots)+uu^\top$, the input transforms to $\kappa u$ and $J'e_x=\sqrt a\,u$, so $h_j=(\kappa/a)e_x^\top(J')^{j+1}e_x$. The spectral decomposition of $J'$ gives nonnegative weights summing to $h_0=\kappa$, and Schur stability puts the eigenvalues in $[0,1)$. The case $a=0$ follows by continuity. A nonnegative mixture of decaying geometric sequences is completely monotone. Since $J$ is nonnegative with row sums at most one, every price row of $J^t$ has $\ell_1$, hence $\ell_2$, norm at most one, attained at $t=0$.

\paragraph{Theorem~\ref{p1:thm:merged}.} With zero initial states, $y_{i,t}=\sum_{s\le t}\eta_ir_i^{t-s}x_s$, so $\sum_iq_iy_{i,t}=\sum_{s\le t}k_{t-s}x_s$ and every response is determined by $(a,\kappa,k)$, hence by $M$. Conversely, equal responses for all probes imply equal $H$, hence equal $G$, and partial fractions of $G$ identify the distinct poles with nonzero residues $w_i$, i.e.\ $M$. After cancellation, $H=\kappa zD(z)/[(z-a)D(z)-\kappa N(z)]$ with $D=\prod(z-r_i)$ and $N(z)=\sum_iw_iz\prod_{j\ne i}(z-r_j)$, as in \eqref{p1:eq:H}; at a retained $r_i$ the denominator equals $-\kappa N(r_i)=-\kappa w_ir_i\prod_{j\ne i}(r_i-r_j)\neq0$ (memory rates are positive because $\eta_i<1$), so no root of $D$ cancels. The remaining numerator factor $z$ cancels if and only if the denominator vanishes at $0$; since $N(0)=0$ the denominator there equals $-aD(0)$, which is nonzero for $a>0$ and zero for $a=0$. Hence the McMillan degree, which is the dimension of a minimal realization, is $1+|{\rm supp}(M)\cap(0,1)|$ for $a'\neq0$ and $|{\rm supp}(M)\cap(0,1)|$ for $a'=0$ (an atom at rate $0$, a cohort with $\eta=1$, contributes the constant $M(\{0\})$ to $G$ and no pole, so it shifts $a$ to $a'=a+\kappa M(\{0\})$ and the argument above applies with $a'$) (numerically, the rank of the Hankel matrix of $(h_j)_{j\ge1}$ is $K+1$ and $K$ respectively for $K=1,2,3$; \texttt{hardening/p1\_checks.py}). For (ii), direct algebra gives
\[
G_\delta(z)-G_0(z)=-\frac{qz(z-1)\delta^2}{(z-1+\bar\eta)\bigl((z-1+\bar\eta)^2-\delta^2\bigr)} .
\]
On $|z|=1$ each factor has modulus at least $\epsilon$. Uniform stability gives $M_H=\sup|H_\theta|<\infty$, and the identity $H_\delta-H_0=z^{-1}H_\delta H_0(G_\delta-G_0)$, whose factor $z^{-1}$ has modulus one on the circle, gives an $\ell_2$ gain at most $C_{\mathcal F}\delta^2$. The mean difference of a power-constrained probe has squared norm at most $C_{\mathcal F}^2\delta^4PT$; the equal-noise KL is at most $C_{\mathcal F}^2\delta^4PT/(2\sigma_\nu^2)$, and Pinsker's inequality with total variation at least $1-2\alpha$ gives \eqref{p1:eq:Tpower}.

\paragraph{Theorem~\ref{p1:thm:quantity}.} In either positive economy $J\ge0$ and $\|J\|_\infty\le1-\epsilon(1-S)<1$, so the response is nonnegative with $\max h=\kappa$ and $\sum h=\kappa/(1-S)$. From the display above, $|G_\delta-G_0|\le2q\delta^2/\epsilon^3$ on the circle; combined with the resolvent identity and positivity this gives $\|H_\delta-H_0\|_\infty\le K_*\delta^2$, and Parseval bounds $\|h^\delta-h^0\|_2$ by the same constant. Young's inequality gives $\|\mu_\delta-\mu_0\|_2\le K_*\delta^2\|v\|_1$, and Gaussian KL with Pinsker gives \eqref{p1:eq:klV}. For (ii), the mean difference is a lower-triangular convolution matrix times $v$; each column is a shifted copy of a prefix of $d_T$, so $\|\mu_\delta-\mu_0\|_2\le V\|d_T\|_2$ with equality for a pulse at $t=0$. For two known equal-covariance Gaussian laws the minimax error is $\Phi(-\|\mu_\delta-\mu_0\|_2/(2\sigma))$, attained by the midpoint likelihood-ratio test. For (iii), the merged measures have equal mass and their first moments differ by $-q\delta^2$, which enters $h_2$ with factor $\kappa^2$ while $h_0,h_1$ coincide; a pulse $A$ separates the third-observation means by $A\kappa^2q\delta^2$. Part (iv) follows from the $\ell_2$ gain bound, and (v) from superposition and $\max_jh_j=\kappa$, which give $\sup_t|\lambda x_t|\le|\lambda|\kappa V$, attained by a single pulse.

\paragraph{Lemma~\ref{p1:lem:toeplitz}.} After the uncontrolled update, the pre-probe displacement at step $j$ is $\sum_{i<j}h_{j-i}v_i$; average linear-book execution adds $\kappa v_j/2$. Summing signed execution costs gives $\tfrac12v^\top T_hv$. Each kernel $\lambda^{|i-j|}$ is a positive-definite AR(1) covariance for $0\le\lambda<1$, and a nonnegative mixture with positive total weight is positive definite. For a round trip the fundamental cash term cancels.

\paragraph{Theorem~\ref{p1:thm:cap}.} Let $U_t=\sum_{j<t}v_j$. Positive responses give $0\le x_t\le\kappa U_t$ and $I_t=-U_t$. Summing cash and marking inventory, $B-E_t=U_tx_t-\sum_{j<t}v_jx^0_{j+1}-\tfrac\kappa2\sum_{j<t}v_j^2+cU_t$. Its quadratic part $v^\top A_tv$ has diagonal $h_{t-1-i}-\kappa/2$ and off-diagonal $\tfrac12(h_{t-1-i}+h_{t-1-j}-h_{j-i})$ for $i<j$; monotonicity of $h$ bounds every entry by $\kappa/2$. Hence $B-E_t\le\kappa U_t^2/2+cU_t\le\kappa V^2/2+cV$, while position value is at most $V(P_*+\kappa V)$; the stated inequality guarantees margin. A single pulse attains both adverse bounds at its first response, so no larger allowance is robust. Solving the quadratic gives \eqref{p1:eq:VB}. The quantity requirement $V_B\gtrsim\delta^{-2}$ then forces $B\gtrsim\delta^{-4}$, and the pulse $A$ funded with $(\mu+\tfrac12)\kappa A^2+(c+\mu P_*)A$ attains it.

\paragraph{Proposition~\ref{p1:prop:disclosure}.} The initial cash receipt is common to both alternatives and the dealer then holds $-A$. At time three $E_{3,r}=B+\kappa A^2/2-cA-A^2h_{r,2}$, so $E_{3,\delta}-E_{3,0}=A^2\kappa^2q\delta^2>0$.

\paragraph{Lemma~\ref{p1:thm:capital}.} Marking residual inventory at $P_*$ gives $E_\infty=B-\tfrac12v^\top T_hv-c\sum|v_t|\ge0$. For each mode $\lambda\le r$ the Toeplitz kernel has spectral density $(1-\lambda^2)/(1-2\lambda\cos\omega+\lambda^2)\ge(1-r)/(1+r)$, so $v^\top T_hv\ge m\|v\|_2^2$ and $\|v\|_2^2\le2B/m$. The transfer bound gives mean-difference squared norm at most $K_*^2\delta^4\|v\|_2^2$, hence the KL bound; Pinsker gives \eqref{p1:eq:Bmin}. For one pulse, cash becomes $B_A-A(P_*+\kappa A/2)-cA$ and inventory $A$; at later marks $E_t=\mu AP_*+A^2h_{t-1}$ and $\mu|I_t|P_t=\mu AP_*+\mu A^2h_{t-1}$, so margin holds for $\mu\le1$.

\paragraph{Proposition~\ref{p1:prop:tri}.} Parts (i) and (ii) follow from \eqref{p1:eq:klconv} with $\Delta h=\Theta(\delta)$ for a rate shift and $\Theta(\delta^2)$ for symmetric splitting, and from additivity of KL over independent pulse cycles. Part (iii) solves the session constraint for $\delta$; substituting $d_T(\delta)\approx\delta^2\tilde d_T$ gives the fourth-root approximation.

\paragraph{Theorems~\ref{p1:thm:spectral} and~\ref{p1:thm:adaptive}.} For a flat round trip, telescoping the cash account gives $L_i(v)=\tfrac12v^\top T_{h_i}v=\frac1{2\pi}\int g_i(\omega)|V(e^{i\omega})|^2d\omega$, with $g_i>0$ by Lemma~\ref{p1:lem:cm}. Parseval gives $\|(H_1-H_0)v\|_2^2=\frac1{2\pi}\int|\Delta\hat H|^2|V|^2\le\sup_\omega(|\Delta\hat H|^2/g_i)L_i(v)$, and the Neyman--Pearson requirement $\|(H_1-H_0)v\|_2\ge2\sigma z_{1-\alpha}$ gives \eqref{p1:eq:Cdet}. For adaptive policies, Doob's optional stopping theorem applied to the log-likelihood-ratio martingale and the data-processing inequality give $\E_0[\Lambda_\tau]\ge{\rm kl}(\alpha,1-\beta)$, while $\E_0[\Lambda_\tau]=\E_0\sum_t(\Delta\mu_t)^2/(2\sigma^2)\le\sup(|\Delta\hat H|^2/g_0)\E_0[L_0]/(2\sigma^2)$, which gives \eqref{p1:eq:adaptive}.

\emph{Part (iv).} Let $G_\theta$ be the $T\times T$ lower-triangular Toeplitz section of $H_\theta/\kappa$, so the observation covariance is $\Sigma_\theta=\sigma_w^2G_\theta G_\theta^\top+\sigma_\nu^2I$. A finite section of a bounded analytic symbol has norm at most its supremum on the circle, so $\|G_\theta\|\le\sup_\omega|H_\theta|/\kappa\le\bar M$ because $|H_\theta(e^{i\omega})|\le\sum_jh_j=\kappa/(1-S)$. The inverse of a lower-triangular Toeplitz section is the section of the reciprocal series; $\kappa/H_\theta$ is analytic on and outside the unit circle because $H_\theta$ is minimum phase (Theorem~\ref{p1:thm:minphase}), and Lemma~\ref{p1:lem:cm} gives ${\rm Re}\,H_\theta(e^{i\omega})=\sum_\ell\omega_\ell(1-\lambda_\ell\cos\omega)/(1-2\lambda_\ell\cos\omega+\lambda_\ell^2)\ge\sum_\ell\omega_\ell/(1+\lambda_\ell)\ge\kappa/2$, so $\|G_\theta^{-1}\|\le2$ and $\Sigma_\theta\succeq\sigma_{\rm eff}^2I$ (the smallest singular value of $G_\theta$ is at least $0.522$ in 762 random positive economies and sections; \texttt{hardening/p1\_checks.py}). The Gaussian divergence splits into a mean and a covariance part. The mean part is $\tfrac12\Delta\mu^\top\Sigma_1^{-1}\Delta\mu\le\|(H_1-H_0)v\|_2^2/(2\sigma_{\rm eff}^2)$. For the covariance part, the eigenvalues $\lambda_k$ of $A=\Sigma_1^{-1/2}\Sigma_0\Sigma_1^{-1/2}$ are at least $c=\sigma_\nu^2/(\sigma_\nu^2+\sigma_w^2\bar M^2)$, and $\lambda-1-\log\lambda\le(\lambda-1)^2/\lambda$, so it is at most $\|A-I\|_F^2/(2c)\le\|\Sigma_0-\Sigma_1\|_F^2/(2c\,\sigma_{\rm eff}^4)$. Writing $\Sigma_0-\Sigma_1=\sigma_w^2(\Delta G\,G_0^\top+G_1\Delta G^\top)$ and $\|\Delta G\|_F^2\le T\|h^\delta-h^0\|_2^2/\kappa^2\le TK_*^2\delta^4/\kappa^2$ (proof of Theorem~\ref{p1:thm:quantity}) gives a covariance part of at most $C_wT\delta^4$. The data-processing inequality turns errors $(\alpha,\beta)$ into $\KL\ge{\rm kl}(\alpha,1-\beta)$, and part (ii) bounds $\|(H_1-H_0)v\|_2^2$ by $\sup_\omega(|\Delta\hat H|^2/g_i)L_i(v)$, which gives \eqref{p1:eq:Cnoise}.

\paragraph{Theorem~\ref{p1:thm:rh}.} Write a feasible flat probe as $v=Eu$ and expand $u=\sum_jc_je_j$ in the $G_i$-orthonormal eigenbasis, so $L_i(v)=\sum_jc_j^2$ and the separation constraint reads $\sum_j\lambda_jc_j^2\ge c_{\rm req}$. Write the top-two component as $(c_{L-2},c_{L-1})=r(\cos\theta,\sin\theta)$ with $r\ge0$; replacing $\theta$ by $\theta+\pi$ flips $d$ and $s_\theta$ and leaves \eqref{p1:eq:rhdefs} unchanged, so $\theta\in[0,\pi)$ suffices. Put $L_0(\theta)=c_{\rm req}/\lambda(\theta)$ and $D=L_i(v)-L_0(\theta)$. The constraint gives $r^2\lambda(\theta)\ge c_{\rm req}-\sum_{j\le L-3}\lambda_jc_j^2$, hence
\[
D=r^2+\sum_{j\le L-3}c_j^2-L_0(\theta)\;\ge\;\sum_{j\le L-3}\Bigl(1-\frac{\lambda_j}{\lambda(\theta)}\Bigr)c_j^2\;\ge\;\Bigl(1-\frac{\lambda_{L-3}}{\lambda(\theta)}\Bigr)\sum_{j\le L-3}c_j^2\;\ge0,
\]
and $r^2\ge L_0(\theta)-\frac{\lambda_{L-3}}{\lambda(\theta)}\sum_{j\le L-3}c_j^2\ge L_0(\theta)\bigl(1-\gamma(\theta)D\bigr)$. By H\"older's inequality with the sign vector $s_\theta$, $c_{\rm fee}\|Eu\|_1\ge c_{\rm fee}r\|Ed(\theta)\|_1+c_{\rm fee}\sum_{j\le L-3}c_j\,s_\theta^\top Ee_j$, and Cauchy--Schwarz with weights $1-\lambda_j/\lambda(\theta)$ bounds the second term below by $-\sqrt{D\,W_\perp(\theta)}$. Using $r\ge\sqrt{L_0(1-\gamma D)_+}\ge\sqrt{L_0}\,(1-\gamma D)_+$ and the fact that the fee is nonnegative, with $x=\sqrt D$,
\[
L_i(v)+c_{\rm fee}\|v\|_1\;\ge\;L_0(\theta)+x^2+\Bigl[F_{\rm dom}(\theta)(1-\gamma(\theta)x^2)_+-\sqrt{W_\perp(\theta)}\,x\Bigr]_+\;\ge\;L_0(\theta)+\Psi(\theta).
\]
Minimizing over the unknown $\theta$ gives \eqref{p1:eq:rh}, and $\Psi\ge0$, $L_0(\theta)\ge c_{\rm req}/\lambda_{L-1}$ give the last inequality. On the piece $\gamma x^2\le1$ with $\gamma F_{\rm dom}<1$ the bracket is a convex quadratic in $x$ with minimum $F_{\rm dom}-W_\perp/(4(1-\gamma F_{\rm dom}))$, which gives the closed form.

\paragraph{Corollary~\ref{p1:cor:garch}.} Since $\Delta h_t=O(\delta^2)$, the variance difference $\sigma^2_{1,t}-\sigma^2_{0,t}=\sigma^2\xi[(h_1+h_0)*v]_t[\Delta h*v]_t=O(\delta^2)$, and the variance term of the Gaussian KL, $\tfrac12\sum(x-1-\log x)$ with $x=\sigma^2_{0,t}/\sigma^2_{1,t}$, is quadratic in that difference. Convex impact only increases dissipation.

\paragraph{Proposition~\ref{p1:prop:crb}.} Standard Cram\'er--Rao for a Gaussian mean model with Jacobian $J$. For $K=2$ the columns of $J$ for the two rates become collinear as $\Delta\eta\to0$; expanding to third order, the smallest singular value of $J$ is of order $(\Delta\eta)^3$, so $[F^{-1}]_{kk}=\Theta((\Delta\eta)^{-6})$.

\paragraph{Theorem~\ref{p1:thm:twocohort}.} Expanding $M_m(\theta_1)-M_m(\theta_0)$ under \eqref{p1:eq:pair} gives zero for $m\le2$ and $4\tau p_-p_+(\Delta\eta)^3$ for $m=3$. The map from moments to impulse responses is analytic with nonsingular derivative at the common limit, so $\|\Delta h\|_2=\Theta(\tau(\Delta\eta)^3)$. Le Cam's two-point lemma with $N$ independent Gaussian experiments gives \eqref{p1:eq:lecam2}. For the permutation-invariant loss, the identity matching of $\theta_0$ and $\theta_1$ costs $2\tau+(p_+-p_-)\Delta\eta=2\tau(1+\Delta\eta/q)$ per atom and the swapped matching costs $\max(2p_+,2p_-)\Delta\eta=(1+2\tau/q)\Delta\eta$; when both are at least $2\tau$, any estimate within $\tau$ of one economy in $d_\infty$ is farther than $\tau$ from the other by the triangle inequality, and the two-point argument is unchanged. For (iii), the upper bound is Theorem~\ref{p1:prop:esprit}; the Prony--Hankel determinant identity follows from $g_m=w_1\eta_1r_1^m+w_2\eta_2r_2^m$.

\paragraph{Theorem~\ref{p1:thm:kcohort}.} Let $(s^A,\omega^A)$ and $(s^B,\omega^B)$ be the $K$-point Gauss and Gauss--Radau rules for the uniform weight on $[-1,1]$, normalized to unit mass. Both are positive and exact for polynomials of degree $2K-2$, and their moments of degree $2K-1$ differ by $g_K\neq0$. With $r=r_c+\Delta\eta\,s$ and weights $W\omega$, the raw moments $\sum w_ir_i^m$ are polynomials in $\Delta\eta$ that agree below degree $2K-1$, so the kernels differ by $O((\Delta\eta)^{2K-1})$ with nonzero leading coefficient; analyticity of the kernel-to-response map gives $\|h^A-h^B\|_2\asymp(\Delta\eta)^{2K-1}$, and Le Cam's lemma gives the bound. For the $W_1$ version, the two merged measures are the images of the fixed normalized rules under $s\mapsto r_c+\Delta\eta\,s$ scaled to mass $W$, so their distance is $W\Delta\eta$ times the distance $2c_K>0$ between the rules, and the same two-point argument applies with the metric loss.

\paragraph{Lemma~\ref{p1:lem:altpair}.} Let $L=2n-2$ and let $p^*$ be the best uniform approximation of $|x|$ on $[-1,1]$ by polynomials of degree $L$, with error $E=E_L(|x|)$. By Chebyshev's equioscillation theorem \citep{p1:rivlin1981} there are points $t_0<\cdots<t_{L+1}$ at which $|x|-p^*(x)$ equals $\pm E$ with alternating signs. The divided-difference weights $w_i=1/\prod_{j\ne i}(t_i-t_j)$ alternate in sign and satisfy $\sum_iw_ip(t_i)=0$ for every polynomial $p$ of degree at most $L$; multiplying by $-1$ if necessary, $w_i$ has the sign of $|t_i|-p^*(t_i)$ for every $i$. Exactly $n$ weights are positive and $n$ negative, and the positive and negative parts have equal mass because constants are annihilated. Normalizing each part to unit mass and mapping $x\mapsto\rho(x+1)/2$ gives $n$-atom probability measures $\pi_A$ (positive part) and $\pi_B$ (negative part) on $[0,\rho]$ with equal moments of orders $0,\ldots,L$. For the 1-Lipschitz function $f(r)=|r-\rho/2|=(\rho/2)|x|$,
\[
\int f\,d(\pi_A-\pi_B)=\frac\rho2\int\bigl(|x|-p^*(x)\bigr)\,d(\pi_A-\pi_B)=\frac\rho2(E+E)=\rho E ,
\]
because $p^*$ integrates to the same value under both measures, $|x|-p^*=E$ on the support of $\pi_A$ and $-E$ on that of $\pi_B$. By Kantorovich--Rubinstein duality $W_1(\pi_A,\pi_B)\ge\rho E\ge c_*\rho/L$.

\paragraph{Theorem~\ref{p1:thm:general}.} Write $K(z)=\sum_mk_mz^m$, $D=1-az-\kappa zK$ and $H=\kappa/D$, which is the generating function of the recursion $h_t=ah_{t-1}+\kappa\sum_mk_mh_{t-1-m}$ with $h_0=\kappa$. Since $k_m\ge0$ and $\sum_mk_m=\int(1-r)^{-1}d\mu\le W/(1-\rho)$, the $\ell_1$ norm of $az+\kappa zK$ is at most $s<1$, so $\|H\|_1\le\kappa/(1-s)$ and $\|1/H\|_1=\|D\|_1/\kappa\le(1+s)/\kappa$.

(i) Let $\mu_A=W\pi_A$ and $\mu_B=W\pi_B$ with $\pi_A,\pi_B$ from Lemma~\ref{p1:lem:altpair}; both lie in $\mathcal M(\rho,W)$. With $\nu=\mu_A-\mu_B$, $\sum_m(k^A_m-k^B_m)z^m=\int(1-rz)^{-1}d\nu(r)=z^{2n-1}\int r^{2n-1}(1-rz)^{-1}d\nu(r)$, because the first $2n-1$ terms of the geometric series integrate to zero. On $|z|=1$ the modulus is at most $2W\rho^{2n-1}/(1-\rho)$, so Parseval gives $\|k^A-k^B\|_2\le2W\rho^{2n-1}/(1-\rho)$. Next, $H_A-H_B=z(K_A-K_B)H_AH_B$, so by Young's inequality $\|h^A-h^B\|_2\le\|H_A\|_1\|H_B\|_1\|k^A-k^B\|_2\le C_h\|k^A-k^B\|_2$. The Kullback--Leibler divergence between the $N$-sample Gaussian experiments is $N\|h^A-h^B\|_2^2/(2\sigma^2)\le(N/\sigma^2)\,2C_h^2W^2\rho^{2(2n-1)}/(1-\rho)^2$, which is at most one as soon as $2(2n-1)\log(1/\rho)\ge\Lambda_N-2\log(1/\rho)$. Take $n$ the smallest integer with this property and $n\ge2$; then $2n-2<\Lambda_N/(2\log(1/\rho))$ when $\Lambda_N>4\log(1/\rho)$. Since $d$ is a norm on signed measures and $d(\mu_A,\mu_B)=W\,W_1(\pi_A,\pi_B)\ge Wc_*\rho/(2n-2)>2c_*\rho W\log(1/\rho)/\Lambda_N$, the Bretagnolle--Huber form of Le Cam's lemma \citep[Lemma 2.6 and Theorem 2.2]{p1:tsybakov2009} gives $\max_{j\in\{A,B\}}P_j(d(\hat\mu,\mu_j)\ge c_*\rho W\log(1/\rho)/\Lambda_N)\ge\tfrac14e^{-1}$.

(ii) \emph{Approximation.} Let $f$ be 1-Lipschitz on $[0,\rho]$ with $f(0)=0$, so $|f|\le\rho$. The function $g(x)=f(\rho(x+1)/2)$ is $(\rho/2)$-Lipschitz on $[-1,1]$, and Jackson's theorem \citep{p1:rivlin1981} gives a polynomial $p$ of degree $L$ with $\sup|g-p|\le(\pi/2)(\rho/2)/(L+1)\le\rho$, hence $\sup|p|\le2\rho$. Its Chebyshev coefficients satisfy $|b_j|\le2\sup|p|\le4\rho$. The polynomial $q(r)=p(2r/\rho-1)$ approximates $f$ with the same error, and its monomial coefficients satisfy $\sum_j|q_j|\le\sum_{j\le L}|b_j|\,\|T_j(2r/\rho-1)\|_{\rm coef}\le4\rho\sum_{j\le L}B(\rho)^j\le\rho B(\rho)^{L+1}$, where $\|\cdot\|_{\rm coef}$ is the sum of absolute monomial coefficients, $\|T_j\|_{\rm coef}\le(1+\sqrt2)^j$, substituting $y=\alpha r-1$ multiplies the coefficient sum of $y^i$ by at most $(1+\alpha)^i$, and $B(\rho)-1\ge4$. \emph{Moment comparison.} If $\max_{j\le L}|\int r^jd\hat\mu-k_j|\le2\delta$, then $\hat\mu$ has mass at most $W+2\delta$ and
\[
\int f\,d(\hat\mu-\mu)=\int(f-q)\,d(\hat\mu-\mu)+\sum_{j\le L}q_j\Bigl(\int r^jd\hat\mu-k_j\Bigr)\le\frac{\pi\rho(W+\delta)}{2(L+1)}+2\delta\rho B(\rho)^{L+1}.
\]
\emph{Kernel recovery.} From $\kappa zK=1-az-\kappa/H$, the kernel values $k_0,\ldots,k_L$ are the coefficients of $(1-az-\kappa/H)/(\kappa z)$, which depend on $h_0,\ldots,h_{L+1}$ only. Work in the algebra of power series modulo $z^{L+2}$ with the $\ell_1$ norm, which is submultiplicative. If the averaged response is $\hat H=H+E$ and $\|1/H\|_1\|E\|_1\le\tfrac12$, then $1/\hat H=(1/H)(1+E/H)^{-1}$ and $\|1/\hat H-1/H\|_1\le2\|1/H\|_1^2\|E\|_1$, so $\max_j|\hat k_j-k_j|\le\delta\equiv2((1+s)/\kappa)^2(L+2)\max_{t\le L+1}|E_t|$. Each $E_t$ is Gaussian with variance $\sigma^2/N$, so with probability at least $1-\beta$, $\max_t|E_t|\le\sigma\sqrt{2\log(2(L+2)/\beta)/N}$. The true $\mu$ is feasible for the linear program with value at most $\delta$, so the minimizer satisfies the moment condition above. \emph{Choice of $L$.} With $L+1=\lfloor\log(N/\sigma^2)/(4\log B(\rho))\rfloor$, $B(\rho)^{L+1}\le(N/\sigma^2)^{1/4}$, so $\delta B(\rho)^{L+1}=O((\sigma^2/N)^{1/4}L\sqrt{\log(L/\beta)})$ and $\delta/(L+1)$ are $o(1/\log(N/\sigma^2))$, while $\pi\rho W/(2(L+1))\le2\pi\rho W\log B(\rho)/(\log(N/\sigma^2)-4\log B(\rho))$. For $N/\sigma^2$ larger than a threshold depending on $(\rho,W,a,\kappa,\beta)$ the side condition $\|1/H\|_1\|E\|_1\le\frac12$ holds on the event and the total is below $3\pi\rho W\log B(\rho)/\log(N/\sigma^2)$, uniformly over $\mathcal M(\rho,W)$. The final statement combines (i) and (ii).

\paragraph{Theorem~\ref{p1:prop:passive}.} For stationary Gaussian processes with spectral densities $f_j\propto|H_j(e^{i\omega})|^2$, the divergence rate is
\[
D=\frac1{4\pi}\int_{-\pi}^{\pi}\bigl(\varrho-1-\log\varrho\bigr)\,d\omega,\qquad \varrho=f_1/f_0=|H_1/H_0|^2 .
\] By the resolvent identity, $H_1/H_0-1=z^{-1}H_1\Delta G$ with $\Delta G=G_1-G_0$, and uniform stability bounds $|H_1|$ on the circle, so $u\equiv\varrho-1=2\,{\rm Re}(z^{-1}H_1\Delta G)+O(|\Delta G|^2)$. Write $G(z)=\sum_kw_k\phi(z,\eta_k)$, with $w_k$ the cohort weights (any factor of $\eta_k$ is absorbed in $\phi$), where $\phi$ is analytic in $\eta$ on a neighbourhood of $[\eta_{\min},\eta_{\max}]$ uniformly for $|z|=1$ (its pole $z=1-\eta$ stays inside the disk). Expanding $\phi$ around $\bar\eta$, $\Delta G=\sum_{j\ge0}\frac{\partial_\eta^j\phi(z,\bar\eta)}{j!}\,\Delta\tilde M_j$, where $\tilde M_j=\sum_kw_k(\eta_k-\bar\eta)^j$. For the pair, $\Delta\tilde M_0=\Delta\tilde M_1=0$ and $\Delta\tilde M_2=q\delta^2$; for the moment-matched economies, $\Delta\tilde M_j=0$ for $j\le2$ and $\Delta\tilde M_3=4\tau p_-p_+(\Delta\eta)^3$ by \eqref{p1:eq:dM3}; for the $K$-cohort Gauss and Gauss--Radau economies with rates $\bar\eta+\Delta\eta\,s$ and weights $W\omega$, both rules integrate polynomials of degree $2K-2$ exactly, so $\Delta\tilde M_j=W(\Delta\eta)^j\sum_k(\omega^A_k(s^A_k)^j-\omega^B_k(s^B_k)^j)$ vanishes for $j\le2K-2$ and equals $\pm Wg_K(\Delta\eta)^{2K-1}\neq0$ at $j=2K-1$. Higher terms are smaller by a factor of the spacing: $|\Delta\tilde M_j|\le2W(\Delta\eta)^j$, and analyticity of $\phi$ in $\eta$ on a disk of radius $\varrho>0$ around $\bar\eta$, uniformly on the circle, gives $|\partial_\eta^j\phi|/j!\le C\varrho^{-j}$, so for $\Delta\eta<\varrho/2$ the tail beyond the first nonzero term is $O((\Delta\eta)^{j_0+1})$. Hence $\sup_\omega|u|=O(\delta^2)$, $O(\tau(\Delta\eta)^3)$, respectively $O((\Delta\eta)^{2K-1})$, and $u=c(\omega)\epsilon+o(\epsilon)$ with $\epsilon=\delta^2$, $\tau(\Delta\eta)^3$ or $(\Delta\eta)^{2K-1}$ and $c(\omega)$ proportional to $2\,{\rm Re}\bigl(z^{-1}H_0\,\partial^{j_0}_\eta\phi(\cdot,\bar\eta)\bigr)/j_0!$, $j_0=2$, $3$ or $2K-1$. The function $c$ does not vanish identically: $f=z^{-1}H_0\partial^j_\eta\phi$ is analytic outside the disk including at infinity, where it vanishes, and is not identically zero; an analytic function on the exterior whose real part vanishes on the circle is an imaginary constant, which here would have to be zero. For $|u|\le\frac12$, $u^2/4\le u-\log(1+u)\le u^2$, so $D\asymp\int c(\omega)^2d\omega\,\epsilon^2$, which gives (i)--(iii). The sample requirements need a finite-$T$ bound that is uniform in the separation, which we obtain directly. The first $T$ observations are $N(0,\Sigma_j)$ with $\Sigma_j$ the $T\times T$ Toeplitz matrices of the spectral densities $f_j=\sigma^2|H_j(e^{i\omega})|^2$. On the compact, uniformly stable family the transfer functions have no zeros or poles on the circle (Theorem~\ref{p1:thm:minphase}), so $0<\underline f\le f_j\le\bar f<\infty$ with constants that do not depend on the separation. A Toeplitz matrix has eigenvalues within the range of its symbol, so $\lambda_{\min}(\Sigma_0)\ge\underline f$ and $\|\Sigma_1-\Sigma_0\|\le\sup_\omega|f_1-f_0|=\sup_\omega f_0|u|\le\bar f\sup_\omega|u|$. With $A=\Sigma_0^{-1/2}\Sigma_1\Sigma_0^{-1/2}$, whose eigenvalues lie in $[\underline f/\bar f,\bar f/\underline f]$, the inequality $\lambda-1-\log\lambda\le(\lambda-1)^2/\lambda$ gives
\[
\KL(P_1^T\|P_0^T)=\tfrac12\sum_k(\lambda_k-1-\log\lambda_k)\le\frac{\bar f}{2\underline f}\,\|A-I\|_F^2\le\frac{\bar f}{2\underline f}\,T\Bigl(\frac{\bar f}{\underline f}\Bigr)^2\sup_\omega|u|^2\le C\,T\epsilon^2,
\]
with $C$ independent of $T$ and of the separation, because $\sup_\omega|u|=O(\epsilon)$ uniformly. The Bretagnolle--Huber inequality \citep[Lemma 2.6]{p1:tsybakov2009} then gives the stated lower bounds on $T$ for every $T$, not only asymptotically. For a fixed pair, $T^{-1}\KL(P_1^T\|P_0^T)\to D$ and Stein's lemma (one error fixed) and the Chernoff bound (both errors vanishing) give the error exponents of the likelihood-ratio test, of the same order in the separation; that statement is pointwise in the pair.

\paragraph{Theorem~\ref{p1:thm:minphase}.} In the delay variable $\lambda=z^{-1}$, $H=\kappa\prod_i(1-r_i\lambda)/\bigl[(1-a\lambda)\prod_i(1-r_i\lambda)-\kappa\lambda\sum_iw_i\prod_{j\ne i}(1-r_j\lambda)\bigr]$, which is \eqref{p1:eq:H} divided through by $z^{K+1}$; its nonstructural zeros are the memory rates, which do not cancel under the nondegenerate assumptions of Theorem~\ref{p1:thm:minphase}. In the z-plane there is also the structural timing zero at zero: the numerator is proportional to $z\prod_i(z-r_i)$. Stability puts the poles inside the disk. A rational spectral density $\sigma^2|H(e^{i\omega})|^2$ has a unique minimum-phase factor up to a unimodular constant, which is $H$; Theorem~\ref{p1:thm:merged} then gives the merged measure.

\paragraph{Proposition~\ref{p1:prop:tagged}.} With tagged flow, $u_{i,t}=r_iu_{i,t-1}+w_ix_t$ holds exactly and involves only cohort $i$, so $\tilde u_{i,t}=r_i\tilde u_{i,t-1}+w_ix_t+e_{i,t}$ with the MA(1) error $e_{i,t}$. Prices are independent of the measurement noise, so $\E[Z_te_{i,t}]=0$ and the IV estimator is consistent when $\Sigma_i=\E[Z_t(\tilde u_{i,t-1},x_t)]$ is nonsingular. Since $u_{i,t-1}=w_i\sum_{s\ge0}r_i^sx_{t-1-s}$, $\E[x_{t-1}u_{i,t-1}]=w_iA_0$ and $\E[x_tu_{i,t-1}]=w_iA_1$, which gives $\Sigma_i$. The price is a nonnegative mixture of geometric responses to its innovation (Lemma~\ref{p1:lem:cm}), so $\gamma_x(s)=\sum_\ell c_\ell\lambda_\ell^{s}$ with $c_\ell=\sigma_\xi^2\kappa^{-2}\omega_\ell\sum_m\omega_m/(1-\lambda_\ell\lambda_m)\ge0$; substituting, $\det\Sigma_i=w_i(\gamma_x(0)A_0-\gamma_x(1)A_1)=w_i\sum_{\ell,m}c_\ell c_m(1-\lambda_\ell\lambda_m)/(1-r_i\lambda_m)$, each term positive. The scores $Z_te_{i,t}$ form a martingale-difference-plus-MA(1) array (prices are independent of the measurement noise), so a martingale central limit theorem applies; their covariances are with $\E[Z_tZ_t^\top e_{i,t}^2]=\sigma_f^2(1+r_i^2)\Gamma_0$ and $\E[Z_tZ_{t-1}^\top e_{i,t}e_{i,t-1}]=-r_i\sigma_f^2\Gamma_1$, so their long-run variance is $\Omega_i$, and writing $\sum_tZ_te_{i,t}=\sum_t(Z_t-r_iZ_{t+1})\varepsilon_{i,t}$ plus a boundary term, the summands are martingale differences with respect to $\mathcal F_t=\sigma((x_s)_{s\in\mathbb Z},\varepsilon_{i,s}:s\le t)$, because prices are independent of the measurement noise, and the martingale central limit theorem gives \eqref{p1:eq:ivvar}. On a uniformly stable family the modes and weights lie in a compact set with $\lambda_\ell\le\bar\lambda<1$, so $\det\Sigma_i\ge w_{\min}(1-\bar\lambda^2)\gamma_x(0)^2>0$ and $\|\Omega_i\|$ is bounded, uniformly in the spacing.

\paragraph{Theorem~\ref{p1:thm:blind}.} With $J_\theta$ upper triangular the eigenvalues are fixed. The stationary covariance solves $P=J_\theta PJ_\theta^\top+{\rm diag}(s_1,s_2)$, and $P_{11}$ is affine in $s_2$ with positive slope, so for each $\theta$ a unique $s_2(\theta)\ge0$ gives $P_{11}=V$ when $V$ exceeds the value at $s_2=0$. The one-step price variance is $s_1$ throughout, and $e_1^\top J_\theta e_2=\theta$ gives $\sup_k\|e_1^\top J_\theta^k\|\ge|\theta|$.

\paragraph{Proposition~\ref{p1:prop:choice}.} (i) $\mathcal L_{\rm disc}<\mathcal L_{\rm price}$ is $F+Mc<M\ell G$. (ii) $T_p$ and $T_t$ are decreasing in $\Delta\eta$, with $T_t\ge\bar T$ exactly when $\Delta\eta\le\Delta\eta_t$ and $T_p\ge\bar T$ exactly when $\Delta\eta\le\Delta\eta_p$. On $\Delta\eta\le\Delta\eta_t$ both minima equal $\bar T$ and $G=0$. On $(\Delta\eta_t,\Delta\eta_p]$, $G=\bar T-c_t(\Delta\eta)^{-2}$, which increases to $G^*=\bar T-T_t(\Delta\eta_p)>0$. On $\Delta\eta>\Delta\eta_p$, $G=c_p(\Delta\eta)^{-(4K-2)}-c_t(\Delta\eta)^{-2}=(\Delta\eta)^{-2}\bigl[c_p(\Delta\eta)^{-(4K-4)}-c_t\bigr]$; the bracket is positive exactly when $\Delta\eta<\Delta\eta_0=(c_p/c_t)^{1/(4K-4)}$, and $G'(\Delta\eta)=(\Delta\eta)^{-3}\bigl[2c_t-(4K-2)c_p(\Delta\eta)^{-(4K-4)}\bigr]<0$ whenever $c_p(\Delta\eta)^{-(4K-4)}>c_t$, because $4K-2>2$. So $G$ decreases wherever it is positive beyond $\Delta\eta_p$, and is non-positive beyond $\Delta\eta_0$. Every superlevel set of $G$ at a positive level is therefore an interval (the negative tail, where $G$ rises back toward zero from below, never reaches a positive level), and the level $(F/M+c)/\ell>0$ lies above $G=0$, so the set where (i) holds is an interval contained in $(\Delta\eta_t,\Delta\eta_0)$ that contains $\Delta\eta_p$ whenever it is non-empty. (iii) The interval is non-empty if and only if the level is below the maximum, $F/M+c<\ell G^*$, that is $M>M^*$; the level decreases in $M$, so the superlevel set grows. $\square$

\section{Appendix: Notation}\label{p1:app:notation}

\begin{longtable}{p{0.18\textwidth}p{0.72\textwidth}}
\toprule
Symbol & Meaning\\
\midrule
\endhead
$x_t$ & pricing error (price minus fundamental)\\
$a$ & price persistence (anchor), $0\le a<1$\\
$\kappa$ & price impact per unit of executed quantity\\
$q_i,\ \eta_i,\ r_i$ & cohort weight, adaptation rate, memory rate $r_i=1-\eta_i$\\
$w_i=q_i\eta_i$ & merged weight\\
$M=\sum_iw_i\delta_{r_i}$ & merged rate--weight measure\\
$k_m,\ G(z)$ & feedback kernel and its transform\\
$H(z),\ h_j$ & closed-loop transfer function and impulse response, $h_0=\kappa$\\
$J,\ S$ & state transition matrix and $S=a+\kappa\sum q_i$\\
$v_t,\ V$ & executed probe quantity and total quantity $\sum|v_t|$\\
$d_T$ & impulse-response difference between hypotheses, first $T$ taps\\
$\sigma$ & observation-noise standard deviation\\
$\delta,\ \Delta\eta$ & half-separation of a known pair; spacing of unknown cohorts\\
$K_*$ & transfer-difference constant $2\kappa^2q/[\epsilon^3(1-S)^2]$\\
$T_h$ & Toeplitz matrix $(h_{|i-j|})$ of execution cost\\
$g_i(\omega)$ & dissipation kernel ${\rm Re}\,\hat H_i(e^{i\omega})-\kappa/2$\\
$B,\ \mu,\ c,\ \phi$ & capital, maintenance-margin fraction, proportional cost, taker fee\\
$V_{\rm eff}$ & per-trade execution capacity $\min(V_{\max},D_{\max}/\kappa)$\\
$\hat Q,\ \hat S$ & reduced-form loop gain and persistence indices\\
\bottomrule
\end{longtable}

\section{Appendix: Further results}\label{p1:app:ia}

\subsection{A finite-window price-surveillance bound}

\begin{theorem}[Finite-window abstention]\label{p1:thm:abstain}
Observe $x_t=sx_{t-1}+w_t$, $w_t\sim N(0,\sigma^2)$, $x_0=0$, $t\le T$, with $s\in(0,2)$. Let a possibly randomized rule $\Psi$ output Stable, Unstable or Abstain, with $\sup_{s<1}P_s(\Psi={\rm Unstable})\le\alpha$ and $\sup_{s\ge1}P_s(\Psi={\rm Stable})\le\alpha$, $\alpha<1/2$. Choose $0<c_0<1$ with $c_0e^{c_0}/(2\sqrt2)<1-2\alpha$. On $B_T=(1-c_0/T,1+c_0/T)$,
\[
\inf_{s\in B_T}P_s(\Psi={\rm Abstain})\;\ge\;1-2\alpha-\frac{c_0e^{c_0}}{2\sqrt2}>0 ,
\]
and for every finite $T$ the minimax worst-case error of binary rules over $s\in(0,2)$ is at least $1/2$.
\end{theorem}

\begin{proof}
Pair $s_0\in(1-c_0/T,1)$ with $s_1=s_0+c_0/T$. Then $\KL(P_{s_0}\|P_{s_1})=\frac{(s_1-s_0)^2}{2\sigma^2}\sum_t\E x_{t-1}^2\le c_0^2e^{2c_0}/4$, so Pinsker bounds total variation by $c_0e^{c_0}/(2\sqrt2)$. Error control gives $P_{s_0}({\rm Stable})\le\alpha+{\rm TV}$ and $P_{s_0}({\rm Unstable})\le\alpha$, which yields the abstention bound. For binary rules, continuity in total variation and Le Cam's two-point bound give worst-case error tending to $1/2$.
\end{proof}

The interval $B_T$ is an abstention layer that any error-controlled monitor must respect in the scalar Gaussian case; richer disclosure can remove it.

\subsection{Passive distinction of a known pair in finite samples}

\begin{proposition}[Passive distinction with zero controlled quantity]\label{p1:prop:passive0}
Keep the known pair, but add independent uniform innovations $w_t\in[-b,b]$ of variance $s_w^2=b^2/3$ at each price update, observe $\tilde x_t=x_t+\nu_t$ with $\nu_t\sim N(0,\sigma^2)$, use independent three-step reset blocks and no probes, with $P_*>3b$. Let $b_1=a+\kappa q\eta$, $b_{2,0}=b_1^2+\kappa q\eta(1-\eta)$ and $b_{2,\delta}=b_{2,0}-\kappa q\delta^2$. In each block the covariance of the first and third observations differs by $\Delta=\kappa qs_w^2\delta^2$. With
\[
R=\max_{r\in\{0,\delta\}}\Bigl\{(s_w^2+\sigma^2)\bigl[s_w^2(1+b_1^2+b_{2,r}^2)+\sigma^2\bigr]+s_w^4b_{2,r}^2\Bigr\},
\]
the midpoint test of the average product has worst-case error at most $\alpha$ whenever $n\ge4R/(\alpha\Delta^2)$ blocks are observed: zero controlled quantity and order $\delta^{-4}$ blocks.
\end{proposition}

The Chebyshev bound is conservative but the order is real. With $b=0.03$, observation s.d.\ $0.01$ and $\delta=0.10$, the bound asks for $3.9\times10^7$ blocks at $\alpha=0.1$; in simulation the midpoint test still errs 32--38\% of the time with $2\times10^5$ blocks. Independent resets and known innovation laws are extra conditions. The proposition does not show that anonymous public records identify participants; it shows that passive excitation can accumulate memory information while the probe account stays untouched.

\begin{proof} In each block $x_1=w_1$, $x_2=b_1w_1+w_2$ and $x_3=b_{2,r}w_1+b_1w_2+w_3$, so $|x_t|\le3b$. Observation noise changes variances but not the covariance $s_w^2b_{2,r}$. The product variance is the product of marginal variances plus the squared covariance plus a fourth-cumulant term $b_{2,r}^2(\E w^4-3s_w^4)\le0$ for uniform innovations, so it is at most $R$. Chebyshev bounds each midpoint error by $4R/(n\Delta^2)$.
\end{proof}

\subsection{Blind inversion of the closed loop}\label{p1:app:era}
\begin{enumerate}[itemsep=1pt]
\item \emph{Impulse decoupling:} $g_t=h_{t+1}-ah_t$.
\item \emph{Minimal realization:} form the Hankel matrix $H_{ij}=g_{i+j-1}$ and run Kung's eigensystem realization with a truncated SVD.
\item \emph{Transmission zeros:} compute the zeros of the realization, which under known $(\kappa,a)$ isolate the roots of the denominator of $G$.
\item \emph{Weights:} with recovered rates, project $u_t=(x_t-ax_{t-1})/\kappa$ on the reconstructed filters $y_{k,t-1}$.
\end{enumerate}
Without noise each step is exact: the realization has the order of the true system, its nonstructural transmission zeros are the memory rates (Theorem~\ref{p1:thm:merged}); the known timing zero is discarded, and the projection returns the weights, so the procedure removes the pole-shift bias of fitting poles directly to the price response. With noise, per-cohort recovery degrades as rates move together at the rate given by Theorems~\ref{p1:thm:kcohort} and~\ref{p1:prop:esprit}, which is why we do not report it as an estimator of individual cohorts.

\subsection{Probing protocols}
Table~\ref{p1:tab:protocols} compares four probing protocols at observation horizon $T=4$ per cycle.

\begin{table}[!htbp]\centering\small
\begin{threeparttable}
\caption{Probing protocols at observation horizon $T=4$ per cycle}\label{p1:tab:protocols}
\begin{tabular}{lcrrrr}
\toprule
Protocol & $\delta$ & Executed vol. & Max trade & Error (SE) & Capital\\
\midrule
A: single pulse (unconstrained) & 0.03 & 131.5 & 131.5 & 0.047 (0.003) & 6{,}985.91\\
 & 0.05 & 47.3 & 47.3 & 0.054 (0.004) & 921.73\\
 & 0.08 & 18.5 & 18.5 & 0.053 (0.004) & 146.73\\
 & 0.12 & 8.2 & 8.2 & 0.044 (0.003) & 31.45\\
B: repeated capped pulses & 0.03 & 2{,}768.8 & 6.2 & 0.050 (0.003) & 7{,}035.75\\
 & 0.05 & 362.5 & 6.2 & 0.054 (0.004) & 923.88\\
 & 0.08 & 56.2 & 6.2 & 0.047 (0.003) & 146.00\\
 & 0.12 & 12.5 & 6.2 & 0.038 (0.003) & 34.88\\
C: sequential test, capped & 0.03 & 1{,}410.2 & 6.2 & 0.045 (0.003) & 3{,}585.14\\
 & 0.05 & 195.6 & 6.2 & 0.038 (0.003) & 500.00\\
 & 0.08 & 35.9 & 6.2 & 0.032 (0.003) & 94.30\\
 & 0.12 & 10.5 & 6.2 & 0.011 (0.002) & 29.85\\
D: generalized eigenvector & 0.03 & 2{,}359.7 & 3.7 & 0.050 (0.003) & 2{,}632.77\\
 & 0.05 & 314.8 & 3.7 & 0.048 (0.003) & 353.67\\
 & 0.08 & 49.2 & 3.7 & 0.046 (0.003) & 57.53\\
 & 0.12 & 9.9 & 3.8 & 0.045 (0.003) & 13.56\\
\bottomrule
\end{tabular}
\begin{tablenotes}\footnotesize\item 2{,}000 replications per hypothesis and cell. Log--log capital slopes: $-3.90$, $-3.84$, $-3.47$, $-3.81$. Protocols A and B reproduce the earlier version exactly. C and D are re-implemented because the earlier code was not retained: C is a Wald test with capped pulses; D repeats the top generalized eigenvector of a flat $(T+1)$-step probe, scaled to a maximum trade of 3.75, and is cheaper than the earlier D. Kind: \lab{recomputed} (\texttt{recompute\_p1\_protocols.py}).\end{tablenotes}
\end{threeparttable}
\end{table}

\subsection{Coordinate-weighted output}\label{p1:app:coordw}
For $J=\begin{psmallmatrix}a&\kappa q\\ \eta a&1-\eta+\eta\kappa q\end{psmallmatrix}$ with $a=0.5$, $\kappa=0.8$, $q=0.2$, $\eta=0.3$ ($\rho(J)=0.8224$, non-normal), the price output has $\sup_t\|e_x^\top J^t\|=1$. The weighted output $C_w=(1,\lambda/{\rm MMR})=(1,160)$ with $\lambda=0.8$ and ${\rm MMR}=0.005$ has $\sup_t\|C_wJ^t\|=\|C_w\|=160.003$, attained at $t=0$, then $122.32$, $98.32$, $80.79$ \lab{recomputed in closed form}. The factor 160 is a static coordinate weight; $C_wz$ is not a margin quantity because $z$ contains no inventory. Margin must be computed on a cash-and-inventory ledger.

\section{Appendix: Reproducibility: source of every number}\label{p1:app:repro}

The replication package contains numbered Colab notebooks and stand-alone scripts; its README lists them in run order. ``Part $N$'' below and in table notes refers to notebook $N$ (Part 1: theorem checks and spot tests; Part 2: BTCUSDT futures depth; Parts 4--13: the cross-section, replications, fresh-day tests, second venues, the CFTC test, the certified guard and post-hoc checks). Result identifiers such as \texttt{X6} or \texttt{U6} name the JSON file that the notebook writes. ``Kind'' distinguishes closed-form or simulation results (recomputed), results from public market data (real data), and audits.

\begin{longtable}{p{0.39\textwidth}p{0.37\textwidth}p{0.14\textwidth}}
\toprule
Result & Source & Kind\\
\midrule
\endhead
Complete monotonicity, 892 economies & \texttt{L26\_complete\_monotonicity} & recomputed\\
$d_2=-\kappa^2q\delta^2$ & \texttt{P1.7-T1\_third\_tap} & recomputed\\
Ledger identity; $T_h$ positive definite & \texttt{R2-T1}, \texttt{R2-T3} & recomputed\\
Figure~\ref{p1:fig:knownpair} & \texttt{make\_figs2.py} & recomputed\\
$\delta_{\min}(3)$, Table~\ref{p1:tab:capbounds} recomputed columns, 61.14, 29.94 & \texttt{recompute\_theory.py} $\to$ \texttt{R\_P1\_theory\_constants.json} & recomputed\\
Pulse benchmarks (271, 58, 12, 11 pulses), Figures~\ref{p1:fig:impulse}--\ref{p1:fig:phase} & \texttt{make\_figs3.py} $\to$ \texttt{R\_P1\_pulse\_benchmarks.json} & recomputed\\
Table~\ref{p1:tab:mcerr}, Table~\ref{p1:tab:radau}, worked example, passive simulation & \texttt{recompute\_extra.py} $\to$ \texttt{R\_P1\_extra.json} & recomputed\\
Daily noise (Table~\ref{p1:tab:roll}) & \texttt{P1.3-T1\_sigma\_obs} & real data\\
Table~\ref{p1:tab:lecam} recomputed columns & same & recomputed\\
$K$-cohort slopes (Table~\ref{p1:tab:kcohort}) & \texttt{T1\_K\_cohort\_rate} & recomputed\\
ESPRIT slope $-6.15$ & \texttt{P1.2-T1\_esprit\_rate} & recomputed\\
Minimum phase, 1,733 economies; rank 6 of 6 & \texttt{T2a}, \texttt{T2b} & recomputed\\
Passive slopes $-4.21$, $-5.93$ & \texttt{T3\_passive\_rate} & recomputed\\
Tagged-flow slope 0.011 & \texttt{T4\_disclosure} & recomputed\\
Spectral blindness gains & \texttt{T5\_spectral\_blindness} & recomputed\\
Noise and capital grid (Table~\ref{p1:tab:capital}) & \texttt{P1.3-T1\_sigma\_obs} & real data\\
Loop gain and placebo (Section~\ref{p1:sec:empirics}) & \texttt{P1.6-T1\_ema\_placebo} & real data\\
Estimator table (Section~\ref{p1:sec:empirics}) & \texttt{P1.5-T1}, \texttt{T2}, \texttt{T3} & recomputed\\
Data integrity (111 files) & \texttt{R2-T9\_data\_integrity} & audit\\
Tables~\ref{p1:tab:ledger}, \ref{p1:tab:capbounds} (certified, feasible, adaptive columns) & \texttt{recompute\_p1\_capital.py} $\to$ \texttt{R\_P1\_capital\_ledger.json} & recomputed\\
Table~\ref{p1:tab:contracts}; resolution floors; Table~\ref{p1:tab:protocols} & \texttt{recompute\_p1\_protocols.py} $\to$ \texttt{R\_P1\_protocols.json} & recomputed\\
SPRT capital figures & \texttt{recompute\_p1\_sprt.py} $\to$ \texttt{R\_P1\_sprt.json} & recomputed\\
Table~\ref{p1:tab:crb} ($N^*$, $\sigma_{\rm crit}$) & \texttt{recompute\_crb.py} $\to$ \texttt{R\_P1\_crb.json} & recomputed\\
Table~\ref{p1:tab:crb} (ML columns), Table~\ref{p1:tab:lecam} (estimators) & \texttt{recompute\_p1\_estimators.py} & recomputed\\
Symmetric multi-asset network (Section~\ref{p1:sec:blind}) & \texttt{recompute\_oct20.py} & recomputed\\
Proof audit: $K_*$ bound and $d_2$ identity (Theorem~\ref{p1:thm:quantity}), moment gaps (Theorem~\ref{p1:thm:twocohort}), spectral-blindness example (Theorem~\ref{p1:thm:blind}) & \texttt{recompute\_audit.py} $\to$ \texttt{R\_audit\_oct2026.json} & recomputed\\
Theorem~\ref{p1:prop:esprit}, Table~\ref{p1:tab:kupper} & \texttt{kcohort\_upper.py} $\to$ \texttt{R\_P1\_kcohort\_upper.json} & recomputed\\
Lem.~\ref{p1:lem:altpair}, Thm~\ref{p1:thm:general}, Tab.~\ref{p1:tab:general} & \texttt{general\_minimax.py} & recomputed\\
Theorem~\ref{p1:thm:discvalue} numbers & \texttt{disclosure\_value.py} $\to$ \texttt{R\_P1\_disclosure\_value.json} & recomputed\\
Section~\ref{p1:sec:empirics} & Part 8: \texttt{V2\_spot\_panel}, \texttt{V3\_fee\_events} & real data\\
Table~\ref{p1:tab:disc}, top panel & Part 4: \texttt{X9\_disclosure\_metrics} & real data\\
Table~\ref{p1:tab:disc}, bottom panel; class-position forecasts & Part 6: \texttt{Z6\_C\_identification}, \texttt{Z6\_secondary}; \texttt{P6\_fresh\_days.json} (65acda36\ldots) & real data\\
Table~\ref{p1:tab:cot} & Part 11: \texttt{C3\_cot\_data}, \texttt{C4\_cot\_disclosure}; \texttt{P11\_cot.json} (6f1d6f35\ldots) & real data\\
Basis-trade signature; real-data Cochrane--Orcutt & Part 13: \texttt{R6}, \texttt{R7} & real data (post hoc)\\
Feedback across 30 perpetuals & Part 4: \texttt{X3\_T6\_cross\_section} & real data\\
Theorem~\ref{p1:thm:rh} full bound (Table~\ref{p1:tab:rh}), direct minimization & \texttt{hardening/rh\_bound.py}, \texttt{rh\_check.py} $\to$ \texttt{R\_P1\_rh\_explicit.json} & recomputed\\
Minimal dimension at $a=0$; Toeplitz sections; passive slopes for $K\le4$; IV determinant & \texttt{hardening/p1\_checks.py} $\to$ \texttt{R\_P1\_hardening\_checks.json} & recomputed\\
Stylized nonlinear and trained online rules (Table~\ref{p1:tab:learners}, Figure~\ref{p1:fig:learners}) & \texttt{recompute\_\allowbreak p1\_\allowbreak learners.py} $\to$ \texttt{R\_P1\_\allowbreak learners.json} (Panel A); \texttt{validate\_\allowbreak trained\_\allowbreak learners\_\allowbreak hardened.py} $\to$ \texttt{R\_P1\_\allowbreak trained\_\allowbreak learners\_\allowbreak hardened.json} (Panels B--C); specification \texttt{P1\_\allowbreak trained\_\allowbreak learners.json} (3e8bf30c\ldots) & recomputed\\
Disclosure in dollars (Table~\ref{p1:tab:cat_dollars}) & \texttt{recompute\_\allowbreak p1\_\allowbreak cat\_\allowbreak costs.py} $\to$ \texttt{R\_P1\_\allowbreak cat\_\allowbreak disclosure\_\allowbreak dollars.json} & recomputed\\
Out-of-sample post-shock return test (Section~\ref{p1:sec:disc_data}) & \texttt{recompute\_\allowbreak prereg\_\allowbreak oct21\_\allowbreak 22nd.py} $\to$ \texttt{R\_P1\_\allowbreak disclosure\_\allowbreak oos.json}; \texttt{P1\_\allowbreak disclosure\_\allowbreak oos.json} (f0182ce5\ldots) & real data\\
Out-of-sample cohort-state trajectory test (Table~\ref{p1:tab:disc_state_oos}) & \texttt{recompute\_\allowbreak prereg\_\allowbreak oct21\_\allowbreak followup.py} $\to$ \texttt{R\_P1\_\allowbreak disclosure\_\allowbreak state\_\allowbreak oos.json}; \texttt{P1\_\allowbreak disclosure\_\allowbreak state\_\allowbreak oos.json} (5136cc6b\ldots) & real data\\
\bottomrule
\end{longtable}

\renewcommand{\arraystretch}{1}
\chapter{Safety without identification: transient amplification and a maximal pre-trade guard}\label{ch:p2}

    This chapter, the second paper of the thesis, condenses the paper \emph{From Human Markets to Machine Markets: Transient Amplification, an Exact Depth Condition, and a Maximal Pre-Trade Guard}, joint with Suboh Alkhushayni, Yuhao Chen, Mansi Bhavsar and Prashant Niroula. It keeps the theorems the rest of the thesis uses, with their proofs, and the empirical program in summary, with every registered test and its outcome in Table~\ref{p2:tab:alltests}. The full version, with every table, robustness check and the source of every number, is in the replication package (\texttt{thesis/ch3\_paper2\_full.tex}).

    \begin{quote}\small\textbf{Summary.} Eigenvalues and volatility do not bound the losses a market produces before it settles. In a limit order book with adaptive traders, cumulative transient risk is $S=\tilde\kappa/(m-q\tilde\kappa)$, long-run impact over reversion net of momentum; a book overshoots its direct impact if and only if $\varphi\varepsilon>(1-a)\kappa$; and a pre-trade guard on the summed impulse response is the largest safe admission rule and does not depend on how the traders split into cohorts. The summed response is certified with $\Theta(\sigma^2\epsilon^{-2}\log(1/\epsilon))$ experiments whatever the cohort spacing, while identifying the cohorts costs polynomially more, and the certified guard is minimax-optimal in lost capacity up to a logarithm. Of 25 registered empirical tests, 15 pass: six-hour estimates of the statistic forecast the size of the next hour's adverse excursion in thin Binance perpetuals beyond a full volatility baseline and VPIN, more weakly in thin US stocks, and not on Coinbase or in the deepest book; tests of a liquidation mechanism fail.\end{quote}

    \section{Introduction}

    On October 10, 2025, crypto derivatives markets liquidated an estimated \$19 billion of leveraged positions in a single day \citep{p2:shares2025}; on May 6, 2010, US equity prices fell and recovered by several percent within minutes \citep{p2:kirilenko2017}. In both episodes the losses that mattered were transient: positions were closed out on the path, before prices settled. Algorithmic markets are closed-loop systems, and a common shortcut equates their stability with eigenvalues inside the unit circle. For a leveraged intermediary that is the wrong object: a market can be spectrally stable and still push a dealer through its maintenance margin within seconds.

    The chapter asks three questions. When does a limit order book amplify a flow shock? How can a dealer or a venue control transient loss without identifying who is trading, which Chapter~\ref{ch:p1} shows is expensive from public data? And what information certifies a safe trading decision, at what cost? The answers, in a linear book model, are these.

    \paragraph{Spectral stability does not bound transient loss.} The safe shock radius of a stable market is the loss budget over its peak transient gain, and that gain can be arbitrarily large at any fixed spectral radius (Theorem~\ref{p2:thm:radius}).

    \paragraph{When a book amplifies.} A book with depletion and refill overshoots its direct impact if and only if $\varphi\varepsilon>(1-a)\kappa$; momentum alone cannot make a stable book overshoot; and cumulative risk has the closed form $S=\tilde\kappa/(m-q\tilde\kappa)$, which needs no depletion (Theorems~\ref{p2:thm:peak} and~\ref{p2:thm:sum}).

    \paragraph{The maximal guard.} Under a declared risk contract, worst-case loss is inventory times the shock bound times the summed absolute impulse response, so a guard on that sum admits exactly the safe actions (Lemma~\ref{p2:thm:guard}); a guard on the peak response admits unsafe ones. The guard survives estimation error (Theorem~\ref{p2:thm:robust}) and needs no identification of the traders (Proposition~\ref{p2:prop:invariance}).

    \paragraph{Certifying is cheaper than identifying.} Two markets can agree on their early responses and differ threefold in risk (Example~\ref{p2:ex:witness}). A finite window certifies the summed response under a memory bound and certifies nothing without one (Theorem~\ref{p2:thm:window}); with noisy experiments the cost of certification is $\Theta(\sigma^2\epsilon^{-2}\log(1/\epsilon))$ whatever the cohort structure, against $(\Delta\eta)^{-(4K-2)}$ for identification (Theorem~\ref{p2:thm:certcost}). The guard built on the certificate has a closed form and is minimax-optimal (Theorem~\ref{p2:thm:robustguard}); disclosure is worth most when memory is long, and a noisy proxy can be worth less than nothing (Theorem~\ref{p2:thm:disclosure}).

    \paragraph{Who controls the risk.} Reversion is a public good: dealers under-provide it, momentum traders over-provide feedback, and the corrective rebate is a function of the guard's own statistic (Proposition~\ref{p2:prop:equil}).

    \paragraph{Evidence.} Twenty-five primaries were registered; twenty-four were tested and one was not opened under its fixed sequence; fifteen pass (Table~\ref{p2:tab:alltests}). The deepest book, BTCUSDT, sits on the stable side of the amplification condition and fails a registered forecasting test. In thin books the statistic forecasts the size of the next hour's adverse excursion in every registered forecasting test on Binance, on new symbols and on untouched days, by up to $2.8$ points of $R^2$ beyond a full volatility baseline; about a fifth as much in thin US stocks; negligibly on Bybit; not at all on Coinbase. We have not identified a deeper mechanism: the effect does not run through open-interest liquidations, and it is carried by slow reversion relative to impact, which is what the statistic measures. As a risk control the certified guard lowers tail loss against a static limit but not against a volatility cap at matched exposure.

    \paragraph{What we do not claim.} The information-safety theorems concern declared observation contracts (controlled reset experiments, exact windows, audited data), not finite-sample guarantees for observational order-book records. Cohorts are any traders whose demand is a linear filter of past prices; nothing attributes an effect to algorithms rather than people. The forecasting result is reduced-form, and its size depends on the venue.

    \paragraph{Related work.} Models of transient impact and resilience \citep{p2:obizhaeva2013,p2:bouchaud2004,p2:bouchaud2018} take the decay of impact as given; we derive when depletion, refill and momentum make a book overshoot and give cumulative impact in closed form. Liquidity spirals \citep{p2:brunnermeier2009,p2:shleifer1992} need funding constraints; our amplification does not, and our guard is the control that bounds it. \citet{p2:sornette2023} show that non-normal interactions create transient bubbles. The guard's worst-case identity is the $\ell_1$ gain of $\ell_1$-optimal control \citep{p2:vidyasagar1986,p2:dahleh1987,p2:dahleh1995}; what is new is its use as an admission rule, its invariance to the cohort coordinates and its identification from prices. The early-warning benchmark is VPIN \citep{p2:easley2012,p2:andersen2014}; forecast evaluation uses \citet{p2:clark2007} with the volatility baselines of \citet{p2:corsi2009} and \citet{p2:riskmetrics1996}, and we pre-register following \citet{p2:nosek2018} and \citet{p2:harvey2017}. Rule 15c3-5 requires pre-trade controls on credit and capital \citep{p2:sec2010}; the guard is a candidate with a provable maximality property.

    \section{The closed-loop market}\label{p2:sec:model}

    \subsection{Price formation}
Time is discrete with bars of length $\Delta$. Let $p_t$ be the price, $p^*_t=p^*_{t-1}+\xi_t$ the unobserved fundamental and $x_t=p_t-p^*_t$ the pricing error. Competitive dealers post aggregate depth $D>0$, so price impact is $\kappa=1/D$, and fundamental arbitrage removes a fraction $\kappa\varphi_f$ of the error each period. With algorithmic demand $u_t$ and liquidity noise, the pricing error follows
\begin{equation}\label{p2:eq:price}
x_{t+1}=ax_t+\kappa u_t+w_{t+1},\qquad a=1-\kappa\varphi_f\in[0,1).
\end{equation}
Algorithmic cohorts trade $u_t=\sum_iq_iy_{i,t}$ with exponential memory $y_{i,t+1}=(1-\eta_i)y_{i,t}+\eta_ix_{t+1}$, a reduced form that arises when an adaptive desk updates its belief by steady-state Kalman filtering or regularized online gradient descent and maximizes expected continuation return $\beta_i y_{i,t}u_{i,t}$ minus quadratic execution cost $\tfrac12\lambda_{{\rm tc},i}u_{i,t}^2$ and inventory penalty $\tfrac12\gamma_{{\rm inv},i}I_{i,t}^2$ ($I_{i,t}=(1-\delta_{\rm inv})I_{i,t-1}+u_{i,t}$, $|I_{i,t}|\le\bar I_i$), yielding $u_{i,t}^*=q_iy_{i,t}-\psi_{{\rm inv},i}I_{i,t-1}$ with $q_i=\beta_i/(\lambda_{{\rm tc},i}+\gamma_{{\rm inv},i})$ and $\psi_{{\rm inv},i}=\gamma_{{\rm inv},i}(1-\delta_{\rm inv})/(\lambda_{{\rm tc},i}+\gamma_{{\rm inv},i})$ (Chapter~\ref{ch:p1}; this single-agent filtering-and-control rule is distinct from the joint dealer--trader equilibrium of Section~\ref{p2:sec:equil}). The baseline equations set $\psi_{{\rm inv},i}=0$, the limit in which inventory holding penalties are small relative to execution costs ($\gamma_{{\rm inv},i}\ll\lambda_{{\rm tc},i}$) or inventory is reset at the bar horizon; Proposition~\ref{p2:prop:inventory_S} below relaxes this and gives the closed-form summed response when cohorts carry inventory across bars with $\psi_{{\rm inv},i}>0$. Stacking $z_t=(x_t,y_{1,t},\ldots,y_{K,t})^\top$ gives $z_{t+1}=Jz_t+bw_{t+1}$, where $b=(1,\eta_1,\ldots,\eta_K)^\top$ because cohorts update on the same period's price (a unit of flow enters through $\kappa b$), and
\begin{equation}\label{p2:eq:J}
J=\begin{pmatrix} a & \kappa q_1 & \cdots & \kappa q_K\\ \eta_1 a & 1-\eta_1+\eta_1\kappa q_1 & \cdots & \eta_1\kappa q_K\\ \vdots & & \ddots & \vdots\\ \eta_K a & \eta_K\kappa q_1 & \cdots & 1-\eta_K+\eta_K\kappa q_K\end{pmatrix}.
\end{equation}

    \subsection{Stability and what an observer can use}
\begin{proposition}[Stability with one cohort]\label{p2:prop:jury}
With one cohort $(q,\eta)$ and $a=1-\kappa\varphi_f$, the market is Schur stable if and only if
\begin{equation}\label{p2:eq:jury}
q<\varphi_f,\qquad \kappa<\frac{4-2\eta}{(2-\eta)\varphi_f-\eta q}\ \ \text{when}\ (2-\eta)\varphi_f>\eta q,\qquad |1-\kappa\varphi_f|(1-\eta)<1 .
\end{equation}
Higher impact $\kappa$ (thinner depth) lowers price persistence $a$; feedback is destabilizing only when cohort aggressiveness $q$ reaches the arbitrage strength $\varphi_f$.
\end{proposition}
Chapter~\ref{ch:p1} shows that prices identify only the merged rate--weight measure $M=\sum_iq_i\eta_i\delta_{1-\eta_i}$ of the cohorts. A pre-trade control must therefore depend on the market only through objects that are the same for every cohort representation of $M$.

\begin{proposition}[Input--output invariance]\label{p2:prop:invariance}
Two populations with the same merged measure generate the same transfer function $H(z)=\sum_kg_kz^{-k}=\kappa z/[z-a-\kappa G(z)]$, $G(z)=\sum_iq_i\eta_iz/(z-1+\eta_i)$, and the same impulse coefficients $g_k=\kappa\,c_{\rm price}^\top J^kb$ for every $k\ge0$. Under any invertible change of state coordinates $z\mapsto Tz$, the Euclidean gain $\|CT^{-1}(TJT^{-1})^k\|_2$ depends on $T$, whereas the price path
\[
\Delta p_{t+k}=\sum_{j\ge1}g_{k-1+j}u_{t-j}+\sum_{m=0}^{k-1}g_{k-1-m}u_{t+m}
\]
and its worst case over $\|u\|_\infty\le\bar u$ depend only on $\{g_k\}$ and are invariant to both $T$ and cohort splitting.
\end{proposition}

Numerically, splitting a cohort of 1{,}000 random economies into two with the same memory rate leaves every impulse coefficient unchanged up to rounding ($\max_k|g_k(\mathcal E^A)-g_k(\mathcal E^B)|\le2.5\times10^{-12}$ over 200 lags; \texttt{recompute\_oct20.py}). Proposition~\ref{p2:prop:invariance} is the reason we build the guard on impulse coefficients and avoid state norms in latent coordinates. When a state norm is needed, we fix units by observable calm-market innovation standard deviations, which is what makes the spot statistic of Section~\ref{p2:sec:evidence} unit-free.
\begin{proposition}[Summed response with carried cohort inventory]\label{p2:prop:inventory_S}
Suppose each cohort $i$ carries inventory $I_{i,t}=(1-\delta_{\rm inv})I_{i,t-1}+u_{i,t}$ with $\delta_{\rm inv}\in(0,1]$ and submits order flow $u_{i,t}=q_i y_{i,t}-\psi_i I_{i,t-1}$ with $q_i\ge0$ and $\psi_i\ge0$. Then:
\begin{enumerate}[label=(\roman*)]
\item The closed-loop transfer function from a flow impulse to price is $H_{\rm inv}(z)=\kappa z/\bigl[z-a-\kappa G_{\rm inv}(z)\bigr]$ with
\[
G_{\rm inv}(z)=\sum_{i=1}^K \frac{q_i\eta_i z(z-1+\delta_{\rm inv})}{(z-1+\eta_i)(z-1+\delta_{\rm inv}+\psi_i)},
\]
which is invariant to splitting any cohort that shares $(\eta_i,\psi_i)$. Whenever the augmented state matrix is Schur stable, the signed cumulative price response is
\begin{equation}\label{p2:eq:Sinv}
S(\psi)=\sum_{t=0}^\infty h_t(\psi)=H_{\rm inv}(1)=\frac{\kappa}{1-a-\kappa\Phi_{\rm eff}(\psi)},\qquad \Phi_{\rm eff}(\psi)\equiv\sum_{i=1}^K\frac{\delta_{\rm inv}}{\delta_{\rm inv}+\psi_i}\,q_i\le\Phi.
\end{equation}
\item If $0\le\psi_i\le 1-\delta_{\rm inv}$ and $\delta_{\rm inv}\ge\eta_i$ for each $i$ and $a+\kappa\Phi_{\rm eff}(\psi)<1$, every open-loop feedback tap $k_m(\psi)$ and hence every closed-loop impulse coefficient $h_t(\psi)$ is nonnegative, so the worst-case $\ell_1$ gain equals the signed sum, $\|h(\psi)\|_1=\sum_{t=0}^\infty |h_t(\psi)|=S(\psi)$. When $\psi_i$ is large enough that inventory payback reverses the sign of late taps, $\|h(\psi)\|_1>S(\psi)$.
\end{enumerate}
\end{proposition}

Inventory mean-reversion $\psi_i>0$ acts as an automatic stabilizer that attenuates each cohort's long-run feedback weight from $q_i$ to $\frac{\delta_{\rm inv}}{\delta_{\rm inv}+\psi_i}q_i$ (across 500 random stable economies \eqref{p2:eq:Sinv} matches the simulated state-space impulse sum to $1.78\times10^{-15}$, and in the monotone regime $\|h(\psi)\|_1=S(\psi)$ to $1.78\times10^{-15}$, reducing $S$ from $2.7778$ at $\psi=0$ to $2.5000$ at the calibrated worked example; in the sign-reversal regime $\|h(\psi)\|_1=2.0960>S(\psi)=1.7241$; \texttt{recompute\_p2\_step2\_theory.py}).

    The model is linear: it describes the transient response to flow shocks of ordinary size, which is what a pre-trade control must bound, not regime changes such as exchange outages. Every input to the guard (impulse coefficients, an operator estimate and its confidence radius, inventory and margin) is observable to a dealer or a venue.

    \section{Spectral stability is not solvency}\label{p2:sec:radius}
    
\begin{theorem}[Safe radius, Kreiss bounds and the failure of the spectral radius]\label{p2:thm:radius}
Let $z_t=J^tz_0$ with $\rho(J)<1$, loss budget $B>0$ and loss output $C$.
\begin{enumerate}[label=(\roman*)]
\item The largest initial-state radius with $\sup_t\|CJ^tz_0\|_2\le B$ for every $\|z_0\|_2\le r$ is
\begin{equation}\label{p2:eq:rexact}
r^*_{\rm exact}=\frac{B}{M_C},\qquad M_C\equiv\sup_{t\ge0}\|CJ^t\|_2 .
\end{equation}
\item With $C=I$ and the Kreiss constant $\mathcal K(J)=\sup_{|z|>1}(|z|-1)\|(zI-J)^{-1}\|_2$, the Kreiss matrix theorem gives $\mathcal K(J)\le\sup_t\|J^t\|_2\le e\,n\,\mathcal K(J)$ for $J\in\R^{n\times n}$, hence
\begin{equation}\label{p2:eq:kreiss}
\frac{B}{e\,n\,\mathcal K(J)}\;\le\;r^*_{\rm exact}\;\le\;\frac{B}{\mathcal K(J)} .
\end{equation}
\item For any $\rho\in(0,1)$ and any $B>0$ there is a family $J_\theta$ with $\rho(J_\theta)=\rho$ for all $\theta$ and an output $C=(1,0)$ such that $r^*_{\rm exact}(J_\theta)\to0$ as $\theta\to\infty$.
\end{enumerate}
\end{theorem}

    For $J_\theta=\begin{psmallmatrix}0.8&\theta\\0&0.4\end{psmallmatrix}$ the spectral radius is $0.8$ throughout, while the price-row gain rises from 1 to 120 as $\theta$ goes from 0 to 100 and the safe radius falls by the same factor. Chapter~\ref{ch:p1} shows more: markets that share their eigenvalues, one-step price variance and stationary price variance can differ arbitrarily in worst-case transient gain, so neither eigenvalue nor volatility monitors bound transient loss. In the cohort economy with nonnegative weights the price response is a nonnegative mixture of decaying exponentials and peaks at the direct impact (Chapter~\ref{ch:p1}, positive-geometry lemma): non-normality is not sufficient for scalar price overshoot, and necessity for Euclidean state-norm growth does not extend to arbitrary input-output responses. In this positive cohort model, momentum alone does not supply overshoot. Depth depletion does.

    \section{When does a book amplify? An exact condition}\label{p2:sec:t6}

    Transient path risk takes two forms. \emph{Peak overshoot} is a response that rises above its direct impact; \emph{cumulative amplification} is a summed response much larger than direct impact. The data support the second form, not the first.

    \subsection{A book with depletion and refill}
Let $x_t$ be the price displacement, $d_t$ the signed depth imbalance (bid minus ask depth near the mid, relative to total) and $u_t$ signed order flow. Flow moves the price directly and depletes one side of the book, which refills at rate $r$; depth imbalance moves the price:
\begin{align}
x_t &= a\,x_{t-1} + \varphi\,d_{t-1} + \kappa\,u_t, \label{p2:eq:x}\\
d_t &= (1-r)\,d_{t-1} + \varepsilon\,u_t, \label{p2:eq:d}
\end{align}
with $\kappa>0$, $a\in[0,1)$ and $r\in(0,1]$. Write $b=1-r$. With the sign convention $d=(B-A)/(B+A)$, a buy that consumes ask depth raises $d$, so $\varepsilon>0$ means \emph{depletion} of the side that is hit, and $\varphi>0$ means that a bid-heavy book pushes the price up. The impulse response to a unit flow shock at $t=0$ is $h_0=\kappa$ and
\begin{equation}\label{p2:eq:rec}
h_k=a\,h_{k-1}+\varphi\varepsilon\,b^{k-1},\qquad k\ge1.
\end{equation}

\begin{definition}
The book \emph{amplifies} if $\max_kh_k>h_0$: the price keeps moving in the direction of the shock after the shock has passed.
\end{definition}

    \subsection{Peak overshoot}
Add momentum traders whose flow $f_t=(1-\eta)f_{t-1}+\eta q\,x_{t-1}$ enters both \eqref{p2:eq:x} and \eqref{p2:eq:d} alongside $u_t$; the two-state model is the case $q=0$. Write $C_2\equiv\varphi\varepsilon-(1-a)\kappa$ and $C_3\equiv\varphi\varepsilon+\kappa^2\eta q-(1-a)\kappa$.

\begin{theorem}[Peak overshoot: depletion versus momentum]\label{p2:thm:peak}
\begin{enumerate}[label=(\roman*)]
\item\label{p2:thm:t6} \emph{Two-state condition.} Without momentum ($q=0$) the book amplifies if and only if $C_2>0$; when $C_2\le0$ the response peaks at impact, $\max_kh_k=\kappa$.
\item\label{p2:prop:shape} \emph{Shape.} Let $q=0$, $0<a<1$ and $0<b<1$. If $a\neq b$, $h_k=Aa^k+Bb^k$ with $B=\varphi\varepsilon/(b-a)$ and $A=\kappa-B$; if $a=b$, $h_k=a^{k-1}(a\kappa+\varphi\varepsilon k)$ for $k\ge1$. When the book amplifies, $h$ viewed as a function of real $k\ge0$ has a single maximum, at
\[
k^*=\frac{\log\bigl(-A\log a/(B\log b)\bigr)}{\log(b/a)}\quad(a\neq b),\qquad k^*=\frac1{\log(1/a)}-\frac{a\kappa}{\varphi\varepsilon}\quad(a=b),
\]
the second being the limit of the first as $b\to a$; the peak lag is $\lfloor k^*\rfloor$ or $\lceil k^*\rceil$ (and is $1$ when $k^*<1$, because $h_1>h_0$). The cumulative response is $\sum_kh_k=\kappa/(1-a)+\varphi\varepsilon/[(1-a)r]$ in all cases, and faster refill lowers both the peak and the cumulative response. When $a=0$ ($h_k=\varphi\varepsilon b^{k-1}$) or $r=1$ ($h_k=a^{k-1}(a\kappa+\varphi\varepsilon)$) the response peaks at lag 1.
\item\label{p2:prop:t6m} \emph{One-period margin with momentum.} $h_1-h_0=C_3$, so $C_3>0$ implies that the book amplifies.
\item\label{p2:thm:nomom} \emph{Momentum alone cannot amplify.} If $\varphi\varepsilon=0$, $q\ge0$ and the system is stable, then $h_k=c_1\lambda_1^k+c_2\lambda_2^k$ with residues $c_1,c_2\ge0$, $c_1+c_2=\kappa$ and $|\lambda_2|\le\lambda_1=\rho<1$, so $h_k\le\kappa\rho^k<\kappa$ for every $k\ge1$; in particular $C_3<0$.
\item\label{p2:thm:noover} \emph{A sufficient condition for no overshoot.} For any signs of $\varphi$, $\varepsilon$ and $q$, if
\[
\bar C_3\equiv|\varphi\varepsilon|\Bigl(1+\frac{\kappa|q|}{r}\Bigr)+\kappa^2|q|-(1-a)\kappa\;\le\;0,
\]
then $\sup_k|h_k|=\kappa$. When $\varphi\varepsilon\ge0$ and $q\ge0$, $C_3>0\Rightarrow$ amplification $\Rightarrow\bar C_3>0$.
\end{enumerate}
\end{theorem}

\begin{proof}
(i) \emph{If.} $h_1-h_0=a\kappa+\varphi\varepsilon-\kappa=C_2>0$. \emph{Only if.} Suppose $C_2\le0$. If $\varphi\varepsilon<0$, then by \eqref{p2:eq:rec} $h_k=a^k\kappa+\varphi\varepsilon\sum_{j<k}a^{k-1-j}b^j\le a^k\kappa\le\kappa$. If $\varphi\varepsilon\ge0$, let $\Delta_k=h_k-h_{k-1}$. Then $\Delta_1=C_2\le0$ and, from \eqref{p2:eq:rec}, $\Delta_{k+1}=a\Delta_k+\varphi\varepsilon(b^{k}-b^{k-1})=a\Delta_k-\varphi\varepsilon\,r\,b^{k-1}\le a\Delta_k$, so by induction $\Delta_k\le0$ for all $k$ and $h$ peaks at $h_0=\kappa$.

(ii) Solving \eqref{p2:eq:rec} gives the two forms. For $a\ne b$, $\frac{d}{dk}h=Aa^k\log a+Bb^k\log b$ has at most one zero because $(b/a)^k$ is monotone, which gives $k^*$; for $a=b$, $\frac{d}{dk}a^{k-1}(a\kappa+\varphi\varepsilon k)=a^{k-1}[(a\kappa+\varphi\varepsilon k)\log a+\varphi\varepsilon]$ vanishes once, at the stated $k^*$, and expanding the first formula in $b-a$ gives the same limit. A unimodal function of real $k$ attains its integer maximum at $\lfloor k^*\rfloor$ or $\lceil k^*\rceil$. Summing the geometric series gives the cumulative response, which falls in $r$; the peak falls in $r$ because $h_k$ is decreasing in $r$ term by term for $\varphi\varepsilon>0$.

(iii) With a unit shock at $t=0$, $x_0=\kappa$, $d_0=\varepsilon$ and $f_1=\eta q\kappa$, so $h_1=a\kappa+\varphi\varepsilon+\kappa^2\eta q$.

(iv) Without the depth channel the state $(x_t,f_t)$ evolves by $M=\begin{psmallmatrix}a+\kappa\eta q&\kappa(1-\eta)\\ \eta q&1-\eta\end{psmallmatrix}$, which is entrywise nonnegative, and $h_k=\kappa(M^k)_{11}$. The off-diagonal product is nonnegative, so the eigenvalues are real, and by the Perron--Frobenius theorem $\lambda_2\le m_{11}\le\lambda_1=\rho$. Writing $(M^k)_{11}=c_1'\lambda_1^k+c_2'\lambda_2^k$ with $c_1'=(m_{11}-\lambda_2)/(\lambda_1-\lambda_2)\in[0,1]$ and $c_2'=1-c_1'$ gives the residues $c_j=\kappa c_j'\ge0$ and $(M^k)_{11}\le\rho^k$. Since $m_{11}\le\rho<1$, $C_3=\kappa(m_{11}-1)<0$.

(v) Suppose $|x_j|\le\kappa$ for $j<k$. Then $|f_j|\le|q|\kappa$ for $j\le k$ and $|d_{k-1}|\le|\varepsilon|\bigl(b^{k-1}+\sum_{j\ge1}b^{k-1-j}|q|\kappa\bigr)\le|\varepsilon|(1+|q|\kappa/r)$, so $|x_k|\le a\kappa+|\varphi\varepsilon|(1+\kappa|q|/r)+\kappa^2|q|\le\kappa$. Induction from $x_0=\kappa$ completes the proof.
\end{proof}

\paragraph{Interpretation.} The left side $\varphi\varepsilon$ of $C_2$ is the extra price pressure that a unit of flow creates one period later through the book it leaves behind. The right side $(1-a)\kappa$ is the part of the direct impact that reverts in one period. The book amplifies exactly when the depth channel adds more than mean reversion removes. Refill $r$ does not enter the condition; it governs how long an overshoot lasts.

    The converse of part~\ref{p2:prop:t6m} fails with two slow modes: near the stability boundary, or with slow refill and slow momentum, a book can have $C_3<0$ and still peak far above its impact, because the depth channel creates a negative residue on an intermediate mode. No condition on the one-period margin alone is exact.

    \subsection{Cumulative transient risk}
\begin{theorem}[Summed response in closed form]\label{p2:thm:sum}
Let $\tilde\kappa=\kappa+\varphi\varepsilon/r$ (long-run impact: direct plus depth) and $m=1-a$ (reversion). In the stable three-state model,
\[
S\equiv\sum_{k\ge0}h_k=\frac{\tilde\kappa}{m-q\tilde\kappa},
\]
and when $\varphi\varepsilon\ge0$ and $q\ge0$ the system is stable only if $m>q\tilde\kappa$. When all $h_k\ge0$, $S$ is the summed absolute response that enters the guard of Section~\ref{p2:sec:guard}.
\end{theorem}

\begin{proof}
$S=e_1^\top(I-M_3)^{-1}(\kappa,\varepsilon,0)^\top$ with $M_3$ the transition matrix of $(x,d,f)$. Direct computation gives $\det(I-M_3)=\eta\,[\,r(m-\kappa q)-\varphi\varepsilon q\,]$ and the stated ratio; the formula matches the simulated sums to a relative error of $7\times10^{-14}$ in 3{,}000 stable draws. If $\varphi\varepsilon\ge0$ and $q\ge0$, $M_3$ is entrywise nonnegative and stability requires $\det(I-M_3)>0$, i.e.\ $m>q\tilde\kappa$.
\end{proof}

Cumulative transient risk is long-run impact divided by reversion net of momentum. Two practical consequences follow. First, the statistic that matters for losses is the summed response, which grows without bound as reversion falls to $q\tilde\kappa$, even when no single lag overshoots (Theorem~\ref{p2:thm:peak}\ref{p2:thm:nomom}). Second, $C_3$ is best read as a one-period amplification margin: it moves with $m$, $q$ and $\tilde\kappa$ in the same direction as $S$ and is what we estimate in Sections~\ref{p2:sec:evidence} and~\ref{p2:sec:thin}. Both statistics are invariant in sign to rescaling flow.

    \section{Identification and financial safety}\label{p2:sec:infosafety}

    Chapter~\ref{ch:p1} shows that identifying the cohorts from prices is expensive. A pre-trade control does not need them; it needs a bound on the risk they create. This section asks what information certifies a safe decision, what it costs, and how much safe trading is lost without it.

    \subsection{Observation and risk contracts}\label{p2:sec:contracts}
We work in the positive cohort economy \eqref{p2:eq:J} with $q_i\ge0$ and known $a$ and $\kappa$. Its merged measure $M=\sum_iq_i\eta_i\delta_{r_i}$, $r_i=1-\eta_i$, has kernel $k_m=\int r^m\,dM$, and the price response to a unit of flow obeys $h_0=\kappa$, $h_{t+1}=ah_t+\kappa\sum_{m\le t}k_mh_{t-m}$. All $h_t\ge0$ (Chapter~\ref{ch:p1}, positive-geometry lemma), and the summed response is
\begin{equation}\label{p2:eq:SPhi}
S=\sum_th_t=\frac{\kappa}{1-a-\kappa\Phi},\qquad \Phi=\int\frac{dM(r)}{1-r}=\sum_iq_i ,
\end{equation}
with stability exactly when $a+\kappa\Phi<1$. A \emph{memory bound} $\bar r<1$ states that every memory rate satisfies $r_i\le\bar r$.

\emph{Risk contract.} A position $I$ is held fixed, each period's flow shock is bounded by $\bar u$, losses are marked to market at every date, and the horizon is unbounded. By Lemma~\ref{p2:thm:guard} the worst-case loss is $|I|\bar uS$, so the safe capacity, the largest shock bound with loss at most the budget $\mathcal B$, is $\mathcal B/(|I|S)$. Section~\ref{p2:sec:guard} adds admission, margin and fills. All statements below are deterministic worst cases under this contract; none is a probability-of-ruin statement.

\emph{Observation contracts.} We consider four: (O1) $N$ independent reset experiments, each returning $h_0,\ldots,h_{T-1}$ with independent Gaussian errors of standard deviation $\sigma$ per tap, the controlled-probe contract of Chapter~\ref{ch:p1}; (O2) an exact window $h_0,\ldots,h_n$; (O3) class-tagged flows; and (O4) audited exposure, which reveals $\Phi$. The theorems concern these contracts. Observational public records of a live book are a different contract, for which we offer evidence (Section~\ref{p2:sec:evidence}) but no finite-sample guarantee.

Given data $D$, let $\Theta(D)$ be the set of economies compatible with it. A decision is \emph{certifiably safe} if its worst-case loss is within budget for every economy in $\Theta(D)$, \emph{certifiably unsafe} if it exceeds the budget for every one, and \emph{undecidable} otherwise. This classification is a definition, not a result; it is the finance version of the data-informativity view of control without identification \citep{p2:vanwaarde2020} and of set-membership estimation \citep{p2:milanese1991}. The results below make it explicit for this economy.

    \subsection{Two markets that look the same}
\begin{example}[An exact ambiguity witness]\label{p2:ex:witness}
Let $a=0.4$ and $\kappa=1$. Market A has one cohort ($q=0.08$, $\eta=0.5$); market B has two ($q=0.02/0.95$, $\eta=0.95$; $q=0.4$, $\eta=0.05$). Their merged measures have the same mass and first moment, so in rational arithmetic their price responses coincide at lags 0, 1 and 2 ($1$, $11/25$, $267/1250$) and first differ at lag 3. Their summed responses are $1.92$ and $5.59$, so for a unit position and budget the safe shock bound is $0.52$ in A and $0.18$ in B.
\end{example}

\begin{proposition}[Unavoidable loss of safe capacity]\label{p2:prop:twoworld}
Let two economies have safe capacities $v_0<v_1$ and data laws $P_0,P_1$. If a data-driven rule $A(D)\in[0,v_1]$ admits an unsafe level in economy 0 with probability at most $\beta$, $P_0(A>v_0)\le\beta$, then
\[
\E_1[v_1-A(D)]\;\ge\;(v_1-v_0)\,\bigl[1-\beta-{\rm TV}(P_0,P_1)\bigr]_+ .
\]
The bound is attained when $P_0=P_1$ by the rule that admits $v_1$ with probability $\beta$ and $v_0$ otherwise.
\end{proposition}

\begin{proof}
Let $E=\{A\le v_0\}$. Then $P_0(E)\ge1-\beta$ and $P_1(E)\ge P_0(E)-{\rm TV}$. On $E$ the capacity given up in economy 1 is at least $v_1-v_0$, and it is nonnegative elsewhere.
\end{proof}

The argument is the two-point method of statistical decision theory \citep{p2:lecam1986}; its use here is to price missing information in units of safe trading. In Example~\ref{p2:ex:witness}, a supervisor who sees the first three coefficients and must admit an unsafe level in B with probability at most 5\% gives up at least $0.95\times0.34=0.32$ units of safe shock capacity in A.

    \subsection{What a finite window certifies}
Lags $1$ to $n$ of the response determine the first $n$ moments of $M$ exactly, through $k_t=(h_{t+1}-ah_t-\kappa\sum_{m<t}k_mh_{t-m})/\kappa^2$. The summed response depends on $M$ only through $\Phi=\int f\,dM$ with $f(r)=(1-r)^{-1}$. Certifying it from a window is therefore an extremal moment problem.

\begin{theorem}[Finite-window certification]\label{p2:thm:window}
Observe $h_0,\ldots,h_n$ exactly under a memory bound $\bar r<1$, and suppose at least one compatible economy is stable. Suppose the implied moment vector $(k_0,\ldots,k_{n-1})$ lies in the interior of the moment space $\mathcal M_n([0,\bar r])$ (equivalently, the Hankel matrices of the moment sequence and of its localization by $r(\bar r-r)$ are positive definite). Let $\mu_-$ and $\mu_+$ be its lower and upper principal representations on $[0,\bar r]$ (the extremal measures of the Markov--Krein theorem, unique in the interior; $\mu_+$ has an atom at $\bar r$), and $\Phi_\pm=\int f\,d\mu_\pm$. On the boundary of the moment space the compatible measure is unique and $\Phi_-=\Phi_+$.
\begin{enumerate}[label=(\roman*)]
\item \emph{Sharp interval.} The summed responses of stable economies compatible with the data form the interval $[S(\Phi_-),S(\Phi_+)]$ if $a+\kappa\Phi_+<1$, and $[S(\Phi_-),\infty)$ otherwise. Every finite endpoint in the stable class is attained by an economy with at most $\lfloor n/2\rfloor+1$ cohorts. When the upper endpoint is infinite it is a supremum approached by stable economies, not an attained response.
\item \emph{Geometric convergence.} $\Phi_+-\Phi_-\le2k_0E_{n-1}$, where $E_{n-1}=(2/\bar r)\,\varrho^{\,n-1}/(x_0^2-1)$ is the error of the best uniform approximation of $f$ on $[0,\bar r]$ by polynomials of degree $n-1$, $x_0=2/\bar r-1$ and $\varrho=x_0-\sqrt{x_0^2-1}<1$. A window of length $n=O\bigl(\log(1/\epsilon)/\log(1/\varrho)\bigr)$ certifies $\Phi$ to precision $\epsilon$, and $\log(1/\varrho)\sim2\sqrt{(1-\bar r)/\bar r}$ as $\bar r\to1$.
\item \emph{No certificate without a memory bound.} If memory rates may lie anywhere in $[0,1)$, then for every $n$ and every window whose moment vector is interior to the moment space, compatible stable economies have arbitrarily large summed responses: no finite window certifies any finite risk bound.
\end{enumerate}
\end{theorem}


    Part (i) applies the Markov--Krein theorem \citep{p2:krein1977,p2:karlin1966}, part (ii) Chebyshev's best approximation of $1/(x_0-x)$ \citep{p2:achieser1956}; the proofs are in Appendix~\ref{p2:app:proofs}. In an eight-cohort economy with memory bound $0.95$ and true summed response $5.56$, lags $0$--$3$ certify nothing, lags $0$--$6$ give $[5.01,7.20]$ and lags $0$--$8$ give $[5.30,6.00]$; for fixed data the upper end grows without limit as the memory bound approaches one, as part (iii) says. The endpoints are proved, not only computed: polynomial bounds whose sign conditions are verified in interval arithmetic with outward rounding turn them into certified intervals \citep{p2:rump2010}.

    \subsection{What certification costs, and what identification costs}
\begin{theorem}[Certification cost of the summed response]\label{p2:thm:certcost}
Consider the class $\mathcal K(\bar r,\bar s)$ of positive economies with memory bound $\bar r$ and $a+\kappa\Phi\le\bar s<1$, observed under contract (O1). Let $\bar\lambda=1-(1-\bar r)(1-\bar s)$ and $\bar h_t$ the average of the $N$ responses.
\begin{enumerate}[label=(\roman*)]
\item \emph{A valid certificate.} For every economy in the class, $h_t\le\kappa\bar\lambda^t$, and
\[
\hat S^+_T=\sum_{t<T}\bar h_t+z_{1-\gamma}\,\sigma\sqrt{T/N}+\frac{\kappa\bar\lambda^T}{1-\bar\lambda}
\]
satisfies $P(S\le\hat S^+_T)\ge1-\gamma$. With $T_\epsilon=\lceil\log(2\kappa/((1-\bar\lambda)\epsilon))/\log(1/\bar\lambda)\rceil$ and $N\ge16z_{1-\gamma}^2\sigma^2T_\epsilon/\epsilon^2$, the excess $\hat S^+-S$ is at most $\epsilon$ with probability at least $1-2\gamma$. Certification to precision $\epsilon$ therefore costs $N=O(\sigma^2\epsilon^{-2}\log(1/\epsilon))$.
\item \emph{A matching lower bound.} Fix an interior single-cohort economy $\theta_0$ and let $C_0=\|\partial h/\partial q\|_2/(\partial S/\partial q)$ at $\theta_0$. Any procedure that is valid at level $1-\gamma$ on the class and overshoots $S$ by less than $2\epsilon$ with probability at least $1/2$ at $\theta_0$ needs $N\ge\bigl((1-2\gamma)\sigma/(2C_0\epsilon)\bigr)^2(1+o(1))$ as $\epsilon\to0$.
\item \emph{The logarithm is necessary.} There is $c>0$ depending only on $(a,\kappa,\bar r,\bar s)$ such that any procedure that is valid at level $1-\gamma$, $\gamma<1/4$, on the class and overshoots $S$ by at most $\epsilon$ with probability at least $1-\gamma$ at every economy of the class needs $N\ge c\,\sigma^2\epsilon^{-2}\log(1/\epsilon)$ for all small $\epsilon$, however many responses each experiment records. With part (i), the cost of certifying $S$ is $\Theta(\sigma^2\epsilon^{-2}\log(1/\epsilon))$.
\item \emph{Independence of the cohort structure.} None of these bounds depends on the number of cohorts or on the spacing of their memories, whereas identifying them costs $N\asymp\sigma^2(\Delta\eta)^{-(4K-2)}$ with exponential memory and $\exp(\Theta(1/w))$ experiments for accuracy $w$ with general memory (Chapter~\ref{ch:p1}).
\end{enumerate}
\end{theorem}

The certificate in part (i) is elementary: it is a sum of averaged taps, a Gaussian margin and a tail that the memory bound controls. Part (iii) shows that its logarithm is not an artefact of truncation: no procedure, however it weights the observed responses, certifies the summed response at the parametric rate, because $1/(1-r)$, the kernel of the total feedback, has a pole at infinite memory and any polynomial that approximates it on the memory range must carry large coefficients. Part (ii) is the parametric bound at a fixed economy; part (iii) is the worst case over the class. The content is in the comparison. The compared inverse problems differ: \citet{p2:tu2018} study $\mathcal H_\infty$-norm estimation relative to FIR input-output model identification, whereas the comparison here concerns a positive response functional relative to latent cohort recovery. These are distinct tasks; no general claim that norm certification is easier than identifying the observable response follows. Table~\ref{p2:tab:certvsid} makes this concrete on Chapter~\ref{ch:p1}'s hardest two-cohort pairs. Those pairs have \emph{identical} summed responses, because they share the total weight $\sum_iq_i$: the directions in which prices are least informative are directions in which the risk a guard must bound does not move. Identifying their weights to $\pm0.01$ needs from 7 to 1.5 million reset experiments as the spacing falls from $0.30$ to $0.04$ (log--log slope $-6.06$), while certifying the summed response to 1\% needs 11 to 18 at every spacing (slope $+0.23$).

\begin{table}[t]\centering\small
\begin{threeparttable}
\caption{Identifying the cohorts versus certifying the risk they create}\label{p2:tab:certvsid}
\begin{tabular}{lcccccccc}
\toprule
Rate spacing $\Delta\eta$ & 0.30 & 0.20 & 0.14 & 0.10 & 0.08 & 0.06 & 0.05 & 0.04\\
\midrule
Identify weights ($\pm0.01$), $N$ & 7 & 88 & 775 & 5{,}928 & 22{,}741 & 128{,}333 & 383{,}859 & 1{,}466{,}362\\
Certify $S$ (1\%), $N$ & 18 & 14 & 13 & 12 & 11 & 11 & 11 & 11\\
\bottomrule
\end{tabular}
\begin{tablenotes}\footnotesize\item The two-cohort moment-matched pairs of Chapter~\ref{ch:p1} ($q=0.2$, mean rate $0.4$, $\tau=0.01$, $a=0.5$, $\kappa=0.8$, per-tap noise $1.6\times10^{-3}$). Identification: Le Cam requirement at 95\% (Chapter~\ref{ch:p1}, Table on two-point sample complexity). Certification: sufficient $N$ from Theorem~\ref{p2:thm:certcost}(i) with memory bound equal to the largest rate in the pair and $\bar s=a+\kappa q+0.02$. The two economies of each pair have the same summed response, $2.353$. Kind: recomputed (\texttt{research\_extension/information\_safety}, test T03).\end{tablenotes}
\end{threeparttable}
\end{table}

    Two refinements are in the full version. If the noise in the taps of one experiment is correlated, the certificate of part (i) can under-cover (75\% at nominal 95\% in simulation), while a Student-$t$ certificate on the summed taps and a Hotelling region over the window keep their level \citep{p2:hotelling1931}; and experiments can be run one at a time with an anytime-valid version of the certificate that stops once the data decide.

    \section{An identification-free pre-trade guard}\label{p2:sec:guard}

    \subsection{Maximality}
A dealer with inventory $I$ chooses an admission level $\alpha\in[0,1]$ (the fraction of incoming algorithmic flow it accepts, or the depth it quotes). The closed-loop operator is $J(\alpha)=J_0+\alpha B_F$ with price impulse response $g_k(\alpha)$. Shocks satisfy $|u_t|\le\bar u$. Over horizon $H$ under a reset initial state $z_0=0$ (or after subtracting a known or directly observed initial drift $b_t=CJ(\alpha)^t z_0$ from the remaining budget), the shock-driven price change is $\sum_{k<H}g_k(\alpha)u_{t-k}$, and the dealer has a loss budget $\mathcal B$. The contract is: inventory $I(\alpha)$ is held fixed over the horizon, the initial state is reset ($z_0=0$) or its drift is observed, each period's shock is bounded by $\bar u$, and loss is the mark-to-market loss at any date in the horizon; because the partial sums $\sum_{k<h}|g_k|$ increase in $h$, the worst date is the last. Margin, hedging, fills and non-zero initial states are treated in the full version.

\begin{lemma}[Summed-gain guard is maximal]\label{p2:thm:guard}
The worst-case loss over horizon $H$ is
\[
L_{\max}(\alpha)=|I(\alpha)|\,\bar u\sum_{k<H}|g_k(\alpha)|.
\]
Hence the guard ``admit $\alpha$ iff $L_{\max}(\alpha)\le\mathcal B$'' is necessary and sufficient for bounded-shock solvency: it admits every safe $\alpha$ and no unsafe one. Replacing $\sum_k|g_k|$ by $\max_k|g_k|$ can admit unsafe levels; replacing it by $H\max_k|g_k|$ is safe but conservative.
\end{lemma}

\begin{proof}
By H\"older's inequality $|\sum_kg_ku_{t-k}|\le\bar u\sum_k|g_k|$, with equality for $u_{t-k}=\bar u\,{\rm sign}(g_k)$, which is admissible. Multiplying by $|I(\alpha)|$ gives $L_{\max}$, so the admission set equals the safe set. Since $\max_k|g_k|\le\sum_k|g_k|\le H\max_k|g_k|$, the peak guard under-states loss and the $H$-scaled guard over-states it.
\end{proof}

    The identity is the classical $\ell_1$ gain of a convolution operator \citep{p2:vidyasagar1986,p2:dahleh1987,p2:dahleh1995}; we use it rather than claim it. The finance content is that the dealer's admission level changes the operator, so the guard is a rule over an action set; that its single input is invariant to how the traders split into cohorts and is identified from prices; and that the same statistic prices the corrective rebate in equilibrium (Proposition~\ref{p2:prop:equil}). The guard coincides with a brute-force adversary's maximal safe level in 200 of 200 random states, and the peak-gain guard admitted an unsafe level in 2{,}000 of 2{,}000. An infinite-horizon gain is computed on a finite horizon with a contraction-block tail bound.

    \subsection{Robustness to estimation error}
\begin{theorem}[Robust gain]\label{p2:thm:robust}
Let $M=\sup_k\|J^k\|$ and $S=\sum_k\|J^k\|$. If $\|E\|\le\epsilon$ and $\epsilon S<1$, then $\sup_k\|(J+E)^k\|\le M/(1-\epsilon S)$.
\end{theorem}

\begin{proof}
Let $P_k=(J+E)^k$. Expanding, $P_k=J^k+\sum_{j<k}J^{k-1-j}EP_j$. With $m_k=\|P_k\|$, $m_k\le M+\epsilon\sum_{j<k}\|J^{k-1-j}\|m_j\le M+\epsilon S\max_{j<k}m_j$. Since $m_0=1\le M$, induction gives $m_k\le M/(1-\epsilon S)$ for all $k$.
\end{proof}

Theorem~\ref{p2:thm:robust} is deterministic: it holds whenever $\|E\|\le\epsilon$. A statistical guarantee follows only through a confidence set for the operator, and then holds with that set's coverage probability, not with certainty. The bound held in all 30,000 random perturbations of 100 non-normal operators, with median tightness $0.50$ (realized gain over bound; maximum $0.57$). In a simulated window estimator of the kind used in Section~\ref{p2:sec:evidence} (600 transitions), a residual-bootstrap 95\% ball for the operator covered the true $J$ in $94.7\%$ of 300 replications, with median radius $0.086$; but $\epsilon S<1$ held in only $8\%$ of windows. The certificate is valid but often unavailable at realistic sample sizes; longer windows or pooled estimates are needed to use it in practice.

The guard needs the summed gain, which is invariant to cohort splitting (Proposition~\ref{p2:prop:invariance}) and certified with $\Theta(\sigma^2\epsilon^{-2}\log(1/\epsilon))$ experiments whatever the spacing of the memories (Theorem~\ref{p2:thm:certcost}); it does not need the algorithms' identities, which are expensive (Chapter~\ref{ch:p1}). In a two-cohort economy with rates $0.7\pm\Delta$, the relative error of the estimated summed gain is flat in $\Delta$ (log--log slope $-0.15$), while the error of the cohort weights grows as $\Delta^{-3.07}$.

    \subsection{An uncertainty-aware guard}\label{p2:sec:robustguard}
The guard of Lemma~\ref{p2:thm:guard} assumes the operator is known. With data, the relevant question is which admission levels are safe for every economy the data allow. We answer it for a venue gate (a venue that admits a fraction of incoming algorithmic flow) in the positive cohort economy, where admitting a fraction $\alpha$ of the algorithmic flow scales every cohort's weight, $M\mapsto\alpha M$, so that $S(\alpha)=\kappa/(1-a-\kappa\alpha\Phi)$.

\begin{lemma}[Soundness of compatible-set guards]\label{p2:lem:sound}
If a data-dependent set $\mathcal C(D)$ contains the true economy with probability at least $1-\gamma$, then every action whose worst-case loss is within budget for all economies in $\mathcal C(D)$ is within budget for the true economy with probability at least $1-\gamma$, including an action selected from the data. If no action qualifies, the guard certifies nothing: it halts new admission and flags the state, because existing inventory can be unsafe even with no new trading.
\end{lemma}

\begin{theorem}[Optimal uncertainty-aware admission]\label{p2:thm:robustguard}
Under the reset or observed-drift contract ($z_0=0$, or known drift subtracted from $\mathcal B$), let position $I$, shock bound $\bar u$ and budget $\mathcal B$ be given, and let $\Phi^+$ be an upper bound on $\Phi$ valid with probability at least $1-\gamma$ (for example $\Phi^+=(1-a-\kappa/\hat S^+)/\kappa$ from Theorem~\ref{p2:thm:certcost}, or $\Phi_+$ from Theorem~\ref{p2:thm:window}). Put $c_I=1-a-\kappa|I|\bar u/\mathcal B$.
\begin{enumerate}[label=(\roman*)]
\item \emph{Closed form and interval.} If $c_I\ge0$, the levels that are safe for every $\Phi\le\Phi^+$ form the interval $[0,\alpha^*]$ with $\alpha^*=\min\{1,\ c_I/(\kappa\Phi^+)\}$.
\item \emph{Safety.} With probability at least $1-\gamma$, every admitted level is safe for the true economy.
\item \emph{Maximality.} Every level above $\alpha^*$ is unsafe for an economy with $\Phi=\Phi^+$; when $\Phi^+$ is attained by a compatible economy (as for $\Phi_+$ in Theorem~\ref{p2:thm:window}), no rule using the same information can admit more without admitting an unsafe level.
\item \emph{Infeasibility.} If $c_I<0$, no level is safe, including $\alpha=0$: the guard halts admission and the position must fall to $|I|\le\mathcal B(1-a)/(\kappa\bar u)$.
\item \emph{Minimax capacity loss.} On the class of Theorem~\ref{p2:thm:certcost} with $\Phi$ bounded away from zero, the certificate-based guard gives up $\alpha^\circ-\alpha^*=O\bigl(\sigma\sqrt{\log N/N}\bigr)$ of the oracle capacity $\alpha^\circ=\min\{1,c_I/(\kappa\Phi)\}$ with probability at least $1-2\gamma$, and every rule that admits an unsafe level with probability at most $\beta<1$ gives up at least $c\,\sigma/\sqrt N$ in expectation at some economy of the class. The guard is minimax-optimal in lost capacity up to a logarithmic factor.
\end{enumerate}
\end{theorem}

Monotonicity of $S(\alpha)$ in $\alpha$ and $\Phi$ gives (i)--(iv); (v) combines the certificate's excess with Proposition~\ref{p2:prop:twoworld} applied to two economies whose capacities differ by order $\sigma/\sqrt N$ (Appendix~\ref{p2:app:proofs}). In 500 random economies the closed form matched a brute-force search over 1{,}001 admission levels, computed from the state matrices rather than from \eqref{p2:eq:SPhi}, in every case; every safe set was an interval; and the 28 infeasible states were flagged (test T06). Across 18 designs the certificate covered the true summed response at least at its nominal level (minimum 93.6\% at nominal 90\%), with a median excess of 3--15\% when the memory bound is $0.6$ or $0.9$ and about 230\% when it is $0.97$, where the tail term dominates (test T04).

    When admission also changes the book, deepening it and speeding reversion as well as adding feedback, risk can first rise and then fall with the admission level, and the safe levels need not form an interval: certifying the largest safe level then no longer certifies those below it. The full version certifies the whole safe set by interval branch and bound \citep{p2:hansen2004}, with soundness and completeness; in a book where admission raises feedback linearly and liquidity quadratically the certified safe set is $[0,0.053]\cup[0.178,1]$.

    \paragraph{Why the peak guard is unsafe in a real book.} At the median BTCUSDT estimates of Section~\ref{p2:sec:evidence}, price impact is persistent ($\hat a\approx0.97$ per minute), so the summed gain over 15 minutes is about 13 times the peak. A guard calibrated to the peak would admit inventory about 13 times larger than is safe against a sustained adversarial flow of the declared size, although that book does not amplify. Persistence alone makes peak-based limits unsafe.

    \paragraph{The guard on a dealer's balance sheet.} The full version carries the guard through maintenance margin, hedging and worst-case fills, and proves forward solvency at every step at which an admissible level exists; when none exists the guard halts new fills and flags the state, because existing inventory can still breach margin. In a matched-shock Monte Carlo inside the model, with every rule calibrated to the same trading capacity, the guard has no forced liquidations against 31.6\% with no control, 31.4\% under a trailing volatility limit and 20.3\% under an inventory cap. With a hidden cohort state the drift of future prices is not identification-invariant under a Euclidean bound, whereas the natural interval bound supports the following sufficient invariant box. Cumulative unforced drift still depends on the updating rates:
    
\begin{proposition}[Hidden-state drift and an invariant-box guard]\label{p2:prop:driftguard}
Let $0\le a<1$, $\kappa>0$, $q_i\ge0$, $0<\eta_i\le1$ and $\Phi=\sum_iq_i\le\Phi^+$. Admission $\alpha\in[0,1]$ scales cohort demand, with $a+\alpha\kappa\Phi^+<1$. Use the original price-first coordinates $x_{t+1}=ax_t+\alpha\kappa\sum_iq_iy_{i,t}+\kappa u_{t+1}$ and $y_{i,t+1}=(1-\eta_i)y_{i,t}+\eta_ix_{t+1}$.
\begin{enumerate}[label=(\roman*)]
\item \emph{Cumulative unforced drift.} For fixed weights and rates and $u=0$,
\begin{equation}\label{p2:eq:driftenv}
\sup_{\|y_0\|_\infty\le\bar y_0}\sum_{t=1}^\infty|x_t|
=\frac{a|x_0|+\alpha\kappa\bar y_0\sum_iq_i/\eta_i}{1-a-\alpha\kappa\Phi}.
\end{equation}
Aligned initial memories attain the bound. Thus the cumulative quantity is not determined by $\Phi$ and $\bar y_0$ alone when rates can approach zero.
\item \emph{Exact finite-horizon pointwise envelope.} Write $c^\top=e_x^\top$, $b=(\kappa,\kappa\eta_1,\ldots,\kappa\eta_K)^\top$ and $\bar z_0=(|x_0|,\bar y_0\mathbf1^\top)^\top$. For $|u_t|\le\bar u$, the worst-case peak over $0\le t\le H$ is
\begin{equation}\label{p2:eq:driftpeak}
\max_{0\le t\le H}\left[c^\top J(\alpha)^t\bar z_0+\bar u\sum_{j=0}^{t-1}c^\top J(\alpha)^jb\right].
\end{equation}
The empty sum at $t=0$ is zero. Aligned memories and constant shocks attain every term simultaneously. This formula retains the time index shared by initial drift and forced response.
\item \emph{Identification-free sufficient box.} For fixed inventory $I\ne0$ and budget $\mathcal B>0$, let $L=\mathcal B/|I|$. If $|x_0|\le L$, $\bar y_0\le L$ and $aL+\alpha\kappa\Phi^+L+\kappa\bar u\le L$, then $|I x_t|\le\mathcal B$ for every date and every compatible positive cohort economy. For $\Phi^+>0$, $\bar u>0$ and $(1-a)L\ge\kappa\bar u$, a sufficient admission limit is
\begin{equation}\label{p2:eq:alphabox}
\alpha_{\rm box}=\min\left\{1,\frac{(1-a)L-\kappa\bar u}{\kappa\Phi^+L}\right\}.
\end{equation}
If these initial-box or feasibility conditions fail, this argument is inconclusive. If $I=0$, the fixed-inventory mark-to-market loss is zero; if $\Phi^+=0$, admission does not affect this plant. This is not a theorem about fills, changing inventory, fees or maintenance margin.
\end{enumerate}
\end{proposition}

At the Table~\ref{p2:tab:certvsid} calibration ($a=0.4$, $\kappa=1$, $\Phi^+=0.4$, $|I|=1$, $\bar u=0.48$, $\mathcal B=1$, $x_0=0$, $\bar y_0=0.25$), the invariant box admits $\alpha_{\rm box}=0.3000$, the same as the reset guard. Its capacity reduction in this calibration is $0.0000$. Indeed $0.4+0.30\times0.4+0.48=1$; every memory update is a convex combination of values in $[-1,1]$. For any larger stable admission, an economy with $\Phi=0.4$ under constant positive shocks has limit $0.48/(0.6-0.4\alpha)>1$, proving maximality in this calibration. The earlier $0.2400$ rule was conservative here; its claimed exact $0.0600$ capacity cost is withdrawn. Original result files are retained as historical records. The corrected checks are in \texttt{recomputed/R\_P2\_step2\_theory\_repaired.json} and \texttt{studies/economic\_review\_20261029/}.

    Re-optimization by the traders does not undo the guard: when a venue throttles admission, filtering cohorts lower their adaptation rates, and in the positive cohort economy every finite-horizon summed response falls with them.

    \section{The value of disclosure for safe trading}\label{p2:sec:disclosure}

    Chapter~\ref{ch:p1} measures disclosure in samples needed to identify the cohorts. For a pre-trade control the relevant value is the safe trading capacity it unlocks.
    
\begin{theorem}[Decision value of disclosure]\label{p2:thm:disclosure}
Under the risk contract of Section~\ref{p2:sec:contracts}, compare an observation contract $O$ with a richer contract $O'$ that reveals everything $O$ reveals and more, both with exact data.
\begin{enumerate}[label=(\roman*)]
\item \emph{Monotonicity.} The compatible sets are nested, $\Theta_{O'}\subseteq\Theta_O$, so $\Phi^+_{O'}\le\Phi^+_O$ and the admitted level of Theorem~\ref{p2:thm:robustguard} satisfies $\alpha^*_{O'}\ge\alpha^*_O$.
\item \emph{Strictness.} $\alpha^*_{O'}>\alpha^*_O$ if and only if $\alpha^*_O<1$, $c_I>0$, and $O'$ excludes every compatible economy with $\Phi=\Phi^+_O$. Under the window contract that economy is the upper principal representation of Theorem~\ref{p2:thm:window}: disclosure has decision value exactly when it rules out a cohort at the memory bound with the extremal weight.
\item \emph{Long memory.} Under the window contract with memory bound $\bar r$, $\alpha^*_O\to0$ as $\bar r\to1$ for any fixed window (Theorem~\ref{p2:thm:window}(iii)), whereas audited exposure, which reveals $\Phi$, admits the oracle capacity $\alpha^\circ$. The decision value of audited disclosure therefore tends to the whole oracle capacity as memory lengthens.
\item \emph{Net value.} With value $w$ per unit of admitted capacity and disclosure cost $c(O')$, the net value $w(\alpha^*_{O'}-\alpha^*_O)-c(O')$ is zero or negative whenever $O$ already identifies $\Phi$ (a window whose moment vector lies on the boundary of the moment space) or the budget does not bind.
\item \emph{Noisy proxies.} With noisy data the certificates of two contracts are confidence bounds, not nested sets, and a richer contract can admit less. Under class-tagged flows measured with noise over $T$ periods and estimated as in Chapter~\ref{ch:p1}, class $i$'s exposure $\hat\Phi_i=\hat w_i/(1-\hat r_i)$ has delta-method standard error
\[
{\rm se}(\hat\Phi_i)=\frac{1}{(1-r_i)\sqrt T}\Bigl[V_{ww}+2\Phi_iV_{wr}+\Phi_i^2V_{rr}\Bigr]^{1/2},
\]
with $V$ the asymptotic covariance of $(\hat w_i,\hat r_i)$, so the tagged bound $\sum_i\hat\Phi_i+z_{1-\gamma}{\rm se}$ admits less than the public certificate whenever its margin exceeds $\Phi^+_O-\Phi$. Intersecting the two certificates, each at level $1-\gamma/2$, gives a bound valid at level $1-\gamma$ that admits at least as much as either certificate at level $1-\gamma/2$, which restores monotonicity.
\end{enumerate}
\end{theorem}

The theorem is set inclusion plus the closed form of Theorem~\ref{p2:thm:robustguard}; its content is the location of the value. Disclosure is worth most where public certificates are weakest, in markets whose participants have long memories. In simulation (test T07; Table~\ref{p2:tab:discsim}), with two classes and one of them slow, audited exposure adds 0.027 of admitted capacity to a public certificate when the slow class's memory rate is 0.6 and 0.122 when it is 0.995, out of an oracle capacity of 0.667. Class-tagged flows measured with 10\% noise add 0.02--0.03 at memory rates up to 0.97 but \emph{lower} capacity at 0.995, because a class's exposure $q=w/(1-r)$ is badly conditioned when $r$ is near one: disclosure of a noisy proxy can be worth less than nothing. Part (v) quantifies this. In a second simulation of the same design (test T15), the standard error of the slow class's exposure rises from $0.0004$ to $0.035$ as $1-r$ falls from $0.4$ to $0.005$, roughly in proportion to $(1-r)^{-1}$ over the last decade, while the public certificate's margin grows only with the square root of its horizon. At $r=0.995$ the tagged bound admits $0.506$ against $0.549$ for the public one; the intersection admits $0.537$, close to the better of the two, and covers the true exposure in 99\% of replications (Table~\ref{p2:tab:discsim}). Chapter~\ref{ch:p1}'s evidence that Binance's and the CFTC's class positions pin down memories that prices leave unrestricted concerns class-level positions, which are proxies, not records of individual algorithms.

\begin{table}[t]\centering\small
\begin{threeparttable}
\caption{Admitted capacity under three observation contracts (simulation)}\label{p2:tab:discsim}
\begin{tabular}{lcccc}
\toprule
Memory rate of the slow class & 0.60 & 0.90 & 0.97 & 0.995\\
\midrule
Public responses (certificate, Theorem~\ref{p2:thm:certcost}) & 0.640 & 0.623 & 0.601 & 0.544\\
Class-tagged flows, 10\% noise & 0.662 & 0.655 & 0.631 & 0.507\\
Audited exposure (oracle) & 0.667 & 0.667 & 0.667 & 0.667\\
Public $\cap$ tagged, each at 97.5\% (test T15) & 0.661 & 0.653 & 0.626 & 0.537\\
\midrule
Value of audited exposure over public & 0.027 & 0.044 & 0.066 & 0.122\\
\bottomrule
\end{tabular}
\begin{tablenotes}\footnotesize\item Two classes ($q=0.06$ with the stated memory rate, $q=0.12$ with rate $0.5$), $a=0.4$, $\kappa=1$, position 1, shock bound $0.48$, budget 1; public: 500 reset experiments, per-tap noise $0.05$, certificate at 95\% with the true memory bound; tagged: 2{,}000 periods of each class's flow with 10\% measurement noise, two-stage least squares as in Chapter~\ref{ch:p1}, 95\% delta-method bound; 300 replications. Entries are mean admitted fractions $\alpha^*$. The intersection row comes from an independent run of the same design (seed 15), whose public and tagged rows are $0.639/0.662$, $0.625/0.655$, $0.600/0.631$ and $0.549/0.506$. Kind: recomputed (tests T07 and T15).\end{tablenotes}
\end{threeparttable}
\end{table}

\paragraph{Pricing versus guarding.} A conservative bound is the right input for a guard and the wrong input for a price. The corrective rebate of Proposition~\ref{p2:prop:equil} depends on $S^2$; computed from an upper bound that overstates $S$ by 10\% it over-subsidizes reversion and loses 2.4--2.8\% of first-best welfare, and by 50\%, 70--80\% (test T12). Guards should use certified upper bounds; rebates should use point estimates.

    \section{Who provides reversion? An equilibrium}\label{p2:sec:equil}

    Theorem~\ref{p2:thm:sum} says that transient risk is governed by $S=\tilde\kappa/(m-q\tilde\kappa)$, and its ingredients are chosen by participants: dealers decide how hard to lean against displacements, which sets reversion $m$, and momentum traders how strongly to chase them, which sets $q$.
    
\subsection{Setup}
There are $N$ dealers. Dealer $i$ supplies reversion intensity $s_i\ge0$ by absorbing order flow against the displacement, so $m=m_0+\sum_is_i$, where $m_0$ is reversion supplied by other means. Leaning earns a spread $v$ per unit of intensity at quadratic inventory cost, so the private value is $vs_i-\tfrac c2s_i^2$. By Lemma~\ref{p2:thm:guard} the worst-case transient loss of a position of size $Q$ facing flow shocks bounded by $U$ is $QU\,S$, so we take aggregate expected tail loss to be $\Lambda S(m,q)$ with $\Lambda>0$. Dealers bear a share $\theta<N$ of it, so each bears $\theta/N$; leveraged users and their counterparties bear the rest. Momentum traders choose aggregate intensity $q$ with private value $bq-\tfrac e2q^2$ and bear none of the tail loss. Long-run impact $\tilde\kappa$ is taken as given.

\subsection{Results}

\begin{proposition}[Equilibrium transient risk]\label{p2:prop:equil}
Let $g(m)\equiv-\partial S/\partial m=\tilde\kappa/(m-q\tilde\kappa)^2$.
\begin{enumerate}[label=(\roman*)]
\item \emph{Under-provision of reversion.} The symmetric Nash equilibrium solves $m-m_0=N v/c+\theta\Lambda g(m)/c$, the planner's optimum solves $m-m_0=N v/c+N\Lambda g(m)/c$, and $m^{\rm eq}<m^{\rm SP}$, so $S^{\rm eq}>S^{\rm SP}$.
\item \emph{Fewer dealers, more risk.} $\partial m^{\rm eq}/\partial N=(v/c)/(1-\theta\Lambda g'(m)/c)>0$, so equilibrium reversion rises and $S^{\rm eq}$ falls with the number of dealers.
\item \emph{Over-provision of momentum.} Momentum traders choose $q=b/e$; the planner chooses $q^{\rm SP}<b/e$, solving $b-eq=\Lambda\tilde\kappa^2/(m-q\tilde\kappa)^2$.
\item \emph{The corrective price is identified.} The first best for reversion is implemented by a rebate per unit of reversion
\[
\tau^*=\Bigl(1-\frac\theta N\Bigr)\Lambda\,g(m)=\Bigl(1-\frac\theta N\Bigr)\Lambda\,\frac{S^2}{\tilde\kappa},
\]
which depends on the market only through $S$ and $\tilde\kappa$.
\item \emph{Endogenous depth.} Suppose in addition that each dealer posts depth $D_i\ge0$ at private value $v_DD_i-\tfrac{c_D}2D_i^2$, that long-run impact is $\tilde\kappa=1/D$ with $D=\sum_iD_i$, and that momentum traders size their flow to the book, $q=\psi D$ with $\psi$ fixed. Then $S=\Gamma/D$ with $\Gamma\equiv S/\tilde\kappa=1/(m-\psi)$, which does not depend on depth. Symmetric equilibrium depth solves $D=Nv_D/c_D+\theta\Lambda\Gamma/(c_DD^2)$ and the planner's $D=Nv_D/c_D+N\Lambda\Gamma/(c_DD^2)$, so $D^{\rm eq}<D^{\rm SP}$; $D^{\rm eq}$ increases with $N$, so $S^{\rm eq}=\Gamma/D^{\rm eq}$ falls with the number of dealers while $S/\tilde\kappa$ does not move. The rebate per unit of depth $\tau_D^*=(1-\theta/N)\Lambda S\tilde\kappa$ implements the first-best depth.
\end{enumerate}
\end{proposition}

\begin{proof}
(i) Dealer $i$'s first-order condition is $v-cs_i+(\theta/N)\Lambda g(m)=0$; the planner's is $v-cs+\Lambda g(m)=0$. Summing over dealers gives the two fixed-point equations. The right-hand sides are decreasing in $m$ because $g$ is decreasing on $m>q\tilde\kappa$, and the planner's lies above the dealers' for every $m$ when $\theta<N$, so its crossing with the line $m-m_0$ lies to the right. (ii) follows by implicit differentiation, since $g'<0$. (iii) $\partial S/\partial q=\tilde\kappa^2/(m-q\tilde\kappa)^2>0$ is the marginal external cost of momentum that traders ignore. (iv) The rebate closes the wedge between the two first-order conditions, and $g=S^2/\tilde\kappa$ by Theorem~\ref{p2:thm:sum}. (v) With $q\tilde\kappa=\psi$, Theorem~\ref{p2:thm:sum} gives $S=\tilde\kappa/(m-\psi)=\Gamma/D$ and $\partial S/\partial D_i=-\Gamma/D^2$. Dealer $i$'s first-order condition for depth is $v_D-c_DD_i+(\theta/N)\Lambda\Gamma/D^2=0$ and the planner's $v_D-c_DD_i+\Lambda\Gamma/D^2=0$; summing gives the two fixed points. In each, the left side $D$ increases and the right side decreases in $D$, so each has a unique positive solution, the planner's to the right of the dealers' because $\theta<N$; and $D-\theta\Lambda\Gamma/(c_DD^2)=Nv_D/c_D$ has a left side increasing in $D$, so $D^{\rm eq}$ increases with $N$. The rebate closes the wedge $(1-\theta/N)\Lambda\Gamma/D^2=(1-\theta/N)\Lambda S\tilde\kappa$.
\end{proof}

At $v=0.002$, $c=0.5$, $\theta=0.3$, $\Lambda=10^{-4}$, $m_0=0.005$, $q=0.02$ and $\tilde\kappa=1$, equilibrium $S$ falls from 27.9 with one dealer to 6.8 with forty, and the planner's $S$ is 1.5--2.7 times lower throughout; the rebate computed from $S$ alone coincides with the first-order wedge to machine precision. With $m=0.03$, $b=0.3$ and $e=20$, momentum traders choose $q=0.015$ against a planner's $0.0062$.

\paragraph{Interpretation.} Reversion is a public good: each dealer who leans against a displacement lowers the tail exposure of every leveraged position in the market but is paid only for his own fills. Momentum is the mirror image. Part (iv) closes the loop with the guard: the policy price of the externality is a function of the same summed response that the guard bounds, and Proposition~\ref{p2:prop:invariance} and Theorem~\ref{p2:thm:certcost} show that this quantity is identified from public data without identifying the traders. A venue that rebates reversion in proportion to $S^2/\tilde\kappa$, or a supervisor that sets pre-trade limits by $S$, prices the externality with an observable.

    Across the liquidity terciles of 30 Binance perpetuals the summed response per dollar of flow rises steeply as depth falls (Spearman correlation with log depth $-0.94$), as part (ii) predicts, but the unit-free multiple $S/|h_0|$ is the same in all three terciles (medians $46.4$, $47.3$ and $46.9$): thinner books carry more cumulative risk because their impact is larger, not because they revert more slowly. That is the pattern of part (v), which we added after seeing the comparison, so the match is a consistency check, not a test.

    \section{Evidence}\label{p2:sec:evidence}

    The empirical program tests the theory's content in order-book data. Every test fixed its specification, hashed it and committed it before its test data were opened, and every test set was opened once except where Table~\ref{p2:tab:alltests} says otherwise. Data are public: Binance USD-M futures and spot archives (one-minute klines with taker volume, depth within $\pm1$--$5\%$ of the mid at about 30-second intervals, open interest), verified against published checksums and screened by validation rules fixed in advance; Bybit trade and order-book archives; Coinbase trades; and consolidated US trades. The out-of-sample criterion is a one-sided Clark--West test \emph{and} a date-block bootstrap interval for the gain in $R^2$ that excludes zero.

    \subsection{The deepest book does not amplify}

    In BTCUSDT futures on 718 training days, the estimated one-period margin $\hat C_3$ is negative on 93.5\% of days and the median implied response peaks at impact. Local projections show no overshoot once future flow is held fixed. A registered test of whether six-hour estimates of $\hat C_3$ forecast the next hour's maximum adverse excursion beyond volatility, flow and depth controls fails ($+0.018$ points of $R^2$, interval including zero), and a volatility flag classifies tail hours far better. In spot data, a unit-free measure of transient amplification on five crash days is indistinguishable from one and from a no-feedback null, and a registered forecasting test fails. Binance's depth bands appear to show replenishment of the side being hit; on a fixed price grid built from Bybit's full order books the sign reverses, so the replenishment is an artefact of band re-centring and depletion is present, but it is not large enough to produce an overshoot. Where the theory predicts nothing, we find nothing.

    \subsection{Thin books: the summed response forecasts the size of excursions}\label{p2:sec:thin}

    Thirty Binance perpetuals were fixed from one 2024 training day before any statistic was computed, spanning BTCUSDT to symbols with under \$5M of daily volume. The amplifying regime is more frequent in thinner books (16.1\% of symbol-days with $\hat C_3>0$ in the thin tercile against 9.4\% in the thick one), although the average margin is negative everywhere. Measured feedback is larger in \emph{deeper} books: feedback and amplification are different objects.

    \begin{table}[t]\centering\small
    \begin{threeparttable}
    \caption{Registered forecasting tests in thin Binance books, 2025--2026}\label{p2:tab:thin}
    \setlength{\tabcolsep}{3pt}
    \begin{tabular}{p{0.42\textwidth}cc}
    \toprule
    Test (what changes) & $\Delta R^2$ (pts) [CI] & Clark--West $t$\\
    \midrule
    Ten thin symbols, volatility and flow controls & $+1.216$ [$+0.273$, $+7.987$] & 8.53\\
    Ten new thin symbols drawn by a seeded rule & $+0.845$ [$+0.303$, $+6.371$] & 11.66\\
    Fresh days, 20 thin symbols, price-impact controls in the baseline & $+1.583$ [$+0.588$, $+3.237$] & 9.46\\
    Fresh days, full volatility-forecasting baseline (HAR, EWMA, hour effects) & $+2.787$ [$+1.425$, $+4.944$] & 12.59\\
    \quad the guard's statistic $S_{60}$ instead of $\hat C_3$ (secondary) & $+2.098$ [$+1.165$, $+3.597$] & ---\\
    Untouched days, model frozen at its training estimates & $+4.196$ [$+1.551$, $+8.233$] & ---\\
    Untouched thin symbols & $+0.727$ [$+0.377$, $+1.428$] & ---\\
    Untouched thin symbols, beyond the baseline plus VPIN & $+0.562$ [$+0.274$, $+1.159$] & ---\\
    \bottomrule
    \end{tabular}
    \begin{tablenotes}\footnotesize\item Outcome: the next hour's maximum adverse excursion; candidate estimated from the previous six hours. Date-block bootstrap intervals (97.5\% for the fresh-day and untouched-day primaries, 95\% otherwise). Every row is a registered primary except the $S_{60}$ row. Full tables and secondaries in the full version.\end{tablenotes}
    \end{threeparttable}
    \end{table}

    The result holds in every registered forecasting test on thin Binance books (Table~\ref{p2:tab:thin}). It strengthens as the baseline gets harder: beyond a full volatility-forecasting baseline that already explains 18.6\% of the variance, $\hat C_3$ adds $2.8$ points of $R^2$, the volatility benchmarks themselves add nothing beyond the simpler controls, and the summed response $S$, the guard's own statistic, adds $2.1$ points. The decomposition matches the theory: a larger momentum term predicts larger excursions, a larger reversion term smaller ones, and the depth term nothing; most of the variation is in impact-weighted reversion, as $S=\tilde\kappa/(m-q\tilde\kappa)$ implies. It is not price impact (with four impact controls it keeps about 80\% of its gain, and impact alone forecasts nothing), not tick size (it survives tick controls, the removal of days at a non-training tick and of any single symbol; unit-free versions work on fresh days), not order-flow toxicity (it keeps about 73\% of its gain beyond VPIN, a comparison made post hoc and then confirmed in a registered test), and not order-flow persistence (a placebo with the same flow persistence but broken alignment with prices adds nothing). In economic terms, within a volatility quintile the top quintile of $\hat C_3$ has next-hour excursions about 26 basis points larger than the bottom quintile, enough to roughly double the hourly liquidation probability of a $20\times$ position.

    \subsection{Other venues}

    The result belongs to thin markets, not to one venue, but its size depends on the venue. In 23 thin US stocks a registered test passes at about a fifth of the Binance size ($+0.50$ points beyond the full baseline), carried again by the reversion term; mid-size stocks show a gain too. On Bybit the statistic is absent with trades alone (a registered failure) and, with order-book depth, meets its registered criterion with an effect about one eightieth of Binance's; on Coinbase spot the registered test fails. On both crypto venues where it fails, thin pairs trade sparsely (11\% of Bybit minutes have no trade, against 0.1\% on Binance), which leaves the series the statistic is estimated from much noisier.

    \subsection{The guard as a risk control}

    The guard's theorems concern worst-case loss under a certified operator; whether its statistic helps size positions in practice is a separate question, and the answer is: modestly, and not instead of volatility. A six-hour estimate scaled into a volatility limit does not improve it (a registered failure), and limits built on the estimate alone are worse than a static limit. With an operator certified from three days of data, the scaled limit lowers tail loss per unit of position by 2.4\% (a narrow registered pass). Used as the theorem uses it, as a one-sided throttle, the certified guard lowers tail loss against a static limit by 8\% on untouched days and 23\% on untouched symbols (registered passes), but not against a volatility cap at matched activity. Over 18-hour holding periods, an uncapped version failed its registered capacity test; a capped throttle of the kind Rule~15c3-5 requires, registered after that failure, lowers the forced-liquidation rate of a $20\times$ position by 7.4 percentage points against a volatility cap, mostly because it admits 11\% less exposure; scaled to the same mean exposure the two rules are indistinguishable.\label{p2:sec:p24}

    \subsection{What the effect is not: the mechanism tests}\label{p2:sec:p22}

    Two registered tests asked whether the forecasting content runs through forced liquidations, measured by five-minute open-interest contractions with adverse price moves. Both fail: $\hat C_3$ does not forecast them, and they mediate none of its forecasting gain ($-6.8\%$). A leverage interaction, specified in advance but not hash-registered, also fails. A follow-up registered after these failures finds that $\hat C_3$ forecasts price persistence over the next 15 to 60 minutes and that this persistence accounts for nearly all of its forecasting gain, but its registered placebo clause fails, so the mediation is not established.\label{p2:sec:p25} A zero-fee natural experiment leaves reversion unchanged and lowers impact. On October 10, 2025, reported descriptively, the summed response was elevated all day but did not warn of anything volatility missed in the final hours. We therefore read the evidence as showing that thin books are risky because their prices revert slowly relative to impact, which is what the statistic measures, not as a test of a deeper mechanism.

    \subsection{Every registered test}\label{p2:sec:alltests}
    Of twenty-five registered primaries, fifteen pass, nine fail and one was not tested because the test it depended on failed (Table~\ref{p2:tab:alltests}).
    
\begin{table}[!htbp]\centering\footnotesize\renewcommand{\arraystretch}{0.80}\setlength{\abovecaptionskip}{3pt}
\begin{threeparttable}
\caption{All pre-registered tests}\label{p2:tab:alltests}
\setlength{\tabcolsep}{3pt}
\begin{tabular}{p{0.31\textwidth}lp{0.27\textwidth}cc}
\toprule
Test (specification hash) & Sample & Primary statistic [CI] & Result & Opened\\
\midrule
BTCUSDT futures, MAE $\mid\hat C_3$ (\texttt{d9521218}) & 2025--26 & $\Delta R^2$ $+0.018$ [$-0.100$, $+0.132$] & fail & 1\\
BTCUSDT spot, MAE $\mid$ amplification (\texttt{b1fd7f8e}) & 2024 & $\Delta R^2$ $-0.023$ [$-0.056$, $-0.002$] & fail & 3\\
Thin books, 10 symbols (\texttt{5910b6b4}) & 2025--26 & $\Delta R^2$ $+1.216$ [$+0.273$, $+7.987$] & pass & 2\\
New thin symbols, 10 (\texttt{6aa8e2c0}) & 2025--26 & $\Delta R^2$ $+0.845$ [$+0.303$, $+6.371$] & pass & 1\\
Fresh days, 20 thin, impact controls (\texttt{65acda36}) & 2025--26 & $\Delta R^2$ $+1.583$ [$+0.588$, $+3.237$] & pass & 1\\
Position sizing, same specification & 2025--26 & $\Delta\log{\rm TLR}$ $-0.017$ [$-0.045$, $+0.011$] & fail & 1\\
Bybit, 14 coins, depth-free & 2025--26 & $\Delta R^2$ $-0.000$ [$-0.001$, $+0.001$] & fail & 1\\
Beyond volatility forecasting, 20 thin (\texttt{71983f1a}) & 2025--26 & $\Delta R^2$ $+2.787$ [$+1.425$, $+4.944$] & pass & 1\\
Certified guard, position sizing (\texttt{b9a56af4}) & 2025--26 & $\Delta\log{\rm TLR}$ $-0.024$ [$-0.051$, $-0.0004$] & pass & 1\\
Bybit with order-book depth, 6 coins (\texttt{ec2145fa}) & 2025--26 & $\Delta R^2$ $+0.033$ [$+0.019$, $+0.053$] & pass & 1\\
Thin US stocks, 23, depth-free (\texttt{f5800a69}) & 2025--26 & $\Delta R^2$ $+0.499$ [$+0.272$, $+0.779$] & pass & 1\\
Untouched days, frozen model (\texttt{69849c96}) & 2025--26 & $\Delta R^2$ $+4.196$ [$+1.551$, $+8.233$] & pass & 1\\
Untouched days, guard vs static (same) & 2025--26 & $\Delta\log{\rm TLR}$ $-0.084$ [$-0.173$, $-0.011$] & pass & 1\\
Untouched symbols, 5 thin (\texttt{0523e476}) & 2025--26 & $\Delta R^2$ $+0.727$ [$+0.377$, $+1.428$] & pass & 1\\
Untouched symbols, beyond VPIN (same) & 2025--26 & $\Delta R^2$ $+0.562$ [$+0.274$, $+1.159$] & pass & 1\\
Untouched symbols, guard vs static (same) & 2025--26 & $\Delta\log{\rm TLR}$ $-0.256$ [$-0.497$, $-0.021$] & pass & 1\\
New venue, Coinbase, 6 thin (\texttt{3233c52b}) & 2025--26 & $\Delta R^2$ $-0.092$ [$-0.407$, $+0.158$] & fail & 1\\
Coinbase, real minus placebo (same) & 2025--26 & $+0.013$ [$-0.331$, $+0.298$] & not tested & 1\\
OI liquidation forecast, 19 thin (\texttt{2b3794e3}) & 2025--26 & $\Delta R^2_{\rm LIQ}$ $-0.795$ [$-1.529$, $-0.179$] & fail & 1\\
OI liquidation mediation, 19 thin (same) & 2025--26 & Share $-6.8\%$ [$-10.2\%$, $-3.1\%$] & fail & 1\\
Non-resetting 18h guard capacity (\texttt{d4169377}) & 2025--26 & $\Delta\overline{N}$ (k\$) $+332.8$ [$-20.4$, $+1180.8$] & fail & 1\\
Capped 18h throttle liq.\ rate (\texttt{698f9119}) & 2025--26 & $\Delta{\rm Liq}$ $-7.41$ [$-11.45$, $-4.15$] pp & pass & 1\\
Capped 18h throttle ES$_{95}$ (same) & 2025--26 & $\Delta{\rm ES}_{95}$ (k\$) $-26.4$ [$-53.1$, $-6.3$] & pass & 1\\
Persistence forecast, 19 thin (\texttt{8ebd67f3}) & 2025--26 & $\Delta R^2_{\rm pers}$ $+3.713$ [$+1.884$, $+6.035$] & pass & 1\\
Persistence mediation + placebo (same) & 2025--26 & Share $97.6\%$ [$96.6\%$, $98.3\%$] & fail & 1\\
\bottomrule
\end{tabular}
\begin{tablenotes}\scriptsize\item $\Delta R^2$ in percentage points. ``Opened'' counts executions of the test cell; the repeated openings used unchanged code and specifications (Sections~\ref{p2:sec:thin} and~\ref{p2:sec:evidence}). Intervals are 95\%, except 97.5\% for the two fresh-day primaries and the two Part 15 primaries. \end{tablenotes}
\end{threeparttable}
\end{table}

    \paragraph{How to read the passes.} The program was sequential: each registration protected one test from specification search on its own test data, but which test to run next was chosen after the earlier results were known. The passing Binance tests share training days and are not independent replications; they vary one thing at a time (symbols, days, baseline), and the US test shares none of their data. Every feature uses only minutes before the forecast hour, and a scrambling audit leaves every feature unchanged when later data are permuted. The gains are modest on top of a strong baseline, a 15\% increase in the explained variance of the next hour's excursion at most; they are not a new predictor that replaces volatility.

    \section{Tests of the information-safety results}\label{p2:sec:istests}

    The theorems of Sections~\ref{p2:sec:infosafety}--\ref{p2:sec:disclosure} concern declared contracts, which market data cannot deliver. They are therefore tested in simulations with a known truth (Table~\ref{p2:tab:simtests}).
    
\begin{table}[!tbp]\centering\footnotesize
\begin{threeparttable}
\caption{Simulation tests of the information-safety results}\label{p2:tab:simtests}
\setlength{\tabcolsep}{3pt}
\begin{tabular}{p{0.05\textwidth}p{0.36\textwidth}p{0.5\textwidth}}
\toprule
Test & Question & Result\\
\midrule
T01 & Can markets agree on early responses yet differ in risk? & Example~\ref{p2:ex:witness} in rational arithmetic: lags 0--2 equal, summed responses $1.92$ and $5.59$\\
T02 & Does Proposition~\ref{p2:prop:twoworld} bound the capacity a safe rule gives up? & Attained exactly with identical data ($0.324$ at $\beta=0.05$); respected in 9 of 9 noisy designs by the certificate rule, whose unsafe-admission rate is at most $3.2\%$\\
T03 & Is certifying risk cheaper than identifying cohorts? & Table~\ref{p2:tab:certvsid}: slope $-6.06$ in $\Delta\eta$ for identification, $+0.23$ for certification\\
T04 & Does the certificate cover at its nominal level? & 18 designs: coverage at or above nominal in all (minimum $93.6\%$ at nominal $90\%$); median excess $3$--$15\%$ with memory bound $0.6$--$0.9$, about $230\%$ at $0.97$\\
T05 & Is the corrected gate exact? & 2{,}000 cases, horizons 1--10, every sign path enumerated: per-period worst case equals $V\sum|g|$ and total-budget worst case equals $V\max|g|$ in all; the peak gain understates the per-period worst case in 1{,}800\\
T06 & Is the closed-form guard globally optimal? & 500 economies: equal to brute force over 1{,}001 levels from the state matrices in all; safe sets always intervals; 28 infeasible states flagged\\
T07 & What does disclosure add? & Table~\ref{p2:tab:discsim}\\
T08 & Does the guard depend on the channel? & Momentum learners, fixed-rule traders, depth and slow dealer reversion calibrated to $S=5$: identical guard capacity ($0.20$); one-period margins differ ($-0.34$ to $-0.20$)\\
T09 & What breaks the certificate? & Coverage at nominal 95\%: as assumed $100\%$; delayed feedback $100\%$; contrarian cohorts $100\%$; depth channel $92.8\%$; memory longer than declared $86.3\%$; feedback 20\% stronger after certification $34.7\%$\\
T10 & Are the window endpoints proved? & Certified by verified polynomial sign conditions for every window; outside the 60-digit endpoints by at most $1.5\times10^{-8}$ relative\\
T11 & What if the tap noise is correlated? & (O1) certificate covers 74--76\% (exact value $75.1\%$); $t$-certificate 94--95\%; Hotelling 100\%; dependent experiments 82--85\%\\
T12 & Should rebates use certified bounds? & No: a 10\% overstatement of $S$ costs $2.4$--$2.8\%$ of first-best welfare, a 50\% overstatement $70$--$80\%$\\
T13 & What if admission changes the dynamics? & Safe sets $[0,0.035]\cup[0.386,1]$ (closed form) and $[0,0.053]\cup[0.178,1]$ (book, certified); no certified level unsafe in 433 sampled economies\\
T14 & Does sequential probing keep its error level? & No wrong decision in 7{,}000 runs; 3.7 to 107 probes against 130 for a fixed design\\
\bottomrule
\end{tabular}
\begin{tablenotes}\footnotesize\item All tests are deterministic given their seeds and run without data (\texttt{research\_extension/information\_safety/run\_all.py} $\to$ \texttt{results/T*.json}, about six minutes). Further checks: the joint per-period and total-budget worst case equals its linear-programming value in 300 random cases (T05b), and the balance-sheet step keeps the reserve in all 17{,}551 sampled states that satisfy the guard inequality (T12b). Kind: recomputed. \end{tablenotes}
\end{threeparttable}
\end{table}

The simulations confirm what the theorems claim and locate what they do not (Table~\ref{p2:tab:simtests}). The certificate is valid where its contract holds, including with delayed feedback and contrarian cohorts, and conservative when the declared memory is long. It fails in two ways that matter for practice: when memory is longer than declared, which is the case Theorem~\ref{p2:thm:window}(iii) warns about, and when feedback strengthens after certification, in which case a stale certificate covers the new risk only a third of the time. A guard must therefore be re-certified as conditions change; it cannot be set once. The remaining tests check the extensions: the window endpoints are proved rather than computed (T10), certificates that do not assume independent taps keep their level when the taps are correlated while the (O1) certificate does not (T11), non-monotone admission is certified by branch and bound (T13), and sequential probing keeps its error level (T14).

    \section{Policy implications}\label{p2:sec:policy}

    Rule 15c3-5 requires pre-trade controls that prevent orders exceeding credit or capital thresholds \citep{p2:sec2010}; static caps and daily value-at-risk limits do not see transient dynamics. Lemma~\ref{p2:thm:guard} names the quantity a solvency-motivated control should bound, inventory times the shock envelope times the summed absolute impulse response, and peak- or eigenvalue-based limits are unsafe or blind exactly where they matter. The evidence says where such a control matters (thin books, under about \$0.4M of depth within 1\% of the mid) and what it can do with public data: correct a volatility limit modestly, not replace it. Because the summed response is certified without identifying the traders, a supervisor or venue can set limits, or price a reversion rebate in proportion to $S^2/\tilde\kappa$, from an observable. Guards should use certified upper bounds; rebates should use point estimates.

    \section{Conclusion}\label{p2:sec:conclusion}

    The theory answers four questions. A book amplifies when depletion outweighs impact-weighted reversion; cumulative transient risk is long-run impact over reversion net of momentum; the guard on the summed response is the largest safe admission rule and needs no identification of the traders; and the information that certifies a safe decision is cheap in exactly the directions in which identification is expensive. The evidence is mixed in an informative way: the deepest book does not amplify and the statistic is useless there, while in thin books it forecasts the size of adverse excursions beyond every baseline we tried, at a size that depends on the venue, and as a practical control it adds to volatility only modestly. Message-level data from a regulated market, data on liquidity-provider identities, and a market observed near the stability boundary would settle what remains open.

    \paragraph{Data, code and AI use.} All data are public and every number is produced by archived code from them; specifications, hashes and unlock counters are in the replication package. The authors used Claude (Anthropic) to write and debug analysis code, check derivations and draft and edit text; every number comes from code the authors executed, and the authors take full responsibility for the content.
    
\begin{thebibliography}{99}\setlength{\itemsep}{2pt}\small
\bibitem[21Shares(2025)]{p2:shares2025} 21Shares (2025). Record crypto liquidations amid tariff shock. Research note, October 13, 2025. \url{https://www.21shares.com/research/record-crypto-liquidations-amid-tariff-shock}.
\bibitem[Achieser(1956)]{p2:achieser1956} Achieser, N. I. (1956). \emph{Theory of Approximation}. New York: Frederick Ungar.
\bibitem[Andersen and Bondarenko(2014)]{p2:andersen2014} Andersen, T. G., and O. Bondarenko (2014). VPIN and the flash crash. \emph{Journal of Financial Markets} 17, 1--46.
\bibitem[Bouchaud et~al.(2004)]{p2:bouchaud2004} Bouchaud, J.-P., Y. Gefen, M. Potters, and M. Wyart (2004). Fluctuations and response in financial markets: The subtle nature of ``random'' price changes. \emph{Quantitative Finance} 4, 176--190.
\bibitem[Bouchaud et~al.(2018)]{p2:bouchaud2018} Bouchaud, J.-P., J. Bonart, J. Donier, and M. Gould (2018). \emph{Trades, Quotes and Prices: Financial Markets Under the Microscope}. Cambridge University Press.
\bibitem[Brunnermeier and Pedersen(2009)]{p2:brunnermeier2009} Brunnermeier, M. K., and L. H. Pedersen (2009). Market liquidity and funding liquidity. \emph{Review of Financial Studies} 22, 2201--2238.
\bibitem[Clark and West(2007)]{p2:clark2007} Clark, T. E., and K. D. West (2007). Approximately normal tests for equal predictive accuracy in nested models. \emph{Journal of Econometrics} 138, 291--311.
\bibitem[Corsi(2009)]{p2:corsi2009} Corsi, F. (2009). A simple approximate long-memory model of realized volatility. \emph{Journal of Financial Econometrics} 7, 174--196.
\bibitem[Dahleh and Diaz-Bobillo(1995)]{p2:dahleh1995} Dahleh, M. A., and I. J. Diaz-Bobillo (1995). \emph{Control of Uncertain Systems: A Linear Programming Approach}. Prentice Hall.
\bibitem[Dahleh and Pearson(1987)]{p2:dahleh1987} Dahleh, M. A., and J. B. Pearson (1987). $\ell^1$-optimal feedback controllers for MIMO discrete-time systems. \emph{IEEE Transactions on Automatic Control} 32, 314--322.
\bibitem[Donoho(1994)]{p2:donoho1994} Donoho, D.~L. (1994). Statistical estimation and optimal recovery. \emph{Annals of Statistics} 22, 238--270.
\bibitem[Easley et~al.(2012)]{p2:easley2012} Easley, D., M. M. L\'opez de Prado, and M. O'Hara (2012). Flow toxicity and liquidity in a high-frequency world. \emph{Review of Financial Studies} 25, 1457--1493.
\bibitem[Hansen and Walster(2004)]{p2:hansen2004} Hansen, E., and G. W. Walster (2004). \emph{Global Optimization Using Interval Analysis}, 2nd ed. New York: Marcel Dekker.
\bibitem[Harvey(2017)]{p2:harvey2017} Harvey, C. R. (2017). Presidential address: The scientific outlook in financial economics. \emph{Journal of Finance} 72, 1399--1440.
\bibitem[Hotelling(1931)]{p2:hotelling1931} Hotelling, H. (1931). The generalization of Student's ratio. \emph{Annals of Mathematical Statistics} 2, 360--378.
\bibitem[J.P. Morgan/Reuters(1996)]{p2:riskmetrics1996} J.P. Morgan/Reuters (1996). \emph{RiskMetrics---Technical Document}, 4th ed. New York: J.P. Morgan.
\bibitem[Karlin and Studden(1966)]{p2:karlin1966} Karlin, S., and W. J. Studden (1966). \emph{Tchebycheff Systems: With Applications in Analysis and Statistics}. New York: Interscience.
\bibitem[Kirilenko et~al.(2017)]{p2:kirilenko2017} Kirilenko, A., A. S. Kyle, M. Samadi, and T. Tuzun (2017). The Flash Crash: High-frequency trading in an electronic market. \emph{Journal of Finance} 72, 967--998.
\bibitem[Krein and Nudel'man(1977)]{p2:krein1977} Krein, M. G., and A. A. Nudel'man (1977). \emph{The Markov Moment Problem and Extremal Problems}. Translations of Mathematical Monographs 50. Providence, RI: American Mathematical Society.
\bibitem[Le Cam(1986)]{p2:lecam1986} Le Cam, L. (1986). \emph{Asymptotic Methods in Statistical Decision Theory}. New York: Springer.
\bibitem[Milanese and Vicino(1991)]{p2:milanese1991} Milanese, M., and A. Vicino (1991). Optimal estimation theory for dynamic systems with set membership uncertainty: An overview. \emph{Automatica} 27, 997--1009.
\bibitem[Nosek et~al.(2018)]{p2:nosek2018} Nosek, B. A., C. R. Ebersole, A. C. DeHaven, and D. T. Mellor (2018). The preregistration revolution. \emph{Proceedings of the National Academy of Sciences} 115, 2600--2606.
\bibitem[Obizhaeva and Wang(2013)]{p2:obizhaeva2013} Obizhaeva, A. A., and J. Wang (2013). Optimal trading strategy and supply/demand dynamics. \emph{Journal of Financial Markets} 16, 1--32.
\bibitem[Rump(2010)]{p2:rump2010} Rump, S. M. (2010). Verification methods: Rigorous results using floating-point arithmetic. \emph{Acta Numerica} 19, 287--449.
\bibitem[SEC(2010)]{p2:sec2010} U.S. Securities and Exchange Commission (2010). Risk management controls for brokers or dealers with market access. Release No.\ 34-63241.
\bibitem[Shleifer and Vishny(1992)]{p2:shleifer1992} Shleifer, A., and R. W. Vishny (1992). Liquidation values and debt capacity: A market equilibrium approach. \emph{Journal of Finance} 47, 1343--1366.
\bibitem[Sornette et~al.(2023)]{p2:sornette2023} Sornette, D., S. Lera, J. Lin, and K. Wu (2023). Non-normal interactions create socio-economic bubbles. \emph{Communications Physics} 6, 261.
\bibitem[Tu et~al.(2018)]{p2:tu2018} Tu, S., R. Boczar, and B. Recht (2018). Minimax lower bounds for $\mathcal H_\infty$-norm estimation. Working paper, arXiv:1809.10855.
\bibitem[Vidyasagar(1986)]{p2:vidyasagar1986} Vidyasagar, M. (1986). Optimal rejection of persistent bounded disturbances. \emph{IEEE Transactions on Automatic Control} 31, 527--534.
\bibitem[van Waarde et~al.(2020)]{p2:vanwaarde2020} van Waarde, H. J., J. Eising, H. L. Trentelman, and M. K. Camlibel (2020). Data informativity: A new perspective on data-driven analysis and control. \emph{IEEE Transactions on Automatic Control} 65, 4753--4768.
\end{thebibliography}

    \section{Appendix: Proofs}\label{p2:app:proofs}
    
\paragraph{Proposition~\ref{p2:prop:jury}.} With one cohort, $J=\begin{psmallmatrix}a&\kappa q\\ \eta a&1-\eta+\eta\kappa q\end{psmallmatrix}$, so ${\rm tr}J=a+1-\eta+\eta\kappa q$ and $\det J=a(1-\eta)$. The Jury conditions for a real $2\times2$ matrix are $|\det J|<1$ and $1\pm{\rm tr}J+\det J>0$. Substituting $a=1-\kappa\varphi_f$: $1-{\rm tr}J+\det J=\eta(1-a-\kappa q)=\eta\kappa(\varphi_f-q)$, positive iff $q<\varphi_f$; and $1+{\rm tr}J+\det J=(4-2\eta)-\kappa[(2-\eta)\varphi_f-\eta q]$, positive iff the stated bound on $\kappa$ holds (automatically when $(2-\eta)\varphi_f\le\eta q$).

\paragraph{Proposition~\ref{p2:prop:invariance}.} Equal merged measures give equal $G$, hence equal $H$, and the impulse coefficients are the Laurent coefficients of $H$ (Chapter~\ref{ch:p1}, Theorem on merged equivalence). Under $z\mapsto Tz$, $c^\top J^kb$ is unchanged because $CT^{-1}(TJT^{-1})^kTb=CJ^kb$, while the row norm $\|CT^{-1}(TJT^{-1})^k\|$ depends on $T$. The displayed price path is the convolution of inputs with $\{g_k\}$, and its worst case over the $\ell_\infty$ ball is $\bar u\sum|g|$ by H\"older.

\paragraph{Proposition~\ref{p2:prop:inventory_S}.} (i) Taking $z$-transforms of $y_{i,t}=(1-\eta_i)y_{i,t-1}+\eta_i x_t$, $I_{i,t}=(1-\delta_{\rm inv})I_{i,t-1}+u_{i,t}$ and $u_{i,t}=q_i y_{i,t}-\psi_i I_{i,t-1}$ gives $Y_i(z)=\frac{\eta_i z}{z-1+\eta_i}X(z)$ and $(z-1+\delta_{\rm inv}+\psi_i)I_i(z)=q_i z Y_i(z)$, whence $U_i(z)=\frac{z-1+\delta_{\rm inv}}{z-1+\delta_{\rm inv}+\psi_i}q_i Y_i(z)$ and $G_{\rm inv}(z)=\sum_i U_i(z)/X(z)$ has the stated form. Evaluating at $z=1$ when the augmented operator is Schur stable gives $G_{\rm inv}(1)=\sum_i \frac{\delta_{\rm inv}}{\delta_{\rm inv}+\psi_i}q_i=\Phi_{\rm eff}(\psi)$ and $S(\psi)=H_{\rm inv}(1)=\kappa/(1-a-\kappa\Phi_{\rm eff}(\psi))$. (ii) Write $r_i=1-\eta_i$ and $\rho_i=1-\delta_{\rm inv}-\psi_i$. When $0\le\psi_i\le 1-\delta_{\rm inv}$ and $\delta_{\rm inv}\ge\eta_i$, we have $0\le\rho_i\le 1-\delta_{\rm inv}\le r_i<1$. For $r_i>\rho_i$, the $m$-th open-loop tap is
$k_{i,m}=q_i\eta_i[r_i^m-\psi_i(r_i^m-\rho_i^m)/(r_i-\rho_i)]
=q_i\eta_i[(1-v_i)r_i^m+v_i\rho_i^m]$, where
$v_i=\psi_i/(r_i-\rho_i)=\psi_i/(\delta_{\rm inv}-\eta_i+\psi_i)\in[0,1]$.
It is nonnegative. If $r_i=\rho_i$, the assumptions force $\psi_i=0$ and the tap is $q_i\eta_i r_i^m\ge0$ (with the usual $r_i^0=1$ convention).
Nonnegative open-loop taps $k_m(\psi)\ge0$ with $a+\kappa\sum_m k_m(\psi)=a+\kappa\Phi_{\rm eff}(\psi)<1$ generate a nonnegative closed-loop renewal sequence $h_t(\psi)\ge0$, so $\|h(\psi)\|_1=\sum_{t\ge0}h_t(\psi)=S(\psi)$.

\paragraph{Theorem~\ref{p2:thm:radius}.} (i) $\sup_t\|CJ^tz_0\|\le\sup_t\|CJ^t\|\|z_0\|$, with equality approached by aligning $z_0$ with the top right singular vector of $CJ^{t^*}$ at the maximizing $t^*$. (ii) The Kreiss matrix theorem with Spijker's sharp constant gives $\mathcal K(J)\le\sup_t\|J^t\|\le en\mathcal K(J)$; substitute in (i). (iii) Take $J_\theta=\begin{psmallmatrix}\rho&\theta\\0&\rho/2\end{psmallmatrix}$; then $\rho(J_\theta)=\rho$ and $(1,0)J_\theta e_2=\theta$, so $M_C\ge|\theta|$ and $r^*_{\rm exact}\le B/|\theta|\to0$.

\paragraph{Proposition~\ref{p2:prop:driftguard}.} (i) Positivity gives $|c^\top J^tz_0|\le c^\top J^t\bar z_0$, with equality for memories aligned with $x_0$ (either sign if $x_0=0$). Summing the original memory recursion, without rescaling its initial condition, gives $\sum_{t\ge0}y_{i,t}=y_{i,0}/\eta_i+\sum_{t\ge1}x_t$. Hence $(1-a-\alpha\kappa\Phi)\sum_{t\ge1}x_t=ax_0+\alpha\kappa\sum_iq_iy_{i,0}/\eta_i$, proving \eqref{p2:eq:driftenv}. (ii) Superposition and positivity give the expression in \eqref{p2:eq:driftpeak}; shocks and memories with the common sign of $x_0$ attain it. Taking the maximum over the same time index proves sharpness for fixed parameters. (iii) If $|x_t|,|y_{i,t}|\le L$, then $|x_{t+1}|\le aL+\alpha\kappa\Phi^+L+\kappa\bar u\le L$. Each subsequent memory is a convex combination of two values bounded by $L$. Induction proves the invariant box. Rearranging gives \eqref{p2:eq:alphabox}. For $\bar u>0$, its endpoint automatically has $a+\alpha\kappa\Phi^+<1$. No general maximality or positive capacity cost is inferred by adding separately maximized cumulative and pointwise quantities.

\paragraph{Theorem~\ref{p2:thm:window}.} The recursion $h_{t+1}=ah_t+\kappa\sum_{m\le t}k_mh_{t-m}$ with $h_0=\kappa\neq0$ gives $k_t$ from $h_0,\ldots,h_{t+1}$, so the window determines $k_0,\ldots,k_{n-1}$ and nothing more about $M$: any nonnegative measure on $[0,\bar r]$ with these moments, written as atoms $(r_i,w_i)$ or as a limit of such, is the merged measure of a positive economy with $\eta_i=1-r_i$ and $q_i=w_i/\eta_i$ that produces the same window. (i) The compatible measures form a convex, weak-* compact set, and $\Phi=\int f\,dM$ is linear and continuous on it ($f$ is continuous on $[0,\bar r]$), so its range is a closed interval $[\Phi_-,\Phi_+]$. Since $f^{(n)}(r)=n!/(1-r)^{n+1}>0$ on $[0,\bar r]$, $(1,r,\ldots,r^{n-1},f)$ is an extended complete Chebyshev system, and by the Markov--Krein theorem \citep[Theorems IV.1.1 and IV.6.1]{p2:krein1977} the maximum is attained at the upper and the minimum at the lower principal representation, which have at most $\lfloor n/2\rfloor+1$ atoms. The theorem applies to moment vectors in the interior of the moment space, where both principal representations exist and are unique; a vector on the boundary has exactly one representing measure, so the interval collapses to a point. $S=\kappa/(1-a-\kappa\Phi)$ is increasing in $\Phi$ on $\Phi<(1-a)/\kappa$, which gives the interval of summed responses, unbounded when $\Phi_+\ge(1-a)/\kappa$ because by convexity every $\Phi\in[\Phi_-,\Phi_+]$ is attained. (ii) If $\mu,\nu$ share moments up to order $n-1$, then $\int f\,d(\mu-\nu)=\int(f-p)\,d(\mu-\nu)$ for every polynomial $p$ of degree $n-1$, so $|\int f\,d(\mu-\nu)|\le2k_0\inf_p\|f-p\|_\infty=2k_0E_{n-1}$. With $r=\bar r(1+x)/2$, $f=(2/\bar r)/(x_0-x)$ on $[-1,1]$ with $x_0=2/\bar r-1>1$, and Chebyshev's formula gives $E_m(1/(x_0-x))=\varrho^m/(x_0^2-1)$ \citep{p2:achieser1956}; we checked it against a linear-programming minimax fit for $m\le8$. Expanding $x_0=1+2(1-\bar r)/\bar r$ gives the asymptotics. (iii) Let the moment vector $c$ be interior to the moment cone of $[0,1]$. Interior points have an open neighbourhood in the cone, and $m(r_0)=(1,r_0,\ldots,r_0^{n-1})$ has norm at most $\sqrt n$, so for some $\epsilon>0$, $c-\epsilon m(r_0)$ is interior for every $r_0\in[0,1]$ and is represented by its lower principal representation, which has no atom at $1$. Adding $\epsilon\delta_{r_0}$ gives a measure on $[0,1)$ with moments $c$ and $\Phi\ge\epsilon/(1-r_0)\to\infty$ as $r_0\to1$; by convexity every larger $\Phi$ below $(1-a)/\kappa$ is attained by a stable economy, so $S$ is unbounded.

\paragraph{Theorem~\ref{p2:thm:certcost}.} (i) The flow impulse enters the stacked state as $b=\kappa(1,\eta_1,\ldots,\eta_K)^\top$, and $h_t=e_x^\top J^tb\le\|J^tb\|_\infty\le\|J\|_\infty^t\|b\|_\infty$ because $J\ge0$. The price row of $J$ sums to $a+\kappa\Phi\le\bar s$ and cohort row $i$ to $1-\eta_i(1-a-\kappa\Phi)\le1-(1-\bar r)(1-\bar s)=\bar\lambda$, so $h_t\le\kappa\bar\lambda^t$ and $\sum_{t\ge T}h_t\le\kappa\bar\lambda^T/(1-\bar\lambda)$. $\sum_{t<T}\bar h_t$ is Gaussian with mean $\sum_{t<T}h_t$ and variance $T\sigma^2/N$, so $\sum_{t<T}\bar h_t+z_{1-\gamma}\sigma\sqrt{T/N}\ge\sum_{t<T}h_t$ with probability $1-\gamma$, and adding the tail bound gives validity. On the two-sided event of probability $1-2\gamma$, $\hat S^+-S\le2z_{1-\gamma}\sigma\sqrt{T/N}+\kappa\bar\lambda^T/(1-\bar\lambda)$; $T_\epsilon$ makes the second term at most $\epsilon/2$ and the stated $N$ makes the first at most $\epsilon/2$. (ii) Let $\theta_1$ raise the cohort weight of $\theta_0$ by $\Delta$ chosen so that $S_1-S_0=2\epsilon$; then $\|h_1-h_0\|_2=2C_0\epsilon(1+o(1))$. Validity at $\theta_1$ gives $P_1(\hat S\ge S_1)\ge1-\gamma$, hence $P_0(\hat S-S_0\ge2\epsilon)\ge1-\gamma-{\rm TV}(P_0,P_1)$. For $N$ independent Gaussian responses, ${\rm KL}=N\|h_1-h_0\|_2^2/(2\sigma^2)$, and Pinsker's inequality gives ${\rm TV}\le\sqrt N\,C_0\epsilon(1+o(1))/\sigma$. An overshoot probability below $1/2$ therefore needs ${\rm TV}>1/2-\gamma$, which is the stated bound. (iv) Write $M=\sum_iq_i(1-r_i)\delta_{r_i}$ for the merged measure, $k_m=\int r^m dM$ for the open-loop kernel and $\mathcal M_0$ for the measures on $[0,\bar r]$ of mass at most $B_0=(1-\bar r)(\bar s-a)/\kappa$; economies with merged measure in $\mathcal M_0$ lie in the class, and finite-atom measures are weak-$*$ dense in $\mathcal M_0$. \emph{Step 1 (closed loop to open loop).} The closed-loop transforms satisfy $H=\kappa/(1-az-\kappa zK)$, so $H_1-H_0=\kappa^{-1}zH_1H_0(K_1-K_0)$ and, by Young's inequality, $\|h_1-h_0\|_2\le\kappa^{-1}S_1S_0\|k_1-k_0\|_2\le C_1\|k_1-k_0\|_2$ on the class, for any number of recorded responses. \emph{Step 2 (two points).} If two economies have $S_1>S_0+\epsilon$, validity at the first and the overshoot bound at the second make $\{\hat S^+\le S_0+\epsilon\}$ have probability at least $1-\gamma$ under $P_0$ and at most $\gamma$ under $P_1$, so the total variation is at least $1-2\gamma$; for Gaussian responses this needs $\|h_1-h_0\|_2\ge2z_{1-\gamma}\sigma/\sqrt N$. Hence $\omega_S(2z_{1-\gamma}\sigma/(C_1\sqrt N))\le\epsilon$, where $\omega_S(\delta)$ is the largest difference in $S$ between two economies with $\|k_1-k_0\|_2\le\delta$. Since $\partial S/\partial\Phi=\kappa^2/(1-a-\kappa\Phi)^2\ge\kappa^2$, $\omega_S\ge\kappa^2\omega_\Phi$. \emph{Step 3 (duality).} For a depth $L$, $\omega_\Phi(\delta)$ over $\mathcal M_0$ with the first $L$ kernel values equals $\min_{c\in\mathbb R^L}\{\delta\|c\|_2+B_0\psi(g-p_c)\}$, with $g(r)=1/(1-r)$, $p_c(r)=\sum_{t<L}c_tr^t$ and $\psi(h)=(\sup h)_++(-\inf h)_+$ on $[0,\bar r]$: the inequality $\le$ is direct, and equality follows from Sion's minimax theorem because $\mathcal M_0^2$ is weak-$*$ compact \citep{p2:donoho1994}; the modulus for all kernel values is the limit as $L\to\infty$, by compactness. \emph{Step 4 (a pole forces large coefficients).} If $\|c\|_2\le R$ and $\eta=\sup_{[0,\bar r]}|g-p_c|$, then $\log(1/\eta)\le K_0\max\{R^2,(1-\bar r)^{-1}\}$ with $K_0=8\log6\sqrt{1-\bar r}\log\rho_*$, $\rho_*=(1+\sqrt{1-\bar r})^2/\bar r$. Indeed $e=g-p_c$ is analytic on the disc of radius $\rho=1-s/2$, $s=\min\{1/(4R^2),(1-\bar r)/4\}$, with $|e(1-s)|\ge1/(2s)$ and $|e|\le3/s$ on $|z|=\rho$ because $|p_c(z)|\le R(1-|z|)^{-1/2}$; the two-constants theorem on the disc slit along $[0,\bar r]$, with the harmonic measure of the slit bounded below by an annulus comparison through the Joukowski map, gives the claim. \emph{Step 5.} For every $c$, either $\eta\ge\delta\sqrt{\log(1/\delta)}$, or $\log(1/\eta)>\tfrac12\log(1/\delta)$ and Step 4 forces $\|c\|_2\ge\sqrt{\log(1/\delta)/(2K_0)}$ once $\log(1/\delta)\ge2K_0/(1-\bar r)$; so $\omega_\Phi(\delta)\ge\min\{B_0,(2K_0)^{-1/2}\}\delta\sqrt{\log(1/\delta)}$ for small $\delta$, and Step 2 gives $N\ge c\,\sigma^2\epsilon^{-2}\log(1/\epsilon)$. (v) The identification rates are Chapter~\ref{ch:p1}'s Theorems on the $K$-cohort lower bound, the matching upper bound and general memory.

\paragraph{Lemma~\ref{p2:lem:sound}.} On the event that $\mathcal C(D)$ contains the true economy, the supremum over $\mathcal C(D)$ dominates the true worst-case loss for every action, including one selected from $D$; the event has probability at least $1-\gamma$.

\paragraph{Theorem~\ref{p2:thm:robustguard}.} (i) For $c_I\ge0$, $\alpha\mapsto S(\alpha)=\kappa/(1-a-\kappa\alpha\Phi)$ is increasing, and so is the worst-case loss $|I|\bar uS(\alpha)$, which is also increasing in $\Phi$; the robust constraint is therefore the constraint at $\Phi^+$, and $|I|\bar u\kappa/(1-a-\kappa\alpha\Phi^+)\le\mathcal B$ is equivalent to $\alpha\le c_I/(\kappa\Phi^+)$. (ii) is Lemma~\ref{p2:lem:sound}. (iii) At $\Phi=\Phi^+$ the constraint binds at $\alpha^*$ when $\alpha^*<1$. (iv) At $\alpha=0$, $S=\kappa/(1-a)$. (v) Upper: $\alpha^*$ is Lipschitz in $\Phi^+$ with constant $c_I/(\kappa\Phi_{\min}^2)$ and $\Phi^+$ is Lipschitz in $\hat S^+$ with constant $1/S_{\min}^2$, so Theorem~\ref{p2:thm:certcost}(i) with $T\asymp\log N$ gives the rate. Lower: take two single-cohort economies in the class whose weights differ by $\Delta=c\,\sigma/\sqrt N$; their oracle capacities differ by a constant times $\Delta$, and ${\rm TV}\le\sqrt N\|\partial h/\partial q\|_2\Delta/(2\sigma)$, which is below $(1-\beta)/2$ for small $c$; Proposition~\ref{p2:prop:twoworld} gives an expected loss of at least $(1-\beta)/2$ times the capacity difference.

\paragraph{Theorem~\ref{p2:thm:disclosure}.} (i) If $O'$ reveals everything $O$ reveals, every economy compatible with $O'$ data is compatible with $O$ data, so suprema over $\Theta_{O'}$ are at most those over $\Theta_O$, and Theorem~\ref{p2:thm:robustguard}(i) is monotone in $\Phi^+$. (ii) $\alpha^*$ is strictly decreasing in $\Phi^+$ when $c_I>0$ and $\alpha^*<1$, and $\Phi^+_{O'}<\Phi^+_O$ exactly when no economy with $\Phi=\Phi^+_O$ remains compatible (the supremum is attained, by Theorem~\ref{p2:thm:window}(i)). (iii) follows from Theorem~\ref{p2:thm:window}(iii). (iv) If $O$ identifies $\Phi$, $\Phi^+_O=\Phi$ and no contract can raise $\alpha^*$. (v) The delta method applied to $(w,r)\mapsto w/(1-r)$, whose gradient is $(1-r)^{-1}(1,\Phi_i)$, gives the standard error. If each certificate is valid at level $1-\gamma/2$, their minimum is valid at level $1-\gamma$ by the union bound and is no larger than either, so the admitted level is at least the larger of the two admitted levels.


\renewcommand{\arraystretch}{1}
\chapter{Certify, don't identify: a theory of safety without identification}\label{ch:p3}

This chapter is the third paper of the thesis and the candidate's own work. It places the results of Chapters~\ref{ch:p1} and~\ref{ch:p2} in one framework, proves when a question about a market's unknown agents is cheap or impossible to certify, and extends the safety results to agents that are nonlinear and keep learning. Results proved here are marked \Proved{}; results proved in Chapters~\ref{ch:p1}--\ref{ch:p2} are marked accordingly; the conjecture is marked as such.

\begin{quote}\small\textbf{Summary.} Regulators ask to know which algorithms trade in a market; we ask which questions about those algorithms can be answered from market data at all, and at what cost. Every question an outside observer can ask is a functional of the agents' merged measure of memory rates and weights, observed through noisy moments. For linear questions we determine the number of experiments needed to certify the answer, up to a constant factor in accuracy, by a modulus of continuity whose dual is a problem of polynomial approximation with penalized coefficients. A question is certifiable at the parametric rate exactly when its kernel has square-summable power-series coefficients; questions singular at infinite memory cannot be certified without a memory bound. Safety depends on the agents only through their total feedback weight, which sits on the boundary: it costs exactly $\sigma^2\epsilon^{-2}\log(1/\epsilon)$ experiments, and the logarithm cannot be removed, whereas identifying the agents costs polynomially or logarithmically more. The summed-gain guard remains safe for agents whose rules are nonlinear and whose memories and weights change over time, given only a bound on their total sensitivity, and when agents keep learning a certificate lasts for its safety margin divided by their learning speed. The safety bound extends to several assets with positive cross-impact, sharply. A discontinuous guard can fail to admit a strategic equilibrium.\end{quote}

\section{Markets with adaptive agents}\label{sec:fw_model}

\begin{definition}[Cohort economy]\label{def:cohort}
A state $x_t\in\R$, the displacement of the price from its fundamental value, evolves as
\begin{equation}\label{eq:fw_plant}
x_{t+1}=a\,x_t+\kappa\,(u_t+w_{t+1}),\qquad u_t=\sum_{i=1}^K q_i\,\varphi_i(y_{i,t}),\qquad y_{i,t}=r_i\,y_{i,t-1}+(1-r_i)\,x_t,
\end{equation}
with persistence $a\in[0,1)$, impact $\kappa>0$, exogenous shocks $w_t$, and $K$ cohorts of agents with weights $q_i\ge0$, memory rates $r_i\in[0,1)$ and response rules $\varphi_i$. The economy is \emph{linear} when $\varphi_i(y)=y$ for all $i$.
\end{definition}

The exponential state $y_{i,t}$ is the reduced form of steady-state Kalman filtering and of regularized online gradient descent (Lemma~\ref{p1:lem:micro}); Section~\ref{sec:bl_robust} relaxes linearity and stationarity. In the linear economy the aggregate response is a linear filter of past states,
\[
u_t=\sum_{m\ge0}k_m\,x_{t-m},\qquad k_m=\int r^m\,dM(r),\qquad M=\sum_{i=1}^K q_i(1-r_i)\,\delta_{r_i},
\]
so the agents enter the observable dynamics only through the \emph{merged measure} $M$, a positive measure on $[0,1)$, and through its moments $k_m$. Cohorts that share a memory rate are indistinguishable, whatever their identities (Theorem~\ref{p1:thm:merged}). We write $\mathcal M(\bar r,B)$ for the positive measures on $[0,\bar r]$ of total mass at most $B$; $\bar r<1$ is a \emph{memory bound}.

\section{Questions are functionals of the merged measure}\label{sec:fw_questions}

Every question an outside observer can answer from the observable dynamics is a functional $F$ of $M$ (and of the known plant parameters $a,\kappa$). Three families matter.

\paragraph{Identification.} $F(M)=M$ itself, in a metric such as the Wasserstein distance $W_1$, or the atoms $(r_i,q_i)$ when $M$ has $K$ atoms. Chapter~\ref{ch:p1} studies these.

\paragraph{Safety.} For bounded shocks $|w_t|\le\bar w$, the worst-case transient loss of a position $I$ is $|I|\,\bar w\,\|h\|_1$, where $h$ is the closed-loop impulse response, and a pre-trade guard on $\|h\|_1$ admits every safe action and no unsafe one (Lemma~\ref{p2:thm:guard}). In the linear economy with $a\ge0$ all $h_t$ are non-negative, so $\|h\|_1$ equals the summed response
\begin{equation}\label{eq:fw_S}
S(M)=\frac{\kappa}{1-a-\kappa\,\Phi(M)},\qquad \Phi(M)=\sum_{m\ge0}k_m=\int\frac{dM(r)}{1-r}=\sum_i q_i ,
\end{equation}
whenever $a+\kappa\Phi<1$ (Theorem~\ref{p2:thm:sum}). Safety therefore depends on the agents only through one linear functional of $M$, the total feedback weight $\Phi$.

\paragraph{Other operational questions.} The response at a given lag, $k_m=\int r^m dM$; the discounted feedback $\int (1-\beta r)^{-1}dM$; the peak of the impulse response. The first two are linear in $M$; the peak is not, and a guard built on it admits unsafe actions (Chapter~\ref{ch:p2}).

\section{Observation contracts and the cost of a question}\label{sec:fw_contracts}

\begin{assumption}[Reset contract]\label{ass:reset}
The observer can run $N$ independent reset experiments. Each returns the first $L$ aggregate response taps with independent Gaussian noise, $\hat k_m=k_m+\sigma\xi_m$, $\xi_m\sim N(0,1)$, $m=0,\dots,L-1$.
\end{assumption}

This is the controlled-probe contract with observed aggregate flow (contract $I_1$ of Chapter~\ref{ch:p1}). Chapter~\ref{ch:p2} works with the closed-loop taps $h_t$ instead; for the summed response the two contracts give the same rate (Theorem~\ref{p2:thm:certcost}). Correlated tap noise, passive data and tagged disclosure are other contracts, treated in Chapters~\ref{ch:p1} and~\ref{ch:p2}.

\begin{definition}[Cost of certification]\label{def:ncert}
For a functional $F$, a class $\mathcal M$, accuracy $\epsilon>0$ and level $1-\alpha$, $N_{\rm cert}(\epsilon;F,\mathcal M)$ is the smallest $N$ for which some interval built from the data has half-width at most $\epsilon$ and covers $F(M)$ with probability at least $1-\alpha$ for every $M\in\mathcal M$.
\end{definition}

\section{The certification--identification dichotomy}\label{sec:fw_dich}

Consider a linear question $F_g(M)=\int g\,dM$, where $g$ is continuous on $[0,\bar r]$; in the examples $g(r)=\sum_{t\ge0}a_t r^t$ converges absolutely there. Write $k_L(M)=(k_0,\dots,k_{L-1})$ for the first $L$ taps and $p_c(r)=\sum_{t<L}c_tr^t$ for the polynomial with coefficient vector $c\in\R^L$. Truncating the Taylor series gives a first measure of how hard the question is. Its \emph{tail at depth} $T$ and its \emph{coefficient energy} are
\[
\tau_g(T)=B\sup_{0\le r\le\bar r}\Bigl|\sum_{t\ge T}a_t r^t\Bigr|,\qquad A_g(T)=\sum_{t<T}a_t^2,
\]
and its \emph{truncation depth} is $T_g(\epsilon)=\min\{T\ge1:\tau_g(T)\le\epsilon/2\}$; $E_T(g)=\min_{\deg p<T}\max_{[0,\bar r]}|g-p|$ is the error of the best polynomial approximation of degree below $T$. The quantity that decides the cost exactly is the \emph{modulus of continuity} of the question with respect to the observed taps,
\[
\omega_L(\delta)=\sup\bigl\{F_g(M_1)-F_g(M_0):\ M_0,M_1\in\mathcal M(\bar r,B),\ \|k_L(M_1)-k_L(M_0)\|_2\le\delta\bigr\},
\]
the largest difference in the answer between two economies whose observable taps differ by at most $\delta$. For a function $h$ on $[0,\bar r]$ put $\psi(h)=(\sup h)_++(-\inf h)_+$; $\psi(h)\ge\|h\|_\infty$, with equality when $h$ changes sign.

\begin{theorem}[Certification of linear questions]\label{thm:dich}
Under Assumption~\ref{ass:reset}, with $z_\beta$ the standard normal $\beta$-quantile and $\alpha\in(0,1/2)$:
\begin{enumerate}[label=(\roman*)]
\item \emph{Upper bound.} If $L\ge T_g(\epsilon)$, then $N_{\rm cert}(\epsilon;F_g,\mathcal M(\bar r,B))\le\bigl\lceil4z_{1-\alpha/2}^2\sigma^2A_g(T_g(\epsilon))\,\epsilon^{-2}\bigr\rceil$. \hfill\Proved
\item \emph{Truncation is necessary.} For every $T$ there are $M,M'\in\mathcal M(\bar r,B)$ whose first $T$ taps coincide and with $|F_g(M)-F_g(M')|\ge2B\,E_T(g)$. Hence any certificate that uses only the first $T$ taps has worst-case half-width at least $B\,E_T(g)$. \hfill\Proved
\item \emph{Long memory makes certification impossible.} Without a memory bound, that is on the positive measures on $[0,1)$ of mass at most $B$, if $\limsup_{r\uparrow1}|g(r)|\sqrt{1-r}=\infty$ then for every level $1-\alpha$ with $\alpha<1/4$ no finite $N$ certifies $F_g$ to any finite accuracy. \hfill\Proved
\item \emph{Exact cost.} For every depth $L$, the modulus has the dual form
\[
\omega_L(\delta)=\min_{c\in\R^L}\bigl\{\delta\|c\|_2+B\,\psi(g-p_c)\bigr\}\qquad(\delta>0),
\]
and $\omega_L(0)=2B\,E_L(g)$. If $N$ experiments certify $F_g$ to half-width $\epsilon$ at level $1-\alpha$, then $\omega_L(2z_{1-\alpha}\sigma/\sqrt N)\le2\epsilon$. Conversely, if $\omega_L(2z_{1-\alpha/2}\sigma/\sqrt N)\le2\epsilon$, an affine interval $\sum_tc_t\bar k_t+b\pm\epsilon$ certifies $F_g$ with $N$ experiments. Writing $N_\omega(\epsilon)=\min\{N:\omega_L(2z_{1-\alpha}\sigma/\sqrt N)\le2\epsilon\}$ and $\lambda_\alpha=z_{1-\alpha/2}/z_{1-\alpha}$,
\[
N_\omega(\epsilon)\;\le\;N_{\rm cert}(\epsilon;F_g,\mathcal M(\bar r,B))\;\le\;N_\omega(\epsilon/\lambda_\alpha).
\]
\hfill\Proved
\end{enumerate}
\end{theorem}

Part (iv) replaces the matching lower bound in terms of coefficient energy that earlier drafts stated as a target. That statement is false (Remark~\ref{rem:energy}); the modulus is the right quantity, and it determines the cost of every linear question up to the factor $\lambda_\alpha$ in accuracy ($\lambda_{0.05}\approx1.19$). The structure of part (iv), a modulus that governs both a lower bound and an affine interval, is the classical theory of estimating a linear functional over a convex class \citep{t:ibragimov1984,t:donoho1994}. What is specific to markets is the dual form: the cost of a question about the agents is the value of a problem of polynomial approximation of $g$ in which the polynomial's coefficients are penalized, because each coefficient multiplies a noisy tap.

\begin{proof}[Proof of (i)--(iii)]
(i) Let $\bar k_t$ be the average of the $N$ replicates of tap $t$ and $T=T_g(\epsilon)$. The estimator $\hat F=\sum_{t<T}a_t\bar k_t$ has mean $\sum_{t<T}a_tk_t$ and variance $\sigma^2A_g(T)/N$. Its bias is $|\sum_{t\ge T}a_tk_t|=|\int\sum_{t\ge T}a_tr^t\,dM(r)|\le\tau_g(T)\le\epsilon/2$ for every $M\in\mathcal M(\bar r,B)$. The interval $\hat F\pm(\epsilon/2+z_{1-\alpha/2}\sigma\sqrt{A_g(T)/N})$ therefore covers $F_g(M)$ with probability at least $1-\alpha$, and its half-width is at most $\epsilon$ once $z_{1-\alpha/2}\sigma\sqrt{A_g(T)/N}\le\epsilon/2$.

(ii) By duality for best uniform approximation, $E_T(g)=\max\{\int g\,d\nu\}$ over signed measures $\nu$ on $[0,\bar r]$ with total variation at most one that annihilate every polynomial of degree below $T$. Take a maximizer and write $\nu=\nu^+-\nu^-$. Orthogonality to constants gives $\nu^+([0,\bar r])=\nu^-([0,\bar r])=:c\le1/2$. Set $M=(B/c)\nu^+$ and $M'=(B/c)\nu^-$: both lie in $\mathcal M(\bar r,B)$, their first $T$ moments coincide, and $F_g(M)-F_g(M')=(B/c)E_T(g)\ge2B\,E_T(g)$. Two measures that produce identical data cannot both be covered by an interval of half-width below half their distance.

(iii) Compare $M_0=0$ with $M_1=c\,\delta_r$. Under $M_1$ the taps are $k_t=c\,r^t$, so the Kullback--Leibler divergence between the two data distributions over any number of taps is at most $Nc^2/(2\sigma^2(1-r^2))$. Choose $c=\sigma\sqrt{(1-r^2)/N}$, which is at most $B$ for $r$ near one; the divergence is then at most $1/2$, so by Pinsker's inequality the total variation distance between the two data distributions is at most $1/2$. The functional differs by $c\,|g(r)|=\sigma\sqrt{(1-r^2)/N}\,|g(r)|\ge\sigma\sqrt{(1-r)/N}\,|g(r)|$, which is unbounded along a sequence $r\uparrow1$ by assumption. Suppose an interval of half-width $\epsilon$ covered $F_g$ with probability at least $1-\alpha$ under both economies. Once the two values differ by more than $2\epsilon$, the event that the interval contains $F_g(M_0)$ has probability at least $1-\alpha$ under $M_0$ and at most $\alpha$ under $M_1$, so the total variation distance is at least $1-2\alpha>1/2$, a contradiction.
\end{proof}

\begin{proof}[Proof of (iv)]
\emph{Duality.} For any $c\in\R^L$ and any admissible pair, with $h=g-p_c$,
\[
F_g(M_1)-F_g(M_0)=\int h\,dM_1-\int h\,dM_0+\langle c,k_L(M_1)-k_L(M_0)\rangle\le B\psi(h)+\delta\|c\|_2,
\]
because $\int h\,dM_1\le B(\sup h)_+$ and $-\int h\,dM_0\le B(-\inf h)_+$ for measures of mass at most $B$. So $\omega_L(\delta)$ is at most the minimum. For the reverse inequality write $\mu=(M_0,M_1)$ and
\[
\Lambda(\mu,c)=\int h\,d(M_1-M_0)+\delta\|c\|_2,\qquad h=g-p_c .
\]
Since $\inf_c\{\delta\|c\|_2-\langle c,v\rangle\}$ equals $0$ if $\|v\|_2\le\delta$ and $-\infty$ otherwise, $\omega_L(\delta)=\sup_\mu\inf_c\Lambda(\mu,c)$. The set $\mathcal M(\bar r,B)^2$ is convex and compact in the weak-$*$ topology, $\Lambda$ is affine and weak-$*$ continuous in $\mu$ (because $h$ is continuous) and convex and continuous in $c$, so Sion's minimax theorem \citep{t:sion1958} gives $\omega_L(\delta)=\inf_c\sup_\mu\Lambda(\mu,c)$. The inner supremum is $B\psi(g-p_c)+\delta\|c\|_2$, attained by putting mass $B$ at a maximizer of $h$ (if $\sup h>0$) and at a minimizer (if $\inf h<0$). For $\delta>0$ the objective is convex, continuous and coercive in $c$, so the infimum is a minimum. At $\delta=0$ the same argument gives $\omega_L(0)=\inf_c B\psi(g-p_c)$; since adding a constant to $p_c$ can make $g-p_c$ change sign and then $\psi=\sup-\inf$, the infimum is $B$ times the smallest oscillation of $g-p$, that is $2B\,E_L(g)$, which is part (ii) with equality.

\emph{Lower bound.} The averages $\bar k=(\bar k_0,\dots,\bar k_{L-1})$ are sufficient and distributed as $N(k_L(M),(\sigma^2/N)I_L)$, so the total variation distance between the data distributions of two economies is $2\Phi_{\mathcal N}(\sqrt N\|\Delta k\|_2/(2\sigma))-1$, with $\Phi_{\mathcal N}$ the normal distribution function. Take an admissible pair with $\|k_L(M_1)-k_L(M_0)\|_2<2z_{1-\alpha}\sigma/\sqrt N$; the distance is then below $1-2\alpha$. If the two answers differed by more than $2\epsilon$, an interval of half-width $\epsilon$ covering each with probability $1-\alpha$ would separate the two data distributions with total variation at least $1-2\alpha$, as in (iii). Hence $\omega_L(t)\le2\epsilon$ for every $t<2z_{1-\alpha}\sigma/\sqrt N$. The modulus is concave (the admissible pairs form a convex set and the constraint is a norm), finite and nondecreasing, hence continuous on $(0,\infty)$, so $\omega_L(2z_{1-\alpha}\sigma/\sqrt N)\le2\epsilon$.

\emph{Upper bound.} Put $\delta=2z_{1-\alpha/2}\sigma/\sqrt N$, let $c$ attain the minimum in the dual form and $h=g-p_c$. The interval $\sum_tc_t\bar k_t+b\pm\epsilon$ with $b=B[(\sup h)_+-(-\inf h)_+]/2$ has bias at most $B\psi(h)/2$ for every economy, by the bounds in the duality step, and its noise is normal with standard deviation $\sigma\|c\|_2/\sqrt N$. It covers $F_g$ with probability at least $1-\alpha$ whenever $\epsilon\ge B\psi(h)/2+z_{1-\alpha/2}\sigma\|c\|_2/\sqrt N=\omega_L(\delta)/2$.

\emph{Sandwich.} The left inequality restates the lower bound. Concavity and $\omega_L(0)\ge0$ give $\omega_L(\lambda t)\le\lambda\,\omega_L(t)$ for $\lambda\ge1$. If $\omega_L(2z_{1-\alpha}\sigma/\sqrt N)\le2\epsilon/\lambda_\alpha$, then $\omega_L(2z_{1-\alpha/2}\sigma/\sqrt N)\le2\epsilon$, and the upper bound applies.
\end{proof}

\paragraph{What the theorem says.} The cost of answering a linear question is the cost of approximating $g$, on the memory range $[0,\bar r]$, by a polynomial whose coefficients are cheap to estimate. Truncating the Taylor series is one such polynomial and gives part (i); it is often not the best one. Questions singular at infinite memory, faster than $(1-r)^{-1/2}$, cannot be certified without a memory bound (part (iii)). The number of agents and the spacing of their memories do not enter. Two consequences make the picture precise.

\begin{proposition}[Which questions are parametric]\label{prop:parametric}
Let $\omega_\infty(\delta)=\inf_L\omega_L(\delta)$, the modulus when enough taps are observed. Then $\sup_{0<\delta\le1}\omega_\infty(\delta)/\delta<\infty$, so that $N_{\rm cert}=O(\sigma^2\epsilon^{-2})$ for observation depths $L$ large enough, if and only if $g$ agrees on $[0,\bar r]$ with a power series $\sum_ta_tr^t$ whose coefficients are square summable. When $g$ is given by a power series on $[0,\bar r]$ with $\bar r>0$, this is the condition $\sum_ta_t^2<\infty$ on its Taylor coefficients. \hfill\Proved
\end{proposition}

\begin{proof}
If $a\in\ell^2$, take $c=(a_0,\dots,a_{L-1})$: by Cauchy--Schwarz $\psi(g-p_c)\le2\sup_{r\le\bar r}|\sum_{t\ge L}a_tr^t|\le2\|a\|_2\,\bar r^L/\sqrt{1-\bar r^2}$, so $\omega_L(\delta)\le\delta\|a\|_2+2B\|a\|_2\bar r^L/\sqrt{1-\bar r^2}$ and $\omega_\infty(\delta)\le\delta\|a\|_2$. Conversely, suppose $\omega_\infty(\delta)\le C\delta$ for $\delta\in(0,1]$. For $\delta_n=1/n$ choose $L_n$ and $c^{(n)}\in\R^{L_n}$ with $\delta_n\|c^{(n)}\|_2+B\psi(g-p_{c^{(n)}})\le2C\delta_n$. Then $\|c^{(n)}\|_2\le2C$ and $\|g-p_{c^{(n)}}\|_\infty\le2C\delta_n/B\to0$. Padded with zeros, the $c^{(n)}$ are bounded in $\ell^2$, so a subsequence converges weakly to some $a\in\ell^2$. For each $r\in[0,\bar r]$ the sequence $(r^t)_t$ lies in $\ell^2$, hence $p_{c^{(n)}}(r)=\langle c^{(n)},(r^t)\rangle\to\sum_ta_tr^t$ along the subsequence, while $p_{c^{(n)}}(r)\to g(r)$. Uniqueness of power series that agree on an interval gives the last sentence.
\end{proof}

The square-summable power series are the restrictions to $[0,\bar r]$ of the Hardy space $H^2$ of the unit disc, whose functions grow at most like $(1-r)^{-1/2}$ as $r\uparrow1$. That is exactly the growth beyond which part (iii) rules out certification without a memory bound: below the threshold a memory bound makes a question parametric or nearly so, above it nothing helps.

The single tap $g(r)=r^m$ and the discounted feedback $g(r)=(1-\beta r)^{-1}$ with $\beta<1$ have square-summable coefficients and are parametric. The total feedback $\Phi$, the quantity that safety depends on, has $a_t\equiv1$: it sits exactly on the boundary, and it pays exactly one logarithm.

\begin{theorem}[The safety question costs exactly one logarithm]\label{thm:philog}
Let $g(r)=1/(1-r)$, so that $F_g=\Phi$. There are constants $0<c\le C$ and $\delta_0\in(0,1)$, depending only on $(\bar r,B)$, such that for $\delta\le\delta_0$
\[
c\,\delta\sqrt{\log(1/\delta)}\;\le\;\omega_L(\delta)\quad\text{for every }L,\qquad
\omega_L(\delta)\;\le\;C\,\delta\sqrt{\log(1/\delta)}\quad\text{for }L\ge\log(1/\delta)/\log(1/\bar r).
\]
Consequently $N_{\rm cert}(\epsilon;\Phi,\mathcal M(\bar r,B))\asymp\sigma^2\epsilon^{-2}\log(1/\epsilon)$ as $\epsilon\to0$, for deep enough observation, and the same holds for the summed response $S=\kappa/(1-a-\kappa\Phi)$ whenever $a+\kappa B/(1-\bar r)<1$. The logarithmic factor in the certification cost of Chapter~\ref{ch:p2} cannot be removed. \hfill\Proved
\end{theorem}

The lower bound rests on a lemma about polynomials: a polynomial that approximates $1/(1-r)$ closely on the memory range must have large coefficients, because $1/(1-z)$ has a pole at $z=1$ while a polynomial with coefficient norm $R$ grows at most like $R(1-|z|)^{-1/2}$ inside the disc.

\begin{lemma}\label{lem:hardy}
There is $K<\infty$ depending only on $\bar r$ such that for every polynomial $p(r)=\sum_{t<L}c_tr^t$ with $\|c\|_2\le R$,
\[
\log\frac1\eta\le K\max\{R^2,\,(1-\bar r)^{-1}\},\qquad \eta=\sup_{0\le r\le\bar r}\Bigl|\frac1{1-r}-p(r)\Bigr| .
\]
\end{lemma}

\begin{proof}
Put $e(z)=1/(1-z)-p(z)$, $E=[0,\bar r]$, $\rho_1=(1+\bar r)/2$ and $\rho_2=(3+\bar r)/4$. Let $s=\min\{1/(4R^2),(1-\bar r)/4\}$, $x=1-s$ and $\rho=1-s/2$, so that $\rho_1<x<\rho$ and $\rho\ge\rho_2$. By Cauchy--Schwarz $|p(z)|\le R(1-|z|^2)^{-1/2}\le R(1-|z|)^{-1/2}$ in the disc. At $z=x$ this is at most $R/\sqrt s\le1/(2s)$, so $|e(x)|\ge1/s-1/(2s)=1/(2s)$. On the circle $|z|=\rho$, $|e(z)|\le2/s+R\sqrt{2/s}\le3/s$. The function $\log|e|$ is subharmonic on the slit disc $D_\rho\setminus E$, at most $\log\eta$ on $E$ and at most $\log(3/s)$ on $|z|=\rho$, so by the two-constants theorem \citep[Theorem~4.3.7]{t:ransford1995}
\[
\log\frac1{2s}\le\log|e(x)|\le\varpi\log\eta+(1-\varpi)\log\frac3s,
\]
where $\varpi$ is the harmonic measure of $E$ in $D_\rho\setminus E$ seen from $x$. Rearranging, $\varpi\log(1/\eta)\le\log6$. To bound $\varpi$ from below, let $m_*>0$ be the minimum over $|z|=\rho_1$ of the harmonic measure of $E$ in $D_{\rho_2}\setminus E$; it is positive because a segment has positive capacity and the circle is a compact subset of the slit disc, and it depends only on $\bar r$. Harmonic measure of $E$ increases with the outer domain, so the harmonic measure in $D_\rho\setminus E$ is at least $m_*$ on $|z|=\rho_1$; it vanishes on $|z|=\rho$. The maximum principle on the annulus $\rho_1<|z|<\rho$ then gives $\varpi\ge m_*\log(\rho/x)/\log(\rho/\rho_1)\ge m_*s/(4\log(1/\rho_1))$, since $\log(\rho/x)\ge\log(1+s/2)\ge s/4$. Hence $\log(1/\eta)\le4\log6\,\log(1/\rho_1)/(m_*s)$, and $1/s=\max\{4R^2,4/(1-\bar r)\}$.
\end{proof}

\begin{proof}[Proof of Theorem~\ref{thm:philog}]
\emph{Upper bound.} Take $c=(1,\dots,1)$ of length $T=\lceil\log(1/\delta)/\log(1/\bar r)\rceil\le L$. Then $g-p_c=r^T/(1-r)$, so $\psi(g-p_c)\le\bar r^T/(1-\bar r)\le\delta/(1-\bar r)$ and $\omega_L(\delta)\le\delta\sqrt T+B\delta/(1-\bar r)\le C\delta\sqrt{\log(1/\delta)}$ for small $\delta$.

\emph{Lower bound.} By part (iv) it suffices to bound $\delta\|c\|_2+B\psi(g-p_c)$ below for every $c$. Let $R=\|c\|_2$ and $\eta=\|g-p_c\|_\infty\le\psi(g-p_c)$. If $\eta\ge\delta\sqrt{\log(1/\delta)}$, the objective is at least $B\delta\sqrt{\log(1/\delta)}$. Otherwise $\eta<\sqrt\delta$, so $\log(1/\eta)>\frac12\log(1/\delta)$, and Lemma~\ref{lem:hardy} gives $K\max\{R^2,(1-\bar r)^{-1}\}>\frac12\log(1/\delta)$. Choose $\delta_0$ with $\log(1/\delta_0)\ge2K/(1-\bar r)$; then for $\delta\le\delta_0$ the maximum is $R^2$, so $R\ge\sqrt{\log(1/\delta)/(2K)}$ and the objective is at least $\delta\sqrt{\log(1/\delta)/(2K)}$. Hence $c=\min\{B,(2K)^{-1/2}\}$ works for every $L$.

\emph{Cost.} By part (iv), $N$ certifies $\Phi$ to half-width $\epsilon$ only if $\omega_L(2z_{1-\alpha}\sigma/\sqrt N)\le2\epsilon$, which with the lower bound forces $N\ge c'\sigma^2\epsilon^{-2}\log(1/\epsilon)$ for small $\epsilon$; the upper bound and part (iv) give a matching $N\le C'\sigma^2\epsilon^{-2}\log(1/\epsilon)$. On $\mathcal M(\bar r,B)$, $\Phi\le B/(1-\bar r)$, and when $a+\kappa B/(1-\bar r)<1$ the map $\Phi\mapsto S$ is increasing with derivative $\kappa^2/(1-a-\kappa\Phi)^2$ bounded above and below, so intervals for $\Phi$ and for $S$ convert into each other with constant changes in accuracy.
\end{proof}

\begin{corollary}[The rates of this thesis as cases of one picture]\label{cor:cases}
\begin{enumerate}[label=(\alph*)]
\item \emph{Safety with a memory bound.} $N_{\rm cert}(\epsilon;\Phi)\asymp\sigma^2\epsilon^{-2}\log(1/\epsilon)$, and the same for $S$ (Theorem~\ref{thm:philog}). Part (i) alone gives the upper bound $O(\sigma^2\epsilon^{-2}\log(1/\epsilon)/\log(1/\bar r))$, the rate of Theorem~\ref{p2:thm:certcost}; Theorem~\ref{thm:philog} shows that its logarithm is necessary. \hfill\Proved
\item \emph{Safety without a memory bound.} $g(r)=(1-r)^{-1}$ satisfies the condition of (iii): no finite number of experiments certifies the risk, as Theorem~\ref{p2:thm:window}(iii) shows for finite windows. \hfill\ProvedIn{CH.~3}
\item \emph{A single response tap, and discounted feedback.} For $g(r)=r^m$, $\omega_L(\delta)\le\delta$; for $g(r)=(1-\beta r)^{-1}$ with $\beta<1$, $\omega_\infty(\delta)\le\delta/\sqrt{1-\beta^2}$. Both are parametric (Proposition~\ref{prop:parametric}). \hfill\Proved
\item \emph{Identifying $K$ cohorts.} Recovering the atoms is not a linear question. It requires moments up to order $2K-1$ to be resolved against moment-matched alternatives whose higher moments differ by $O((\Delta\eta)^{2K-1})$, which costs $(\Delta\eta)^{-(4K-2)}$ experiments (Theorems~\ref{p1:thm:kcohort} and~\ref{p1:prop:esprit}). \hfill\ProvedIn{CH.~2}
\item \emph{Identifying a general memory measure.} Recovering $M$ in $W_1$ means certifying $F_g$ uniformly over all $1$-Lipschitz $g$. Their best approximations by degree-$T$ polynomials improve only as $1/T$, and the coefficients of those polynomials grow geometrically in $T$, so accuracy improves only logarithmically in $N$: the rate $1/\log N$ of Theorem~\ref{p1:thm:general}. \hfill\ProvedIn{CH.~2}
\end{enumerate}
\end{corollary}

The link in (e) between coefficient growth and the logarithmic rate is a heuristic reading of Chapter~\ref{ch:p1}'s proof through the dual form of part (iv), not a separate result here. The ladder is now exact at its first rungs: a single tap and any discounted feedback are parametric, the undiscounted total feedback that safety depends on costs exactly one extra logarithm, questions singular faster than $(1-r)^{-1/2}$ cannot be certified without a memory bound, and identification costs polynomially or logarithmically more.

\begin{remark}[Why coefficient energy is the wrong lower-bound quantity]\label{rem:energy}
Earlier drafts stated, as a target, a lower bound $N_{\rm cert}\ge c\,\sigma^2A_g(c'T_g(\epsilon))\epsilon^{-2}$ with constants uniform over a regular coefficient class. It is false. Fix $\epsilon_0>0$, $\beta\in(0,1)$, $\eta>0$ with $\eta/(1-\beta\bar r)<1/2$, and set $g_m(r)=(\epsilon_0/B)[(1-r/\bar r)^m+\eta/(1-\beta r)]$. Its coefficients are eventually positive with polynomial growth; on $[0,\bar r]$ it takes values in an interval of length below $3\epsilon_0/(2B)$, so a constant polynomial gives $\psi(g_m-p_c)<3\epsilon_0/(2B)$ with $\|c\|_2$ bounded and the question is certified to accuracy $\epsilon_0$ with no data. Yet $T_{g_m}(\epsilon_0)=m+1$ and $A_{g_m}(c'T_{g_m}(\epsilon_0))\epsilon_0^{-2}\ge B^{-2}(m/\bar r-\eta\beta)^2\to\infty$, because the alternating binomial tail $\sum_{t=T}^m(-1)^t\binom mt=(-1)^T\binom{m-1}{T-1}$ keeps $\tau_{g_m}(T)>\epsilon_0/2$ for $T\le m$. Taylor coefficients can be large while a different polynomial approximates the question cheaply; part (iv) measures the cheapest one.
\end{remark}

\paragraph{Numerical certificates.} The script \texttt{checks/check\_modulus.py} solves the primal (pairs of measures on a grid) and the dual (coefficient vectors) of part (iv) for the total feedback, for discounted feedback with $\beta=0.8$ and for the tap $r^3$, at $\bar r\in\{0.5,0.9\}$ and $\delta$ from $10^{-1}$ to $10^{-5}$, and stores both. Without any solver, the check rescales the stored measures to exact feasibility (a lower bound on $\omega_L$) and evaluates the stored coefficients over the whole interval with a curvature allowance between grid points (an upper bound, by weak duality). The brackets agree to a relative $4\times10^{-4}$ or better. Over four decades of $\delta$, $\omega_L(\delta)/\delta$ settles at the values $\|a\|_2$ of Proposition~\ref{prop:parametric} for the two parametric questions ($1/\sqrt{1-\beta^2}=5/3$ for the discounted question and $1$ for the tap) and keeps growing for $\Phi$ (from $6.6$ to $10.1$ at $\bar r=0.9$), while $\omega_L(\delta)/(\delta\sqrt{\log(1/\delta)})$ stays between $1.1$ and $4.4$. The check also verifies $\omega_L(0)=2B\,E_L(g)$ at small depths. These computations illustrate the theorems; they are not part of their proofs.

\section{Disclosure changes the contract, not the question}\label{sec:fw_disclosure}

Disclosure regimes are observation contracts. Class-tagged flow reveals each class's own moments, which turns identification into an ordinary estimation problem and is worth a factor $(\Delta\eta)^{-(4K-4)}$ in samples (Theorem~\ref{p1:thm:discvalue}). For safety the value of disclosure is different: it raises safe capacity exactly when it rules out the extremal compatible economy, and it is worth most when memory is long, where Theorem~\ref{thm:dich}(iii) says public data certify nothing (Theorem~\ref{p2:thm:disclosure}). Supervision that wants safety should therefore ask for the disclosures that bound memory or reveal $\Phi$, not for identities.

\section{A robust safety bound: nonlinear and changing agents}\label{sec:bl_robust}

Papers~1 and~2 work with agents whose rules are linear filters with fixed memory and weight. Real agents cap positions, trade only beyond thresholds, adapt their gains, and keep learning while the operator watches. This section and the next show that the safety side survives all of this under one condition: a bound on the agents' total sensitivity.


Return to the economy of Definition~\ref{def:cohort} and let everything about the agents vary over time.

\begin{assumption}[Bounded total sensitivity]\label{ass:sector}
For each cohort $i$ and period $t$, the memory rate $r_i(t)\in[0,1)$ is arbitrary, the weight satisfies $|q_i(t)|\le\bar q_i$, and the response rule satisfies the sector condition $|\varphi_{i,t}(y)|\le s_i|y|$ for all $y$. Write $\bar\Phi=\sum_i\bar q_i s_i$ for the total sensitivity.
\end{assumption}

The sector condition covers saturating momentum ($s_i=1$ for $\tanh$), position limits, threshold and dead-zone rules ($y\mathbf 1\{|y|>\tau\}$), and any rule whose response never exceeds a multiple of its signal. It does not require the rule to be smooth, monotone, stationary or known.

\begin{theorem}[Robust safety bound]\label{thm:sector}
Under Assumption~\ref{ass:sector}, if $|a|+\kappa\bar\Phi<1$ and the economy starts at rest, then for every shock sequence with $|w_t|\le\bar w$,
\[
\sup_{t\ge0}|x_t|\;\le\;\bar w\,\frac{\kappa}{1-|a|-\kappa\bar\Phi}.
\]
When $a\ge0$, this is the worst-case loss bound of the linear stationary economy with weights $\bar q_is_i$, so the summed-gain guard of Chapter~\ref{ch:p2} computed at total sensitivity $\bar\Phi$ is safe for every economy satisfying the assumption. \hfill\Proved
\end{theorem}

\begin{proof}
Let $Z_t=\max_{s\le t}|x_s|$. Each memory state is a convex combination of its previous value and the current state, so $|y_{i,t}|\le r_i(t)|y_{i,t-1}|+(1-r_i(t))|x_t|\le\max(|y_{i,t-1}|,|x_t|)$, and by induction from $y_{i,0}=0$, $|y_{i,t}|\le Z_t$. Then
\[
|x_{t+1}|\le|a|\,|x_t|+\kappa\Bigl(\sum_i|q_i(t)|\,|\varphi_{i,t}(y_{i,t})|+|w_{t+1}|\Bigr)\le(|a|+\kappa\bar\Phi)\,Z_t+\kappa\bar w .
\]
Let $Z^*=\kappa\bar w/(1-|a|-\kappa\bar\Phi)$. If $Z_t\le Z^*$ then $|x_{t+1}|\le(|a|+\kappa\bar\Phi)Z^*+\kappa\bar w=Z^*$, so $Z_{t+1}\le Z^*$; and $Z_0=0$. For the second claim, with $a\ge0$ the linear stationary economy with weights $\bar q_is_i$ is a positive system, its impulse response is non-negative, and its worst-case peak response to shocks bounded by $\bar w$ is $\bar w\sum_th_t=\bar w\,\kappa/(1-a-\kappa\bar\Phi)$ (equation~\eqref{eq:fw_S}). The bound of the theorem therefore coincides with the quantity the guard controls (Lemma~\ref{p2:thm:guard}).
\end{proof}

\paragraph{What the theorem says.} An operator who wants a safe admission level does not need to know how many agents there are, what their memories are, whether their rules are linear, or whether they change over time. It needs one number, an upper bound on the agents' total sensitivity $\bar\Phi$. The bound is tight: when the agents are linear, stationary, positive and at their maximal weights, it is attained. Nonlinearity, memory changes and negative feedback can only make the market safer than the bound says.

\begin{remark}[What the theorem does not cover]
Rules that amplify their signal, $|\varphi(y)|>s|y|$ for the declared $s$, break the bound; so do changes in the market's own persistence $a$ or impact $\kappa$, which the operator must certify separately (Chapter~\ref{ch:p2}, Section~\ref{p2:sec:infosafety}). The bound controls the peak of the price path, which is what a fixed position's loss depends on; dynamic positions are handled by the balance-sheet guard of Chapter~\ref{ch:p2}.
\end{remark}

\begin{proposition}[Several assets]\label{prop:multiasset}
Let $x_t\in\R^d$ evolve as $x_{t+1}=Ax_t+\mathcal K\bigl(\sum_iQ_i(t)\varphi_{i,t}(y_{i,t})+w_{t+1}\bigr)$, with memory states $y_{i,t}=R_i(t)y_{i,t-1}+(I-R_i(t))x_t$, $R_i(t)$ diagonal with entries in $[0,1)$, sector conditions $|\varphi_{i,t}(y)|\le s_i|y|$ componentwise, and $\|Q_i(t)\|_\infty\le\bar q_i$. If $\|A\|_\infty+\|\mathcal K\|_\infty\sum_i\bar q_is_i<1$, then $\sup_t\|x_t\|_\infty\le\bar w\,\|\mathcal K\|_\infty/(1-\|A\|_\infty-\|\mathcal K\|_\infty\sum_i\bar q_is_i)$ for $\|w_t\|_\infty\le\bar w$. \hfill\Proved
\end{proposition}

\begin{proof}
The proof of Theorem~\ref{thm:sector} goes through with absolute values replaced by the sup norm and scalars by induced sup norms: each component of $y_{i,t}$ is a convex combination of its past value and the corresponding component of $x_t$.
\end{proof}

The induced-norm bound discards the pattern of cross-impact. For positive cross-impact, the following comparison gives the sharp replacement.

\begin{proposition}[Positive cross-impact comparison]\label{prop:multiasset_positive}
Assume the economy starts at rest, $A,\mathcal K,\overline Q_i$ are entrywise non-negative, $|Q_i(t)|\le\overline Q_i$ entrywise, and $|\varphi_{i,t}(y)|\le S_i|y|$ for a non-negative diagonal matrix $S_i$. Memory matrices $R_i(t)$ are diagonal with entries in $[0,1)$. Put $G=A+\mathcal K\sum_i\overline Q_i S_i$. If $\rho(G)<1$, then for every $\|w_t\|_\infty\le\bar w$,
\[
\sup_t\|x_t\|_\infty\le
\bar w\,\bigl\|(I-G)^{-1}\mathcal K\bigr\|_\infty.
\]
This bound is sharp over the stated class. \hfill\Proved
\end{proposition}

\begin{proof}
The Neumann series $\sum_{j\ge0}G^j=(I-G)^{-1}$ has non-negative entries. Write $\mathbf1$ for the vector of ones and set $z=\bar w(I-G)^{-1}\mathcal K\mathbf1\ge0$. Then $Gz+\bar w\mathcal K\mathbf1=z$. Suppose $|x_t|\le z$ and $|y_{i,t}|\le z$ componentwise. The plant and the sector bounds give
\[
|x_{t+1}|\le A z+\mathcal K\Bigl(\sum_i\overline Q_i S_i z+\bar w\mathbf1\Bigr)=z.
\]
Each next memory coordinate is a convex combination of its previous value and the corresponding next state coordinate, so its absolute value is also bounded by the corresponding coordinate of $z$. Induction from rest proves the componentwise bound for all times. Because the resolvent times $\mathcal K$ is non-negative, its induced sup norm is its maximum row sum, and $\|z\|_\infty$ equals the displayed bound.

For sharpness, choose $R_i(t)=0$, $Q_i(t)=\overline Q_i$, $\varphi_{i,t}(y)=S_i y$, and $w_t=\bar w\mathbf1$. The state then satisfies $x_{t+1}=Gx_t+\bar w\mathcal K\mathbf1$, so from rest $x_t=\bar w\sum_{j<t}G^j\mathcal K\mathbf1$ increases componentwise to $z$. Its supremum norm therefore approaches the bound. Sharpness means a supremum over the admitted model class, not attainment at a finite time in every economy.
\end{proof}

\paragraph{Registered numerical validation.}
% Generated by checks/export_multiasset_numbers.py from R_multiasset_recovery.json
\newcommand{\MultiassetCases}{128}
\newcommand{\MultiassetPassed}{128}

The specification and its hash were pushed before the seeded instances were generated, and the unlock counter was advanced before the single opening. The joint primary required both componentwise enclosure for nonlinear, time-varying cohorts and convergence to the bound in the maximal positive linear instance. It passed in \MultiassetPassed{} of \MultiassetCases{} registered random instances. The generated matrices, tolerances and diagnostics are recorded in \texttt{checks/R\_multiasset.json}. This is a numerical check of the proved statement under the registered construction, not evidence for it.

\section{Certificates expire: learning agents}\label{sec:bl_learning}

Theorem~\ref{thm:sector} allows weights to change, provided they stay below their bounds. When agents learn, the bounds themselves are the issue: a certificate issued today says nothing about weights next month. The theorem turns this into a computable lifetime.

\begin{corollary}[Certificate lifetime]\label{thm:lifetime}
Suppose a certificate at time $0$ establishes $|q_i(0)|\le q_i^+$ for each cohort, and the guard's loss budget is met for every total sensitivity up to $\Phi_{\max}$, that is $\kappa\bar w/(1-a-\kappa\Phi_{\max})$ times the position equals the budget. If each weight changes by at most $\gamma_i$ per period, $|q_i(t+1)-q_i(t)|\le\gamma_i$, and $\sum_iq_i^+s_i\le\Phi_{\max}$ with $\sum_i\gamma_is_i>0$, then the guard remains safe for at least
\[
T_{\rm valid}=\Bigl\lfloor\frac{\Phi_{\max}-\sum_iq_i^+s_i}{\sum_i\gamma_is_i}\Bigr\rfloor
\]
periods under the sector bounds. \hfill\Proved
\end{corollary}

\begin{proof}
On $[0,T_{\rm valid}]$, $|q_i(t)|\le q_i^++\gamma_iT_{\rm valid}=:\bar q_i$ and $\sum_i\bar q_is_i\le\Phi_{\max}$. Theorem~\ref{thm:sector} applies on that horizon, since its induction uses only periods up to the current one.
\end{proof}

The lifetime is the certificate's safety margin divided by the agents' speed of learning. This is a guaranteed lower bound, not a claim that loss first crosses the budget at this time. Realized safety can last much longer if learning slows, changes sign, or stops below the sensitivity budget. When $\sum_i\gamma_i s_i=0$, the sensitivity envelope does not grow and the same certificate remains valid for every horizon under the stated sector and plant assumptions; the displayed quotient is not used. If the initial sensitivity bound already exceeds $\Phi_{\max}$, this argument supplies no valid initial certificate. Recertification frequency can therefore be based on the certified sensitivity margin and the learning-speed bound, while additional information about the realized learning path may justify a longer horizon.

\paragraph{Formal scope.}
The finite-cohort, finite-horizon safety induction is checked by \nolinkurl{CurrentMM.thesis_sector_safety}. The drift, floor and sensitivity calculations are checked by \nolinkurl{CurrentMM.thesis_lifetime_envelope} and its helper lemmas. The deterministic certification arithmetic is checked by \nolinkurl{CurrentMM.thesis_variance_radius} and \nolinkurl{CurrentMM.thesis_interval_half_width}; the new bias lemma covers finite positive atomic measures. For several assets, \nolinkurl{CurrentMM.thesis_multiasset_comparison} checks the comparison step of Proposition~\ref{prop:multiasset_positive} for any number of assets and cohorts, nonlinear time-varying sector-bounded responses, entrywise weight bounds and all times: every non-negative supersolution bounds every state and memory coordinate; \nolinkurl{CurrentMM.thesis_multiasset_linear_monotone} checks that the maximal positive linear economy rises from rest and stays below every supersolution. The Neumann series that produces the resolvent supersolution and the limit in the sharpness argument are written proofs. For finite atomic measures, \nolinkurl{CurrentMM.thesis_modulus_weak_duality} and \nolinkurl{CurrentMM.thesis_modulus_affine_bias} check the deterministic half of Theorem~\ref{thm:dich}(iv): the weak-duality bound on the modulus and the bias of the affine interval. These declarations do not formalize Gaussian coverage, a general-measure integration theorem, the minimax step and the lower bound of part (iv), or Theorem~\ref{thm:philog}, which rest on Sion's minimax theorem and on potential theory. The coverage register records these boundaries separately.

\paragraph{Identification under learning, by contrast.} Chapter~\ref{ch:p1}, Section~\ref{p1:sec:learners}, shows that trained online learners obey the same $\delta^{-4}$ law for separating memory rates when the competing populations' weights match, and that populations whose learned weights differ are told apart by those weights instead. Learning therefore leaves identification as hard as before while making certification only perishable, not harder per renewal.

\section{Markets of trading algorithms that learn}\label{sec:bl_ai}

Markets in which learning algorithms, including reinforcement-learning and language-model agents, trade with each other are where the results of this chapter matter most. Two of them apply directly. Certification of safety needs only a bound on total sensitivity (Theorem~\ref{thm:sector}); and the certificate lasts for its safety margin divided by the agents' learning speed (Corollary~\ref{thm:lifetime}). The second result gives a supervisory rule with no counterpart in identity-based regimes: recertify at a frequency set by how fast the agents update.

\paragraph{A registered simulation with learning agents.}
% generated by check_multiagent_v2.py; do not edit
\newcommand{\MAvTwoRuns}{200}
\newcommand{\MAvTwoPonePasses}{200}
\newcommand{\MAvTwoAdvRuns}{100}
\newcommand{\MAvTwoPtwoHits}{98}
\newcommand{\MAvTwoPtwoRate}{98.0\%}
\newcommand{\MAvTwoPtwoVerdict}{passed}

% generated by export_multiagent_v2_table.py; do not edit
\newcommand{\MAvTwoTableRows}{%
Adversarial & $0.001$ & $300$ & 25/25 & 1.27 & 25/25 & 3.97 & 2.7\\
Adversarial & $0.003$ & $100$ & 25/25 & 1.30 & 25/25 & 3.99 & 2.7\\
Adversarial & $0.010$ & $30$ & 25/25 & 1.33 & 25/25 & 4.03 & 2.9\\
Adversarial & $0.030$ & $10$ & 25/25 & 1.80 & 17/25 & 4.70 & 3.7\\
Random & $0.001$ & $300$ & 0/25 & --- & 0/25 & --- & 5.4\\
Random & $0.003$ & $100$ & 0/25 & --- & 0/25 & --- & 5.8\\
Random & $0.010$ & $30$ & 0/25 & --- & 0/25 & --- & 6.2\\
Random & $0.030$ & $10$ & 0/25 & --- & 0/25 & --- & 7.1\\}

Corollary~\ref{thm:lifetime} is a guarantee, so a simulation cannot confirm it; what a simulation can show is whether it is informative when agents learn from their own trading. We registered one (specification hash \texttt{58c32f9b\ldots}, pushed with its runner in a commit of its own before the run). The economy is that of Definition~\ref{def:cohort} with $a=0.3$, $\kappa=0.4$ and four cohorts with memory rates $0.2,0.5,0.7,0.85$, each following the saturating momentum rule $\varphi(y)=\tanh y$ (sector bound $1$). Each cohort learns its weight online by a gradient step on its own trading profit, clipped to the declared learning rate $\gamma$, so weights rise or fall but never move by more than $\gamma$ per period. Weights start at $0.05$ (total sensitivity $0.2$), the budget is $\Phi_{\max}=1.4$ (so the loss budget is $\kappa\bar w/(1-a-\kappa\Phi_{\max})=2.86$ for $\bar w=1$), and $T_{\rm valid}=\lfloor1.2/(4\gamma)\rfloor$ is $300$, $100$, $30$ and $10$ periods for $\gamma=0.001,0.003,0.01,0.03$. Shocks are bounded by $\bar w$: in the adversarial regime they push in the direction of the price displacement nine times in ten; in the random regime they are uniform. There are 25 runs for each regime and learning rate, each lasting five lifetimes.

The registered check (P1) held in \MAvTwoPonePasses{} of \MAvTwoRuns{} runs: up to $T_{\rm valid}$ every price path stayed within the loss budget and total sensitivity within $\Phi_{\max}$, as the corollary guarantees. The informative primary (P2) asked whether learning actually exhausts the sensitivity budget within three lifetimes in at least half of the adversarial runs. It \MAvTwoPtwoVerdict{}: it did in \MAvTwoPtwoHits{} of \MAvTwoAdvRuns{} runs. Table~\ref{tab:multiagent_v2} shows that under adversarial shocks the budget is exhausted after 1.3 to 1.8 lifetimes, and the realized loss first exceeds its budget after four to five. So the lifetime bound is close to tight when trend-following pays, and a certificate that is not renewed does eventually admit losses beyond its budget. Under random shocks trend-following loses money on average, weights decay towards zero and no run comes near the budget: there the bound is conservative, because it charges every period for the fastest possible learning.

\begin{table}[ht]
\centering\small
\begin{threeparttable}
\caption{Registered simulation with learning agents: certificate lifetime and realized loss}\label{tab:multiagent_v2}
\begin{tabular}{lccccccc}
\toprule
Shocks & $\gamma$ & $T_{\rm valid}$ & Budget exhausted & Median time & Loss breached & Median time & Median margin\\
\midrule
\MAvTwoTableRows
\bottomrule
\end{tabular}
\begin{tablenotes}\footnotesize\item 25 runs per row, horizon $5T_{\rm valid}$. ``Budget exhausted'': runs in which total sensitivity exceeds $\Phi_{\max}$; ``Loss breached'': runs in which $|x_t|$ exceeds the loss budget. Median times are in units of $T_{\rm valid}$, over the runs in which the event occurs. ``Median margin'': loss budget divided by the largest $|x_t|$ up to $T_{\rm valid}$. Source: \texttt{checks/R\_multiagent\_v2.json}. An earlier simulation, whose weight updates could only increase and so could not test anything, is withdrawn (Chapter~\ref{ch:conclusion}).\end{tablenotes}
\end{threeparttable}
\end{table}

\section{Agents who respond to the guard}\label{sec:bl_strategic}

A guard changes the environment it certifies. If agents learn that the operator throttles admission when measured sensitivity is high, they may change how they trade.

\begin{conjecture}[Guard-consistent equilibrium]\label{conj:strategic}
Suppose each cohort chooses its weight $q_i\in[0,\bar q_i]$ to maximize expected trading profit net of a convex cost, given an admission rule that throttles the dealer's position when certified total sensitivity exceeds $\Phi_{\max}$, and given the equilibrium reversion supplied by dealers (Proposition~\ref{p2:prop:equil}). Then (i) an equilibrium exists; (ii) in every equilibrium the guard is safe; and (iii) equilibrium total sensitivity is weakly lower than without the guard. \hfill\Conj
\end{conjecture}

\begin{remark}[Counterexample to existence under a discontinuous guard]\label{rem:strategic_counterexample}
Part (i) of the conjecture is false under the stated assumptions. Consider one cohort with $q\in[0,1]$, a perfect certificate $\Phi=q$, and admitted dealer position $I(q)=1$ for $q\le1/2$ and $I(q)=0$ for $q>1/2$. This is an admission rule that throttles above the certified threshold. Let expected trading profit, as a function of weight and admitted position, be
\[
\pi(q,I)=(1-I)(1-q),\qquad c(q)=0.
\]
The profit primitive is continuous in $(q,I)$ and the zero cost is convex. After composition with the guard, utility is zero at $q\le1/2$ and equals $1-q$ at $q>1/2$. If $q\le1/2$, the admissible deviation $q'=3/4$ strictly increases utility. If $q>1/2$, the admissible deviation $q'=(q+1/2)/2$ lies strictly between $1/2$ and $q$ and also strictly increases utility. Thus no weight maximizes utility: the one-player game has no pure equilibrium despite its compact strategy set, convex cost, and exact sensitivity certificate.

Randomization does not restore existence. Utility is strictly below $1/2$ everywhere, so its expectation under any probability distribution is strictly below $1/2$; a pure weight sufficiently close to the threshold from above earns more. The universal profitable-deviation statement and nonexistence of a pure maximizer are formalized as \nolinkurl{CurrentMM.thesis_strategic_profitable_deviation} and \nolinkurl{CurrentMM.thesis_strategic_no_equilibrium}. The probability-integral step for arbitrary mixed strategies is a written argument, not part of those Lean declarations. Exact rational diagnostics are recorded separately in \texttt{checks/R\_strategic.json}.

The failure concerns equilibrium existence, not the safety invariant. For example take $a=0$, $\kappa=1/4$, and $\bar w=1$. On admitted profiles $q\le1/2$, unit position loss is bounded by $2/7$; on throttled profiles the position is zero. Every admitted decision therefore respects this budget, although the agents' game has no equilibrium. A revised strategic theorem needs assumptions that rule out the best-response discontinuity. Convex cost alone also supplies no ordering of equilibrium sensitivity with and without a guard; that comparative-static claim remains open under suitable additional payoff restrictions.
\end{remark}

\section{What the framework does not cover}\label{sec:fw_scope}

The framework is linear in the observable layer and assumes the reset contract for its sharpest statements. Sections~\ref{sec:bl_robust} and~\ref{sec:bl_learning} extend the safety results to sector-bounded nonlinear and to learning agents; strategic equilibrium existence fails without additional assumptions (Remark~\ref{rem:strategic_counterexample}); existence for suitably regularized guards remains open. Passive observational data are slower but obey the same polynomial order for identification (Theorem~\ref{p1:prop:passive}); a finite-sample theory of certification from observational order-book data, with dependence and non-stationarity, is not developed here, and the empirical certificates of Chapter~\ref{ch:p2} are bootstrap diagnostics rather than guarantees.

\chapter{AI agents in the loop: supervising the market, not the models}\label{ch:ai}

Supervisors are now asking how markets and banks should be governed when trading and other decisions are delegated to AI agents. In October 2026 the Financial Stability Board issued sound practices for AI adoption that centre on governance inside each firm, risk-based triggers for human approval before an agent acts, and concentration risk from shared vendors \citep{t:fsb2026}; US banking agencies asked banks how they monitor and audit their agents and how they would shut them down \citep{t:reuters2026}; and the Bank for International Settlements warned that institutions using similar models, data or decision rules may act at the same time and amplify price moves, and that supervisors should look beyond individual institutions to common exposures \citep{t:bis2026}. The instruments on the table act on one firm or one agent at a time. The risk that the last warning names is a property of the loop that many agents form with prices. This chapter applies the thesis's central result to that gap: the stability of the loop can be certified from public data without inspecting any agent, and the one thing the certificate needs from the agents' owners is a bound on their memory.

\section{What the theory already says about AI agents}\label{sec:ai_map}

Four results of the earlier chapters apply to AI agents as they stand.

\paragraph{Their internals are not needed, and cannot be learned from the market.} Identifying the agents from market data costs a number of observations that grows polynomially in the inverse of their similarity (Chapter~\ref{ch:p1}), and passive observation does not help (Theorem~\ref{p1:prop:passive}). An AI agent is harder to identify than the exponential-memory cohorts of the theorems, not easier: its rule is high-dimensional, it changes when the model is updated, and, in the Financial Stability Board's words, it does not reason like a human.

\paragraph{Nonlinear and changing agents are covered.} The robust bound of Section~\ref{sec:bl_robust} holds for agents whose rules are nonlinear and whose memories and weights change over time, given a bound on their total sensitivity, and when agents keep learning a certificate lasts for its margin divided by their learning speed (Corollary~\ref{thm:lifetime}). These are the properties that distinguish an AI agent from a fixed algorithm.

\paragraph{The system-level quantity is certifiable from public data.} The margin is decided from prices and anonymous flow with a record of order $1/M$, for any number of cohorts with finite memory, uniformly as the margin vanishes, and for the pooled estimator used in practice (Theorems~\ref{thm:passivemulti}, \ref{thm:uniformcert} and~\ref{thm:pooled}).

\paragraph{The owners must disclose a memory bound.} Without one, no passive certificate exists: a slow momentum cohort can make a market unstable without leaving a trace in a finite record (Theorem~\ref{thm:hidden}).

\section{Monoculture}\label{sec:ai_mono}

Agents built on one model and fed the same data respond alike. In the cohort model they share a memory rate, so they merge into a single cohort whose weight is the sum of theirs (Theorem~\ref{p1:thm:merged}). The next proposition makes the consequence for stability explicit.

\begin{proposition}[Monoculture]\label{prop:monoculture}
Let $N$ agents trade on a market with persistence $a$ and impact $\kappa$. A share $s\in[0,1]$ of them run a common model with response $q_D$; the others have mean response $\bar q$.
\begin{enumerate}[label=(\roman*)]
\item The margin is linear in the adoption share: $M(s)=M(0)-s\,\kappa N(q_D-\bar q)$, with $M(0)=1-a-\kappa N\bar q$.
\item If $q_D>\bar q$ and $M(0)>0$, the market is stable for $s<s^*$ and unstable for $s>s^*$, where $s^*=M(0)/(\kappa N(q_D-\bar q))$.
\item The adopters form one cohort of the merged measure: adoption lowers the number of distinct cohorts, which makes identification easier, while it lowers the margin, which makes certification urgent.
\item If each adopter's decisions carry noise of variance $\sigma^2$ and adopters of one model share it, the adopters' noise contributes variance $(sN)^2\sigma^2$ rather than $sN\sigma^2$: correlated decisions enlarge the shocks that the loop amplifies by a factor equal to the number of adopters.
\end{enumerate}
\hfill\Proved
\end{proposition}

\begin{proof}
(i) The total feedback is $\Phi=N((1-s)\bar q+sq_D)$ and $M=1-a-\kappa\Phi$. (ii) $M(s)$ is affine and decreasing in $s$ with zero at $s^*$. (iii) Agents with a common model have a common memory rate, and cohorts with a common rate merge (Theorem~\ref{p1:thm:merged}). (iv) The sum of $n$ copies of one noise has variance $n^2\sigma^2$; the sum of $n$ independent ones has variance $n\sigma^2$.
\end{proof}

Part~(ii) is static: it freezes the other agents. When they learn, they respond to the price path that the adopters create, and the threshold can move. Whether it moves up, because the others learn to trade against the adopters, or down, because they learn to follow them, is a question about learning dynamics that the proposition does not answer; Test~B1 answers it in simulation. The linear threshold, the sign changes and the noise identity are checked in Lean (\nolinkurl{CurrentMM.thesis_mono_margin}, \nolinkurl{CurrentMM.thesis_mono_threshold}, \nolinkurl{CurrentMM.thesis_mono_below}, \nolinkurl{CurrentMM.thesis_mono_above}, \nolinkurl{CurrentMM.thesis_mono_common_noise}).

\section{A two-number disclosure and a loop-level trigger}\label{sec:ai_policy}

The results suggest a supervisory instrument with two parts.

\paragraph{What deployers disclose.} The owner of an AI trading agent discloses two numbers: a bound $s_i$ on the agent's sensitivity to the price displacement (the sector constant of Assumption~\ref{ass:sector}) and a bound $\bar r_i$ on its memory, the longest horizon over which its response depends on past prices. Neither reveals the model, its data or its strategy. The sum of the sensitivity bounds certifies safety by itself through the robust bound of Section~\ref{sec:bl_robust}, conservatively; the memory bound is what makes the public-data certificate possible (Theorem~\ref{thm:hidden}) and what the disclosed estimator of Test~B1 uses.

\paragraph{What the supervisor computes.} From public prices and anonymous flow the supervisor estimates the margin daily on a rolling window, with the pooled estimator augmented by the disclosed memory bound, and acts on two signals: the loss of the certificate that the margin exceeds a buffer, and a projected time to the boundary, the current margin divided by its recent rate of decline. A pause or a review is triggered by the loop, not by any one firm. Firm-level kill switches stop one agent; a trigger on the margin addresses the risk that many agents create together.

\paragraph{How it maps onto current guidance.} The Financial Stability Board's practice on human oversight asks for risk-based triggers before an agent acts; the trigger here is a market-level input to those triggers. Its practice on third-party dependencies asks firms to oversee concentration on vendors; Proposition~\ref{prop:monoculture} says what that concentration does to stability and at what adoption share it becomes dangerous. The US agencies' question of how a bank would shut its agents down has a market-level counterpart: when should all of them slow down?

\section{Test B1: learning agents and a vendor model}\label{sec:ai_b1}

% generated by b1_ailoop.py; do not edit
\newcommand{\BoneControl}{100}\newcommand{\BoneRamp}{200}
\newcommand{\BoneFA}{0}
\newcommand{\BoneEarly}{193}\newcommand{\BoneEarlyShare}{96}\newcommand{\BoneLead}{13.0}
\newcommand{\BoneGain}{17.0}\newcommand{\BonePos}{100}\newcommand{\BoneNeg}{0}
\newcommand{\BoneBelow}{200}\newcommand{\BoneBelowN}{200}
\newcommand{\BoneShareCrisis}{0.46}\newcommand{\BoneStatic}{0.82}
\newcommand{\BonePone}{passed}\newcommand{\BonePtwo}{passed}\newcommand{\BonePthree}{passed}\newcommand{\BonePfour}{passed}

\paragraph{Design.} The market is the closed loop of the thesis on five-minute bars, with 50 agents. Independent learners each have their own memory and learn their response online; with reverting prices they learn to trade against moves. A vendor model is a pre-trained momentum product with one memory rate; every agent that adopts it responds identically. Adoption ramps from none to 80\% of the agents between days 24 and 96 of a 120-day run, in two versions: a fast vendor model (memory rate $0.9$ per bar) and a slow one ($0.995$, a half-life of about eleven hours); in the control nobody adopts. The observer sees only the price displacement and the anonymous flow and runs the pooled estimator of Theorem~\ref{thm:pooled} daily on a ten-day window, alone (standard) or with the flow equation augmented by the signal at a disclosed memory bound of $0.995$ (disclosed). A crisis is the first day on which the displacement exceeds ten times its standard deviation before adoption. The warning rules were fixed in an exploratory pilot on other seeds: the certificate of a positive margin is lost; the certificate of a margin above a buffer of $0.1$ is lost; or the projected time to the boundary, the estimate divided by its rate of decline over five days, falls below ten days. Test~B1 was registered, with its code and four primaries, before its seeds (1001--1100, three scenarios, 300 runs) were run, and it ran once.

\paragraph{Result.} All four primaries passed. In the \BoneControl{} control runs the standard estimator raised no false alarm under any rule (\BoneFA{} for the time-to-boundary rule). Every one of the \BoneRamp{} runs with adoption ended in a crisis, and the time-to-boundary rule warned before it in \BoneEarly{} ($\BoneEarlyShare\%$), with a median lead of \BoneLead{} days (\BoneLeadFast{} for the fast vendor model, \BoneLeadSlow{} for the slow one). The level of the margin alone gave no warning: the standard certificate was lost on the day of the crisis (median lead \BoneLostLead{} days). With the slow vendor model, the disclosed estimator's buffer warning came earlier than the standard one in \BonePos{} of 100 runs and later in \BoneNeg{}, a median gain of \BoneGain{} days (median leads \BoneBufDiscSlow{} and \BoneBufStdSlow{} days). And in all \BoneBelowN{} runs the crisis came at an adoption share below the static threshold of Proposition~\ref{prop:monoculture}, computed with the learners frozen when adoption began: a median share of \BoneShareCrisis{} against a median threshold of \BoneStatic{}.

\paragraph{What the secondaries show.} Disclosure has a price in false alarms. The disclosed estimator is noisier: in the controls its buffer rule fired in \BoneFAdiscBuf{} of 100 runs, its lost-certificate rule in \BoneFAdiscLost{}, and its time-to-boundary rule in \BoneFAdiscTtz{}, so only the standard estimator's trend and the disclosed estimator's buffer are usable together, and the gain of Primary~3 is earned partly by a rule that would false-alarm about one time in six. The standard estimator's daily estimate differed from the true margin by a median of \BoneTrack{}, mostly as a lag of a few days behind a margin that falls quickly.

\paragraph{Reading.} Three findings matter for supervision. The certificate's level is not an early warning, but its trend is: a margin falling toward zero is visible in public data about two weeks before the crisis, with no false alarms in markets where nothing changes. A disclosed memory bound buys more warning when the shared model is slow, which is the case Theorem~\ref{thm:hidden} says prices alone cannot see, at a cost in false alarms that a supervisor would have to weigh. And learning makes monoculture dangerous earlier than static analysis predicts: as the adopters push prices, the independent learners stop trading against the move and start following it, so the market crosses the boundary at about half the adoption share at which a vendor-by-vendor assessment, freezing everyone else, would put it. That last finding is the strongest argument for supervising the loop rather than the models: the danger is created by the interaction, and an assessment of any one model, however careful, misses it.


\section{What this chapter does not show}\label{sec:ai_limits}

The agents of Test~B1 are simple learning rules and a fixed momentum product, not language-model agents or the systems banks deploy; adoption is imposed rather than chosen; the observer sees the true displacement rather than a moving-average proxy; and the warning rules and the crisis definition were chosen in an exploratory pilot before registration. The empirical chapters show that the certificate behaves as the theory says in real crypto markets, but those markets were far from the boundary and nothing in their data says which traders are AI agents. The chapter therefore offers a framework and an instrument with simulated evidence, not evidence about AI agents in real markets; Test~H1, with participant-level data, is the first step toward the latter.

\chapter{Conclusion and open problems}\label{ch:conclusion}

This thesis asked whether safety in a market with adaptive agents requires knowing who the agents are. For the markets studied here, it does not. Prices reveal the agents only through a merged measure of their memories and weights, and recovering the agents from that measure costs trading capital of order $\delta^{-4}$, experiments of order $(\Delta\eta)^{-(4K-2)}$, and with general memory an accuracy that improves only logarithmically (Chapter~\ref{ch:p1}). Safety depends on the agents only through their total feedback weight, which is certified at the parametric rate under a memory bound and cannot be certified without one; a guard built on that certificate is the largest safe admission rule (Chapter~\ref{ch:p2}). The theory of Chapter~\ref{ch:p3} places these results in one picture, in which the cost of a question is set by its behaviour at long memory, and shows that the safety bound needs only the agents' total sensitivity, even when their rules are nonlinear and their memories and weights change, and that learning makes certificates perishable at a computable rate. In data, disclosed trader classes reveal what prices do not, and the safety statistic forecasts adverse excursions in thin books, while the mechanism behind that forecast remains unidentified.

The policy message is direct. Supervisors who want to control transient risk should ask for the information that certifies safety, a bound on memory and on total sensitivity, renewed as often as the agents learn, rather than for the identities of the agents. Identity-based regimes may serve other purposes, such as enforcement and attribution, but this thesis shows they are neither necessary nor efficient for safety.

\section{Evidence that would settle the thesis}\label{sec:evidence}

The empirical work of Chapters~\ref{ch:p1} and~\ref{ch:p2} tests each paper's own predictions. This section states what that evidence establishes for the thesis as a whole, what it does not, and two tests designed to test the thesis directly. Both are registered: each specification was pushed to the public repository in a commit of its own before any of its data were downloaded. Neither has been run.

\subsection{What the existing evidence establishes}\label{sec:ev_existing}

\paragraph{Prices do not reveal the agents.} In 30 Binance perpetuals, trader-class positions pin down memory rates 28 to 124 times more precisely than prices and flow, and in a pre-registered test the CFTC's trader-class reports do the same in all eight regulated futures markets examined (Sections~\ref{p1:sec:disc_data} and~\ref{p1:sec:cot}). This is the identification half of the thesis in data: the merged measure leaves class memories essentially unrestricted.

\paragraph{The safety statistic carries information.} In thin Binance books, six-hour estimates of the amplification margin forecast the size of the next hour's adverse excursion out of sample in every pre-registered forecasting test, beyond a full volatility-forecasting baseline and VPIN; the effect is about a fifth as large in thin US stocks, negligible on Bybit and absent on Coinbase (Section~\ref{p2:sec:thin} and Table~\ref{p2:tab:alltests}).

\paragraph{What it does not establish.} The forecasting result is reduced form. Tests of a leverage or liquidation mechanism fail (Section~\ref{p2:sec:p22}), and the follow-up that points to slow reversion is close to a check of how the statistic is built (Section~\ref{p2:sec:p25}). As a risk control, the certified guard improves on a static limit but not, at matched exposure, on a volatility cap (Section~\ref{p2:sec:p24}). And from passive data the certificate itself is not always precise: in Paper~1's registered test the summed response was significantly different from zero in only 5 of 30 thin markets (Section~\ref{p1:sec:disc_data}). None of the existing tests checks the thesis's central claim directly: that, on the same data, certifying is much cheaper than identifying.

\subsection{Test E1: certify versus identify on fresh data}\label{sec:ev_e1}

\paragraph{Idea.} The theory says that the summed response is a linear functional of the merged measure and is cheap to certify, while the agents' memories are expensive to identify. In passive market data both can be estimated from the same regression of order flow on lagged returns: the sum of its coefficients is the certificate quantity $\hat\Phi$, and two exponentials fitted to the same coefficients give the memory rates. Published trader-class positions give a third, class-tagged estimate of memory. The test compares the precision of the three on data not used before.

\paragraph{Design.} Thirty Binance USD-M perpetuals (those of Section~\ref{p1:sec:disc_data}) on 105 fresh dates: the 4th, 5th, 18th, 19th and 23rd of every month from January 2025 to September 2026. A scan of every result file, notebook and script in the repository finds no earlier use of these symbol-dates. From five-minute returns, signed taker flow and the three class ratios, with Paper~1's estimators and day-bootstrap intervals, each symbol gives $I_{\rm cert}$, the width of the interval for $\hat\Phi$ relative to $|\hat\Phi|$, and $I_{\rm id}$, the width of the interval for the separation $\hat r_{\rm slow}-\hat r_{\rm fast}$ of the price-only memory rates relative to the separation. Primary~1: $I_{\rm id}\ge10\,I_{\rm cert}$, that is, certifying is at least ten times cheaper than identifying, in at least 75\% of symbols. Primary~2: the price-only memory intervals are at least ten times wider than the class-tagged ones in at least 75\% of symbols, a replication of Paper~1 on new days. Each primary must also pass a one-sided binomial test, Holm-corrected across the two. Primary~1 is a risky prediction: Paper~1's earlier test suggests that $\hat\Phi$ may be imprecise in thin markets, in which case it fails, and the failure will be reported. It would not contradict the theorems, which concern controlled experiments, but it would limit the empirical claim to what passive data can deliver.

\paragraph{Status.} Registered on 24 October 2026 (specification hash \texttt{0b6c31c0\ldots}; pipeline, download code and a Colab notebook committed with it). The notebook downloads the data, opens the test once and writes the results, the numbers printed here and a log of the SHA-256 of every downloaded file.

\subsection{Test E2: a prospective test of the mechanism}\label{sec:ev_e2}

\paragraph{Idea.} Every pre-registered test so far used archived data. A prospective test removes any doubt about the order of registration and data: the data do not exist when the specification is written.

\paragraph{Design.} The original specification (pushed on 23 October 2026) asks, for 20 thin perpetuals over at least three months after registration, whether six-hour $\hat C_3$ forecasts next-hour liquidation volume beyond volatility, impact and lagged liquidations (Primary~1), and whether adding realized next-hour liquidations to the excursion model removes at least half of $\hat C_3$'s forecasting gain (Primary~2). Binance's forced-liquidation stream, which it named as the source, cannot be recorded from the United States, and the public archive contains no liquidation files. An addendum, pushed on its own before any post-registration file was downloaded, keeps the hypotheses and replaces liquidation volume by forced deleveraging measured from open interest: the open interest destroyed during five-minute moves larger than twice their recent standard deviation. It also fixes the window (24 October 2026 to 31 January 2027), the split into 50 training and 50 test dates, and the estimators, which are those of the thin-book tests of Section~\ref{p2:sec:thin}. The price of the substitution is stated in advance: E2 can no longer test the liquidation channel with direct liquidation data, which was its advantage over the open-interest proxy of Section~\ref{p2:sec:p22}.

\paragraph{Status.} The analysis code and its notebook are committed with the addendum and refuse to run before 2 February 2027.

\subsection{Withdrawn results}\label{sec:ev_withdrawn}

Two results reported in an earlier draft of this thesis are withdrawn. The first was presented as a run of Test~E1 in 38 markets. It did not estimate anything from fresh data: it read results archived from earlier tests and used hard-coded and invented intervals, and its design compared quantities in different units. The specification it ran under is superseded by the one above, and none of its numbers appears in this thesis. The second was a simulation of learning agents whose weights, by construction, could only increase; its ``realized lifetimes'' measured that construction and could not have fallen below the bound they were said to test. Section~\ref{sec:bl_ai} replaces it with a registered simulation in which weights follow the agents' trading profits. The code and output of both are kept in the repository, marked as withdrawn, so that the record is complete.

\subsection{Rules for every new test}\label{sec:ev_rules}

\begin{enumerate}
\item The specification, its hash and an unlock counter are pushed to the public repository in a commit of their own, so the platform records when they became public, before any test data are downloaded or recorded.
\item Each test has at most two primaries with a pre-stated pass rule and a pre-stated correction across them.
\item A failed primary is reported as failed. An explanation found afterwards is labelled post hoc, and any test that follows from it is labelled as registered after the failure.
\item Every number in the thesis is produced by code that writes a result file, and an automated check compares the printed numbers with those files.
\item A result found to be wrong after it was reported is withdrawn in the text, with the reason, and its files are kept and marked rather than deleted.
\end{enumerate}

\section{Open problems}

\begin{enumerate}
\item \emph{Nonlinear questions.} Theorem~\ref{thm:dich}(iv) gives the exact cost of every linear question. Extend it to questions that are not linear in the merged measure, beginning with the peak of the impulse response, and make the link to the $1/\log N$ rate of identification (Corollary~\ref{cor:cases}(e)) a theorem rather than a reading of Chapter~\ref{ch:p1}'s proof.
\item \emph{Constants.} Theorem~\ref{thm:philog} fixes the order $\sigma^2\epsilon^{-2}\log(1/\epsilon)$ of the cost of certifying safety but not its constant, which depends on a harmonic measure; compute it, or the limit of $\omega_L(\delta)/(\delta\sqrt{\log(1/\delta)})$, as a function of the memory bound.
\item \emph{Minimal disclosure for safety.} Characterize the least informative disclosure under which a given safety question becomes certifiable from finite data. Theorem~\ref{thm:dich}(iii) says a memory bound is necessary for the summed response; is it sufficient in a sense that can be made operational?
\item \emph{Strategic agents.} Remark~\ref{rem:strategic_counterexample} refutes existence under a discontinuous admission rule. Determine sufficient assumptions for existence under a regularized guard and for the claimed comparison of equilibrium sensitivity.
\item \emph{Observational certification.} Develop finite-sample certificates from observational order-book data with dependence and non-stationarity, replacing the bootstrap diagnostics of Chapter~\ref{ch:p2}; Test~E1 measures how far passive data are from the reset contract of Theorem~\ref{thm:dich}.
\item \emph{The mechanism of thin-book risk.} Run Tests E1 and E2 of Section~\ref{sec:evidence}.
\end{enumerate}

\chapter{Conclusion}\label{ch:conclusion}

This thesis asked whether safety in a market with adaptive agents requires knowing who the agents are. For the markets studied here it does not, and the reason is structural rather than empirical. Prices reveal the agents only through a merged measure of their memories and weights, and recovering the agents from that measure costs capital of order $\delta^{-4}$, experiments of order $(\Delta\eta)^{-(4K-2)}$, and with general memory an accuracy that improves only logarithmically (Chapter~\ref{ch:p1}). Safety depends on the agents only through their total feedback weight. That one number is certified at the parametric rate up to a single logarithm under a memory bound, from passive data at a cost of order one over the stability margin, and not at all without a memory bound; a guard built on the certificate is the largest safe admission rule (Chapters~\ref{ch:p2} and~\ref{ch:p3}). The directions in which prices are least informative about the agents are directions in which the risk does not move. Certify, don't identify.

\paragraph{What the evidence supports.} The identification half holds in data: prices and order flow leave the agents' memories essentially unrestricted, on new symbols and new days, while disclosed trader-class positions restrict them. The certification half holds in a narrower form. With enough days, public data decide the sign of the stability margin in new markets whose agents they do not identify; a lower margin forecasts next-day variance beyond the standard measure of self-excitation; and in thin books the summed response forecasts the size of the next hour's adverse excursion beyond a full volatility baseline. But the decided margins rest mainly on price reversion, the markets studied were far from the boundary, the margin does not forecast adverse moves or rank damage on stress days, and a guard estimated from public data does no better than a volatility limit. Of the 23 primary hypotheses of the thesis's registered tests, 11 passed, and of the 25 of Paper~2's program, 15 passed. These are the results of one researcher's program on public data, reported with their failures, not a settled empirical case.

\paragraph{AI agents.} The question has become practical. When trading and lending decisions are delegated to AI agents built on a few shared models, the risk supervisors fear is that the agents act alike, and the monoculture proposition makes that precise: a shared model lowers the margin in proportion to its adoption share and multiplies the shocks the loop amplifies. In simulation with learning agents, the margin's trend, computed from public data alone, warned about two weeks before instability with no false alarm, while a threshold computed from a model of the agents erred by nearly a factor of two; with language-model desks it again warned before every crisis with no false alarm, but only three to four days ahead. A supervisor does not need to read the agents' code or predict how they learn; it needs a bound on their memory, renewed as often as they change, and a trigger on the loop. Test~B2 showed that the trigger's specificity survives language-model desks and its lead does not; Test~H1 will say whether real accounts on a live venue have the distinct memories the theory assumes.

\paragraph{The policy message.} Supervisors who want to control transient risk should ask for the information that certifies safety, a bound on memory and on total sensitivity, renewed as often as the agents learn, rather than for the identities or internals of the agents. Identity-based regimes serve other purposes, such as enforcement and attribution, but in the model of this thesis they are neither necessary nor efficient for safety. Passive observation of prices delivers neither kind of knowledge reliably when markets are far from the boundary; disclosure, or an outside shock that acts as an experiment, is what makes certification attainable in practice.

\section{Open problems}\label{sec:open}

\begin{enumerate}
\item \emph{Nonlinear questions.} Theorem~\ref{thm:dich}(iv) gives the exact cost of every linear question, and Proposition~\ref{prop:nonlinear} settles questions that depend on finitely many taps, such as the peak. Extend the exact theory to questions that depend nonlinearly on infinitely many taps, and make the link to the $1/\log N$ rate of identification (Corollary~\ref{cor:cases}(e)) a theorem rather than a reading of Chapter~\ref{ch:p1}'s proof.
\item \emph{Constants.} Theorem~\ref{thm:philog} gives explicit constants, but its lower constant is far below its upper one ($\PhilogLowB$ against $\PhilogUpB$ at $\bar r=0.9$). Find the limit of $\omega_L(\delta)/(\delta\sqrt{\log(1/\delta)})$; the computed brackets sit near the upper constant.
\item \emph{Minimal disclosure for safety.} Characterize the least informative disclosure under which a given safety question becomes certifiable from finite data. Theorem~\ref{thm:dich}(iii) says a memory bound is necessary for the summed response; is it sufficient in a sense that can be made operational?
\item \emph{Strategic agents.} A continuous throttle with a buffer makes every strategy profile safe and has equilibria (Proposition~\ref{prop:throttle}), and agents cannot move the certificate within the memory bound without moving the margin (Theorem~\ref{thm:nogaming}). Open: whether equilibrium sensitivity is lower with the guard than without it, and how to supervise agents that can switch their response abruptly at a time of their choosing.
\item \emph{Observational certification.} Develop finite-sample certificates from observational order-book data with dependence and non-stationarity, replacing the bootstrap diagnostics of Chapter~\ref{ch:p2}. Test~E1 found passive five-minute data far from the reset contract of Theorem~\ref{thm:dich}; natural experiments (Proposition~\ref{prop:natural}) are the nearest substitute, and their sampling theory under dependent shocks is open. Theorems~\ref{thm:passivecert}, \ref{thm:passivemulti} and~\ref{thm:uniformcert} settle the cost of passive certification for one closed-loop root and, uniformly as the margin vanishes, for any number of cohorts with finite memory; Theorem~\ref{thm:hidden} shows that without a memory bound no passive certificate exists. The pooled price-and-flow estimator that the tests use attains the same order (Theorem~\ref{thm:pooled}). What remains is a version uniform over geometric memory structures and over the pooled estimator's nuisance parameters, and the gap between the moving-average displacement the tests use and the distance from the fundamental value the theory uses.
\item \emph{The evidence.} Run Test~E2 after 2 February 2027; measure the feedback term itself, not only the margin, from disclosure or outside shocks; and extend the cohort model to contrarian agents, which the exploration finds to be common.
\end{enumerate}


\appendix
\chapter{Status of results and notation across chapters}\label{app:status}

\section{Status register}

\begin{longtable}{p{0.47\textwidth}p{0.22\textwidth}p{0.2\textwidth}}
\toprule
Result & Location & Status\\
\midrule
\endhead
\multicolumn{3}{l}{\emph{Identification (Chapter~\ref{ch:p1})}}\\
Merged rate--weight equivalence & Thm~\ref{p1:thm:merged} & proved\\
Executed quantity $\Theta(\delta^{-2})$ & Thm~\ref{p1:thm:quantity} & proved\\
Capital $\Theta(\delta^{-4})$ & Thm~\ref{p1:thm:cap} & proved\\
Spectral round-trip loss bound; adaptive probing; finite-horizon dual & Thms~\ref{p1:thm:spectral}, \ref{p1:thm:adaptive}, \ref{p1:thm:rh} & proved\\
$K$ cohorts: $(\Delta\eta)^{-(4K-2)}$, lower and upper & Thms~\ref{p1:thm:kcohort}, \ref{p1:prop:esprit} & proved\\
General memory: $\Theta(1/\log N)$ & Thm~\ref{p1:thm:general} & proved\\
Passive data: same polynomial order & Thm~\ref{p1:prop:passive} & proved\\
Tagged flow removes the barrier; value of disclosure & Prop.~\ref{p1:prop:tagged}, Thm~\ref{p1:thm:discvalue} & proved\\
Spectral blind spot & Thm~\ref{p1:thm:blind} & proved\\
Trained learners: $\delta^{-4}$ when weights match & Sec.~\ref{p1:sec:learners} & simulation\\
Class disclosure pins memories (Binance, CFTC) & Secs.~\ref{p1:sec:disc_data}, \ref{p1:sec:cot} & data, pre-registered\\
\midrule
\multicolumn{3}{l}{\emph{Safety (Chapter~\ref{ch:p2})}}\\
Closed form of $S$ & Thm~\ref{p2:thm:sum} & proved\\
Exact overshoot condition & Thm~\ref{p2:thm:peak} & proved\\
Spectral stability is not solvency & Thm~\ref{p2:thm:radius} & proved\\
Summed-gain guard is maximal & Lemma~\ref{p2:thm:guard} & proved\\
Invariance to cohort splits & Prop.~\ref{p2:prop:invariance} & proved\\
Finite-window certification and its limit & Thm~\ref{p2:thm:window} & proved\\
Certification at the parametric rate & Thm~\ref{p2:thm:certcost} & proved\\
Optimal uncertainty-aware guard & Thm~\ref{p2:thm:robustguard} & proved\\
Value of disclosure for safe capacity & Thm~\ref{p2:thm:disclosure} & proved\\
Equilibrium under-supply of reversion & Prop.~\ref{p2:prop:equil} & proved\\
Carried inventory; hidden drift & Props.~\ref{p2:prop:inventory_S}, \ref{p2:prop:driftguard} & proved\\
Thin-book forecasting; guard backtests; mechanism tests & Secs.~\ref{p2:sec:thin}, \ref{p2:sec:istests} & data, pre-registered\\
\midrule
\multicolumn{3}{l}{\emph{Paper 3 (Chapter~\ref{ch:p3})}}\\
Upper bound on certification cost for linear questions & Thm~\ref{thm:dich}(i) & proved here\\
Truncation necessity; impossibility at infinite memory & Thm~\ref{thm:dich}(ii)--(iii) & proved here\\
Matching lower bound & Thm~\ref{thm:dich}(iv) & target; uniform version refuted\\
Fixed-atom parametric obstruction & Sec.~\ref{sec:dich_partial} & written argument; Lean pending\\
Robust safety bound (nonlinear, time-varying agents) & Thm~\ref{thm:sector} & proved here\\
Several assets (induced-norm bound) & Prop.~\ref{prop:multiasset} & proved here\\
Certificate lifetime & Cor.~\ref{thm:lifetime} & proved here\\
Guard-consistent equilibrium & Conj.~\ref{conj:strategic} & conjecture\\
Learning-agent simulation & Sec.~\ref{sec:bl_ai} & pending data\\
Ground-truth test E1; prospective test E2 & Sec.~\ref{sec:evidence} & pending data\\
\bottomrule
\end{longtable}

\section{Notation across chapters}

\begin{longtable}{p{0.3\textwidth}p{0.3\textwidth}p{0.3\textwidth}}
\toprule
Chapter~\ref{ch:p3} & Chapter~\ref{ch:p1} & Chapter~\ref{ch:p2}\\
\midrule
memory rate $r_i$ & $1-\eta_i$ & $1-\eta_i$\\
weight $q_i$ & $q_i$ & $q_i$\\
persistence $a$ & $a$ & $a=1-m$ (reversion $m$)\\
impact $\kappa$ & $\kappa$ & $\kappa$, $\tilde\kappa$ with depth\\
merged measure $M=\sum_iq_i(1-r_i)\delta_{r_i}$ & $M=\sum_iq_i\eta_i\delta_{1-\eta_i}$ & taps $k_m$ of $M$\\
total weight $\Phi=\int(1-r)^{-1}dM$ & $\sum_iq_i$ & $q$ (aggregate)\\
summed response $S=\kappa/(1-a-\kappa\Phi)$ & --- & $S=\tilde\kappa/(m-q\tilde\kappa)$\\
\bottomrule
\end{longtable}




% BEGIN ECONOMIC REVIEW 20261029
\chapter{Development revision and validation status}
This is a working draft for author and specialist review. The revision corrects
hidden-state drift accounting, distinguishes cumulative displacement from peak
mark-to-market loss, and restricts claims to their stated model and observation
contracts. It does not declare every theorem independently verified or every
research objective complete.

\paragraph{Returned execution record.} The author returned a completed full-mode
notebook archive using source commit \texttt{19805d7e}, seed 20261029,
Python 3.13.18 and NumPy 2.1.3. All 10 exported-file hashes and the four analytical
source modules were checked. The run completed E1, E2, E3, E4 and E6; its log
reports 18 targeted tests passing. This is an execution reproduction of the same
seeded code, not independent scientific validation. No real-market CSV was
provided, and the GPU, known-source-feedback and prospective studies were not run.

\paragraph{What the checks establish.} E1 cross-checks the corrected drift identity
against a resolvent and exercises a sufficient invariant-box guard. E2 demonstrates
filter-induced persistence under martingale nulls. E3 compares known alternatives
with equal aggregate risk; it is an illustrative witness, not a uniform rate
theorem. E4 matches baseline capacity on training episodes and reports test
capacity differences; the guard and volatility control both have zero observed
budget breaches in the three bounded-shock synthetic regimes. This does not
establish universal guard superiority or economic welfare.

\paragraph{Coverage and the observation contract.} E6 uses 1,000 replications,
32 reset responses and 512 taps. The nominal 95\% interval assumes independent
Gaussian tap errors and a valid memory envelope. The intervals in the table
measure Monte Carlo uncertainty about coverage, not financial risk.
\begin{center}\small
\begin{tabular}{lrrc}
\toprule
Setting & Coverage & MC 95\% interval & Contract holds\\
\midrule
Independent Gaussian & 94.6\% & [93.0, 95.8] & yes\\
Correlated tap errors & 48.3\% & [45.2, 51.4] & no\\
Heavy-tailed readings & 95.1\% & [93.6, 96.3] & no\\
Biased readings & 25.6\% & [23.0, 28.4] & no\\
Unseen slow memory & 0.0\% & [0.0, 0.4] & no\\
\bottomrule
\end{tabular}
\end{center}
Failures under correlated errors, biased readings and a false memory floor expose
the contract's limits. The heavy-tail outcome in this configuration is not a
general robustness result. All outcomes, including negative ones, are retained.

\paragraph{Remaining obligations.} Independent review of the load-bearing
identification and certification proofs, continuous numerical extrema, economic
primitives, the public-observation bridge and policy counterfactuals remains
necessary. The current two-paper name-map audit finds 54 of 57 formal environments
with existing declaration mappings, three unmapped statements and seven revised
statements requiring semantic review. A label mapping is not a new Lean proof.
The experiment design and full results are retained under
\texttt{studies/economic\_review\_20261029/}.

% END ECONOMIC REVIEW 20261029

% BEGIN INFORMATION DECISION REVISION 20261029
\chapter{Economic decision and returned information-decision experiment}

\subsection*{What this revision adds}
The information question is distinct from choosing a preferred amount of risk.
For a specified financial decision, the relevant issue is which observations
support an exposure constraint, what remains unidentified, and whether obtaining
additional information changes the attainable decision after costs. This does
not introduce portfolio choice or control without full identification as new
ideas. The general identification-versus-decision-value result still requires
independent proof and a comparison with the closest existing results.

\subsection*{Time conventions and the institution}
The new application is a price-taking, unlevered own-account investor, not a
venue-wide supervisor. It begins with cash, buys once, holds a fixed position
for 60 one-minute model periods, and closes fully. Round-trip linear and
quadratic charges enter cash and wealth. The account cannot change the market's
feedback, and the experiment does not measure customer liquidity or welfare.
This is separate from the paper's dealer/venue admission contracts.

\begin{center}\small
\begin{tabular}{p{0.26\textwidth}p{0.65\textwidth}}
\toprule
Quantity & Interpretation\\
\midrule
Sampling interval & One minute in the new synthetic account; other historical studies retain their stated sampling clocks.\\
Historical lookback & Often six hours in the existing forecasting studies; not a holding period.\\
Forecast horizon & Often the next hour's adverse excursion; not an executed round-trip trade.\\
New holding period & Exactly 60 model minutes, with explicit entry and exit charges.\\
Other theory horizons & The existing finite-horizon theorems retain their declared $H$; no universal holding period is imposed.\\
\bottomrule
\end{tabular}
\end{center}

\subsection*{Information and decision design}
The new comparison uses controlled independent reset observations of 12 impulse
taps, a known plant and reading variance, and a declared finite candidate class.
It does not treat public market trading as a reset experiment. Common training
and validation paths supply 768 and 384 additional reset episodes per repetition;
512 further paths evaluate the selected actions. These are in addition to the
reported tap-reading budget, not evidence of learning from 32 observations in
total. Verified disclosure restricts
the minimum-memory rate to an audited band. The oracle knows the true response,
not future shocks. All certificate arms use the same training objective.

The account selects a position by empirical expected logarithmic terminal
wealth, subject to its information-supported exposure restriction. This is a
finite-horizon application of logarithmic utility, not a newly proved Kelly
growth-optimality theorem. Equal response bounds alone do not imply equal
expected returns or identical utility-optimal portfolios.

The new account has initial wealth of 1,000 dollars and an initial price of
100 dollars; its position grid runs from zero to 9.9 units. Its round-trip
dollar charge is $c(I)=0.002I+0.0002I^2$. For bounded external flow, the
constraint adds this charge to position times the largest admissible absolute
price response over the holding horizon. The resulting bound concerns loss
from initial wealth at each date. It is not peak-to-trough drawdown, protection
against arbitrary unbounded shocks, or a rolling deployment guarantee.

\subsection*{Returned execution and proof scope}
The returned full-mode archive uses seed 20261029. All 35 manifest entries and
eight core source files were checked. The run contains 20 independent
repetitions per regime, not 16,800 independent experiments: the 16,800 policy
rows reuse tapes across policies, budgets and observation counts. Twelve
targeted checks passed. Eight conditional decision-logic lemmas compiled with
Lean 4.19.0 and their axiom printouts report no axioms. This is a seeded
execution reproduction. It neither completes the older manuscript's Lean
coverage nor formalizes the Gaussian confidence argument, market dynamics,
general minimax separation, portfolio optimality or welfare.

The displayed slice uses 32 reset episodes and a one-dollar loss budget from
initial wealth of 1,000 dollars. It is an illustrative presentation of an
exploratory grid, not a primary selected before the results were seen. All
regimes, budgets, policies and negative outcomes remain in the result archive.
The inputs and complete outputs are retained with the verification record in
the new information-decision study directory.

\subsection*{Aggregate information and unresolved cohorts: returned simulation}
\begin{center}\small
\begin{tabular}{lrrr}
\toprule
Regime & Unique candidate (\%) & Certified units & Oracle units\\
\midrule
No feedback & 0.0 & 9.405 & 9.900\\
Moderate feedback & 5.0 & 5.000 & 5.000\\
Near boundary & 0.0 & 1.200 & 1.200\\
\bottomrule
\end{tabular}
\end{center}
Here ``unique'' means one remaining candidate in the declared finite class,
not identification of arbitrary real-world algorithms. In the moderate and
near-boundary configurations, the aggregate certificate reaches the oracle's
grid capacity in this slice even though the candidate decomposition is usually
unresolved. This is an illustrative decision comparison; it does not prove a
uniform rate separation for continuous classes.

\subsection*{Financial decisions: returned simulation}
{\footnotesize\setlength{\tabcolsep}{4pt}\renewcommand{\arraystretch}{0.9}
\begin{longtable}{p{0.36\textwidth}rrrr}
\toprule
Policy & Units & Log growth (bp) & ES95 (\$) & Breach (\%)\\
\midrule
\endfirsthead
\toprule
Policy & Units & Log growth (bp) & ES95 (\$) & Breach (\%)\\
\midrule
\endhead
\multicolumn{5}{l}{\textit{No feedback}}\\
Aggregate certificate & 9.405 & 0.578 & 0.394 & 0.000\\
Unpaid disclosure & 9.405 & 0.578 & 0.394 & 0.000\\
Paid disclosure & 9.405 & 0.478 & 0.404 & 0.000\\
Oracle response & 9.900 & 0.614 & 0.415 & 0.000\\
Volatility and persistence & 9.900 & 0.614 & 0.415 & 0.000\\
Point response & 9.900 & 0.614 & 0.415 & 0.000\\
Expected-log objective & 9.900 & 0.614 & 0.415 & 0.000\\
Fractional log, validation & 9.900 & 0.614 & 0.415 & 0.000\\
Empirical loss constraint & 9.900 & 0.614 & 0.415 & 0.000\\
Robust empirical log & 0.000 & 0.000 & 0.000 & 0.000\\
\midrule
\multicolumn{5}{l}{\textit{Moderate feedback}}\\
Aggregate certificate & 5.000 & 1.330 & 0.234 & 0.000\\
Unpaid disclosure & 5.000 & 1.330 & 0.234 & 0.000\\
Paid disclosure & 4.900 & 1.205 & 0.240 & 0.000\\
Oracle response & 5.000 & 1.330 & 0.234 & 0.000\\
Volatility and persistence & 8.550 & 2.214 & 0.404 & 0.000\\
Point response & 5.000 & 1.330 & 0.234 & 0.000\\
Expected-log objective & 9.900 & 2.537 & 0.469 & 0.000\\
Fractional log, validation & 9.900 & 2.537 & 0.469 & 0.000\\
Empirical loss constraint & 9.900 & 2.537 & 0.469 & 0.000\\
Robust empirical log & 5.195 & 1.374 & 0.244 & 0.000\\
\midrule
\multicolumn{5}{l}{\textit{Near boundary}}\\
Aggregate certificate & 1.200 & 1.386 & 0.084 & 0.000\\
Unpaid disclosure & 1.200 & 1.386 & 0.084 & 0.000\\
Paid disclosure & 1.200 & 1.286 & 0.094 & 0.000\\
Oracle response & 1.200 & 1.386 & 0.084 & 0.000\\
Volatility and persistence & 2.555 & 2.943 & 0.180 & 0.000\\
Point response & 1.200 & 1.386 & 0.084 & 0.000\\
Expected-log objective & 9.900 & 11.252 & 0.705 & 0.195\\
Fractional log, validation & 9.900 & 11.252 & 0.705 & 0.195\\
Empirical loss constraint & 9.900 & 11.252 & 0.705 & 0.195\\
Robust empirical log & 8.730 & 9.941 & 0.620 & 0.078\\
\midrule
\bottomrule
\end{longtable}}
Values average 20 repetitions. Log growth is $10^4$ times the mean logarithm
of terminal wealth divided by initial wealth; ES95 is the mean of each
repetition's worst five percent of initial-wealth losses. It is not a return
forecast or a peak-to-trough drawdown guarantee. Position sizes differ across
policies; the table reports a risk--growth tradeoff rather than equal-exposure
dominance. A zero observed breach is not proof of zero breach probability.

\subsection*{Disclosure cost and the negative result}
\begin{center}\small
\begin{tabular}{lrrr}
\toprule
Paid minus aggregate & Mean (bp) & Lower (bp) & Upper (bp)\\
\midrule
No feedback & -0.100 & -0.100 & -0.100\\
Moderate feedback & -0.126 & -0.126 & -0.125\\
Near boundary & -0.100 & -0.100 & -0.100\\
\bottomrule
\end{tabular}
\end{center}
The intervals are exploratory percentile bootstraps over repetition-level
paired means, not simultaneous intervals across the grid. In this slice the
unpaid disclosure does not expand capacity, while the assumed one-cent charge
reduces net log wealth. This run therefore does not establish a positive net
economic value of disclosure. The price is hypothetical, not a calibrated
reporting cost. Stronger disclosure need not be worth purchasing when aggregate
information already supports the relevant decision.

\subsection*{Assumption failures and calibration diagnostics}
The table reports finite-sample frequencies for all seven regimes, combining
the three reading budgets. Each band outcome is counted once per repetition
and reading budget (60 outcomes per regime); repeating the same outcome across
four financial loss budgets would not create further observations. The three
reading budgets within a repetition are not an independent replication.
\begin{center}\small
\begin{tabular}{lcrr}
\toprule
Regime & Contract valid & Tap-band hit (\%) & Empty set (\%)\\
\midrule
No feedback & Yes & 96.7 & 3.3\\
Moderate feedback & Yes & 93.3 & 5.0\\
Near boundary & Yes & 96.7 & 3.3\\
Hidden slow memory & No & 98.3 & 11.7\\
Correlated readings & No & 98.3 & 1.7\\
Biased readings & No & 0.0 & 100.0\\
Behaviour change & No & 95.0 & 5.0\\
\bottomrule
\end{tabular}
\end{center}
The nominal one-decision band coverage is 95 percent under the valid contract;
these small, dependent-budget frequencies are diagnostics, not a new coverage
theorem. A tap-band hit concerns the observed response, not membership in the
finite candidate class or stability of future behaviour. In particular, hidden
slow memory falls outside the declared class even when its first 12 taps fit
the bands. Under behaviour change the readings describe the old response.
Biased readings empty the set throughout this run and force the certificate
policy to abstain. Correlated readings use an unverified standard-error model.
None of these four regimes retains the stated guarantee; a high observed band
hit or zero observed financial breaches does not make an invalid contract valid.

\subsection*{Pending evidence: deliberately empty result cells}
The following numerical cells remain empty. The returned archive contained no
market CSV, dynamic-inventory experiment, institution-level intervention or
external replication. The general theorem and independent novelty review also
remain open; a table placeholder is not a theorem statement or a proof.

\begin{center}\small
\begin{tabular}{p{0.62\textwidth}cc}
\toprule
Pending study & Estimate & Uncertainty\\
\midrule
Untouched market-data test beyond order-flow imbalance and depth & & \\[1.8ex]
Full changing-inventory, execution and funding comparison & & \\[1.8ex]
Dealer/venue intervention with endogenous market response & & \\[1.8ex]
Net value under empirically calibrated disclosure costs & & \\[1.8ex]
Independent implementation and replication & & \\[1.8ex]
\bottomrule
\end{tabular}
\end{center}
\noindent Further data cannot by themselves prove a general identification or
decision-value theorem. Its exact statement, proof and literature priority
require separate review. No publication rating follows from the returned run.

% END INFORMATION DECISION REVISION 20261029

\chapter*{Bibliography for Chapters~1, 4 and 5}
\addcontentsline{toc}{chapter}{Bibliography for Chapters 1, 4 and 5}
\begingroup
\renewcommand{\bibsection}{}
\begin{thebibliography}{99}\setlength{\itemsep}{2pt}\small
\bibitem[Brunnermeier and Pedersen(2009)]{t:brunnermeier2009} Brunnermeier, M.~K. and Pedersen, L.~H. (2009). Market liquidity and funding liquidity. \emph{Review of Financial Studies} 22(6), 2201--2238.
\bibitem[Cai and Low(2011)]{t:cai2011} Cai, T.~T. and Low, M.~G. (2011). Testing composite hypotheses, Hermite polynomials and optimal estimation of a nonsmooth functional. \emph{Annals of Statistics} 39(2), 1012--1041.
\bibitem[Cand\`es and Fernandez-Granda(2014)]{t:candes2014} Cand\`es, E.~J. and Fernandez-Granda, C. (2014). Towards a mathematical theory of super-resolution. \emph{Communications on Pure and Applied Mathematics} 67(6), 906--956.
\bibitem[Farina and Rinaldi(2000)]{t:farina2000} Farina, L. and Rinaldi, S. (2000). \emph{Positive Linear Systems: Theory and Applications}. Wiley.
\bibitem[Heinrich and Kahn(2018)]{t:heinrich2018} Heinrich, P. and Kahn, J. (2018). Strong identifiability and optimal minimax rates for finite mixture estimation. \emph{Annals of Statistics} 46(6A), 2844--2870.
\bibitem[Kirilenko et~al.(2017)]{t:kirilenko2017} Kirilenko, A., Kyle, A.~S., Samadi, M. and Tuzun, T. (2017). The Flash Crash: High-frequency trading in an electronic market. \emph{Journal of Finance} 72(3), 967--998.
\bibitem[Lepski et~al.(1999)]{t:lepski1999} Lepski, O., Nemirovski, A. and Spokoiny, V. (1999). On estimation of the $L_r$ norm of a regression function. \emph{Probability Theory and Related Fields} 113(2), 221--253.
\bibitem[Ljung(1999)]{t:ljung1999} Ljung, L. (1999). \emph{System Identification: Theory for the User}, 2nd edition. Prentice Hall.
\bibitem[Manski(2003)]{t:manski2003} Manski, C.~F. (2003). \emph{Partial Identification of Probability Distributions}. Springer.
\bibitem[Megretski and Rantzer(1997)]{t:megretski1997} Megretski, A. and Rantzer, A. (1997). System analysis via integral quadratic constraints. \emph{IEEE Transactions on Automatic Control} 42(6), 819--830.
\bibitem[Moitra(2015)]{t:moitra2015} Moitra, A. (2015). Super-resolution, extremal functions and the condition number of Vandermonde matrices. In \emph{Proceedings of the 47th Annual ACM Symposium on Theory of Computing}, 821--830.
\bibitem[Rantzer(2015)]{t:rantzer2015} Rantzer, A. (2015). Scalable control of positive systems. \emph{European Journal of Control} 24, 72--80.
\bibitem[Wu and Yang(2020)]{t:wu2020} Wu, Y. and Yang, P. (2020). Optimal estimation of Gaussian mixtures via denoised method of moments. \emph{Annals of Statistics} 48(4), 1981--2007.
\end{thebibliography}

\endgroup

\end{document}
